% ============================================================
%  MODUL PEMBELAJARAN MATEMATIKA SMA
%  Kurikulum Merdeka
%  BAGIAN 2: Fase F (Kelas 11 & 12) + Matematika Tingkat Lanjut (Fase F+)
%  Kompilasi: pdfLaTeX / XeLaTeX (disarankan Overleaf)
% ============================================================
\documentclass[11pt,a4paper]{report}

% ---------- PAKET ----------
\usepackage[utf8]{inputenc}
\usepackage[T1]{fontenc}
\usepackage[bahasa]{babel}
\usepackage{amsmath,amssymb,amsthm,mathtools}
\usepackage{geometry}
\usepackage{xcolor}
\usepackage{graphicx}
\usepackage{enumitem}
\usepackage{tikz}
\usetikzlibrary{calc,arrows.meta,angles,quotes,positioning}
\usepackage{pgfplots}
\pgfplotsset{compat=1.18}
\usepackage[most]{tcolorbox}
\usepackage{fancyhdr}
\usepackage{titlesec}
\usepackage{array,booktabs}
\usepackage[colorlinks=true,linkcolor=biru,urlcolor=biru]{hyperref}

\geometry{a4paper,margin=2.5cm,headheight=15pt}

% ---------- WARNA ----------
\definecolor{biru}{HTML}{1F4E79}
\definecolor{birumuda}{HTML}{D6E4F0}
\definecolor{hijau}{HTML}{2E7D32}
\definecolor{hijaumuda}{HTML}{E3F2E1}
\definecolor{oranye}{HTML}{C75B12}
\definecolor{oranyemuda}{HTML}{FBE8DA}
\definecolor{abu}{HTML}{F2F2F2}
\definecolor{ungu}{HTML}{6A1B9A}
\definecolor{ungumuda}{HTML}{EDE2F3}

% ---------- HEADER / FOOTER ----------
\pagestyle{fancy}
\fancyhf{}
\renewcommand{\headrulewidth}{0.8pt}
\renewcommand{\footrulewidth}{0pt}
\fancyhead[L]{\small\textcolor{biru}{\leftmark}}
\fancyhead[R]{\small\textcolor{biru}{Matematika SMA — Fase F \& F+}}
\fancyfoot[C]{\thepage}

% ---------- GAYA JUDUL BAB ----------
\titleformat{\chapter}[display]
  {\normalfont\huge\bfseries\color{biru}}
  {\filright\large BAB \thechapter}{8pt}{\Huge}
\titlespacing*{\chapter}{0pt}{0pt}{20pt}

% ---------- LINGKUNGAN TEOREMA ----------
\newtheoremstyle{gaya}{6pt}{6pt}{}{}{\bfseries\color{biru}}{.}{.5em}{}
\theoremstyle{gaya}
\newtheorem{definisi}{Definisi}[chapter]
\newtheorem{sifat}{Sifat}[chapter]
\newtheorem{teorema}{Teorema}[chapter]

% ---------- KOTAK KHUSUS ----------
\newtcolorbox{tujuan}{
  enhanced, colback=birumuda, colframe=biru, boxrule=0.6pt,
  arc=3pt, left=10pt, right=10pt, top=8pt, bottom=8pt,
  title=\textbf{Tujuan Pembelajaran}, fonttitle=\bfseries\color{white},
  coltitle=white, attach boxed title to top left={xshift=10pt,yshift=-10pt},
  boxed title style={colback=biru,arc=2pt}}

\newtcolorbox{konsep}[1]{
  enhanced, breakable, colback=hijaumuda, colframe=hijau, boxrule=0.6pt,
  arc=3pt, left=10pt, right=10pt, top=6pt, bottom=6pt,
  title={#1}, fonttitle=\bfseries\color{white}, coltitle=white,
  attach boxed title to top left={xshift=10pt,yshift=-10pt},
  boxed title style={colback=hijau,arc=2pt}}

\newtcolorbox{contoh}[1]{
  enhanced, breakable, colback=oranyemuda, colframe=oranye, boxrule=0.6pt,
  arc=3pt, left=10pt, right=10pt, top=6pt, bottom=6pt,
  title={\textbf{Contoh: #1}}, fonttitle=\bfseries\color{white}, coltitle=white,
  attach boxed title to top left={xshift=10pt,yshift=-10pt},
  boxed title style={colback=oranye,arc=2pt}}

\newtcolorbox{latihan}{
  enhanced, breakable, colback=ungumuda, colframe=ungu, boxrule=0.6pt,
  arc=3pt, left=10pt, right=10pt, top=8pt, bottom=8pt,
  title=\textbf{Latihan Soal}, fonttitle=\bfseries\color{white}, coltitle=white,
  attach boxed title to top left={xshift=10pt,yshift=-10pt},
  boxed title style={colback=ungu,arc=2pt}}

\newcommand{\jawab}{\par\smallskip\noindent\textit{\textcolor{oranye}{Penyelesaian:}}\ }

% ---------- WARNA TAMBAHAN ----------
\definecolor{merah}{HTML}{C62828}
\definecolor{merahmuda}{HTML}{FCE4E4}
\definecolor{tosca}{HTML}{00838F}
\definecolor{toscamuda}{HTML}{E0F2F1}
\definecolor{emas}{HTML}{8D6E00}
\definecolor{emasmuda}{HTML}{FFF6D6}

\usepackage{multicol}
\usepackage{pifont} % untuk \ding (checkbox & bintang)

% ---------- KOTAK PEDAGOGIS TAMBAHAN ----------
% Pertanyaan Pemantik
\newtcolorbox{pemantik}{
  enhanced, breakable, colback=emasmuda, colframe=emas, boxrule=0.6pt,
  arc=3pt, left=10pt, right=10pt, top=8pt, bottom=8pt,
  title=\textbf{Pertanyaan Pemantik}, fonttitle=\bfseries\color{white}, coltitle=white,
  attach boxed title to top left={xshift=10pt,yshift=-10pt},
  boxed title style={colback=emas,arc=2pt}}

% Studi Kasus Kontekstual
\newtcolorbox{studikasus}[1]{
  enhanced, breakable, colback=abu, colframe=biru, boxrule=0.6pt,
  arc=3pt, left=10pt, right=10pt, top=8pt, bottom=8pt,
  title={\textbf{Studi Kasus: #1}}, fonttitle=\bfseries\color{white}, coltitle=white,
  attach boxed title to top left={xshift=10pt,yshift=-10pt},
  boxed title style={colback=biru,arc=2pt}}

% Kekeliruan yang Sering Terjadi
\newtcolorbox{kesalahan}{
  enhanced, breakable, colback=merahmuda, colframe=merah, boxrule=0.6pt,
  arc=3pt, left=10pt, right=10pt, top=8pt, bottom=8pt,
  title=\textbf{Kekeliruan yang Sering Terjadi}, fonttitle=\bfseries\color{white}, coltitle=white,
  attach boxed title to top left={xshift=10pt,yshift=-10pt},
  boxed title style={colback=merah,arc=2pt}}

% Matematika dalam Dunia Nyata
\newtcolorbox{dunianyata}{
  enhanced, breakable, colback=toscamuda, colframe=tosca, boxrule=0.6pt,
  arc=3pt, left=10pt, right=10pt, top=8pt, bottom=8pt,
  title=\textbf{Matematika dalam Dunia Nyata}, fonttitle=\bfseries\color{white}, coltitle=white,
  attach boxed title to top left={xshift=10pt,yshift=-10pt},
  boxed title style={colback=tosca,arc=2pt}}

% Refleksi Diri
\newtcolorbox{refleksi}{
  enhanced, breakable, colback=toscamuda, colframe=tosca, boxrule=0.6pt,
  arc=3pt, left=10pt, right=10pt, top=8pt, bottom=8pt,
  title=\textbf{Refleksi Diri}, fonttitle=\bfseries\color{white}, coltitle=white,
  attach boxed title to top left={xshift=10pt,yshift=-10pt},
  boxed title style={colback=tosca,arc=2pt}}

% Challenge Corner
\newtcolorbox{challenge}{
  enhanced, breakable, colback=ungumuda, colframe=ungu, boxrule=1pt,
  arc=3pt, left=10pt, right=10pt, top=8pt, bottom=8pt,
  title=\textbf{\ding{72}\ Challenge Corner}, fonttitle=\bfseries\color{white}, coltitle=white,
  attach boxed title to top left={xshift=10pt,yshift=-10pt},
  boxed title style={colback=ungu,arc=2pt}}

% Catatan Penting (sidenote ringkas)
\newtcolorbox{catatan}{
  enhanced, breakable, colback=white, colframe=biru, boxrule=0.6pt,
  arc=2pt, left=8pt, right=8pt, top=5pt, bottom=5pt,
  borderline west={3pt}{0pt}{biru}}

% Penanda bagian soal akhir bab
% (argumen kedua sengaja tidak ditampilkan; cukup "Bagian A", "Bagian B", dst.)
\newcommand{\bagiansoal}[2]{%
  \par\medskip\noindent
  {\large\bfseries\color{biru}Bagian #1}\par\smallskip
  \addcontentsline{toc}{subsection}{Bagian #1}}

% Checklist refleksi
\newcommand{\cek}{\ding{113}\ } % kotak kosong

% ---------- FORMAT KUNCI JAWABAN ----------
% Nomor soal pada kunci jawaban
\newcommand{\knomor}[1]{\par\medskip\noindent{\bfseries\color{ungu}#1}\par\nopagebreak}
% Tiga bagian wajib: Soal -- Analisis -- Pembahasan
\newcommand{\ksoal}{\noindent\textbf{\color{biru}Soal.}\space}
\newcommand{\kanalisis}{\par\noindent\textbf{\color{oranye}Analisis.}\space}
\newcommand{\kbahas}{\par\noindent\textbf{\color{hijau}Pembahasan.}\space}
% Sub-judul bagian (A/B/C/D/Challenge) pada kunci
\newcommand{\kbagian}[1]{\par\medskip\noindent{\large\bfseries\color{biru}#1}\par\smallskip}

% ============================================================
\begin{document}

% ---------- HALAMAN JUDUL ----------
\begin{titlepage}
\centering
\vspace*{2cm}
{\color{biru}\rule{\textwidth}{2pt}}\\[0.6cm]
{\Huge\bfseries\color{biru} MODUL PEMBELAJARAN\\[0.3cm] MATEMATIKA SMA}\\[0.4cm]
{\Large Kurikulum Merdeka — Bagian 2}\\[0.2cm]
{\color{biru}\rule{\textwidth}{2pt}}\\[1.2cm]

{\LARGE\bfseries FASE F \& MATEMATIKA TINGKAT LANJUT}\\[0.3cm]
{\Large Kelas 11 \& 12}\\[2cm]

\vfill
{\large Tahun Pelajaran 2026/2027}
\end{titlepage}

\tableofcontents
\thispagestyle{empty}
\clearpage

% ============================================================
\part*{\centering FASE F --- KELAS 11}
\addcontentsline{toc}{part}{FASE F --- KELAS 11}

% ============================================================
\chapter{Komposisi Fungsi dan Fungsi Invers}
% ============================================================
\begin{tujuan}
Setelah mempelajari bab ini, peserta didik mampu:
\begin{itemize}[leftmargin=*,topsep=2pt]
  \item menentukan hasil komposisi dua fungsi atau lebih beserta \emph{domain}-nya;
  \item menyelesaikan komposisi yang melibatkan fungsi pecahan, akar, eksponen, dan logaritma;
  \item menentukan fungsi yang belum diketahui jika komposisinya diberikan;
  \item menjelaskan syarat suatu fungsi memiliki invers (satu-satu/bijektif) dan mengujinya secara aljabar maupun grafik;
  \item menentukan fungsi invers serta memahami kesimetrian grafik $f$ dan $f^{-1}$ terhadap garis $y=x$;
  \item memodelkan situasi nyata (konversi, diskon bertingkat, enkode--dekode) dengan komposisi dan invers.
\end{itemize}
\textbf{Capaian Pembelajaran (Fase F).} Peserta didik dapat menyatakan dan menganalisis fungsi melalui operasi komposisi dan invers, serta menggunakannya untuk menyelesaikan masalah kontekstual.
\end{tujuan}

\begin{pemantik}
\begin{itemize}[leftmargin=*,topsep=2pt]
  \item Saat kamu menukar rupiah ke dolar lalu dolar ke yen, kamu melakukan dua proses berurutan. Bisakah dua proses itu digabung menjadi satu "mesin" tunggal?
  \item Jika sebuah mesin mengubah masukan menjadi keluaran, adakah "mesin kebalikan" yang mengembalikan keluaran menjadi masukan semula?
  \item Apakah setiap fungsi punya kebalikan? Mengapa fungsi $f(x)=x^2$ menimbulkan masalah?
\end{itemize}
\end{pemantik}

\begin{studikasus}{Diskon Bertingkat di Toko Daring}
Sebuah toko daring memberi dua potongan harga berurutan untuk satu produk: pertama \emph{diskon} $20\%$, lalu \emph{voucher} potong Rp25.000.
\begin{center}
\begin{tabular}{c|c|c}
\toprule
Harga awal $x$ & Setelah diskon $g(x)=0{,}8x$ & Setelah voucher $f(g(x))=0{,}8x-25000$ \\
\midrule
100.000 & 80.000 & 55.000 \\
200.000 & 160.000 & 135.000 \\
300.000 & 240.000 & 215.000 \\
$\vdots$ & $\vdots$ & $\vdots$ \\
\bottomrule
\end{tabular}
\end{center}
\textbf{Amati dan duga.} (a) Tuliskan satu rumus tunggal yang langsung memetakan harga awal $x$ ke harga akhir. (b) Apakah urutannya penting --- samakah hasilnya jika voucher dipakai \emph{sebelum} diskon? (c) Jika kasir hanya mencatat harga akhir, bagaimana cara "membalik" untuk menemukan harga awal? Pertanyaan (a) dan (b) adalah inti \textbf{komposisi fungsi}, sedangkan (c) adalah inti \textbf{fungsi invers}.
\end{studikasus}

\section*{Peta Konsep Bab}
\addcontentsline{toc}{section}{Peta Konsep Bab}
\begin{center}
\begin{tikzpicture}[
  node distance=8mm,
  konsepbox/.style={draw=biru,fill=birumuda,rounded corners,align=center,inner sep=5pt,font=\small\bfseries},
  subbox/.style={draw=hijau,fill=hijaumuda,rounded corners,align=center,inner sep=4pt,font=\footnotesize},
  garis/.style={-Stealth,biru,thick}]
  \node[konsepbox] (akar) {FUNGSI: KOMPOSISI \&\ INVERS};
  \node[subbox,below left=12mm and 10mm of akar] (komp) {Komposisi\\$(f\circ g)(x)$};
  \node[subbox,below right=12mm and 10mm of akar] (inv) {Invers\\$f^{-1}(x)$};
  \node[subbox,below=10mm of komp] (dom) {Domain \&\\urutan};
  \node[subbox,below=10mm of inv] (syarat) {Syarat:\\satu-satu};
  \draw[garis] (akar)--(komp); \draw[garis] (akar)--(inv);
  \draw[garis] (komp)--(dom); \draw[garis] (inv)--(syarat);
  \draw[garis,dashed,gray] (komp)--(inv) node[midway,above,font=\tiny\itshape,black]{$f\circ f^{-1}=I$};
\end{tikzpicture}\\
\small \textbf{Peta konsep.} Komposisi menggabungkan fungsi secara berurutan; invers membatalkannya. Keduanya bertemu pada hubungan $(f\circ f^{-1})(x)=x$.
\end{center}

\section{Komposisi Fungsi}

\subsection*{Membangun konsep: fungsi sebagai mesin}
Bayangkan fungsi sebagai \emph{mesin}: masukkan $x$, keluar $g(x)$. Jika keluaran mesin $g$ langsung dialirkan ke mesin $f$, kita memperoleh satu mesin gabungan. Inilah komposisi: $g$ dikerjakan dahulu, hasilnya menjadi masukan bagi $f$.

\begin{center}
\begin{tikzpicture}[every node/.style={font=\small},>=Stealth]
  \node[draw=hijau,fill=hijaumuda,rounded corners,minimum height=9mm] (g) {$g$};
  \node[draw=oranye,fill=oranyemuda,rounded corners,minimum height=9mm,right=22mm of g] (f) {$f$};
  \draw[->,thick] ($(g.west)-(18mm,0)$) -- (g.west) node[midway,above]{$x$};
  \draw[->,thick] (g) -- (f) node[midway,above]{$g(x)$};
  \draw[->,thick] (f.east) -- ($(f.east)+(20mm,0)$) node[midway,above]{$f(g(x))$};
\end{tikzpicture}\\
\small Diagram mesin untuk $(f\circ g)(x)=f\big(g(x)\big)$.
\end{center}

\begin{definisi}
\emph{Komposisi fungsi} $f$ dan $g$ adalah fungsi baru
\[
(f \circ g)(x) = f\big(g(x)\big),
\]
yang terdefinisi untuk semua $x$ pada domain $g$ dengan $g(x)$ termasuk domain $f$. Fungsi $g$ dikerjakan lebih dahulu, kemudian hasilnya menjadi masukan bagi $f$.
\end{definisi}

\begin{contoh}{Komposisi Dua Fungsi}
Diketahui $f(x) = 2x + 1$ dan $g(x) = x^2$. Tentukan $(f\circ g)(x)$ dan $(g\circ f)(x)$.

\jawab
$(f\circ g)(x) = f(x^2) = 2x^2 + 1.$\\
$(g\circ f)(x) = g(2x+1) = (2x+1)^2 = 4x^2 + 4x + 1.$\\
Perhatikan bahwa $f\circ g \neq g\circ f$: \textbf{urutan komposisi penting}.
\end{contoh}

\begin{konsep}{Domain Komposisi}
Agar $(f\circ g)(x)$ bermakna, dua syarat harus dipenuhi:
\begin{enumerate}[leftmargin=*,topsep=2pt,label=(\arabic*)]
  \item $x$ berada pada domain $g$;
  \item $g(x)$ berada pada domain $f$.
\end{enumerate}
Karena itu domain komposisi bisa \emph{lebih sempit} daripada domain $g$ saja.
\end{konsep}

\begin{contoh}{Komposisi dengan Fungsi Akar}
Diketahui $f(x)=\sqrt{x}$ dan $g(x)=x-3$. Tentukan $(f\circ g)(x)$ beserta domainnya.

\jawab $(f\circ g)(x)=\sqrt{x-3}$. Syaratnya $x-3\ge 0$, sehingga domainnya $x\ge 3$.
\end{contoh}

\begin{contoh}{Komposisi Tiga Fungsi}
Diketahui $f(x)=x+1$, $g(x)=2x$, dan $h(x)=x^2$. Tentukan $(f\circ g\circ h)(x)$.

\jawab Kerjakan dari dalam: $h(x)=x^2$, lalu $g(x^2)=2x^2$, lalu $f(2x^2)=2x^2+1$.
Jadi $(f\circ g\circ h)(x)=2x^2+1$.
\end{contoh}

\begin{kesalahan}
Banyak peserta didik menukar urutan dan menghitung $(f\circ g)(x)=g\big(f(x)\big)$. Ingat: pada $f\circ g$, fungsi \emph{paling kanan} ($g$) dikerjakan \emph{lebih dahulu}. Selain itu, $f\circ g$ \emph{tidak sama} dengan perkalian $f\cdot g$.
\end{kesalahan}

\begin{dunianyata}
Pada grafika komputer, sebuah objek sering melewati rangkaian transformasi: \emph{skala} lalu \emph{rotasi} lalu \emph{translasi}. Setiap langkah adalah fungsi, dan rangkaiannya adalah komposisi. Itulah mengapa mengubah urutan langkah dapat menghasilkan tampilan yang berbeda.
\end{dunianyata}

\begin{latihan}
\begin{enumerate}[leftmargin=*]
  \item Diketahui $f(x)=x+3$, $g(x)=2x$. Tentukan $(f\circ g)(x)$ dan $(g\circ f)(2)$.
  \item Diketahui $f(x)=\sqrt{x}$ dan $g(x)=4-x^2$. Tentukan $(f\circ g)(x)$ dan domainnya.
\end{enumerate}
\end{latihan}

\section{Fungsi Invers}

\subsection*{Membangun konsep: membatalkan sebuah proses}
Jika $f$ mengubah $x$ menjadi $y$, maka inversnya $f^{-1}$ mengubah $y$ kembali menjadi $x$. Invers adalah "mesin pembatal": $(f\circ f^{-1})(x)=x$ dan $(f^{-1}\circ f)(x)=x$.

\begin{konsep}{Syarat Memiliki Invers}
Suatu fungsi memiliki invers (sebagai fungsi) \textbf{jika dan hanya jika fungsi itu satu-satu (injektif)}: setiap keluaran berasal dari tepat satu masukan. Secara grafik, fungsi satu-satu lolos \emph{uji garis mendatar} --- tidak ada garis horizontal yang memotong grafik lebih dari satu kali. Langkah mencari invers: tulis $y=f(x)$, tukar peran $x$ dan $y$, lalu nyatakan $y$ dalam $x$.
\end{konsep}

\begin{contoh}{Menentukan Invers}
Tentukan invers dari $f(x) = 2x + 1$.

\jawab $y = 2x+1 \Rightarrow x = 2y+1 \Rightarrow y = \dfrac{x-1}{2}$.
Jadi $f^{-1}(x) = \dfrac{x-1}{2}.$ Periksa: $f\big(f^{-1}(x)\big)=2\cdot\tfrac{x-1}{2}+1=x.$
\end{contoh}

\begin{center}
\begin{tikzpicture}
\begin{axis}[axis lines=middle,width=7.5cm,height=6cm,
  xlabel=$x$,ylabel=$y$,xmin=-4,xmax=5,ymin=-4,ymax=5,
  legend pos=north west,legend style={font=\small}]
\addplot[very thick,biru,domain=-3:2,samples=2]{2*x+1};
\addlegendentry{$f(x)=2x+1$}
\addplot[very thick,hijau,domain=-4:4,samples=2]{(x-1)/2};
\addlegendentry{$f^{-1}(x)=\frac{x-1}{2}$}
\addplot[dashed,gray,domain=-4:4,samples=2]{x};
\addlegendentry{$y=x$}
\end{axis}
\end{tikzpicture}\\
\small Grafik fungsi dan inversnya saling mencerminkan terhadap garis $y=x$.
\end{center}

\begin{contoh}{Invers Fungsi Pecahan}
Tentukan invers dari $f(x)=\dfrac{2x-1}{x+3}$.

\jawab $y=\dfrac{2x-1}{x+3}\Rightarrow y(x+3)=2x-1\Rightarrow yx+3y=2x-1$.
Kelompokkan: $x(y-2)=-1-3y$, sehingga $x=\dfrac{-1-3y}{y-2}=\dfrac{3y+1}{2-y}$.
Jadi $f^{-1}(x)=\dfrac{3x+1}{2-x}$ dengan $x\neq 2$.
\end{contoh}

\begin{konsep}{Invers Eksponen dan Logaritma}
Fungsi eksponen dan logaritma adalah pasangan invers:
\[
f(x)=a^x \ \Longleftrightarrow\ f^{-1}(x)={}^{a}\!\log x, \qquad a>0,\ a\neq1.
\]
Itulah sebabnya grafik $y=a^x$ dan $y={}^{a}\!\log x$ simetris terhadap garis $y=x$.
\end{konsep}

\begin{kesalahan}
Notasi $f^{-1}(x)$ berarti \emph{fungsi invers}, \textbf{bukan} $\dfrac{1}{f(x)}$. Sebagai contoh, jika $f(x)=2x+1$ maka $f^{-1}(x)=\tfrac{x-1}{2}$, sedangkan $\tfrac{1}{f(x)}=\tfrac{1}{2x+1}$ --- keduanya berbeda total.
\end{kesalahan}

\begin{dunianyata}
Pada keamanan data, proses \emph{enkripsi} mengubah pesan menjadi sandi, dan \emph{dekripsi} mengembalikannya. Dekripsi hanya mungkin bila enkripsinya satu-satu --- persis syarat keberadaan invers. Skala desibel, pH, dan Richter juga memakai logaritma, yaitu invers dari eksponen.
\end{dunianyata}

\begin{latihan}
\begin{enumerate}[leftmargin=*]
  \item Tentukan invers dari $f(x)=\dfrac{3x-2}{x+1}$.
  \item Tentukan invers dari $f(x)=3^{x}+2$ dan sebutkan domain inversnya.
\end{enumerate}
\end{latihan}

\section*{Ringkasan Bab}
\addcontentsline{toc}{section}{Ringkasan Bab}
\begin{itemize}[leftmargin=*,topsep=2pt]
  \item \textbf{Komposisi:} $(f\circ g)(x)=f\big(g(x)\big)$; kerjakan fungsi terdalam (paling kanan) lebih dahulu. Umumnya $f\circ g\neq g\circ f$.
  \item \textbf{Domain komposisi:} $x$ di domain $g$ \emph{dan} $g(x)$ di domain $f$.
  \item \textbf{Invers ada} jika fungsi satu-satu (lolos uji garis mendatar); berlaku $(f\circ f^{-1})(x)=(f^{-1}\circ f)(x)=x$.
  \item \textbf{Mencari invers:} tukar $x\leftrightarrow y$ lalu selesaikan untuk $y$.
  \item \textbf{Simetri:} grafik $f$ dan $f^{-1}$ mencerminkan terhadap $y=x$; eksponen dan logaritma adalah pasangan invers.
\end{itemize}

\begin{refleksi}
\begin{itemize}[leftmargin=*,topsep=2pt]
  \item[\cek] Saya dapat menghitung $(f\circ g)(x)$ dan menjelaskan mengapa urutannya penting.
  \item[\cek] Saya dapat menentukan domain sebuah komposisi.
  \item[\cek] Saya dapat menguji apakah suatu fungsi memiliki invers.
  \item[\cek] Saya dapat menentukan $f^{-1}(x)$ dan memverifikasinya melalui $f\circ f^{-1}=I$.
\end{itemize}
\end{refleksi}

\section*{Soal Akhir Bab \thechapter}
\addcontentsline{toc}{section}{Soal Akhir Bab \thechapter}

\bagiansoal{A}{(Fundamental --- setara SNBT Dasar)}
\begin{enumerate}[leftmargin=*]
  \item Diketahui $f(x)=2x+1$ dan $g(x)=x-3$. Tentukan $(f\circ g)(x)$.
  \item Diketahui $f(x)=x^2$ dan $g(x)=x+2$. Hitung $(g\circ f)(2)$.
  \item Tentukan invers dari $f(x)=3x-6$.
  \item Diketahui $f(x)=\dfrac{x+4}{2}$. Tentukan $f^{-1}(x)$.
  \item Diketahui $f(x)=2x$ dan $g(x)=x+5$. Hitung $(f\circ g)(3)$.
\end{enumerate}

\bagiansoal{B}{(Intermediate --- setara SNBT Sulit)}
\begin{enumerate}[leftmargin=*]
  \item Diketahui $f(x)=2x+1$ dan $g(x)=x^2-1$. Tentukan $(f\circ g)(x)$ dan $(g\circ f)(x)$, lalu tunjukkan keduanya tidak sama.
  \item Diketahui $(f\circ g)(x)=4x+5$ dan $g(x)=2x+1$. Tentukan $f(x)$.
  \item Tentukan invers dari $f(x)=\dfrac{2x-1}{x+3}$ beserta nilai $x$ yang tidak diperbolehkan pada inversnya.
  \item Diketahui $f(x)=\sqrt{x-2}$. Tentukan domain $f$ dan rumus $f^{-1}(x)$.
  \item Diketahui $f(x)=3x-2$ dan $g(x)=x+a$. Jika $(f\circ g)(2)=10$, tentukan $a$.
  \item Diketahui $f(x)=x+1$ dan $(f\circ g)(x)=x^2+2$. Tentukan $g(x)$.
  \item Tentukan invers dari $f(x)=2^{x+1}$.
\end{enumerate}

\bagiansoal{C}{(Advanced --- setara Olimpiade SMA, gabungan antarbab)}
\begin{enumerate}[leftmargin=*]
  \item \textbf{(Komposisi + Eksponen).} Diketahui $f(x)=2^{x}$ dan $g(x)=x+3$. Selesaikan $(f\circ g)(x)=32$.
  \item \textbf{(Invers + Logaritma).} Diketahui $f(x)={}^{10}\!\log(x-1)$. Tentukan $f^{-1}(x)$ dan domainnya.
  \item \textbf{(Komposisi + Fungsi Kuadrat).} Diketahui $(f\circ g)(x)=x^2-4x+7$ dan $f(x)=x+3$. Tentukan $g(x)$ serta titik minimumnya.
  \item \textbf{(Komposisi Berulang).} Diketahui $f(x)=\dfrac{x}{x+1}$. Tentukan $(f\circ f\circ f)(x)$.
  \item \textbf{(Invers --- Involusi).} Diketahui $f(x)=\dfrac{x+1}{x-1}$ ($x\neq1$). Buktikan bahwa $f^{-1}(x)=f(x)$.
  \item \textbf{(Komposisi + Trigonometri).} Diketahui $f(x)=\sin x$ dan $g(x)=2x$. Bandingkan $(f\circ g)(x)$ dan $(g\circ f)(x)$ serta sebutkan periodenya.
  \item \textbf{(Komposisi + Barisan).} Barisan didefinisikan $x_0=5$ dan $x_{n+1}=f(x_n)$ dengan $f(x)=2x-3$. Tentukan rumus $x_n$ dalam $n$.
\end{enumerate}

\bagiansoal{D}{(Expert --- setara tahun pertama universitas, sintesis \& pembuktian)}
\begin{enumerate}[leftmargin=*]
  \item \textbf{(Komposisi + Pembuktian).} Buktikan bahwa komposisi dua fungsi injektif juga injektif.
  \item \textbf{(Invers + Eksponen).} Diketahui $f(x)=\dfrac{e^{x}-e^{-x}}{2}$. Tentukan $f^{-1}(x)$ dengan menyelesaikan persamaan kuadrat dalam $e^{x}$.
  \item \textbf{(Komposisi + Fungsi Rasional).} Diketahui $f(x)=\dfrac{1}{1-x}$. Tunjukkan bahwa $(f\circ f\circ f)(x)=x$ untuk setiap $x$ pada domain.
  \item \textbf{(Invers + Statistika).} Skor baku didefinisikan $z=\dfrac{x-\mu}{\sigma}$ dengan $\mu,\sigma$ tetap ($\sigma>0$). Tentukan transformasi inversnya dan jelaskan maknanya.
  \item \textbf{(Komposisi + Paritas).} Buktikan: jika $g$ fungsi genap, maka $f\circ g$ genap untuk sembarang $f$; dan jika $f,g$ keduanya ganjil, maka $f\circ g$ ganjil.
\end{enumerate}

\begin{challenge}
\begin{enumerate}[leftmargin=*,topsep=2pt]
  \item Tentukan semua fungsi linear $f(x)=ax+b$ yang memenuhi $(f\circ f)(x)=x$ untuk setiap $x$.
  \item Fungsi $f$ memenuhi $f(x)+2f\!\left(\dfrac1x\right)=3x$ untuk setiap $x\neq0$. Tentukan $f(x)$.
  \item Diketahui $f(x)=\dfrac{x}{\sqrt{1+x^2}}$. Tunjukkan dengan induksi bahwa komposisi ke-$n$ memenuhi $f^{(n)}(x)=\dfrac{x}{\sqrt{1+nx^2}}$.
\end{enumerate}
\end{challenge}

% ============================================================
\chapter{Lingkaran}
% ============================================================
\begin{tujuan}
Setelah mempelajari bab ini, peserta didik mampu:
\begin{itemize}[leftmargin=*,topsep=2pt]
  \item menyusun persamaan lingkaran dari bentuk baku maupun bentuk umum, serta menentukan pusat dan jari-jari;
  \item menentukan kedudukan titik dan kedudukan garis terhadap lingkaran;
  \item menentukan persamaan garis singgung lingkaran di suatu titik, dengan gradien tertentu, dan dari titik di luar lingkaran;
  \item menghitung panjang garis singgung dan \emph{daya} (power) sebuah titik;
  \item menghitung panjang busur dan luas juring;
  \item menerapkan konsep lingkaran pada masalah nyata (radar, GPS, desain taman).
\end{itemize}
\textbf{Capaian Pembelajaran (Fase F).} Peserta didik dapat menyatakan lingkaran secara aljabar dan geometris serta menggunakannya untuk menyelesaikan masalah kedudukan, garis singgung, dan pengukuran.
\end{tujuan}

\begin{pemantik}
\begin{itemize}[leftmargin=*,topsep=2pt]
  \item Sebuah stasiun radar mendeteksi semua objek dalam jarak $50$ km. Bagaimana menulis "daerah jangkauan" itu sebagai persamaan?
  \item Mengapa garis singgung sebuah roda selalu tegak lurus jari-jari di titik sentuhnya dengan jalan?
  \item Dari sebuah titik di luar lingkaran, mengapa kita selalu bisa menarik tepat \emph{dua} garis singgung yang sama panjang?
\end{itemize}
\end{pemantik}

\begin{studikasus}{Jangkauan Tiga Menara Pemancar}
Sebuah operator seluler memasang tiga menara. Tiap menara menjangkau pelanggan dalam radius tertentu, sehingga daerah layanannya berupa cakram (lingkaran).
\begin{center}
\begin{tabular}{c|c|c}
\toprule
Menara & Pusat $(a,b)$ km & Radius $r$ km \\
\midrule
P & $(0,0)$ & $5$ \\
Q & $(8,0)$ & $5$ \\
R & $(4,6)$ & $4$ \\
\bottomrule
\end{tabular}
\end{center}
\textbf{Amati dan duga.} (a) Tuliskan persamaan daerah layanan tiap menara. (b) Apakah pelanggan di titik $(4,0)$ terlayani oleh P? oleh Q? (c) Di mana dua daerah layanan saling tumpang tindih? Pertanyaan ini menuntut kita menyatakan lingkaran sebagai persamaan dan menentukan \textbf{kedudukan titik} serta \textbf{kedudukan dua lingkaran}.
\end{studikasus}

\section*{Peta Konsep Bab}
\addcontentsline{toc}{section}{Peta Konsep Bab}
\begin{center}
\begin{tikzpicture}[
  node distance=8mm,
  konsepbox/.style={draw=biru,fill=birumuda,rounded corners,align=center,inner sep=5pt,font=\small\bfseries},
  subbox/.style={draw=hijau,fill=hijaumuda,rounded corners,align=center,inner sep=4pt,font=\footnotesize},
  garis/.style={-Stealth,biru,thick}]
  \node[konsepbox] (akar) {LINGKARAN};
  \node[subbox,below left=12mm and 14mm of akar] (pers) {Persamaan\\(baku \&\ umum)};
  \node[subbox,below=12mm of akar] (ked) {Kedudukan\\titik \&\ garis};
  \node[subbox,below right=12mm and 14mm of akar] (sing) {Garis\\singgung};
  \node[subbox,below=10mm of pers] (buja) {Busur \&\\juring};
  \node[subbox,below=10mm of sing] (daya) {Daya titik};
  \draw[garis] (akar)--(pers); \draw[garis] (akar)--(ked); \draw[garis] (akar)--(sing);
  \draw[garis] (pers)--(buja); \draw[garis] (sing)--(daya);
\end{tikzpicture}\\
\small \textbf{Peta konsep.} Dari persamaan lingkaran tumbuh tiga cabang: kedudukan (titik/garis), garis singgung beserta daya titik, dan pengukuran busur/juring.
\end{center}

\section{Persamaan Lingkaran}

\subsection*{Membangun konsep dari definisi}
Lingkaran adalah himpunan semua titik yang berjarak sama (yaitu $r$) dari sebuah titik tetap (pusat). Jika pusatnya $(a,b)$ dan sebuah titik $(x,y)$ berjarak $r$ darinya, maka menurut rumus jarak $\sqrt{(x-a)^2+(y-b)^2}=r$. Mengkuadratkan kedua ruas langsung memberi persamaan baku --- jadi persamaan lingkaran hanyalah "rumus jarak yang dikuadratkan".

\begin{konsep}{Bentuk Persamaan}
\textbf{Pusat $(a,b)$, jari-jari $r$ (bentuk baku):}
\[
(x-a)^2 + (y-b)^2 = r^2.
\]
\textbf{Bentuk umum:} $x^2 + y^2 + Ax + By + C = 0$, dengan
\[
\text{pusat}=\left(-\tfrac{A}{2}, -\tfrac{B}{2}\right), \qquad
r = \sqrt{\tfrac{A^2}{4}+\tfrac{B^2}{4}-C}.
\]
Persamaan berupa lingkaran sejati hanya jika $\tfrac{A^2}{4}+\tfrac{B^2}{4}-C>0$.
\end{konsep}

\begin{center}
\begin{tikzpicture}[scale=0.9]
  \draw[->] (-1,0)--(5,0) node[right]{$x$};
  \draw[->] (0,-1)--(0,5) node[above]{$y$};
  \draw[very thick,biru] (2,2) circle (1.8);
  \fill (2,2) circle (1.5pt) node[below left]{\small $(a,b)$};
  \draw[oranye,thick,-{Latex}] (2,2)--(3.27,3.27) node[midway,above left]{\small $r$};
\end{tikzpicture}
\end{center}

\smallskip\noindent\textbf{Mengenal anatomi lingkaran.} Sebelum menghitung, kenali unsur-unsurnya:
\begin{center}
\begin{tikzpicture}[scale=1.25,>=Stealth,line join=round]
\draw[very thick,biru] (0,0) circle (2);
\fill (0,0) circle (1.6pt);
\node[font=\scriptsize,above right=-1pt] at (0,0){$O$};
% diameter (horizontal)
\draw[hijau!60!black,thick] (180:2)--(0:2);
\node[hijau!60!black,font=\scriptsize] at (1.05,0.2){diameter $=2r$};
% jari-jari
\draw[oranye,very thick] (0,0)--(72:2);
\node[oranye,font=\scriptsize] at (54:1.15){$r$};
% tali busur + apotema
\coordinate (P) at (-65:2); \coordinate (Q) at (235:2);
\draw[ungu,thick] (P)--(Q);
\coordinate (M) at ($(P)!0.5!(Q)$);
\draw[merah,thick] (0,0)--(M);
\draw[gray] ($(M)!2.4mm!(0,0)$)--($($(M)!2.4mm!(0,0)$)!2.4mm!-90:(0,0)$)--($(M)!2.4mm!-90:(0,0)$);
\node[ungu,font=\scriptsize] at (0,-2.3){tali busur};
\node[merah,font=\scriptsize,left] at ($(0,0)!0.5!(M)$){apotema};
\end{tikzpicture}
\end{center}

\begin{konsep}{Unsur-unsur Lingkaran (Istilah Penting)}
\begin{itemize}[leftmargin=*,topsep=2pt]
  \item \textbf{Jari-jari} ($r$): ruas dari pusat $O$ ke titik pada lingkaran.
  \item \textbf{Diameter} ($d=2r$): tali busur \emph{terpanjang}, yaitu yang melalui pusat.
  \item \textbf{Tali busur}: ruas yang menghubungkan dua titik pada lingkaran (tidak harus lewat pusat).
  \item \textbf{Apotema}: jarak (tegak lurus) dari pusat ke sebuah tali busur. Apotema selalu \emph{membagi dua} tali busur.
  \item \textbf{Busur}: bagian lengkung lingkaran (busur kecil/minor dan busur besar/mayor).
  \item \textbf{Juring}: daerah di antara dua jari-jari dan busur ("potongan pizza").
  \item \textbf{Tembereng}: daerah di antara tali busur dan busur.
  \item \textbf{Garis singgung} menyentuh lingkaran di \emph{satu} titik; \textbf{garis sekan} memotongnya di \emph{dua} titik.
\end{itemize}
\end{konsep}

\begin{contoh}{Persamaan dan Pusat}
Tentukan pusat dan jari-jari dari $x^2+y^2-4x+2y-4=0$.

\jawab $A=-4,\ B=2,\ C=-4$. Pusat $\left(-\tfrac{A}{2},-\tfrac{B}{2}\right)=(2,-1)$,
$r=\sqrt{4+1+4}=3$.
\end{contoh}

\begin{kesalahan}
Pada bentuk baku $(x-a)^2+(y-b)^2=r^2$, tandanya \emph{berlawanan}: pusat $(2,-1)$ berasal dari $(x-2)^2+(y+1)^2$. Selain itu, ruas kanan adalah $r^2$, \textbf{bukan} $r$ --- jika tertulis $=9$ maka $r=3$, bukan $9$.
\end{kesalahan}

\section{Kedudukan Titik dan Garis}
\begin{konsep}{Kedudukan Titik terhadap Lingkaran}
Untuk lingkaran dengan kuasa $K(x,y)=(x-a)^2+(y-b)^2-r^2$, titik $(x_0,y_0)$:
\[
K(x_0,y_0)<0 \ \text{(di dalam)}, \quad
=0 \ \text{(pada)}, \quad
>0 \ \text{(di luar)}.
\]
\end{konsep}

\begin{center}
\begin{tikzpicture}[scale=1.0,line join=round]
\draw[very thick,biru] (0,0) circle (2);
\fill (0,0) circle (1.2pt); \node[font=\scriptsize,below=1pt] at (0,0){$O$};
\fill[hijau!60!black] (0.5,0.3) circle (2.2pt);
\node[hijau!60!black,font=\scriptsize,right] at (0.62,0.32){di dalam: $K<0$};
\fill[biru] (90:2) circle (2.2pt);
\node[biru,font=\scriptsize,above] at (90:2){pada: $K=0$};
\fill[merah] (2.5,1.5) circle (2.2pt);
\node[merah,font=\scriptsize,right] at (2.5,1.5){di luar: $K>0$};
\draw[gray,dashed] (0,0)--(2.5,1.5);
\end{tikzpicture}\\
\small Tanda kuasa $K(x_0,y_0)=(x_0-a)^2+(y_0-b)^2-r^2$ menentukan kedudukan titik: bandingkan jaraknya ke pusat dengan $r$.
\end{center}

\begin{contoh}{Kedudukan Titik}
Tentukan kedudukan titik $(5,1)$ terhadap lingkaran $x^2+y^2-4x-2y-4=0$ (pusat $(2,1)$, $r=3$).

\jawab Jarak pusat ke titik $=\sqrt{(5-2)^2+(1-1)^2}=3=r$. Karena sama dengan $r$, titik \textbf{terletak pada} lingkaran.
\end{contoh}

\begin{konsep}{Kedudukan Garis terhadap Lingkaran}
Substitusikan garis ke persamaan lingkaran sehingga diperoleh persamaan kuadrat dengan diskriminan $D$:
\[
D>0 \ \text{garis memotong (2 titik)}, \quad
D=0 \ \text{menyinggung}, \quad
D<0 \ \text{tidak memotong}.
\]
Alternatif: bandingkan jarak pusat ke garis $d$ dengan $r$ ($d<r$, $d=r$, $d>r$).
\end{konsep}

\smallskip\noindent\textbf{Visualisasi kedudukan garis terhadap lingkaran.}
\begin{center}
\begin{tabular}{ccc}
% memotong
\begin{tikzpicture}[scale=0.95,line join=round]
\draw[very thick,biru] (0,0) circle (1.2); \fill (0,0) circle (1pt);
\draw[hijau!60!black,thick] (-1.7,0.5)--(1.7,0.5);
\fill[oranye] (-1.09,0.5) circle (1.6pt); \fill[oranye] (1.09,0.5) circle (1.6pt);
\draw[gray,dashed] (0,0)--(0,0.5); \node[font=\scriptsize,right] at (0,0.27){$d$};
\end{tikzpicture}
&
% menyinggung
\begin{tikzpicture}[scale=0.95,line join=round]
\draw[very thick,biru] (0,0) circle (1.2); \fill (0,0) circle (1pt);
\draw[hijau!60!black,thick] (-1.7,1.2)--(1.7,1.2);
\fill[merah] (0,1.2) circle (1.6pt);
\draw[gray,dashed] (0,0)--(0,1.2); \node[font=\scriptsize,right] at (0,0.6){$d$};
\end{tikzpicture}
&
% tidak memotong
\begin{tikzpicture}[scale=0.95,line join=round]
\draw[very thick,biru] (0,0) circle (1.2); \fill (0,0) circle (1pt);
\draw[hijau!60!black,thick] (-1.7,1.75)--(1.7,1.75);
\draw[gray,dashed] (0,0)--(0,1.75); \node[font=\scriptsize,right] at (0,0.9){$d$};
\end{tikzpicture}
\\[-1pt]
\footnotesize \textbf{Memotong} ($d<r$, $D>0$) & \footnotesize \textbf{Menyinggung} ($d=r$, $D=0$) & \footnotesize \textbf{Tidak memotong} ($d>r$, $D<0$)\\
\footnotesize 2 titik potong & \footnotesize 1 titik singgung & \footnotesize tak ada titik\\
\end{tabular}
\end{center}

\smallskip\noindent\textbf{Visualisasi kedudukan dua lingkaran} (jarak pusat $d$, jari-jari $r_1,r_2$).
\begin{center}
\begin{tabular}{ccccc}
\begin{tikzpicture}[scale=0.6]
\draw[very thick,biru] (-0.85,0) circle (0.55); \draw[very thick,oranye] (0.85,0) circle (0.55);
\fill (-0.85,0) circle (0.7pt); \fill (0.85,0) circle (0.7pt);
\end{tikzpicture}
&
\begin{tikzpicture}[scale=0.6]
\draw[very thick,biru] (-0.6,0) circle (0.6); \draw[very thick,oranye] (0.6,0) circle (0.6);
\fill (-0.6,0) circle (0.7pt); \fill (0.6,0) circle (0.7pt); \fill[merah] (0,0) circle (1.2pt);
\end{tikzpicture}
&
\begin{tikzpicture}[scale=0.6]
\draw[very thick,biru] (-0.4,0) circle (0.6); \draw[very thick,oranye] (0.4,0) circle (0.6);
\fill (-0.4,0) circle (0.7pt); \fill (0.4,0) circle (0.7pt);
\end{tikzpicture}
&
\begin{tikzpicture}[scale=0.6]
\draw[very thick,biru] (0,0) circle (1.0); \draw[very thick,oranye] (0.6,0) circle (0.4);
\fill (0,0) circle (0.7pt); \fill (0.6,0) circle (0.7pt); \fill[merah] (1.0,0) circle (1.2pt);
\end{tikzpicture}
&
\begin{tikzpicture}[scale=0.6]
\draw[very thick,biru] (0,0) circle (1.0); \draw[very thick,oranye] (0.25,0) circle (0.3);
\fill (0,0) circle (0.7pt); \fill (0.25,0) circle (0.7pt);
\end{tikzpicture}
\\[-1pt]
\footnotesize Lepas (luar) & \footnotesize Singgung luar & \footnotesize Berpotongan & \footnotesize Singgung dalam & \footnotesize Di dalam\\
\footnotesize $d>r_1+r_2$ & \footnotesize $d=r_1+r_2$ & \footnotesize $|r_1-r_2|<d<r_1+r_2$ & \footnotesize $d=|r_1-r_2|$ & \footnotesize $d<|r_1-r_2|$\\
\end{tabular}
\end{center}

\section{Garis Singgung Lingkaran}
\begin{konsep}{Tiga Kasus Garis Singgung}
\begin{itemize}[leftmargin=*,topsep=2pt]
  \item \textbf{Di titik $(x_1,y_1)$ pada lingkaran $x^2+y^2=r^2$:} \ $x_1x+y_1y=r^2$.
  \item \textbf{Dengan gradien $m$ pada $x^2+y^2=r^2$:} \ $y=mx\pm r\sqrt{1+m^2}$.
  \item \textbf{Dari titik luar $P$:} panjang garis singgung $=\sqrt{K(P)}=\sqrt{x_P^2+y_P^2+Ax_P+By_P+C}$.
\end{itemize}
Garis singgung selalu \emph{tegak lurus} jari-jari di titik singgung.
\end{konsep}

\begin{center}
\begin{tikzpicture}[scale=0.8,>=Stealth]
  \draw[very thick,biru] (0,0) circle (2);
  \fill (0,0) circle (1.5pt) node[below left]{\small $O$};
  \coordinate (T) at (1.2,1.6);
  \fill[oranye] (T) circle (1.8pt) node[below right]{\small $T$};
  \draw[oranye,thick] (0,0)--(T);
  \draw[hijau,very thick] ($(T)!-1.7cm!90:(0,0)$)--($(T)!1.7cm!90:(0,0)$);
  % penanda sudut siku-siku di T
  \draw[gray] ($(T)!3mm!(0,0)$) -- ($($(T)!3mm!(0,0)$)!3mm!90:(0,0)$) -- ($(T)!3mm!90:(0,0)$);
  \node[hijau] at (2.9,2.0){\small garis singgung};
\end{tikzpicture}\\
\small Garis singgung di $T$ tegak lurus jari-jari $OT$.
\end{center}

\smallskip\noindent\textbf{Dua garis singgung dari sebuah titik di luar.}
\begin{center}
\begin{tikzpicture}[scale=1.05,>=Stealth,line join=round]
\draw[very thick,biru] (0,0) circle (1.5);
\fill (0,0) circle (1.2pt) node[below left,font=\scriptsize]{$O$};
\coordinate (P) at (4,0); \fill (P) circle (1.2pt) node[right,font=\scriptsize]{$P$};
\coordinate (A) at (0.5625,1.39); \coordinate (B) at (0.5625,-1.39);
\draw[oranye,thick] (0,0)--(A); \draw[oranye,thick] (0,0)--(B);
\draw[hijau!60!black,very thick] (P)--(A); \draw[hijau!60!black,very thick] (P)--(B);
\draw[gray,dashed] (0,0)--(P);
\fill[merah] (A) circle (1.6pt) node[above left,font=\scriptsize]{$A$};
\fill[merah] (B) circle (1.6pt) node[below left,font=\scriptsize]{$B$};
\draw[gray] ($(A)!2.4mm!(0,0)$)--($($(A)!2.4mm!(0,0)$)!2.4mm!-90:(0,0)$)--($(A)!2.4mm!-90:(0,0)$);
\draw[gray] ($(B)!2.4mm!(0,0)$)--($($(B)!2.4mm!(0,0)$)!2.4mm!90:(0,0)$)--($(B)!2.4mm!90:(0,0)$);
\node[hijau!60!black,font=\scriptsize] at (2.6,1.05){$PA=PB=\sqrt{K(P)}$};
\end{tikzpicture}\\
\small Dari titik luar $P$ selalu ada \emph{dua} garis singgung yang \emph{sama panjang}; panjangnya $\sqrt{K(P)}$ dengan $OA\perp PA$ dan $OB\perp PB$.
\end{center}

\begin{contoh}{Garis Singgung di Suatu Titik}
Tentukan persamaan garis singgung lingkaran $x^2+y^2=25$ di titik $(3,4)$.

\jawab Gunakan $x_1x+y_1y=r^2$: $\ 3x+4y=25.$
\end{contoh}

\begin{dunianyata}
Pada sistem \textbf{GPS}, posisi penerima ditentukan dari jaraknya ke beberapa satelit. Setiap jarak mendefinisikan sebuah lingkaran (di bidang) atau bola (di ruang); perpotongan lingkaran-lingkaran itulah lokasi penerima. Stasiun \textbf{radar} bekerja serupa: batas jangkauannya adalah sebuah lingkaran.
\end{dunianyata}

\section{Panjang Busur dan Luas Juring}
\begin{konsep}{Rumus}
Untuk sudut pusat $\theta$ (derajat) dan jari-jari $r$:
\[
\text{Panjang busur} = \frac{\theta}{360^\circ}\cdot 2\pi r, \qquad
\text{Luas juring} = \frac{\theta}{360^\circ}\cdot \pi r^2.
\]
Jika $\theta$ dalam radian: busur $=r\theta$ dan juring $=\tfrac12 r^2\theta$.
\end{konsep}

\begin{center}
\begin{tabular}{cc}
% juring + busur
\begin{tikzpicture}[scale=1.0,line join=round]
\draw[very thick,biru] (0,0) circle (1.6);
\fill[oranye,opacity=0.25] (0,0)--(30:1.6) arc(30:110:1.6)--cycle;
\draw[hijau!60!black,very thick] (30:1.6) arc(30:110:1.6);
\draw (0,0)--(30:1.6); \draw (0,0)--(110:1.6);
\fill (0,0) circle (1pt);
\node[font=\scriptsize] at (70:0.55){$\theta$};
\node[hijau!60!black,font=\scriptsize] at (70:1.95){busur};
\node[oranye!75!black,font=\scriptsize] at (62:1.05){juring};
\end{tikzpicture}
&
% tembereng
\begin{tikzpicture}[scale=1.0,line join=round]
\draw[very thick,biru] (0,0) circle (1.6);
\fill[ungu,opacity=0.22] (30:1.6) arc(30:110:1.6)--cycle;
\draw[ungu,thick] (30:1.6)--(110:1.6);
\fill (0,0) circle (1pt);
\node[ungu!80!black,font=\scriptsize] at (70:1.18){tembereng};
\node[font=\scriptsize] at (70:1.95){tali busur};
\end{tikzpicture}
\\[-1pt]
\footnotesize \textbf{Juring} (antara 2 jari-jari \&\ busur) & \footnotesize \textbf{Tembereng} (antara tali busur \&\ busur)\\
\end{tabular}
\end{center}

\smallskip\noindent\textbf{Sudut pusat dan sudut keliling.}
\begin{center}
\begin{tikzpicture}[scale=1.15,line join=round]
\draw[very thick,biru] (0,0) circle (1.6);
\coordinate (A) at (210:1.6); \coordinate (B) at (330:1.6); \coordinate (C) at (95:1.6);
\fill (0,0) circle (1pt) node[font=\scriptsize,above right=-1pt]{$O$};
\fill (A) circle (1.2pt) node[font=\scriptsize,left]{$A$};
\fill (B) circle (1.2pt) node[font=\scriptsize,right]{$B$};
\fill (C) circle (1.2pt) node[font=\scriptsize,above]{$C$};
\draw[oranye,thick] (0,0)--(A); \draw[oranye,thick] (0,0)--(B);
\draw[hijau!60!black,thick] (C)--(A); \draw[hijau!60!black,thick] (C)--(B);
\node[oranye,font=\scriptsize] at (270:0.5){$2\beta$};
\node[hijau!60!black,font=\scriptsize] at (95:1.02){$\beta$};
\end{tikzpicture}\\
\small Pada busur yang sama ($AB$): sudut pusat $\angle AOB$ tepat \emph{dua kali} sudut keliling $\angle ACB$.
\end{center}

\begin{konsep}{Sudut Pusat dan Sudut Keliling}
Untuk busur yang sama, \textbf{sudut pusat} $=2\times$ \textbf{sudut keliling}. Akibatnya, semua sudut keliling yang menghadap busur sama besarnya, dan sudut keliling yang menghadap \emph{diameter} selalu $90^\circ$ (sudut dalam setengah lingkaran).
\end{konsep}

\begin{contoh}{Busur dan Juring}
Sebuah lingkaran berjari-jari 6 cm dengan sudut pusat $60^\circ$. Tentukan panjang busur dan luas juringnya.

\jawab
Busur $= \tfrac{60}{360}\cdot 2\pi(6) = 2\pi$ cm.
Juring $= \tfrac{60}{360}\cdot \pi(6)^2 = 6\pi$ cm$^2$.
\end{contoh}

\begin{konsep}{Bahasa Alternatif yang Sering Mengecoh di Soal}
\begin{itemize}[leftmargin=*,topsep=2pt]
  \item \textbf{jari-jari} ($r$) vs \textbf{diameter} ($2r$) vs \textbf{tali busur}: tali busur tidak harus melalui pusat; diameter adalah tali busur khusus yang lewat pusat.
  \item \textbf{busur} (garis lengkung) vs \textbf{juring} (daerah "potongan pizza") vs \textbf{tembereng} (daerah antara tali busur dan busur). Jangan tertukar saat soal meminta "luas".
  \item \textbf{garis singgung} (\emph{tangen}, satu titik) vs \textbf{garis sekan} (\emph{secant}, dua titik). "Garis menyinggung" $\Rightarrow$ syarat $D=0$ atau $d=r$.
  \item \textbf{apotema} = jarak \emph{tegak lurus} pusat ke tali busur (bukan ke titik sembarang).
  \item \textbf{"panjang garis singgung dari titik $P$"} = jarak $P$ ke titik singgung $=\sqrt{K(P)}$, \emph{bukan} jarak $P$ ke pusat.
  \item \textbf{"daya/kuasa titik"} $K(P)$: bernilai negatif (di dalam), nol (pada), positif (di luar).
  \item \textbf{"busur minor"} (lebih pendek) vs \textbf{"busur mayor"} (lebih panjang) untuk tali busur yang sama.
  \item \textbf{"lingkaran satuan"} = berpusat $O$ dengan jari-jari $1$.
\end{itemize}
\end{konsep}

\section*{Ringkasan Bab}
\addcontentsline{toc}{section}{Ringkasan Bab}
\begin{itemize}[leftmargin=*,topsep=2pt]
  \item \textbf{Baku:} $(x-a)^2+(y-b)^2=r^2$; \textbf{umum:} $x^2+y^2+Ax+By+C=0$ dengan pusat $\left(-\tfrac A2,-\tfrac B2\right)$.
  \item \textbf{Kedudukan titik:} bandingkan $K(x_0,y_0)$ dengan $0$ (atau jarak ke pusat dengan $r$).
  \item \textbf{Kedudukan garis:} diskriminan substitusi, atau jarak pusat ke garis $d$ versus $r$.
  \item \textbf{Garis singgung:} di titik $x_1x+y_1y=r^2$; gradien $m$: $y=mx\pm r\sqrt{1+m^2}$; panjang dari titik luar $=\sqrt{K(P)}$.
  \item \textbf{Busur} $=\tfrac{\theta}{360^\circ}2\pi r$, \textbf{juring} $=\tfrac{\theta}{360^\circ}\pi r^2$.
\end{itemize}

\begin{refleksi}
\begin{itemize}[leftmargin=*,topsep=2pt]
  \item[\cek] Saya dapat mengubah bentuk umum ke bentuk baku dan menemukan pusat serta jari-jari.
  \item[\cek] Saya dapat menentukan kedudukan titik dan garis terhadap lingkaran.
  \item[\cek] Saya dapat menyusun garis singgung untuk ketiga kasus.
  \item[\cek] Saya dapat menghitung panjang busur dan luas juring.
\end{itemize}
\end{refleksi}

\section*{Soal Akhir Bab \thechapter}
\addcontentsline{toc}{section}{Soal Akhir Bab \thechapter}

\bagiansoal{A}{(Fundamental --- setara SNBT Dasar)}
\begin{enumerate}[leftmargin=*]
  \item Tentukan persamaan lingkaran berpusat $(0,0)$ dan berjari-jari $5$.
  \item Tentukan persamaan lingkaran berpusat $(3,-2)$ dan berjari-jari $4$.
  \item Tentukan pusat dan jari-jari dari $x^2+y^2-6x+4y-12=0$.
  \item Hitung panjang busur lingkaran berjari-jari $6$ cm dengan sudut pusat $90^\circ$.
  \item Hitung luas juring dengan $r=10$ cm dan sudut pusat $72^\circ$.
\end{enumerate}

\bagiansoal{B}{(Intermediate --- setara SNBT Sulit)}
\begin{enumerate}[leftmargin=*]
  \item Tentukan kedudukan titik $(1,2)$ terhadap lingkaran $x^2+y^2=9$.
  \item Sebuah lingkaran berpusat $(2,3)$ dan melalui titik $(5,7)$. Tentukan persamaannya.
  \item Tentukan persamaan garis singgung lingkaran $x^2+y^2=25$ di titik $(3,4)$.
  \item Tentukan persamaan garis singgung lingkaran $x^2+y^2=20$ yang bergradien $2$.
  \item Tentukan kedudukan garis $y=x+1$ terhadap lingkaran $x^2+y^2=4$.
  \item Tentukan kedudukan titik $(5,1)$ terhadap lingkaran $x^2+y^2-4x-2y-4=0$.
  \item Sebuah tali busur lingkaran berjari-jari $5$ berjarak $3$ dari pusat. Tentukan panjang tali busur itu.
\end{enumerate}

\bagiansoal{C}{(Advanced --- setara Olimpiade SMA, gabungan antarbab)}
\begin{enumerate}[leftmargin=*]
  \item \textbf{(Garis singgung dari titik luar).} Tentukan panjang garis singgung dari titik $(7,1)$ ke lingkaran $x^2+y^2=9$.
  \item \textbf{(Daya titik).} Hitung daya titik $P(6,0)$ terhadap lingkaran $x^2+y^2=4$, lalu tafsirkan tandanya.
  \item \textbf{(Lingkaran + Garis).} Tunjukkan bahwa garis $3x+4y=25$ menyinggung lingkaran $x^2+y^2=25$, dan tentukan titik singgungnya.
  \item \textbf{(Lingkaran + Fungsi Kuadrat).} Garis $y=mx$ menyinggung lingkaran $(x-4)^2+y^2=4$. Tentukan semua nilai $m$.
  \item \textbf{(Lingkaran + Geometri).} Tentukan persamaan lingkaran yang melalui $A(0,0)$, $B(6,0)$, dan $C(0,8)$.
  \item \textbf{(Lingkaran + Trigonometri).} Titik $T$ pada lingkaran $x^2+y^2=16$ membentuk sudut $60^\circ$ terhadap sumbu-$x$ positif. Tentukan koordinat $T$ dan panjang busur dari $(4,0)$ ke $T$.
\end{enumerate}

\bagiansoal{D}{(Expert --- setara tahun pertama universitas, sintesis \& pembuktian)}
\begin{enumerate}[leftmargin=*]
  \item \textbf{(Pembuktian).} Buktikan bahwa garis singgung $x_1x+y_1y=r^2$ pada lingkaran $x^2+y^2=r^2$ tegak lurus jari-jari di titik singgung $(x_1,y_1)$.
  \item \textbf{(Lingkaran + Optimasi).} Tentukan titik pada lingkaran $x^2+y^2=1$ yang terdekat dan terjauh dari titik $(3,4)$, beserta jaraknya.
  \item \textbf{(Dua lingkaran).} Tentukan persamaan garis kuasa (sumbu radikal) dari $x^2+y^2=4$ dan $x^2+y^2-6x+5=0$, lalu jelaskan maknanya.
  \item \textbf{(Lingkaran + Sistem).} Tentukan persamaan lingkaran yang melalui $(0,0)$, $(4,0)$, dan $(0,6)$ menggunakan bentuk umum.
  \item \textbf{(Lingkaran + Trigonometri/Pembuktian).} Buktikan bahwa untuk setiap $\theta$, titik $\big(a+r\cos\theta,\ b+r\sin\theta\big)$ selalu terletak pada lingkaran $(x-a)^2+(y-b)^2=r^2$.
\end{enumerate}

\begin{challenge}
\begin{enumerate}[leftmargin=*,topsep=2pt]
  \item Dari titik $P(7,1)$ ditarik dua garis singgung ke lingkaran $x^2+y^2=25$, menyentuhnya di $A$ dan $B$. Tentukan panjang tali busur singgung $AB$.
  \item Tentukan tempat kedudukan (lokus) titik-titik yang jaraknya ke $A(0,0)$ dua kali jaraknya ke $B(6,0)$. Tunjukkan bahwa lokusnya berupa lingkaran dan tentukan pusat serta jari-jarinya.
\end{enumerate}
\end{challenge}

% ============================================================
\chapter{Statistika Lanjutan (Data Bivariat)}
% ============================================================
\begin{tujuan}
Setelah mempelajari bab ini, peserta didik mampu:
\begin{itemize}[leftmargin=*,topsep=2pt]
  \item menyajikan data dua variabel dalam diagram pencar serta mengidentifikasi pola hubungannya;
  \item menghitung kovarians dan koefisien korelasi Pearson, lalu menafsirkan arah dan kekuatannya;
  \item menyusun persamaan garis regresi linear dan menggunakannya untuk prediksi;
  \item menghitung dan memaknai residual serta mengenali outlier;
  \item membedakan korelasi dari kausalitas dan mengenali peran variabel perancu.
\end{itemize}
\textbf{Capaian Pembelajaran (Fase F).} Peserta didik dapat menganalisis hubungan dua variabel secara kuantitatif dan menggunakannya untuk membuat prediksi yang bertanggung jawab.
\end{tujuan}

\begin{pemantik}
\begin{itemize}[leftmargin=*,topsep=2pt]
  \item Apakah makin lama belajar selalu berarti makin tinggi nilai? Seberapa erat keduanya terkait?
  \item Bagaimana menarik satu garis "terbaik" yang mewakili awan titik pada scatter plot?
  \item Penjualan es krim dan jumlah orang tenggelam naik bersamaan setiap musim panas. Apakah es krim menyebabkan orang tenggelam?
\end{itemize}
\end{pemantik}

\begin{studikasus}{Jam Belajar dan Nilai Ulangan}
Seorang guru mencatat jam belajar ($x$) dan nilai ulangan ($y$) lima siswa.
\begin{center}
\begin{tabular}{c|ccccc}
\toprule
$x$ (jam) & 1 & 2 & 3 & 4 & 5 \\
$y$ (nilai) & 2 & 3 & 5 & 4 & 6 \\
\bottomrule
\end{tabular}
\end{center}
\textbf{Amati dan duga.} (a) Apakah ada kecenderungan nilai naik seiring jam belajar? (b) Bisakah kita menarik satu garis untuk \emph{memprediksi} nilai siswa yang belajar $6$ jam? (c) Seberapa "percaya diri" prediksi itu? Bab ini menjawab ketiganya melalui \textbf{korelasi} (seberapa erat) dan \textbf{regresi} (garis prediksi). Data ini akan kita pakai berulang sebagai contoh.
\end{studikasus}

\section*{Peta Konsep Bab}
\addcontentsline{toc}{section}{Peta Konsep Bab}
\begin{center}
\begin{tikzpicture}[
  node distance=8mm,
  konsepbox/.style={draw=biru,fill=birumuda,rounded corners,align=center,inner sep=5pt,font=\small\bfseries},
  subbox/.style={draw=hijau,fill=hijaumuda,rounded corners,align=center,inner sep=4pt,font=\footnotesize},
  garis/.style={-Stealth,biru,thick}]
  \node[konsepbox] (akar) {DATA BIVARIAT};
  \node[subbox,below left=12mm and 12mm of akar] (sc) {Diagram\\pencar};
  \node[subbox,below=12mm of akar] (kor) {Kovarians \&\\korelasi $r$};
  \node[subbox,below right=12mm and 12mm of akar] (reg) {Regresi\\$\hat y=a+bx$};
  \node[subbox,below=10mm of kor] (kaus) {Korelasi $\neq$\\kausalitas};
  \node[subbox,below=10mm of reg] (res) {Prediksi,\\residual, outlier};
  \draw[garis] (akar)--(sc); \draw[garis] (akar)--(kor); \draw[garis] (akar)--(reg);
  \draw[garis] (kor)--(kaus); \draw[garis] (reg)--(res);
\end{tikzpicture}\\
\small \textbf{Peta konsep.} Dari awan titik kita ukur keeratan ($r$) dan tarik garis prediksi (regresi); keduanya harus ditafsirkan hati-hati: korelasi bukan kausalitas.
\end{center}

\begin{catatan}
\textbf{Pengingat Materi Kelas 10 (MAE Bab 7 --- Statistika).} Bab ini membangun di atas statistika deskriptif \emph{satu variabel}. Pastikan kamu ingat:
\begin{itemize}[leftmargin=*,topsep=2pt]
  \item Mean, median, modus --- cara menghitung dan menafsirkannya.
  \item Ukuran penyebaran: jangkauan, varians $s^2$, simpangan baku $s$.
  \item Penyajian data: histogram, poligon frekuensi, diagram batang dan lingkaran.
  \item Membaca dan membandingkan distribusi data dari grafik.
\end{itemize}
\end{catatan}

\section{Diagram Pencar dan Korelasi}

\subsection*{Membangun konsep: dari awan titik ke angka}
\emph{Data bivariat} mencatat dua variabel sekaligus. Scatter plot memperlihatkan polanya secara visual, tetapi kita ingin satu \emph{angka} yang mengukur seberapa erat dan ke arah mana hubungannya. Angka itu lahir dari gagasan: ketika $x$ di atas rata-ratanya, apakah $y$ juga cenderung di atas rata-ratanya?

\begin{center}
\begin{tikzpicture}
\begin{axis}[width=8.5cm,height=5.5cm,xlabel=Jam belajar ($x$),ylabel=Nilai ($y$),
  xmin=0,xmax=6,ymin=0,ymax=8,grid=both,grid style={abu}]
\addplot[only marks,mark=*,biru,mark size=2pt] coordinates {(1,2)(2,3)(3,5)(4,4)(5,6)};
\addplot[very thick,oranye,domain=0.5:5.5]{1.3+0.9*x};
\end{axis}
\end{tikzpicture}\\
\small \textbf{Korelasi positif:} titik cenderung naik. Garis oranye $\hat y=1{,}3+0{,}9x$ adalah garis regresi (akan dihitung).
\end{center}

\begin{konsep}{Kovarians dan Korelasi Pearson}
Dengan $\bar x,\bar y$ rata-rata, definisikan
\[
S_{xx}=\sum (x_i-\bar x)^2,\quad
S_{yy}=\sum (y_i-\bar y)^2,\quad
S_{xy}=\sum (x_i-\bar x)(y_i-\bar y).
\]
\textbf{Kovarians:} $\operatorname{cov}(x,y)=\dfrac{1}{n}S_{xy}$. \quad
\textbf{Koefisien korelasi Pearson:}
\[
r=\frac{S_{xy}}{\sqrt{S_{xx}\,S_{yy}}},\qquad -1\le r\le 1.
\]
Rumus hitung praktis: $S_{xy}=\sum x_iy_i-n\bar x\bar y$ dan $S_{xx}=\sum x_i^2-n\bar x^2$.
\end{konsep}

\begin{konsep}{Menafsirkan $r$}
\begin{itemize}[leftmargin=*,topsep=2pt]
  \item \textbf{Tanda:} $r>0$ korelasi positif (searah), $r<0$ negatif (berlawanan).
  \item \textbf{Kekuatan:} $|r|$ dekat $1$ = kuat, dekat $0$ = lemah. Patokan kasar: $|r|\ge0{,}7$ kuat, $0{,}4$--$0{,}7$ sedang, $<0{,}4$ lemah.
\end{itemize}
\end{konsep}

\begin{contoh}{Menghitung Korelasi}
Untuk data studi kasus ($x:1,2,3,4,5$; $y:2,3,5,4,6$), hitung $r$.

\jawab $\bar x=3,\ \bar y=4$. $\sum x_iy_i=2+6+15+16+30=69$, $\sum x_i^2=55$, $\sum y_i^2=90$.
$S_{xy}=69-5(3)(4)=9$, $S_{xx}=55-45=10$, $S_{yy}=90-80=10$.
$r=\dfrac{9}{\sqrt{10\cdot10}}=0{,}9$ --- korelasi positif yang kuat.
\end{contoh}

\section{Regresi Linear dan Prediksi}
\begin{konsep}{Garis Regresi Kuadrat Terkecil}
Garis $\hat y=a+bx$ yang meminimalkan jumlah kuadrat residual memiliki
\[
b=\frac{S_{xy}}{S_{xx}},\qquad a=\bar y-b\bar x.
\]
Garis ini selalu melewati $(\bar x,\bar y)$. \textbf{Residual} titik ke-$i$: $e_i=y_i-\hat y_i$ (selisih nilai nyata dan prediksi).
\end{konsep}

\begin{contoh}{Menyusun Garis Regresi}
Dengan data yang sama, tentukan garis regresi dan prediksi nilai untuk $x=6$.

\jawab $b=\dfrac{S_{xy}}{S_{xx}}=\dfrac{9}{10}=0{,}9$, $a=4-0{,}9(3)=1{,}3$.
Jadi $\hat y=1{,}3+0{,}9x$. Untuk $x=6$: $\hat y=1{,}3+5{,}4=6{,}7$.
\end{contoh}

\begin{kesalahan}
Prediksi regresi paling tepercaya \emph{di dalam} rentang data (interpolasi). Memakai garis untuk $x$ jauh di luar data (\emph{ekstrapolasi}), misalnya memprediksi nilai untuk $50$ jam belajar, sering menyesatkan karena pola bisa berubah.
\end{kesalahan}

\begin{dunianyata}
Sebuah \emph{outlier} --- titik yang menyimpang jauh dari pola --- dapat "menarik" garis regresi dan menggeser nilai $r$ secara dramatis. Karena itu analis data selalu memeriksa scatter plot sebelum memercayai satu angka korelasi.
\end{dunianyata}

\section{Korelasi dan Kausalitas}
\begin{konsep}{Korelasi $\neq$ Kausalitas}
Korelasi kuat \textbf{tidak} membuktikan bahwa satu variabel menyebabkan yang lain. Bisa jadi ada \emph{variabel perancu} (confounding) yang memengaruhi keduanya, atau hubungannya kebetulan. Contoh: penjualan es krim dan kasus tenggelam berkorelasi karena keduanya dipicu oleh \emph{cuaca panas}, bukan saling menyebabkan.
\end{konsep}

\begin{center}
\begin{tikzpicture}[
  kotak/.style={draw=biru,fill=birumuda,rounded corners,align=center,font=\footnotesize\bfseries,inner sep=6pt,minimum width=2.6cm},
  perancu/.style={draw=oranye,fill=oranyemuda,rounded corners,align=center,font=\footnotesize\bfseries,inner sep=6pt,minimum width=2.6cm},
  panah/.style={-Stealth,oranye,very thick}]
\node[perancu] (cuaca) at (3.5,2.2) {Cuaca Panas\\(variabel perancu)};
\node[kotak] (eskrim) at (0,0) {Es Krim $\uparrow$};
\node[kotak] (tenggelam) at (7,0) {Kasus Tenggelam $\uparrow$};
\draw[panah] (cuaca.south west)--(eskrim.north east) node[midway,above left,font=\small]{menyebabkan};
\draw[panah] (cuaca.south east)--(tenggelam.north west) node[midway,above right,font=\small]{menyebabkan};
\draw[merah,very thick,dashed] (eskrim.east)--(tenggelam.west);
\node[font=\small\bfseries,merah,fill=white,inner sep=1pt] at (3.5,0.42) {$\times$\ korelasi semu};
\node[font=\footnotesize,merah,fill=white,inner sep=1pt] at (3.5,-0.15) {BUKAN hubungan sebab-akibat};
\end{tikzpicture}\\
\small \textbf{Korelasi semu (spurious correlation).} Es krim dan kasus tenggelam berkorelasi karena keduanya naik saat cuaca panas --- bukan karena saling menyebabkan. Variabel perancu tersembunyi menciptakan korelasi palsu.
\end{center}

\section*{Ringkasan Bab}
\addcontentsline{toc}{section}{Ringkasan Bab}
\begin{itemize}[leftmargin=*,topsep=2pt]
  \item \textbf{Kovarians} $=\tfrac1n S_{xy}$; \textbf{korelasi} $r=\dfrac{S_{xy}}{\sqrt{S_{xx}S_{yy}}}\in[-1,1]$.
  \item Tanda $r$ menunjukkan arah; $|r|$ menunjukkan kekuatan hubungan linear.
  \item \textbf{Regresi:} $\hat y=a+bx$ dengan $b=\dfrac{S_{xy}}{S_{xx}}$, $a=\bar y-b\bar x$; garis melewati $(\bar x,\bar y)$.
  \item \textbf{Residual} $e_i=y_i-\hat y_i$; outlier dapat mengubah garis dan $r$.
  \item \textbf{Korelasi bukan kausalitas} --- waspadai variabel perancu.
\end{itemize}

\begin{refleksi}
\begin{itemize}[leftmargin=*,topsep=2pt]
  \item[\cek] Saya dapat menghitung $r$ dan menafsirkan arah serta kekuatannya.
  \item[\cek] Saya dapat menyusun garis regresi dan memakainya untuk prediksi.
  \item[\cek] Saya dapat menghitung residual dan menjelaskan dampak outlier.
  \item[\cek] Saya dapat menjelaskan mengapa korelasi tidak membuktikan sebab-akibat.
\end{itemize}
\end{refleksi}

\section*{Soal Akhir Bab \thechapter}
\addcontentsline{toc}{section}{Soal Akhir Bab \thechapter}

\bagiansoal{A}{(Fundamental --- setara SNBT Dasar)}
\begin{enumerate}[leftmargin=*]
  \item Diberikan data $x:2,4,6$ dan $y:1,3,5$. Hitung $\bar x$ dan $\bar y$.
  \item Pada scatter plot, titik cenderung turun dari kiri-atas ke kanan-bawah. Sebutkan jenis korelasinya.
  \item Koefisien korelasi suatu data adalah $r=0{,}85$. Tafsirkan arah dan kekuatannya.
  \item Garis regresi suatu data adalah $\hat y=3+2x$. Prediksikan $y$ saat $x=5$.
  \item Diketahui $\hat y=10-0{,}5x$. Berapa perubahan prediksi $y$ jika $x$ bertambah $1$ satuan?
\end{enumerate}

\bagiansoal{B}{(Intermediate --- setara SNBT Sulit)}
\textit{Soal B.1--B.5 memakai data} $x:1,2,3,4,5$ \textit{dan} $y:2,3,5,4,6$.
\begin{enumerate}[leftmargin=*]
  \item Hitung $S_{xx}$ dan $S_{xy}$.
  \item Tentukan persamaan garis regresi $\hat y=a+bx$.
  \item Hitung koefisien korelasi Pearson $r$.
  \item Hitung kovarians (populasi) data tersebut.
  \item Prediksikan $y$ saat $x=6$, lalu hitung residualnya jika nilai nyata $y=7$.
  \item Koefisien korelasi antara jam bermain game dan nilai ujian adalah $r=-0{,}95$. Tafsirkan, lalu jelaskan apakah ini membuktikan game menurunkan nilai.
  \item Sebuah titik $(10,2)$ menyimpang jauh dari pola data lain yang naik. Sebutkan istilahnya dan jelaskan pengaruhnya terhadap garis regresi.
\end{enumerate}

\bagiansoal{C}{(Advanced --- setara Olimpiade SMA, gabungan antarbab)}
\begin{enumerate}[leftmargin=*]
  \item \textbf{(Pembuktian).} Buktikan bahwa garis regresi $\hat y=a+bx$ selalu melalui titik $(\bar x,\bar y)$.
  \item \textbf{(Korelasi + Transformasi).} Jika setiap nilai $x$ digandakan menjadi $2x$ (nilai $y$ tetap), bagaimana koefisien korelasi $r$ berubah? Jelaskan.
  \item \textbf{(Regresi).} Diketahui $S_{xy}=24$, $S_{xx}=16$, $\bar x=5$, $\bar y=20$. Tentukan persamaan regresi dan prediksi pada $x=8$.
  \item \textbf{(Korelasi).} Diketahui $S_{xy}=18$, $S_{xx}=25$, $S_{yy}=16$. Hitung $r$ dan gradien regresi $b$.
  \item \textbf{(Regresi ganda).} Garis regresi $y$ atas $x$ adalah $\hat y=2+0{,}5x$ dengan $r=0{,}6$. Tentukan gradien garis regresi $x$ atas $y$. (Petunjuk: $b_{yx}\cdot b_{xy}=r^2$.)
  \item \textbf{(Korelasi vs Kausalitas).} Penjualan es krim dan kasus tenggelam berkorelasi positif kuat. Jelaskan dengan konsep variabel perancu mengapa hal ini bukan bukti sebab-akibat.
\end{enumerate}

\bagiansoal{D}{(Expert --- setara tahun pertama universitas, sintesis \& pembuktian)}
\begin{enumerate}[leftmargin=*]
  \item \textbf{(Pembuktian).} Buktikan bahwa $b=r\cdot\dfrac{s_y}{s_x}$ dengan $s_x,s_y$ simpangan baku $x$ dan $y$.
  \item \textbf{(Korelasi + Ketaksamaan).} Dengan ketaksamaan Cauchy--Schwarz, buktikan $|r|\le 1$.
  \item \textbf{(Regresi + Kalkulus).} Turunkan rumus $b=\dfrac{S_{xy}}{S_{xx}}$ dan $a=\bar y-b\bar x$ dengan meminimalkan $\sum (y_i-a-bx_i)^2$ memakai turunan parsial.
  \item \textbf{(Statistika + Logaritma).} Data mengikuti model $y=c\,x^{k}$. Tunjukkan bahwa $\log y$ linear terhadap $\log x$, dan jelaskan bagaimana regresi linear pada $(\log x,\log y)$ menaksir $k$.
  \item \textbf{(Korelasi).} Jika $r=0$, apakah pasti tidak ada hubungan apa pun antara $x$ dan $y$? Berikan satu contoh tandingan dengan hubungan nonlinear.
\end{enumerate}

\begin{challenge}
\begin{enumerate}[leftmargin=*,topsep=2pt]
  \item Data $x:1,2,3,4$ dan $y:1,2,3,10$. Hitung $r$. Lalu hitung ulang $r$ tanpa titik terakhir, dan jelaskan dampak outlier terhadap korelasi.
  \item Buktikan bahwa untuk garis regresi kuadrat terkecil, jumlah seluruh residual selalu nol: $\sum (y_i-\hat y_i)=0$.
\end{enumerate}
\end{challenge}

% ============================================================
\chapter{Bilangan Kompleks}
% ============================================================
\begin{tujuan}
Setelah mempelajari bab ini, peserta didik mampu:
\begin{itemize}[leftmargin=*,topsep=2pt]
  \item memahami satuan imajiner $i$ dan pola pangkatnya;
  \item melakukan penjumlahan, pengurangan, perkalian, dan pembagian bilangan kompleks;
  \item menentukan konjugat, modulus, dan argumen serta menggambarkannya pada diagram Argand;
  \item mengubah bentuk kartesius ke bentuk polar dan sebaliknya;
  \item menerapkan teorema De Moivre untuk pemangkatan dan penarikan akar;
  \item menentukan akar pangkat-$n$ suatu bilangan kompleks dan memaknainya secara geometris.
\end{itemize}
\textbf{Capaian Pembelajaran (Fase F+).} Peserta didik dapat bekerja dengan bilangan kompleks dalam bentuk aljabar dan polar serta menggunakannya untuk menyelesaikan persamaan dan masalah trigonometris.
\end{tujuan}

\begin{pemantik}
\begin{itemize}[leftmargin=*,topsep=2pt]
  \item Persamaan $x^2+1=0$ tak punya solusi real. Apa jadinya jika kita "menciptakan" bilangan yang kuadratnya $-1$?
  \item Mengapa mengalikan dengan $i$ ternyata sama dengan memutar $90^\circ$ pada bidang?
  \item Bisakah sebuah bilangan memiliki \emph{lebih dari dua} akar pangkat tiga?
\end{itemize}
\end{pemantik}

\begin{studikasus}{Putaran pada Bidang dan Sinyal AC}
Insinyur listrik menyatakan arus bolak-balik (AC) sebagai bilangan kompleks: besarnya menjadi \emph{modulus}, dan fasenya menjadi \emph{argumen}. Mengalikan dua bilangan kompleks berarti \textbf{mengalikan modulus} dan \textbf{menjumlahkan sudut} --- persis seperti memutar dan menskalakan sebuah vektor.
\begin{center}
\begin{tabular}{c|c|c}
\toprule
Operasi & Modulus & Argumen (sudut) \\
\midrule
$z_1 z_2$ & $|z_1|\,|z_2|$ & $\theta_1+\theta_2$ \\
$z_1/z_2$ & $|z_1|/|z_2|$ & $\theta_1-\theta_2$ \\
$z^n$ & $|z|^n$ & $n\theta$ \\
\bottomrule
\end{tabular}
\end{center}
\textbf{Amati dan duga.} Mengapa memangkatkan bilangan kompleks ($z^n$) cukup dengan memangkatkan modulus dan \emph{mengalikan} sudutnya? Itulah inti teorema De Moivre yang akan kita buktikan.
\end{studikasus}

\section*{Peta Konsep Bab}
\addcontentsline{toc}{section}{Peta Konsep Bab}
\begin{center}
\begin{tikzpicture}[
  node distance=8mm,
  konsepbox/.style={draw=biru,fill=birumuda,rounded corners,align=center,inner sep=5pt,font=\small\bfseries},
  subbox/.style={draw=hijau,fill=hijaumuda,rounded corners,align=center,inner sep=4pt,font=\footnotesize},
  garis/.style={-Stealth,biru,thick}]
  \node[konsepbox] (akar) {BILANGAN KOMPLEKS};
  \node[subbox,below left=12mm and 12mm of akar] (kar) {Bentuk aljabar\\$a+bi$};
  \node[subbox,below=12mm of akar] (geo) {Diagram Argand\\modulus, argumen};
  \node[subbox,below right=12mm and 12mm of akar] (pol) {Bentuk polar\\$r\,\mathrm{cis}\,\theta$};
  \node[subbox,below=10mm of kar] (op) {Operasi \&\\konjugat};
  \node[subbox,below=10mm of pol] (dm) {De Moivre \&\\akar pangkat-$n$};
  \draw[garis] (akar)--(kar); \draw[garis] (akar)--(geo); \draw[garis] (akar)--(pol);
  \draw[garis] (kar)--(op); \draw[garis] (pol)--(dm);
\end{tikzpicture}\\
\small \textbf{Peta konsep.} Satu bilangan kompleks punya dua "wajah": aljabar ($a+bi$) untuk operasi dasar, dan polar ($r\,\mathrm{cis}\,\theta$) untuk pemangkatan serta penarikan akar.
\end{center}

\section{Satuan Imajiner dan Operasi Dasar}

\subsection*{Membangun konsep: memperluas sistem bilangan}
Agar persamaan seperti $x^2=-1$ punya solusi, kita perkenalkan satuan imajiner $i$ dengan $i^2=-1$. Pangkat $i$ lalu berulang dengan periode $4$:
\[
i^1=i,\quad i^2=-1,\quad i^3=-i,\quad i^4=1,\quad i^5=i,\ \dots
\]
Untuk $i^n$, cukup tinjau sisa $n$ dibagi $4$.

\begin{definisi}
Satuan imajiner $i$ didefinisikan oleh $i^2 = -1$. \emph{Bilangan kompleks} berbentuk
$z = a + bi$, dengan $a=\operatorname{Re}(z)$ bagian real dan $b=\operatorname{Im}(z)$ bagian imajiner.
\end{definisi}

\begin{konsep}{Operasi Aljabar}
\[
(a+bi)\pm(c+di) = (a\pm c)+(b\pm d)i,
\]
\[
(a+bi)(c+di) = (ac-bd)+(ad+bc)i.
\]
\textbf{Konjugat:} $\bar z = a - bi$. \textbf{Modulus:} $|z| = \sqrt{a^2+b^2}$. \\
\textbf{Pembagian} dilakukan dengan mengalikan konjugat penyebut:
\[
\frac{z_1}{z_2}=\frac{z_1\,\bar z_2}{z_2\,\bar z_2}=\frac{z_1\,\bar z_2}{|z_2|^2}.
\]
Sifat kunci: $z\bar z=|z|^2$.
\end{konsep}

\begin{contoh}{Perkalian dan Pembagian}
Hitung $(2+i)(3-i)$ dan $\dfrac{3+2i}{1-i}$.

\jawab $(2+i)(3-i)=6-2i+3i-i^2=6+i+1=7+i$.\\
$\dfrac{3+2i}{1-i}=\dfrac{(3+2i)(1+i)}{(1-i)(1+i)}=\dfrac{3+3i+2i+2i^2}{2}=\dfrac{1+5i}{2}=\dfrac12+\dfrac52 i.$
\end{contoh}

\begin{kesalahan}
Jangan menulis $\sqrt{-4}=\sqrt{-1}\cdot\sqrt{4}=2i$ lalu menerapkan $\sqrt{a}\sqrt{b}=\sqrt{ab}$ pada bilangan negatif secara sembarangan: $\sqrt{-1}\cdot\sqrt{-1}\neq\sqrt{1}=1$, melainkan $i\cdot i=-1$. Aturan akar perkalian hanya berlaku untuk bilangan non-negatif.
\end{kesalahan}

\section{Diagram Argand, Modulus, dan Argumen}
\begin{center}
\begin{tikzpicture}[scale=0.9]
  \coordinate (O) at (0,0);
  \coordinate (Ra) at (3,0);
  \coordinate (Z) at (3,2);
  \draw[->] (-0.5,0)--(4,0) node[right]{Re};
  \draw[->] (0,-0.5)--(0,3.5) node[above]{Im};
  \draw[very thick,biru,-{Latex}] (O)--(Z) node[above right]{$z=a+bi$};
  \draw[dashed,gray] (Ra)--(Z) (0,2)--(Z);
  \node[below] at (Ra){\small $a$}; \node[left] at (0,2){\small $b$};
  \node[biru] at (1.4,1.4){\small $r=|z|$};
  \pic[draw,angle radius=0.7cm,"$\theta$"] {angle=Ra--O--Z};
\end{tikzpicture}\\
\small \textbf{Diagram Argand:} $z$ digambar sebagai titik/vektor. Panjangnya $r=|z|$ (modulus), sudutnya $\theta=\arg z$ (argumen).
\end{center}

\begin{center}
\begin{tikzpicture}[scale=0.82]
\draw[->] (-0.4,0)--(4.2,0) node[right]{\small Re};
\draw[->] (0,-0.4)--(0,5.2) node[above]{\small Im};
% Grid dots
\foreach \x in {1,2,3} \foreach \y in {1,2,3,4}
  \fill[gray!25] (\x,\y) circle(1.5pt);
% Tick labels
\foreach \x in {1,2,3} \node[below,font=\tiny] at (\x,-0.1) {\x};
\foreach \y in {1,2,3,4} \node[left,font=\tiny] at (-0.1,\y) {\y};
% z1 = 1+i
\draw[-Stealth,biru,very thick] (0,0)--(1,1);
\node[biru,font=\small,above left] at (0.5,0.5) {$z_1=1{+}i$};
% z2 = 2+3i
\draw[-Stealth,hijau,very thick] (0,0)--(2,3);
\node[hijau,font=\small,left] at (1,1.8) {$z_2=2{+}3i$};
% z1+z2 = 3+4i (resultant)
\draw[-Stealth,oranye,very thick] (0,0)--(3,4);
\node[oranye,font=\small,right] at (3,4) {$z_1{+}z_2=3{+}4i$};
% Parallelogram completion (dashed)
\draw[biru,dashed,thin] (1,1)--(3,4);
\draw[hijau,dashed,thin] (2,3)--(3,4);
\fill[oranye] (3,4) circle(2.5pt);
\fill[biru] (1,1) circle(2pt);
\fill[hijau] (2,3) circle(2pt);
\end{tikzpicture}\\
\small \textbf{Penjumlahan = penjumlahan vektor.} $(1{+}i)+(2{+}3i)=3{+}4i$; titik ujung diperoleh via aturan jajaran genjang (garis putus-putus). Bagian real dan imajiner dijumlahkan terpisah.
\end{center}

\begin{konsep}{Argumen}
Argumen $\theta=\arg z$ memenuhi $\tan\theta=\dfrac{b}{a}$, dengan \emph{kuadran $z$ menentukan} $\theta$ yang tepat. Biasanya dipilih argumen utama $-180^\circ<\theta\le180^\circ$ (atau $0\le\theta<360^\circ$).
\end{konsep}

\section{Bentuk Polar dan Teorema De Moivre}
\begin{konsep}{Bentuk Polar}
\[
z = r(\cos\theta + i\sin\theta)=r\,\mathrm{cis}\,\theta, \quad
r = \sqrt{a^2+b^2}, \quad \tan\theta = \frac{b}{a}.
\]
Perkalian/pembagian menjadi mudah: $r_1\mathrm{cis}\,\theta_1\cdot r_2\mathrm{cis}\,\theta_2=r_1r_2\,\mathrm{cis}(\theta_1+\theta_2)$.
\end{konsep}

\begin{contoh}{Kartesius ke Polar}
Ubah $z = 1 + i$ ke bentuk polar.

\jawab $r=\sqrt{1+1}=\sqrt2$, $\tan\theta = 1$ dan $z$ di kuadran I $\Rightarrow \theta = 45^\circ$.
Jadi $z = \sqrt2(\cos45^\circ + i\sin45^\circ).$
\end{contoh}

\begin{teorema}[De Moivre]
Untuk setiap bilangan bulat $n$,
\[
\big(\cos\theta+i\sin\theta\big)^n=\cos(n\theta)+i\sin(n\theta).
\]
Akibatnya $z^n=r^n\big(\cos n\theta+i\sin n\theta\big)$.
\end{teorema}

\begin{contoh}{Pemangkatan dengan De Moivre}
Hitung $(1+i)^8$.

\jawab $1+i=\sqrt2\,\mathrm{cis}\,45^\circ$. Maka $(1+i)^8=(\sqrt2)^8\,\mathrm{cis}(8\cdot45^\circ)=16\,\mathrm{cis}\,360^\circ=16(1+0i)=16.$
\end{contoh}

\section{Akar Pangkat-$n$ Bilangan Kompleks}
\begin{konsep}{Rumus Akar ke-$n$}
Akar pangkat-$n$ dari $z=r\,\mathrm{cis}\,\theta$ ada \emph{tepat $n$ buah}:
\[
w_k=\sqrt[n]{r}\,\Big(\cos\tfrac{\theta+360^\circ k}{n}+i\sin\tfrac{\theta+360^\circ k}{n}\Big),\quad k=0,1,\dots,n-1.
\]
Secara geometris, akar-akar ini adalah titik-titik sudut sebuah \textbf{segi-$n$ beraturan} pada lingkaran berjari-jari $\sqrt[n]{r}$.
\end{konsep}

\begin{contoh}{Akar Pangkat Tiga dari 8}
Tentukan semua akar pangkat tiga dari $8$.

\jawab $8=8\,\mathrm{cis}\,0^\circ$, $\sqrt[3]{8}=2$, sudut $\dfrac{0+360^\circ k}{3}$ untuk $k=0,1,2$ yaitu $0^\circ,120^\circ,240^\circ$:
\[
w_0=2,\quad w_1=2\,\mathrm{cis}\,120^\circ=-1+\sqrt3\,i,\quad w_2=2\,\mathrm{cis}\,240^\circ=-1-\sqrt3\,i.
\]
\end{contoh}

\section*{Ringkasan Bab}
\addcontentsline{toc}{section}{Ringkasan Bab}
\begin{itemize}[leftmargin=*,topsep=2pt]
  \item $i^2=-1$; pangkat $i$ berulang periode $4$.
  \item Operasi aljabar pada $a+bi$; pembagian memakai konjugat; $z\bar z=|z|^2$.
  \item \textbf{Polar:} $z=r\,\mathrm{cis}\,\theta$ dengan $r=|z|$, $\theta=\arg z$.
  \item \textbf{De Moivre:} $(\cos\theta+i\sin\theta)^n=\cos n\theta+i\sin n\theta$.
  \item \textbf{Akar ke-$n$:} ada $n$ akar, membentuk segi-$n$ beraturan pada lingkaran $\sqrt[n]{r}$.
\end{itemize}

\begin{refleksi}
\begin{itemize}[leftmargin=*,topsep=2pt]
  \item[\cek] Saya dapat melakukan keempat operasi pada bilangan kompleks, termasuk pembagian.
  \item[\cek] Saya dapat menentukan modulus, konjugat, dan argumen serta menggambarnya di diagram Argand.
  \item[\cek] Saya dapat mengubah ke bentuk polar dan memangkatkan dengan De Moivre.
  \item[\cek] Saya dapat menentukan semua akar pangkat-$n$ suatu bilangan kompleks.
\end{itemize}
\end{refleksi}

\section*{Soal Akhir Bab \thechapter}
\addcontentsline{toc}{section}{Soal Akhir Bab \thechapter}

\bagiansoal{A}{(Fundamental --- setara SNBT Dasar)}
\begin{enumerate}[leftmargin=*]
  \item Hitung $i^2$, $i^3$, $i^4$, dan $i^{10}$.
  \item Hitung $(3+2i)+(1-4i)$.
  \item Hitung $(2+i)(3-i)$.
  \item Tentukan konjugat dan modulus dari $z=3-4i$.
  \item Ubah $z=1+i$ ke bentuk polar.
\end{enumerate}

\bagiansoal{B}{(Intermediate --- setara SNBT Sulit)}
\begin{enumerate}[leftmargin=*]
  \item Hitung $(5+3i)-(2-i)$.
  \item Hitung $\dfrac{3+2i}{1-i}$ dalam bentuk $a+bi$.
  \item Tentukan modulus dan argumen dari $z=-1+\sqrt3\,i$.
  \item Hitung $i^{2025}$.
  \item Jika $z=2+i$, hitung $z\bar z$ dan tunjukkan hasilnya sama dengan $|z|^2$.
  \item Selesaikan $z^2=-9$ dalam bilangan kompleks.
  \item Ubah $z=2(\cos 60^\circ+i\sin 60^\circ)$ ke bentuk kartesius.
\end{enumerate}

\bagiansoal{C}{(Advanced --- setara Olimpiade SMA, gabungan antarbab)}
\begin{enumerate}[leftmargin=*]
  \item \textbf{(De Moivre).} Hitung $(1+i)^8$ menggunakan bentuk polar.
  \item \textbf{(Kompleks + Persamaan Kuadrat).} Selesaikan $x^2-4x+13=0$.
  \item \textbf{(Akar Pangkat-$n$).} Tentukan semua akar pangkat tiga dari $8$.
  \item \textbf{(Kompleks + Trigonometri).} Dengan $(\cos\theta+i\sin\theta)^2$, turunkan rumus $\cos2\theta$ dan $\sin2\theta$.
  \item \textbf{(Kompleks + Geometri).} Titik $z=3+4i$ dirotasi $90^\circ$ berlawanan jarum jam terhadap $O$ (yaitu dikali $i$). Tentukan hasilnya.
  \item \textbf{(Sifat Modulus).} Buktikan $|z_1 z_2|=|z_1|\,|z_2|$ menggunakan bentuk polar.
\end{enumerate}

\bagiansoal{D}{(Expert --- setara tahun pertama universitas, sintesis \& pembuktian)}
\begin{enumerate}[leftmargin=*]
  \item \textbf{(Pembuktian).} Buktikan teorema De Moivre untuk setiap bilangan asli $n$ dengan induksi.
  \item \textbf{(Akar Persatuan).} Tentukan semua akar $z^4=1$ dan tunjukkan jumlahnya nol.
  \item \textbf{(Kompleks + Barisan/Deret).} Hitung jumlah $1+i+i^2+\cdots+i^{2025}$.
  \item \textbf{(Kompleks + Trigonometri).} Dengan $(\cos\theta+i\sin\theta)^3$, buktikan $\cos3\theta=4\cos^3\theta-3\cos\theta$.
  \item \textbf{(Akar Pangkat-$n$ + Pembuktian).} Tunjukkan bahwa akar-akar pangkat-$n$ dari $1$ membentuk segi-$n$ beraturan, dan jumlah seluruhnya nol untuk $n\ge2$.
\end{enumerate}

\begin{challenge}
\begin{enumerate}[leftmargin=*,topsep=2pt]
  \item Hitung $\big(1+i\sqrt3\big)^{12}$.
  \item Tentukan semua akar pangkat enam dari $-64$ dalam bentuk polar.
  \item Buktikan: jika $z+\dfrac1z=2\cos\theta$, maka $z^n+\dfrac{1}{z^n}=2\cos(n\theta)$ untuk setiap bilangan asli $n$.
\end{enumerate}
\end{challenge}

% ============================================================
\chapter{Polinomial (Suku Banyak)}
% ============================================================
\begin{tujuan}
Setelah mempelajari bab ini, peserta didik mampu:
\begin{itemize}[leftmargin=*,topsep=2pt]
  \item melakukan operasi aljabar dan pembagian suku banyak (bersusun maupun Horner);
  \item menerapkan teorema sisa dan teorema faktor;
  \item menentukan akar-akar rasional melalui teorema akar rasional dan memahami multiplisitas akar;
  \item menggunakan hubungan akar dan koefisien (rumus Vieta);
  \item menggambarkan perilaku grafik polinomial dari derajat, akar, dan koefisien utama.
\end{itemize}
\textbf{Capaian Pembelajaran (Fase F+).} Peserta didik dapat menganalisis suku banyak melalui pembagian, pemfaktoran, dan keterkaitan akar--koefisien untuk menyelesaikan persamaan polinomial.
\end{tujuan}

\begin{pemantik}
\begin{itemize}[leftmargin=*,topsep=2pt]
  \item Bagaimana mengetahui sisa pembagian $P(x)$ oleh $(x-2)$ \emph{tanpa} membagi panjang?
  \item Jika sebuah persamaan kubik punya akar bulat, di mana kita "menebaknya"?
  \item Tanpa mencari akarnya, bisakah kita tahu jumlah dan hasil kali akar sebuah polinomial?
\end{itemize}
\end{pemantik}

\begin{studikasus}{Mendesain Lintasan Roller Coaster}
Seorang perancang wahana memodelkan ketinggian rel dengan polinomial $h(x)$. Titik saat rel menyentuh tanah adalah \emph{akar} $h(x)=0$; bentuk naik-turunnya ditentukan derajat dan koefisien utama.
\begin{center}
\begin{tikzpicture}
\begin{axis}[axis lines=middle,width=9cm,height=4.6cm,
  xlabel=$x$,ylabel=$h$,xmin=-2.6,xmax=3.6,ymin=-6,ymax=8,xtick={-2,1,3},ytick=\empty,samples=80]
\addplot[very thick,biru,domain=-2.4:3.4]{(x+2)*(x-1)*(x-3)*0.6};
\addplot[only marks,mark=*,oranye,mark size=1.6pt] coordinates {(-2,0)(1,0)(3,0)};
\end{axis}
\end{tikzpicture}\\
\small $h(x)=0{,}6(x+2)(x-1)(x-3)$: tiga akar $x=-2,1,3$ adalah titik sentuh tanah.
\end{center}
\textbf{Amati dan duga.} (a) Berapa kali kurva memotong sumbu-$x$, dan apa kaitannya dengan derajat? (b) Apa yang terjadi pada ujung kiri dan kanan kurva? (c) Jika dua akar berimpit, bagaimana bentuk kurvanya? Pertanyaan ini menghubungkan \textbf{akar}, \textbf{multiplisitas}, dan \textbf{grafik}.
\end{studikasus}

\section*{Peta Konsep Bab}
\addcontentsline{toc}{section}{Peta Konsep Bab}
\begin{center}
\begin{tikzpicture}[
  node distance=8mm,
  konsepbox/.style={draw=biru,fill=birumuda,rounded corners,align=center,inner sep=5pt,font=\small\bfseries},
  subbox/.style={draw=hijau,fill=hijaumuda,rounded corners,align=center,inner sep=4pt,font=\footnotesize},
  garis/.style={-Stealth,biru,thick}]
  \node[konsepbox] (akar) {POLINOMIAL};
  \node[subbox,below left=12mm and 14mm of akar] (bagi) {Pembagian\\\&\ Horner};
  \node[subbox,below=12mm of akar] (teo) {Teorema sisa\\\&\ faktor};
  \node[subbox,below right=12mm and 14mm of akar] (akr) {Akar rasional\\\&\ multiplisitas};
  \node[subbox,below=10mm of teo] (vieta) {Vieta\\(akar--koefisien)};
  \node[subbox,below=10mm of akr] (graf) {Grafik\\polinomial};
  \draw[garis] (akar)--(bagi); \draw[garis] (akar)--(teo); \draw[garis] (akar)--(akr);
  \draw[garis] (teo)--(vieta); \draw[garis] (akr)--(graf);
\end{tikzpicture}\\
\small \textbf{Peta konsep.} Pembagian melahirkan teorema sisa/faktor, yang membuka jalan ke akar rasional, multiplisitas, relasi Vieta, dan bentuk grafik.
\end{center}

\section{Pembagian, Teorema Sisa, dan Faktor}

\subsection*{Membangun konsep: mengapa sisa = $P(k)$}
Bila $P(x)$ dibagi $(x-k)$ menghasilkan hasil bagi $H(x)$ dan sisa $S$ (konstanta), maka $P(x)=(x-k)H(x)+S$. Substitusi $x=k$ menghapus suku pertama, menyisakan $P(k)=S$. Inilah dasar teorema sisa --- mengevaluasi jauh lebih cepat daripada membagi panjang.

\begin{konsep}{Dua Teorema Penting}
\begin{itemize}[leftmargin=*,topsep=2pt]
  \item \textbf{Teorema Sisa:} jika $P(x)$ dibagi $(x-k)$, sisanya adalah $P(k)$.
  \item \textbf{Teorema Faktor:} $(x-k)$ adalah faktor $P(x)$ jika dan hanya jika $P(k)=0$.
\end{itemize}
\end{konsep}

\begin{contoh}{Horner dan Teorema Faktor}
Bagi $P(x) = x^3 - 2x^2 - 5x + 6$ oleh $(x-1)$ dengan skema Horner.

\jawab Susun koefisien $1,\ -2,\ -5,\ 6$ dengan pembagi $k=1$:
\[
\begin{array}{r|rrrr}
1 & 1 & -2 & -5 & 6\\
  &   & 1  & -1 & -6\\\hline
  & 1 & -1 & -6 & \boxed{0}
\end{array}
\]
Sisa $=0$, maka $(x-1)$ faktor dari $P(x)$, dan
$P(x) = (x-1)(x^2 - x - 6) = (x-1)(x-3)(x+2).$
\end{contoh}

\section{Akar Rasional dan Multiplisitas}
\begin{konsep}{Teorema Akar Rasional}
Jika polinomial berkoefisien bulat $a_nx^n+\cdots+a_0$ memiliki akar rasional $\dfrac{p}{q}$ (bentuk paling sederhana), maka $p\mid a_0$ (membagi konstanta) dan $q\mid a_n$ (membagi koefisien utama). Ini mempersempit daftar "tebakan" akar.
\end{konsep}

\begin{konsep}{Multiplisitas Akar}
Jika $(x-k)^m$ membagi $P(x)$ tetapi $(x-k)^{m+1}$ tidak, maka $k$ adalah akar bermultiplisitas $m$. Pada grafik: multiplisitas \emph{ganjil} membuat kurva \emph{memotong} sumbu-$x$, sedangkan multiplisitas \emph{genap} membuat kurva \emph{menyinggung} lalu berbalik.
\end{konsep}

\begin{contoh}{Mencari Akar Rasional}
Tentukan semua akar $P(x)=2x^3-3x^2-3x+2$.

\jawab Kandidat $\dfrac pq$ dengan $p\mid2$, $q\mid2$: $\pm1,\pm2,\pm\tfrac12$. Uji $x=2$: $16-12-6+2=0$, jadi $(x-2)$ faktor. Horner memberi $2x^3-3x^2-3x+2=(x-2)(2x^2+x-1)=(x-2)(2x-1)(x+1)$. Akar: $x=2,\ \tfrac12,\ -1$.
\end{contoh}

\section{Hubungan Akar dan Koefisien (Vieta)}
\begin{konsep}{Rumus Vieta}
Untuk $ax^2+bx+c=0$ dengan akar $x_1,x_2$:
\[
x_1+x_2=-\frac{b}{a},\qquad x_1x_2=\frac{c}{a}.
\]
Untuk $ax^3+bx^2+cx+d=0$ dengan akar $x_1,x_2,x_3$:
\[
\sum x_i=-\frac{b}{a},\quad \sum_{i<j}x_ix_j=\frac{c}{a},\quad x_1x_2x_3=-\frac{d}{a}.
\]
\end{konsep}

\begin{contoh}{Memakai Vieta}
Akar-akar $x^2-5x+6=0$ adalah $x_1,x_2$. Tanpa mencari akarnya, hitung $x_1^2+x_2^2$.

\jawab $x_1+x_2=5$, $x_1x_2=6$. Maka $x_1^2+x_2^2=(x_1+x_2)^2-2x_1x_2=25-12=13.$
\end{contoh}

\begin{kesalahan}
Pada Vieta untuk kubik, perhatikan tanda berselang-seling: $\sum x_i=-\tfrac ba$ (negatif) tetapi $x_1x_2x_3=-\tfrac da$ (juga negatif untuk derajat ganjil). Salah tanda adalah kekeliruan tersering.
\end{kesalahan}

\section*{Ringkasan Bab}
\addcontentsline{toc}{section}{Ringkasan Bab}
\begin{itemize}[leftmargin=*,topsep=2pt]
  \item \textbf{Teorema sisa:} sisa $P(x)\div(x-k)$ adalah $P(k)$; \textbf{faktor:} $(x-k)\mid P \Leftrightarrow P(k)=0$.
  \item \textbf{Horner} mempercepat pembagian oleh $(x-k)$.
  \item \textbf{Akar rasional} $\tfrac pq$: $p\mid a_0$, $q\mid a_n$.
  \item \textbf{Multiplisitas} ganjil memotong, genap menyinggung sumbu-$x$.
  \item \textbf{Vieta} menghubungkan jumlah/hasil kali akar dengan koefisien.
\end{itemize}

\begin{refleksi}
\begin{itemize}[leftmargin=*,topsep=2pt]
  \item[\cek] Saya dapat membagi polinomial dengan Horner dan memakai teorema sisa/faktor.
  \item[\cek] Saya dapat mendaftar dan menguji kandidat akar rasional.
  \item[\cek] Saya dapat menjelaskan multiplisitas dan kaitannya dengan grafik.
  \item[\cek] Saya dapat memakai Vieta untuk menghitung ekspresi simetri akar.
\end{itemize}
\end{refleksi}

\section*{Soal Akhir Bab \thechapter}
\addcontentsline{toc}{section}{Soal Akhir Bab \thechapter}

\bagiansoal{A}{(Fundamental --- setara SNBT Dasar)}
\begin{enumerate}[leftmargin=*]
  \item Tentukan sisa pembagian $P(x)=2x^3+x^2-7$ oleh $(x-2)$.
  \item Tunjukkan bahwa $(x-1)$ adalah faktor $P(x)=x^3-1$.
  \item Faktorkan $x^3+2x^2-5x-6$ sepenuhnya (diketahui $x=-1$ adalah akar).
  \item Akar-akar $x^2-7x+10=0$ adalah $x_1,x_2$. Tentukan $x_1+x_2$ dan $x_1x_2$.
  \item Daftarkan semua kandidat akar rasional dari $P(x)=x^3-4x+1$.
\end{enumerate}

\bagiansoal{B}{(Intermediate --- setara SNBT Sulit)}
\begin{enumerate}[leftmargin=*]
  \item Tentukan $a$ jika $(x+2)$ adalah faktor dari $P(x)=x^3+ax^2-x+6$.
  \item Tentukan semua akar $P(x)=2x^3-3x^2-3x+2$.
  \item Akar-akar $x^2-5x+6=0$ adalah $x_1,x_2$. Hitung $x_1^2+x_2^2$ dan $\dfrac{1}{x_1}+\dfrac{1}{x_2}$.
  \item Polinomial $P(x)=x^3-6x^2+11x-6$. Tentukan ketiga akarnya menggunakan teorema akar rasional.
  \item Jika $P(x)$ dibagi $(x-1)$ bersisa $3$ dan dibagi $(x+2)$ bersisa $-3$, tentukan sisa pembagian $P(x)$ oleh $(x-1)(x+2)$.
  \item Tentukan nilai $k$ agar $x=3$ menjadi akar bermultiplisitas $\ge2$ dari $P(x)=x^3-9x^2+kx-27$. (Petunjuk: $P(3)=0$ dan $P'(3)=0$.)
\end{enumerate}

\bagiansoal{C}{(Advanced --- setara Olimpiade SMA, gabungan antarbab)}
\begin{enumerate}[leftmargin=*]
  \item \textbf{(Vieta + Barisan).} Akar-akar $x^3-6x^2+11x-6=0$ membentuk barisan aritmetika. Verifikasi dan tentukan bedanya.
  \item \textbf{(Polinomial + Eksponen).} Dengan substitusi $u=2^x$, selesaikan $8^x-7\cdot4^x+14\cdot2^x-8=0$ untuk $x$.
  \item \textbf{(Vieta).} Jika $x_1,x_2,x_3$ akar $x^3+px+q=0$, nyatakan $x_1^2+x_2^2+x_3^2$ dalam $p,q$.
  \item \textbf{(Polinomial + Pembuktian).} Buktikan bahwa jika polinomial berkoefisien bulat bernilai ganjil di $x=0$ dan di $x=1$, maka ia tak memiliki akar bulat.
  \item \textbf{(Polinomial + Grafik).} Sketsa perilaku $P(x)=(x-1)^2(x+2)$: tentukan akar, multiplisitasnya, serta perilaku ujung kurva.
  \item \textbf{(Vieta + Simetri).} Akar $x^3-3x^2+kx-1=0$ adalah $a,b,c$. Jika $a+b+c=ab+bc+ca$, tentukan $k$.
\end{enumerate}

\bagiansoal{D}{(Expert --- setara tahun pertama universitas, sintesis \& pembuktian)}
\begin{enumerate}[leftmargin=*]
  \item \textbf{(Pembuktian Vieta).} Untuk $x^2+bx+c=(x-x_1)(x-x_2)$, buktikan $x_1+x_2=-b$ dan $x_1x_2=c$ dengan menjabarkan ruas kanan.
  \item \textbf{(Polinomial + Bilangan Kompleks).} Tunjukkan bahwa akar tak-real polinomial berkoefisien real selalu hadir berpasangan konjugat. Lalu tentukan polinomial kubik berkoefisien real bermonik dengan akar $2$ dan $1+i$.
  \item \textbf{(Identitas Newton).} Dengan $x_1,x_2,x_3$ akar $x^3+px+q=0$, hitung $x_1^3+x_2^3+x_3^3$ dalam $p,q$.
  \item \textbf{(Polinomial + Kalkulus).} Buktikan bahwa $k$ adalah akar bermultiplisitas $\ge2$ dari $P$ jika dan hanya jika $P(k)=0$ dan $P'(k)=0$.
  \item \textbf{(Polinomial + Pembuktian).} Buktikan bahwa polinomial berderajat $n$ memiliki paling banyak $n$ akar berbeda.
\end{enumerate}

\begin{challenge}
\begin{enumerate}[leftmargin=*,topsep=2pt]
  \item Diketahui $x+\dfrac1x=3$. Tentukan $x^3+\dfrac1{x^3}$ tanpa mencari $x$.
  \item Polinomial monik berderajat $4$ memenuhi $P(1)=1$, $P(2)=4$, $P(3)=9$, $P(4)=16$. Tentukan $P(5)$. (Petunjuk: tinjau $P(x)-x^2$.)
  \item Buktikan bahwa jumlah kuadrat akar-akar persamaan $x^n+a_{n-1}x^{n-1}+\cdots+a_0=0$ sama dengan $a_{n-1}^2-2a_{n-2}$.
\end{enumerate}
\end{challenge}

% ============================================================
\chapter{Matriks}
% ============================================================
\begin{tujuan}
Setelah mempelajari bab ini, peserta didik mampu:
\begin{itemize}[leftmargin=*,topsep=2pt]
  \item mengenali jenis-jenis matriks (identitas, diagonal, segitiga, simetris) dan melakukan operasinya;
  \item menghitung determinan ordo $2$ dan ordo $3$ (aturan Sarrus dan ekspansi kofaktor);
  \item menentukan minor, kofaktor, adjoin, dan invers matriks;
  \item menyelesaikan sistem persamaan linear dengan invers dan aturan Cramer;
  \item menggunakan matriks untuk memodelkan data dan transformasi (kaitan dengan AI \& Data Science).
\end{itemize}
\textbf{Capaian Pembelajaran (Fase F+).} Peserta didik dapat menggunakan matriks dan determinan untuk menyelesaikan sistem persamaan linear serta merepresentasikan data dan transformasi.
\end{tujuan}

\begin{pemantik}
\begin{itemize}[leftmargin=*,topsep=2pt]
  \item Bagaimana satu "kotak angka" bisa menyimpan seluruh sistem persamaan sekaligus?
  \item Mengapa beberapa matriks bisa dibalik (punya invers) dan sebagian lain tidak?
  \item Apa makna geometris dari determinan bernilai nol?
\end{itemize}
\end{pemantik}

\begin{studikasus}{Tabel Penjualan sebagai Matriks}
Sebuah toko mencatat penjualan dua produk di tiga cabang. Data alami berbentuk \emph{matriks}: baris = produk, kolom = cabang. Operasi matriks lalu menjadi bahasa untuk mengolahnya --- penjumlahan menggabungkan dua periode, perkalian dengan vektor harga menghitung pendapatan.
\begin{center}
\begin{tabular}{c|ccc}
\toprule
& Cabang 1 & Cabang 2 & Cabang 3 \\
\midrule
Produk A & 12 & 9 & 15 \\
Produk B & 7 & 11 & 8 \\
\bottomrule
\end{tabular}
\end{center}
\textbf{Amati dan duga.} Jika harga A dan B adalah Rp$20$rb dan Rp$35$rb, bagaimana perkalian matriks $\times$ vektor menghasilkan pendapatan tiap cabang sekaligus? Representasi inilah yang dipakai mesin dalam pemrosesan data dan jaringan saraf tiruan.
\end{studikasus}

\section*{Peta Konsep Bab}
\addcontentsline{toc}{section}{Peta Konsep Bab}
\begin{center}
\begin{tikzpicture}[
  node distance=8mm,
  konsepbox/.style={draw=biru,fill=birumuda,rounded corners,align=center,inner sep=5pt,font=\small\bfseries},
  subbox/.style={draw=hijau,fill=hijaumuda,rounded corners,align=center,inner sep=4pt,font=\footnotesize},
  garis/.style={-Stealth,biru,thick}]
  \node[konsepbox] (akar) {MATRIKS};
  \node[subbox,below left=12mm and 14mm of akar] (op) {Jenis \&\\operasi};
  \node[subbox,below=12mm of akar] (det) {Determinan\\(Sarrus)};
  \node[subbox,below right=12mm and 14mm of akar] (inv) {Kofaktor,\\adjoin, invers};
  \node[subbox,below=10mm of det] (spl) {SPL \&\\Cramer};
  \node[subbox,below=10mm of inv] (apl) {Transformasi\\\&\ data};
  \draw[garis] (akar)--(op); \draw[garis] (akar)--(det); \draw[garis] (akar)--(inv);
  \draw[garis] (det)--(spl); \draw[garis] (inv)--(apl);
\end{tikzpicture}\\
\small \textbf{Peta konsep.} Dari operasi dan jenis matriks tumbuh determinan, lalu invers (via kofaktor/adjoin), yang menyelesaikan SPL dan merepresentasikan transformasi serta data.
\end{center}

\section{Jenis Matriks dan Operasi}
\begin{konsep}{Beberapa Jenis Matriks}
\begin{itemize}[leftmargin=*,topsep=2pt]
  \item \textbf{Identitas} $I$: diagonal $1$, lainnya $0$; $AI=IA=A$.
  \item \textbf{Diagonal:} elemen di luar diagonal utama nol.
  \item \textbf{Segitiga} atas/bawah: nol di bawah/atas diagonal.
  \item \textbf{Simetris:} $A^\top=A$ (elemen $a_{ij}=a_{ji}$).
\end{itemize}
Perkalian matriks \emph{tidak komutatif}: umumnya $AB\neq BA$.
\end{konsep}

\section{Determinan Ordo 2 dan Ordo 3}
\begin{konsep}{Determinan}
\textbf{Ordo $2$:} untuk $A=\begin{pmatrix}a&b\\c&d\end{pmatrix}$, $\det A=ad-bc$.\\[4pt]
\textbf{Ordo $3$ (aturan Sarrus):}
\[
\det\begin{pmatrix}a&b&c\\d&e&f\\g&h&i\end{pmatrix}
=aei+bfg+cdh-ceg-afh-bdi.
\]
\end{konsep}

\begin{contoh}{Determinan Ordo 3 (Sarrus)}
Hitung $\det\begin{pmatrix}1&2&3\\0&1&4\\5&6&0\end{pmatrix}$.

\jawab $=(1)(1)(0)+(2)(4)(5)+(3)(0)(6)-(3)(1)(5)-(1)(4)(6)-(2)(0)(0)$
$=0+40+0-15-24-0=1.$
\end{contoh}

\section{Kofaktor, Adjoin, dan Invers}
\begin{konsep}{Invers via Adjoin}
\textbf{Minor} $M_{ij}$ = determinan setelah baris $i$ dan kolom $j$ dibuang. \textbf{Kofaktor} $C_{ij}=(-1)^{i+j}M_{ij}$. \textbf{Adjoin} = transpos matriks kofaktor. Maka
\[
A^{-1}=\frac{1}{\det A}\,\operatorname{adj}(A),\qquad \det A\neq0.
\]
Untuk ordo $2$: $A^{-1}=\dfrac{1}{ad-bc}\begin{pmatrix}d&-b\\-c&a\end{pmatrix}$.
\end{konsep}

\begin{contoh}{Determinan dan Invers Ordo 2}
Diketahui $A = \begin{pmatrix} 2 & 1\\ 3 & 4 \end{pmatrix}$. Tentukan $\det A$ dan $A^{-1}$.

\jawab $\det A = (2)(4)-(1)(3) = 5$.
$A^{-1} = \dfrac{1}{5}\begin{pmatrix} 4 & -1\\ -3 & 2 \end{pmatrix}.$
\end{contoh}

\begin{kesalahan}
Matriks dengan $\det A=0$ disebut \emph{singular} dan \textbf{tak punya invers}. Selalu cek determinan sebelum menulis $A^{-1}$; jangan pula menukar urutan pada $\operatorname{adj}$ ordo $2$ (yang ditukar diagonal utamanya, yang dinegasikan diagonal lainnya).
\end{kesalahan}

\section{Penyelesaian SPL dengan Matriks}
\begin{konsep}{Dua Cara}
Untuk $A\mathbf x=\mathbf b$ dengan $\det A\neq0$:
\[
\textbf{Invers: } \mathbf x=A^{-1}\mathbf b, \qquad
\textbf{Cramer: } x_i=\frac{\det A_i}{\det A},
\]
dengan $A_i$ = $A$ yang kolom ke-$i$-nya diganti $\mathbf b$.
\end{konsep}

\begin{contoh}{Sistem Persamaan}
Selesaikan $\begin{cases} 2x + y = 5\\ 3x + 4y = 10 \end{cases}$ dengan matriks.

\jawab Bentuk $A\mathbf{x}=\mathbf{b}$ dengan $A=\begin{pmatrix}2&1\\3&4\end{pmatrix}$,
$\mathbf{b}=\begin{pmatrix}5\\10\end{pmatrix}$. Maka
\[
\mathbf{x}=A^{-1}\mathbf{b}=\tfrac15\begin{pmatrix}4&-1\\-3&2\end{pmatrix}\begin{pmatrix}5\\10\end{pmatrix}
=\tfrac15\begin{pmatrix}10\\5\end{pmatrix}=\begin{pmatrix}2\\1\end{pmatrix}.
\]
Jadi $x=2,\ y=1$.
\end{contoh}

\begin{dunianyata}
Pada \emph{Data Science} dan \emph{AI}, gambar, teks, dan tabel direpresentasikan sebagai matriks; jaringan saraf tiruan pada dasarnya adalah rangkaian perkalian matriks. Determinan dan invers menjadi alat untuk menyelesaikan sistem dan menstabilkan perhitungan.
\end{dunianyata}

\begin{center}
\begin{tikzpicture}[scale=1.15]
% Original unit square
\fill[birumuda] (0,0) rectangle (1,1);
\draw[biru,very thick] (0,0) rectangle (1,1);
\node[font=\scriptsize,biru] at (0.5,0.5) {luas $=1$};
\node[font=\footnotesize,biru,above] at (0.5,1.08) {Persegi Satuan};
% Arrow with matrix label
\draw[-Stealth,very thick,oranye] (1.3,0.5)--(2.3,0.5);
\node[font=\tiny,oranye,above] at (1.8,0.6)
  {$M=\bigl(\begin{smallmatrix}2&1\\0&1\end{smallmatrix}\bigr)$};
\node[font=\tiny,oranye,below] at (1.8,0.4){$\det M=2$};
% Transformed parallelogram: corners (0,0)→(0,0),(1,0)→(2,0),(1,1)→(3,1),(0,1)→(1,1)
% Shift all by +2.6 on x-axis
\fill[oranyemuda] (2.6,0)--(4.6,0)--(5.6,1)--(3.6,1)--cycle;
\draw[oranye,very thick] (2.6,0)--(4.6,0)--(5.6,1)--(3.6,1)--cycle;
\node[font=\scriptsize,oranye] at (4.1,0.5) {luas $=2$};
\node[font=\footnotesize,oranye,above] at (4.1,1.08) {Jajar Genjang};
% Det annotation
\node[font=\tiny,merah,align=center] at (4.1,-0.35) {$\det M=2$ $\Rightarrow$ luas berlipat 2};
\end{tikzpicture}\\
\small \textbf{Determinan = faktor skala luas.} Matriks $M$ mengubah persegi satuan (luas 1) menjadi jajar genjang (luas 2); $|\det M|$ mengukur perubahan luas. Jika $\det M=0$, bidang "runtuh" menjadi garis --- tak punya invers.
\end{center}

\section*{Ringkasan Bab}
\addcontentsline{toc}{section}{Ringkasan Bab}
\begin{itemize}[leftmargin=*,topsep=2pt]
  \item Jenis matriks: identitas, diagonal, segitiga, simetris; perkalian tak komutatif.
  \item \textbf{Determinan} ordo $2$: $ad-bc$; ordo $3$: aturan Sarrus.
  \item \textbf{Invers:} $A^{-1}=\tfrac{1}{\det A}\operatorname{adj}(A)$; ada hanya jika $\det A\neq0$.
  \item \textbf{SPL:} $\mathbf x=A^{-1}\mathbf b$ atau aturan Cramer $x_i=\tfrac{\det A_i}{\det A}$.
  \item Matriks merepresentasikan data dan transformasi.
\end{itemize}

\begin{refleksi}
\begin{itemize}[leftmargin=*,topsep=2pt]
  \item[\cek] Saya dapat mengenali jenis matriks dan melakukan operasinya.
  \item[\cek] Saya dapat menghitung determinan ordo $2$ dan ordo $3$.
  \item[\cek] Saya dapat menentukan invers via adjoin dan menyelesaikan SPL.
  \item[\cek] Saya dapat menjelaskan peran matriks dalam data dan transformasi.
\end{itemize}
\end{refleksi}

\section*{Soal Akhir Bab \thechapter}
\addcontentsline{toc}{section}{Soal Akhir Bab \thechapter}

\bagiansoal{A}{(Fundamental --- setara SNBT Dasar)}
\begin{enumerate}[leftmargin=*]
  \item Diketahui $A=\begin{pmatrix}1&2\\3&4\end{pmatrix}$, $B=\begin{pmatrix}0&1\\-1&2\end{pmatrix}$. Hitung $A+B$.
  \item Hitung $\det\begin{pmatrix}3&2\\1&4\end{pmatrix}$.
  \item Tentukan invers dari $\begin{pmatrix}1&2\\3&5\end{pmatrix}$.
  \item Hitung hasil kali $\begin{pmatrix}1&2\\3&4\end{pmatrix}\begin{pmatrix}1\\1\end{pmatrix}$.
  \item Sebutkan jenis matriks $\begin{pmatrix}2&0&0\\0&5&0\\0&0&3\end{pmatrix}$.
\end{enumerate}

\bagiansoal{B}{(Intermediate --- setara SNBT Sulit)}
\begin{enumerate}[leftmargin=*]
  \item Hitung $\det\begin{pmatrix}1&2&3\\0&1&4\\5&6&0\end{pmatrix}$ dengan aturan Sarrus.
  \item Diketahui $A=\begin{pmatrix}2&1\\1&3\end{pmatrix}$. Tunjukkan $A$ simetris dan hitung $A^2$.
  \item Selesaikan $\begin{cases} x+2y=4\\ 3x+5y=7 \end{cases}$ dengan invers matriks.
  \item Selesaikan sistem yang sama dengan aturan Cramer.
  \item Tentukan nilai $k$ agar $\begin{pmatrix}k&2\\8&k\end{pmatrix}$ singular (tak punya invers).
  \item Jika $\det A=3$ dan $A$ berordo $2$, tentukan $\det(2A)$ dan $\det(A^{-1})$.
  \item Diketahui $AB=\begin{pmatrix}1&0\\0&1\end{pmatrix}$ dengan $A=\begin{pmatrix}2&1\\3&2\end{pmatrix}$. Tentukan $B$.
\end{enumerate}

\bagiansoal{C}{(Advanced --- setara Olimpiade SMA, gabungan antarbab)}
\begin{enumerate}[leftmargin=*]
  \item \textbf{(Matriks + Transformasi).} Matriks $R=\begin{pmatrix}0&-1\\1&0\end{pmatrix}$ merotasi $90^\circ$. Hitung $R^4$ dan jelaskan maknanya secara geometris.
  \item \textbf{(Matriks + SPL 3 variabel).} Hitung determinan koefisien sistem $\begin{cases}x+y+z=6\\ 2x-y+z=3\\ x+2y-z=2\end{cases}$ dan tentukan $x$ dengan aturan Cramer.
  \item \textbf{(Determinan + Aljabar).} Buktikan $\det(AB)=\det A\cdot\det B$ untuk matriks ordo $2$ sembarang.
  \item \textbf{(Matriks + Barisan).} Untuk $A=\begin{pmatrix}1&1\\0&1\end{pmatrix}$, tentukan rumus $A^n$ dan buktikan dengan induksi.
  \item \textbf{(Matriks + Data).} Tabel penjualan $\begin{pmatrix}12&9&15\\7&11&8\end{pmatrix}$ (baris=produk, kolom=cabang). Dengan harga $\begin{pmatrix}20\\35\end{pmatrix}$ ribu, tentukan pendapatan tiap cabang melalui perkalian matriks yang sesuai.
  \item \textbf{(Determinan + Geometri).} Tunjukkan bahwa luas segitiga bertitik sudut $(x_1,y_1)$, $(x_2,y_2)$, $(x_3,y_3)$ adalah $\tfrac12|D|$, dengan
  \[
  D=\det\begin{pmatrix}x_2-x_1&x_3-x_1\\y_2-y_1&y_3-y_1\end{pmatrix},
  \]
  lalu hitung luas segitiga $(0,0),(4,0),(1,3)$.
\end{enumerate}

\bagiansoal{D}{(Expert --- setara tahun pertama universitas, sintesis \& pembuktian)}
\begin{enumerate}[leftmargin=*]
  \item \textbf{(Cayley--Hamilton ordo 2).} Buktikan bahwa setiap matriks $A$ ordo $2$ memenuhi $A^2-(\operatorname{tr}A)A+(\det A)I=O$.
  \item \textbf{(Invers + Pembuktian).} Buktikan bahwa jika $A$ dan $B$ dapat dibalik, maka $(AB)^{-1}=B^{-1}A^{-1}$.
  \item \textbf{(Determinan + Sifat).} Buktikan bahwa $\det(A^{-1})=\dfrac{1}{\det A}$ menggunakan sifat $\det(AB)=\det A\det B$.
  \item \textbf{(Matriks + Diagonalisasi sederhana).} Diketahui $A=\begin{pmatrix}2&0\\0&3\end{pmatrix}$. Tentukan rumus umum $A^n$ dan jelaskan mengapa matriks diagonal mudah dipangkatkan.
  \item \textbf{(Matriks + Sistem tak-tunggal).} Tunjukkan bahwa jika $\det A=0$, maka $A\mathbf x=\mathbf b$ tidak memiliki solusi tunggal, dan beri contoh kasus tanpa solusi serta kasus tak hingga solusi.
\end{enumerate}

\begin{challenge}
\begin{enumerate}[leftmargin=*,topsep=2pt]
  \item Untuk $F=\begin{pmatrix}1&1\\1&0\end{pmatrix}$, buktikan bahwa $F^n=\begin{pmatrix}F_{n+1}&F_n\\F_n&F_{n-1}\end{pmatrix}$ dengan $F_n$ bilangan Fibonacci, lalu simpulkan identitas $F_{n+1}F_{n-1}-F_n^2=(-1)^n$.
  \item Tentukan semua matriks $A$ ordo $2$ yang memenuhi $A^2=I$ dengan $A\neq\pm I$.
\end{enumerate}
\end{challenge}

% ============================================================
\chapter{Geometri Transformasi}
% ============================================================
\begin{tujuan}
Setelah mempelajari bab ini, peserta didik mampu:
\begin{itemize}[leftmargin=*,topsep=2pt]
  \item menentukan bayangan titik/bangun oleh translasi, refleksi, rotasi, dan dilatasi;
  \item menyatakan transformasi sebagai operasi matriks dan menghitung bayangan dengan perkalian matriks;
  \item menyusun transformasi berurutan (komposisi) dan menentukan matriks gabungannya;
  \item menentukan luas bayangan melalui determinan matriks transformasi;
  \item menerapkan transformasi pada grafika komputer dan animasi.
\end{itemize}
\textbf{Capaian Pembelajaran (Fase F+).} Peserta didik dapat memodelkan transformasi geometri dengan matriks dan menggunakannya untuk menganalisis komposisi serta dampaknya pada bangun.
\end{tujuan}

\begin{pemantik}
\begin{itemize}[leftmargin=*,topsep=2pt]
  \item Saat sebuah karakter game berputar lalu bergeser, dalam urutan apa komputer menghitungnya?
  \item Mengapa rotasi dan refleksi tidak mengubah ukuran, tetapi dilatasi mengubahnya?
  \item Bisakah dua refleksi berturut-turut menghasilkan sebuah rotasi?
\end{itemize}
\end{pemantik}

\begin{studikasus}{Menggerakkan Objek di Layar}
Pada grafika komputer, posisi tiap titik objek disimpan sebagai vektor, dan setiap gerakan (geser, putar, perbesar, cermin) adalah sebuah \emph{transformasi} yang diwakili matriks. Menggabungkan beberapa gerakan = mengalikan matriks-matriksnya.
\begin{center}
\begin{tikzpicture}[scale=0.7]
  \draw[->] (-0.5,0)--(5,0) node[right]{\small $x$};
  \draw[->] (0,-0.5)--(0,4) node[above]{\small $y$};
  \draw[biru,thick,fill=birumuda] (1,0.5)--(2,0.5)--(2,1.5)--(1,1.5)--cycle;
  \node[biru] at (1.5,1){\scriptsize objek};
  \draw[oranye,thick,fill=oranyemuda] (2.5,1)--(4.5,1)--(4.5,3)--(2.5,3)--cycle;
  \node[oranye] at (3.5,2){\scriptsize bayangan};
  \draw[->,dashed,gray] (2.1,1)--(2.4,1.5);
\end{tikzpicture}\\
\small Sebuah persegi mengalami dilatasi (diperbesar) lalu translasi (digeser).
\end{center}
\textbf{Amati dan duga.} Jika objek diperbesar $2\times$ lalu digeser, samakah hasilnya dengan digeser dulu baru diperbesar? Inilah pentingnya \textbf{urutan komposisi} transformasi.
\end{studikasus}

\section*{Peta Konsep Bab}
\addcontentsline{toc}{section}{Peta Konsep Bab}
\begin{center}
\begin{tikzpicture}[
  node distance=8mm,
  konsepbox/.style={draw=biru,fill=birumuda,rounded corners,align=center,inner sep=5pt,font=\small\bfseries},
  subbox/.style={draw=hijau,fill=hijaumuda,rounded corners,align=center,inner sep=4pt,font=\footnotesize},
  garis/.style={-Stealth,biru,thick}]
  \node[konsepbox] (akar) {TRANSFORMASI GEOMETRI};
  \node[subbox,below left=12mm and 18mm of akar] (iso) {Isometri:\\translasi, refleksi, rotasi};
  \node[subbox,below right=12mm and 14mm of akar] (dil) {Dilatasi\\(ubah ukuran)};
  \node[subbox,below=11mm of iso] (mat) {Matriks\\transformasi};
  \node[subbox,below=11mm of dil] (komp) {Komposisi \&\\grafika komputer};
  \draw[garis] (akar)--(iso); \draw[garis] (akar)--(dil);
  \draw[garis] (iso)--(mat); \draw[garis] (dil)--(komp); \draw[garis] (mat)--(komp);
\end{tikzpicture}\\
\small \textbf{Peta konsep.} Translasi, refleksi, rotasi (isometri --- tak ubah ukuran) dan dilatasi semuanya dapat ditulis sebagai matriks; komposisinya menjadi perkalian matriks yang dipakai grafika komputer.
\end{center}

\section{Jenis-Jenis Transformasi}
\begin{konsep}{Matriks Transformasi (terhadap titik asal)}
\[
\begin{array}{ll}
\textbf{Refleksi sumbu-}x: \begin{pmatrix}1&0\\0&-1\end{pmatrix} &
\textbf{Refleksi sumbu-}y: \begin{pmatrix}-1&0\\0&1\end{pmatrix}\\[10pt]
\textbf{Refleksi garis }y=x: \begin{pmatrix}0&1\\1&0\end{pmatrix} &
\textbf{Dilatasi } [O,k]: \begin{pmatrix}k&0\\0&k\end{pmatrix}\\[10pt]
\multicolumn{2}{l}{\textbf{Rotasi sudut }\alpha\textbf{ terhadap }O: \begin{pmatrix}\cos\alpha&-\sin\alpha\\ \sin\alpha&\cos\alpha\end{pmatrix}}
\end{array}
\]
\textbf{Translasi} oleh $\begin{pmatrix}a\\b\end{pmatrix}$: $(x',y') = (x+a,\ y+b)$ (bukan perkalian matriks $2\times2$, melainkan penjumlahan).
\end{konsep}

\smallskip\noindent\textbf{Galeri empat transformasi dasar} (segitiga \textcolor{biru}{biru} $=$ objek, \textcolor{oranye}{oranye} $=$ bayangan).
\begin{center}
\begin{tabular}{cc}
% Translasi
\begin{tikzpicture}[scale=0.55,>=Stealth,line join=round]
\draw[gray,->](-3.4,0)--(3,0); \draw[gray,->](0,-1)--(0,3);
\coordinate (A) at (1,0.5);\coordinate (B) at (2.2,0.5);\coordinate (C) at (1,1.6);
\coordinate (A2) at (-2,1.1);\coordinate (B2) at (-0.8,1.1);\coordinate (C2) at (-2,2.2);
\fill[biru,opacity=0.25](A)--(B)--(C)--cycle;\draw[biru,thick](A)--(B)--(C)--cycle;
\fill[oranye,opacity=0.25](A2)--(B2)--(C2)--cycle;\draw[oranye,thick](A2)--(B2)--(C2)--cycle;
\draw[->,dashed,gray](A)--(A2);
\end{tikzpicture}
&
% Refleksi sumbu-y
\begin{tikzpicture}[scale=0.55,>=Stealth,line join=round]
\draw[gray,->](-3,0)--(3,0); \draw[hijau!60!black,very thick](0,-1)--(0,3);
\coordinate (A) at (1,0.5);\coordinate (B) at (2.2,0.5);\coordinate (C) at (1,1.6);
\fill[biru,opacity=0.25](A)--(B)--(C)--cycle;\draw[biru,thick](A)--(B)--(C)--cycle;
\fill[oranye,opacity=0.25](-1,0.5)--(-2.2,0.5)--(-1,1.6)--cycle;\draw[oranye,thick](-1,0.5)--(-2.2,0.5)--(-1,1.6)--cycle;
\node[hijau!60!black,font=\scriptsize,above] at (0,3){sumbu-$y$};
\end{tikzpicture}
\\[-1pt]
\footnotesize \textbf{Translasi} (geser; bentuk \&\ ukuran tetap) & \footnotesize \textbf{Refleksi} (cermin; menghadap terbalik)\\[3pt]
% Rotasi
\begin{tikzpicture}[scale=0.55,>=Stealth,line join=round]
\draw[gray,->](-3,0)--(3,0); \draw[gray,->](0,-1)--(0,3);
\coordinate (A) at (1,0.5);\coordinate (B) at (2.2,0.5);\coordinate (C) at (1,1.6);
\fill[biru,opacity=0.25](A)--(B)--(C)--cycle;\draw[biru,thick](A)--(B)--(C)--cycle;
\fill[oranye,opacity=0.25](-0.5,1)--(-0.5,2.2)--(-1.6,1)--cycle;\draw[oranye,thick](-0.5,1)--(-0.5,2.2)--(-1.6,1)--cycle;
\draw[->,dashed,gray](1,0.5) to[bend left=40] (-0.5,1);
\fill (0,0) circle (2.6pt) node[font=\scriptsize,below right]{$O$};
\end{tikzpicture}
&
% Dilatasi
\begin{tikzpicture}[scale=0.55,>=Stealth,line join=round]
\draw[gray,->](-1,0)--(4.2,0); \draw[gray,->](0,-1)--(0,3.2);
\coordinate (A) at (1,0.5);\coordinate (B) at (2.2,0.5);\coordinate (C) at (1,1.6);
\draw[dashed,gray](0,0)--(3.3,0.75); \draw[dashed,gray](0,0)--(1.5,2.4);
\fill[oranye,opacity=0.18](1.5,0.75)--(3.3,0.75)--(1.5,2.4)--cycle;\draw[oranye,thick](1.5,0.75)--(3.3,0.75)--(1.5,2.4)--cycle;
\fill[biru,opacity=0.3](A)--(B)--(C)--cycle;\draw[biru,thick](A)--(B)--(C)--cycle;
\fill (0,0) circle (2.6pt) node[font=\scriptsize,below left]{$O$};
\end{tikzpicture}
\\[-1pt]
\footnotesize \textbf{Rotasi} $90^\circ$ terhadap $O$ (ukuran tetap) & \footnotesize \textbf{Dilatasi} $[O,1{,}5]$ (ukuran berubah)\\
\end{tabular}
\end{center}

\begin{center}
\begin{tikzpicture}[scale=0.8]
  \draw[->] (-3,0)--(4,0) node[right]{$x$};
  \draw[->] (0,-1)--(0,4) node[above]{$y$};
  \fill[biru] (3,1) circle (2pt) node[right]{\small $A(3,1)$};
  \fill[oranye] (-1,3) circle (2pt) node[above left]{\small $A'(-1,3)$};
  \draw[dashed,gray,-{Latex}] (3,1) to[bend left=30] (-1,3);
\end{tikzpicture}\\
\small Rotasi $90^\circ$ berlawanan arah jarum jam: $A(3,1)\mapsto A'(-1,3)$.
\end{center}

\begin{contoh}{Rotasi Titik}
Tentukan bayangan titik $A(3,1)$ oleh rotasi $90^\circ$ berlawanan arah jarum jam terhadap $O$.

\jawab Untuk $\alpha=90^\circ$: $\begin{pmatrix}\cos90^\circ&-\sin90^\circ\\ \sin90^\circ&\cos90^\circ\end{pmatrix}=\begin{pmatrix}0&-1\\1&0\end{pmatrix}$. Maka
$\begin{pmatrix}0&-1\\1&0\end{pmatrix}\begin{pmatrix}3\\1\end{pmatrix}
= \begin{pmatrix}-1\\3\end{pmatrix}.$ Jadi $A'(-1,3)$.
\end{contoh}

\smallskip\noindent\textbf{Refleksi terhadap berbagai garis \&\ rotasi istimewa} (titik $P(3,1)$).
\begin{center}
\begin{tabular}{cc}
% Refleksi 4 garis
\begin{tikzpicture}[scale=0.5,>=Stealth,line join=round]
\draw[gray,->](-4,0)--(4,0) node[font=\scriptsize,right]{$x$}; \draw[gray,->](0,-4)--(0,4) node[font=\scriptsize,above]{$y$};
\draw[gray,dashed](-3.5,-3.5)--(3.5,3.5) node[font=\scriptsize,right]{$y{=}x$};
\draw[gray,dashed](-3.5,3.5)--(3.5,-3.5) node[font=\scriptsize,right]{$y{=}{-}x$};
\fill[biru](3,1) circle (3pt) node[font=\scriptsize,above right]{$P$};
\fill[oranye](3,-1) circle (2.6pt) node[font=\scriptsize,below right]{$x$};
\fill[oranye](-3,1) circle (2.6pt) node[font=\scriptsize,above left]{$y$};
\fill[oranye](1,3) circle (2.6pt) node[font=\scriptsize,above]{$y{=}x$};
\fill[oranye](-1,-3) circle (2.6pt) node[font=\scriptsize,below]{$y{=}{-}x$};
\end{tikzpicture}
&
% Rotasi 90/180/270
\begin{tikzpicture}[scale=0.5,>=Stealth,line join=round]
\draw[gray,->](-4,0)--(4,0) node[font=\scriptsize,right]{$x$}; \draw[gray,->](0,-4)--(0,4) node[font=\scriptsize,above]{$y$};
\fill (0,0) circle (2.6pt) node[font=\scriptsize,below left]{$O$};
\fill[biru](3,1) circle (3pt) node[font=\scriptsize,above right]{$P$};
\fill[oranye](-1,3) circle (2.6pt) node[font=\scriptsize,above]{$90^\circ$};
\fill[oranye](-3,-1) circle (2.6pt) node[font=\scriptsize,left]{$180^\circ$};
\fill[oranye](1,-3) circle (2.6pt) node[font=\scriptsize,below]{$270^\circ$};
\draw[->,dashed,gray](3,1) to[bend left=25] (-1,3);
\end{tikzpicture}
\\[-1pt]
\footnotesize \textbf{Refleksi}: $(3,1)\!\to\!(3,{-}1),({-}3,1),(1,3),({-}1,{-}3)$ & \footnotesize \textbf{Rotasi} terhadap $O$: $90^\circ,180^\circ,270^\circ$ (berlawanan jarum jam)\\
\end{tabular}
\end{center}

\section{Komposisi Transformasi}
\begin{konsep}{Transformasi Berurutan}
Jika titik dikenai $T_1$ (matriks $M_1$) lalu $T_2$ (matriks $M_2$), maka bayangan akhir
\[
\mathbf x'=M_2\big(M_1\mathbf x\big)=(M_2M_1)\mathbf x.
\]
\textbf{Matriks yang diterapkan lebih dahulu ditulis paling kanan.} Karena perkalian matriks tak komutatif, urutan menentukan hasil.
\end{konsep}

\begin{contoh}{Komposisi Dua Transformasi}
Titik $P(2,3)$ direfleksikan terhadap sumbu-$x$, lalu dirotasi $90^\circ$ (berlawanan jarum jam). Tentukan bayangannya.

\jawab Refleksi sumbu-$x$: $M_1=\begin{pmatrix}1&0\\0&-1\end{pmatrix}$; rotasi: $M_2=\begin{pmatrix}0&-1\\1&0\end{pmatrix}$. Gabungan $M_2M_1=\begin{pmatrix}0&1\\1&0\end{pmatrix}$. Maka $\begin{pmatrix}0&1\\1&0\end{pmatrix}\begin{pmatrix}2\\3\end{pmatrix}=\begin{pmatrix}3\\2\end{pmatrix}$, jadi $P'(3,2)$.
\end{contoh}

\smallskip\noindent\textbf{Dua refleksi $=$ satu rotasi.} Mencerminkan terhadap dua garis yang berpotongan di $O$ dengan sudut $\varphi$ setara dengan \emph{rotasi} sebesar $2\varphi$ terhadap $O$.
\begin{center}
\begin{tikzpicture}[scale=0.95,>=Stealth,line join=round]
\draw[gray,dashed](-2.7,0)--(2.7,0) node[font=\scriptsize,right]{$\ell_1$};
\draw[gray,dashed](-2.34,-1.35)--(2.34,1.35) node[font=\scriptsize,right]{$\ell_2$};
\fill (0,0) circle (1.5pt) node[font=\scriptsize,below]{$O$};
\draw[biru,very thick,->](0,0)--(1.607,1.915) node[font=\scriptsize,above]{objek};
\draw[gray,thick,dashed,->](0,0)--(1.607,-1.915) node[font=\scriptsize,below]{cermin $\ell_1$};
\draw[oranye,very thick,->](0,0)--(-0.855,2.349) node[font=\scriptsize,above left]{hasil akhir};
\draw[hijau!60!black](0.9,0) arc(0:30:0.9); \node[hijau!60!black,font=\scriptsize] at (1.18,0.27){$\varphi$};
\draw[merah](50:1.35) arc(50:110:1.35); \node[merah,font=\scriptsize] at (80:1.62){$2\varphi$};
\end{tikzpicture}\\
\small Refleksi terhadap $\ell_1$ lalu $\ell_2$ (bersudut $\varphi$) menghasilkan rotasi $2\varphi$ terhadap $O$.
\end{center}

\begin{konsep}{Luas Bayangan dan Determinan}
Sebuah transformasi bermatriks $M$ mengalikan setiap luas dengan faktor $|\det M|$. Karena $|\det|=1$ untuk rotasi dan refleksi (isometri), luas terjaga; dilatasi $[O,k]$ memiliki $\det=k^2$, sehingga luas menjadi $k^2$ kali.
\end{konsep}

\begin{center}
\begin{tikzpicture}[scale=0.55,line join=round]
\draw[gray,->](-0.5,0)--(2.3,0);\draw[gray,->](0,-0.5)--(0,2.3);
\fill[biru,opacity=0.25](0,0)--(1,0)--(1,1)--(0,1)--cycle;\draw[biru,thick](0,0)--(1,0)--(1,1)--(0,1)--cycle;
\node[font=\scriptsize] at (0.5,0.5){luas $1$};
\end{tikzpicture}
\quad $\xrightarrow{\ M=\left(\begin{smallmatrix}2&1\\0&3\end{smallmatrix}\right)\ }$ \quad
\begin{tikzpicture}[scale=0.55,line join=round]
\draw[gray,->](-0.5,0)--(4,0);\draw[gray,->](0,-0.5)--(0,3.6);
\fill[oranye,opacity=0.22](0,0)--(2,0)--(3,3)--(1,3)--cycle;\draw[oranye,thick](0,0)--(2,0)--(3,3)--(1,3)--cycle;
\node[font=\scriptsize] at (1.6,1.4){luas $6$};
\end{tikzpicture}\\
\small Persegi satuan (luas $1$) menjadi jajar genjang berluas $|\det M|=|2\cdot3-1\cdot0|=6$. Faktor pengali luas $=|\det M|$.
\end{center}

\begin{kesalahan}
Translasi \textbf{bukan} transformasi linear $2\times2$ (tak ada matriks $2\times2$ yang menggesernya), karena ia tidak memetakan titik asal ke titik asal. Pada grafika komputer, translasi ditangani dengan koordinat homogen ($3\times3$). Selain itu, pada komposisi, jangan terbalik: matriks pertama diterapkan ada di sebelah \emph{kanan}.
\end{kesalahan}

\begin{dunianyata}
Mesin permainan dan perangkat lunak animasi merepresentasikan setiap objek dengan vektor titik dan menggerakkannya melalui rantai perkalian matriks transformasi setiap \emph{frame}. Itulah sebabnya kartu grafis (GPU) dioptimalkan untuk perkalian matriks besar-besaran.
\end{dunianyata}

\begin{konsep}{Bahasa Alternatif yang Sering Mengecoh di Soal}
\begin{itemize}[leftmargin=*,topsep=2pt]
  \item \textbf{Arah rotasi}: sudut \emph{positif} berarti \textbf{berlawanan} arah jarum jam (CCW). "Searah jarum jam $\theta$" $=$ rotasi $-\theta$.
  \item \textbf{Dilatasi} $[P,k]$: \emph{pusat} $P$ penting --- $[P,k]\neq[O,k]$ kecuali $P=O$. Jika $|k|>1$ membesar, $|k|<1$ mengecil, dan $k<0$ sekaligus memutar $180^\circ$.
  \item \textbf{"Refleksi terhadap garis $x=h$"} (garis tegak) berbeda dari \textbf{"refleksi terhadap sumbu-$y$"} (yaitu $x=0$); begitu pula $y=k$ vs sumbu-$x$.
  \item \textbf{Komposisi}: "$T_1$ \emph{dilanjutkan} $T_2$" $\Rightarrow$ matriks $M_2M_1$ (yang dikerjakan lebih dulu ditulis di \emph{kanan}). Urutan berpengaruh.
  \item \textbf{Isometri} $=$ transformasi yang mempertahankan jarak: translasi, refleksi, rotasi. Dilatasi \emph{bukan} isometri (kecuali $k=\pm1$).
  \item \textbf{Bayangan} (\emph{image}) vs \textbf{objek/prapeta} (\emph{preimage}): bila soal memberi bayangan dan meminta objek, gunakan transformasi \emph{invers}.
  \item \textbf{Titik tetap}: translasi tak punya titik tetap; pada rotasi/dilatasi, titik tetapnya adalah \emph{pusat}.
  \item \textbf{"Setengah putaran"} $=180^\circ$, \textbf{"seperempat putaran"} $=90^\circ$.
\end{itemize}
\end{konsep}

\section*{Ringkasan Bab}
\addcontentsline{toc}{section}{Ringkasan Bab}
\begin{itemize}[leftmargin=*,topsep=2pt]
  \item Translasi menggeser ($+\mathbf t$); refleksi, rotasi, dilatasi dapat ditulis sebagai matriks.
  \item Rotasi sudut $\alpha$: $\begin{psmallmatrix}\cos\alpha&-\sin\alpha\\ \sin\alpha&\cos\alpha\end{psmallmatrix}$.
  \item \textbf{Komposisi:} matriks gabungan $M_2M_1$ (yang pertama di kanan); urutan penting.
  \item \textbf{Luas} bayangan $=|\det M|\times$ luas asli; isometri menjaga luas.
  \item Transformasi adalah tulang punggung grafika komputer.
\end{itemize}

\begin{refleksi}
\begin{itemize}[leftmargin=*,topsep=2pt]
  \item[\cek] Saya dapat menentukan bayangan titik oleh keempat transformasi dasar.
  \item[\cek] Saya dapat menulis transformasi sebagai matriks dan menghitung bayangan.
  \item[\cek] Saya dapat menyusun komposisi transformasi dengan urutan yang benar.
  \item[\cek] Saya dapat menghubungkan determinan dengan perubahan luas.
\end{itemize}
\end{refleksi}

\section*{Soal Akhir Bab \thechapter}
\addcontentsline{toc}{section}{Soal Akhir Bab \thechapter}

\bagiansoal{A}{(Fundamental --- setara SNBT Dasar)}
\begin{enumerate}[leftmargin=*]
  \item Tentukan bayangan $A(2,5)$ oleh translasi $\begin{pmatrix}3\\-1\end{pmatrix}$.
  \item Tentukan bayangan $B(4,-2)$ oleh refleksi terhadap sumbu-$x$.
  \item Tentukan bayangan $C(1,3)$ oleh rotasi $90^\circ$ berlawanan jarum jam terhadap $O$.
  \item Tentukan bayangan $D(2,-3)$ oleh dilatasi $[O,2]$.
  \item Tentukan bayangan $E(3,4)$ oleh refleksi terhadap garis $y=x$.
\end{enumerate}

\bagiansoal{B}{(Intermediate --- setara SNBT Sulit)}
\begin{enumerate}[leftmargin=*]
  \item Tentukan bayangan $P(2,-3)$ oleh refleksi terhadap sumbu-$x$ dilanjutkan dilatasi $[O,2]$.
  \item Tentukan bayangan $Q(1,2)$ oleh rotasi $180^\circ$ terhadap $O$.
  \item Tentukan matriks gabungan untuk "refleksi sumbu-$y$ dilanjutkan rotasi $90^\circ$", lalu terapkan pada $(3,1)$.
  \item Titik $(5,2)$ ditranslasi $\begin{pmatrix}-2\\4\end{pmatrix}$ lalu direfleksikan terhadap sumbu-$x$. Tentukan bayangannya.
  \item Sebuah segitiga berluas $6$ didilatasi $[O,3]$. Tentukan luas bayangannya.
  \item Tentukan bayangan $(2,0)$ oleh rotasi $60^\circ$ terhadap $O$.
  \item Garis $y=2x+1$ direfleksikan terhadap sumbu-$x$. Tentukan persamaan bayangannya.
\end{enumerate}

\bagiansoal{C}{(Advanced --- setara Olimpiade SMA, gabungan antarbab)}
\begin{enumerate}[leftmargin=*]
  \item \textbf{(Transformasi + Matriks).} Tunjukkan bahwa rotasi $90^\circ$ dilanjutkan rotasi $90^\circ$ sama dengan rotasi $180^\circ$, melalui perkalian matriks.
  \item \textbf{(Transformasi + Trigonometri).} Buktikan bahwa komposisi rotasi sudut $\alpha$ dan rotasi sudut $\beta$ (sepusat) adalah rotasi sudut $\alpha+\beta$, dan turunkan rumus jumlah sudut $\cos(\alpha+\beta)$, $\sin(\alpha+\beta)$.
  \item \textbf{(Transformasi + Lingkaran).} Lingkaran $x^2+y^2=4$ didilatasi $[O,3]$. Tentukan persamaan bayangannya dan bandingkan luasnya.
  \item \textbf{(Komposisi).} Dua refleksi terhadap sumbu-$x$ lalu sumbu-$y$ menghasilkan transformasi apa? Buktikan dengan matriks.
  \item \textbf{(Transformasi + Determinan).} Sebuah transformasi bermatriks $M=\begin{pmatrix}2&1\\0&3\end{pmatrix}$ dikenakan pada persegi satuan. Tentukan luas bayangannya.
  \item \textbf{(Transformasi + Fungsi).} Grafik $y=x^2$ didilatasi $[O,2]$. Tentukan persamaan kurva bayangannya.
\end{enumerate}

\bagiansoal{D}{(Expert --- setara tahun pertama universitas, sintesis \& pembuktian)}
\begin{enumerate}[leftmargin=*]
  \item \textbf{(Pembuktian).} Buktikan bahwa matriks rotasi $R(\alpha)$ memenuhi $R(\alpha)R(\beta)=R(\alpha+\beta)$ dan $R(\alpha)^{-1}=R(-\alpha)$.
  \item \textbf{(Refleksi sebagai matriks).} Turunkan matriks refleksi terhadap garis $y=(\tan\theta)x$ melalui $O$, dan tunjukkan determinannya $-1$.
  \item \textbf{(Isometri).} Buktikan bahwa rotasi dan refleksi mempertahankan jarak antartitik (isometri), sedangkan dilatasi $[O,k]$ mengalikan jarak dengan $|k|$.
  \item \textbf{(Komposisi refleksi).} Buktikan bahwa komposisi dua refleksi terhadap dua garis yang berpotongan di $O$ dengan sudut $\varphi$ adalah rotasi sebesar $2\varphi$.
  \item \textbf{(Transformasi + Koordinat homogen).} Jelaskan bagaimana translasi dapat dinyatakan sebagai perkalian matriks $3\times3$ memakai koordinat homogen, dan tuliskan matriksnya.
\end{enumerate}

\begin{challenge}
\begin{enumerate}[leftmargin=*,topsep=2pt]
  \item Sebuah titik dirotasi $40^\circ$ sebanyak $9$ kali berturut-turut terhadap $O$. Tentukan transformasi tunggal yang setara, dan jelaskan.
  \item Tentukan matriks transformasi tunggal yang setara dengan "dilatasi $[O,2]$ dilanjutkan rotasi $90^\circ$", lalu tentukan apakah hasil ini sama dengan urutan kebalikannya.
\end{enumerate}
\end{challenge}

% ============================================================
\chapter{Analisis Grafik Fungsi}
% ============================================================
\begin{tujuan}
Setelah mempelajari bab ini, peserta didik mampu:
\begin{itemize}[leftmargin=*,topsep=2pt]
  \item mengenali bentuk dan ciri grafik fungsi linear, kuadrat, rasional, akar, eksponensial, logaritma, dan piecewise;
  \item menentukan domain, range, dan asimtot suatu fungsi;
  \item menganalisis transformasi grafik (geseran, regangan, pencerminan);
  \item menentukan kesimetrian (fungsi genap/ganjil);
  \item menafsirkan grafik dalam konteks nyata.
\end{itemize}
\textbf{Capaian Pembelajaran (Fase F+).} Peserta didik dapat menganalisis berbagai jenis fungsi melalui grafik, ciri-ciri, dan transformasinya.
\end{tujuan}

\begin{pemantik}
\begin{itemize}[leftmargin=*,topsep=2pt]
  \item Hanya dengan melihat bentuk grafik, dapatkah kamu menebak jenis fungsinya?
  \item Mengapa grafik $y=\tfrac1x$ "menghindari" garis tertentu tetapi tak pernah menyentuhnya?
  \item Bagaimana satu perubahan kecil pada rumus ($+3$, $-x$, $2\times$) menggeser atau mencerminkan seluruh grafik?
\end{itemize}
\end{pemantik}

\begin{studikasus}{Membaca Grafik di Berita}
Grafik muncul di mana-mana: harga saham, kurva kasus penyakit, laju unduh. Memahami \emph{bentuk} fungsi membantu menafsirkannya --- linear berarti laju tetap, eksponensial berarti melipat ganda, dan asimtot menandai batas yang tak terlampaui.
\begin{center}
\begin{tikzpicture}
\begin{axis}[axis lines=middle,width=5.6cm,height=3.8cm,title={\scriptsize Linear},
  xmin=0,xmax=4,ymin=0,ymax=5,xtick=\empty,ytick=\empty,samples=2]
\addplot[very thick,biru,domain=0:4]{1+x};
\end{axis}
\end{tikzpicture}\hspace{2mm}
\begin{tikzpicture}
\begin{axis}[axis lines=middle,width=5.6cm,height=3.8cm,title={\scriptsize Kuadrat},
  xmin=-2.2,xmax=2.2,ymin=-0.5,ymax=4.5,xtick=\empty,ytick=\empty,samples=50]
\addplot[very thick,hijau,domain=-2:2]{x^2};
\end{axis}
\end{tikzpicture}\hspace{2mm}
\begin{tikzpicture}
\begin{axis}[axis lines=middle,width=5.6cm,height=3.8cm,title={\scriptsize Eksponensial},
  xmin=-2.2,xmax=2.2,ymin=-0.3,ymax=4.5,xtick=\empty,ytick=\empty,samples=50]
\addplot[very thick,oranye,domain=-2:2]{2^x};
\end{axis}
\end{tikzpicture}
\end{center}
\textbf{Amati dan duga.} Mana yang "naik dengan laju tetap", mana yang "melengkung", dan mana yang "meledak"? Membedakan bentuk-bentuk ini adalah inti analisis grafik.
\end{studikasus}

\section*{Peta Konsep Bab}
\addcontentsline{toc}{section}{Peta Konsep Bab}
\begin{center}
\begin{tikzpicture}[
  node distance=8mm,
  konsepbox/.style={draw=biru,fill=birumuda,rounded corners,align=center,inner sep=5pt,font=\small\bfseries},
  subbox/.style={draw=hijau,fill=hijaumuda,rounded corners,align=center,inner sep=4pt,font=\footnotesize},
  garis/.style={-Stealth,biru,thick}]
  \node[konsepbox] (akar) {ANALISIS GRAFIK FUNGSI};
  \node[subbox,below left=12mm and 16mm of akar] (gal) {Galeri \&\\ciri fungsi};
  \node[subbox,below=12mm of akar] (dom) {Domain, range,\\asimtot};
  \node[subbox,below right=12mm and 16mm of akar] (tr) {Transformasi\\grafik};
  \node[subbox,below=10mm of dom] (sim) {Simetri:\\genap/ganjil};
  \draw[garis] (akar)--(gal); \draw[garis] (akar)--(dom); \draw[garis] (akar)--(tr);
  \draw[garis] (dom)--(sim);
\end{tikzpicture}\\
\small \textbf{Peta konsep.} Kenali jenis fungsi, lalu telaah domain/range/asimtot, transformasi, dan kesimetriannya.
\end{center}

\section{Galeri dan Ciri Fungsi}
\begin{center}
\begin{tikzpicture}
\begin{axis}[axis lines=middle,width=6.5cm,height=4.8cm,
  title={\small Fungsi Rasional $f(x)=\frac{1}{x-1}$},
  xmin=-3,xmax=5,ymin=-4,ymax=4,xtick=\empty,ytick=\empty]
\addplot[very thick,biru,domain=-3:0.85,samples=60]{1/(x-1)};
\addplot[very thick,biru,domain=1.15:5,samples=60]{1/(x-1)};
\draw[dashed,oranye] (axis cs:1,-4)--(axis cs:1,4);
\end{axis}
\end{tikzpicture}
\hspace{0.5cm}
\begin{tikzpicture}
\begin{axis}[axis lines=middle,width=6.5cm,height=4.8cm,
  title={\small Fungsi Akar $f(x)=\sqrt{x}$},
  xmin=-1,xmax=6,ymin=-1,ymax=3,xtick=\empty,ytick=\empty]
\addplot[very thick,hijau,domain=0:6,samples=60]{sqrt(x)};
\end{axis}
\end{tikzpicture}
\end{center}

\begin{konsep}{Domain, Range, dan Asimtot}
\textbf{Domain} = himpunan masukan yang sah; \textbf{range} = himpunan keluaran. Hindari pembagian nol dan akar bilangan negatif. \textbf{Asimtot} adalah garis yang didekati grafik tanpa disentuh: rasional $\tfrac{1}{x-1}$ punya asimtot tegak $x=1$ dan datar $y=0$.
\end{konsep}

\begin{konsep}{Fungsi Piecewise dan Tangga}
\textbf{Fungsi piecewise} didefinisikan dengan aturan berbeda pada interval berbeda, contoh:
\[
f(x) = \begin{cases} x+2, & x < 0\\ x^2, & x \ge 0 \end{cases}
\]
\textbf{Fungsi tangga} (mis. fungsi lantai $\lfloor x \rfloor$) bernilai konstan pada tiap interval lalu "melompat".
\end{konsep}

\begin{center}
\begin{tikzpicture}
\begin{axis}[axis lines=middle,width=7.5cm,height=5cm,
  xlabel=$x$,ylabel=$y$,xmin=-3,xmax=3,ymin=-2,ymax=4,samples=2]
\addplot[very thick,biru,domain=-3:0]{x+2};
\addplot[very thick,oranye,domain=0:2,samples=40]{x^2};
\fill[biru](0,2)circle(2pt); \draw[oranye,fill=white](0,0)circle(2pt);
\end{axis}
\end{tikzpicture}\\
\small Grafik fungsi piecewise: garis $x+2$ untuk $x<0$ dan parabola $x^2$ untuk $x\ge0$.
\end{center}

\section{Transformasi dan Simetri Grafik}
\begin{konsep}{Aturan Transformasi Grafik}
Dari grafik $y=f(x)$:
\[
\begin{array}{ll}
f(x)+k:\ \text{geser } k \text{ ke atas} & f(x-h):\ \text{geser } h \text{ ke kanan}\\
a\,f(x):\ \text{regang vertikal } a\times & -f(x):\ \text{cermin terhadap sumbu-}x\\
f(-x):\ \text{cermin terhadap sumbu-}y & f(ax):\ \text{rapat horizontal } \tfrac1a\times
\end{array}
\]
\end{konsep}

\begin{center}
\begin{tikzpicture}
\begin{axis}[axis lines=middle,width=11.5cm,height=6.2cm,
  xlabel=$x$,ylabel=$y$,xmin=-3.2,xmax=5.5,ymin=-3.5,ymax=6.5,
  xtick=\empty,ytick=\empty,samples=80,clip=false,
  legend style={font=\tiny,at={(1.01,1)},anchor=north west,draw=abu}]
\addplot[very thick,biru,domain=-2.5:2.5]{x^2};
\addlegendentry{$f(x)=x^2$ (acuan)}
\addplot[very thick,hijau,domain=-2.5:2.5]{x^2+2};
\addlegendentry{$f(x)+2$: naik 2 satuan}
\addplot[very thick,oranye,domain=-0.3:4.5]{(x-2)^2};
\addlegendentry{$f(x-2)$: geser kanan 2}
\addplot[very thick,merah,dashed,domain=-2.5:2.5]{-x^2};
\addlegendentry{$-f(x)$: cermin sb.-$x$}
% Annotation: vertical shift arrow
\draw[-Stealth,hijau,thin] (axis cs:2,4)--(axis cs:2,6) node[right,font=\tiny]{$+2$};
% Annotation: horizontal shift arrow
\draw[-Stealth,oranye,thin] (axis cs:0,0.2)--(axis cs:2,0.2) node[above,font=\tiny]{$+2$};
\end{axis}
\end{tikzpicture}\\
\small \textbf{Empat transformasi dari $f(x)=x^2$.} Perubahan di luar $f(\cdot)$ menggeser/mencerminkan secara vertikal; perubahan di argumen menggeser secara horizontal (ingat: $f(x-2)$ geser \emph{kanan}, bukan kiri!).
\end{center}

\begin{konsep}{Fungsi Genap dan Ganjil}
\textbf{Genap:} $f(-x)=f(x)$ (simetris terhadap sumbu-$y$, mis. $x^2$). \textbf{Ganjil:} $f(-x)=-f(x)$ (simetris terhadap titik asal, mis. $x^3$).
\end{konsep}

\begin{contoh}{Menganalisis Transformasi}
Jelaskan bagaimana grafik $y=-2(x-1)^2+3$ diperoleh dari $y=x^2$.

\jawab Dari $x^2$: geser $1$ ke kanan ($\to(x-1)^2$), regang vertikal $2\times$ lalu cermin terhadap sumbu-$x$ ($\to-2(x-1)^2$), lalu geser $3$ ke atas. Puncaknya $(1,3)$, terbuka ke bawah.
\end{contoh}

\begin{kesalahan}
Pada $f(x-h)$, geseran ke \emph{kanan} sebesar $h$ (bukan ke kiri) meski tandanya minus --- ini sering tertukar. Demikian pula $f(x+3)$ menggeser ke \emph{kiri} $3$.
\end{kesalahan}

\section*{Ringkasan Bab}
\addcontentsline{toc}{section}{Ringkasan Bab}
\begin{itemize}[leftmargin=*,topsep=2pt]
  \item Kenali ciri grafik: linear (garis), kuadrat (parabola), rasional (asimtot), akar, eksponen, logaritma, piecewise.
  \item \textbf{Domain} hindari nol penyebut dan akar negatif; \textbf{asimtot} = garis batas.
  \item \textbf{Transformasi:} $f(x)+k$, $f(x-h)$, $a f(x)$, $f(ax)$, $-f(x)$, $f(-x)$.
  \item \textbf{Simetri:} genap $f(-x)=f(x)$; ganjil $f(-x)=-f(x)$.
\end{itemize}

\begin{refleksi}
\begin{itemize}[leftmargin=*,topsep=2pt]
  \item[\cek] Saya dapat mengenali jenis fungsi dari grafiknya.
  \item[\cek] Saya dapat menentukan domain, range, dan asimtot.
  \item[\cek] Saya dapat menjelaskan transformasi sebuah grafik.
  \item[\cek] Saya dapat menguji kesimetrian fungsi.
\end{itemize}
\end{refleksi}

\section*{Soal Akhir Bab \thechapter}
\addcontentsline{toc}{section}{Soal Akhir Bab \thechapter}

\bagiansoal{A}{(Fundamental --- setara SNBT Dasar)}
\begin{enumerate}[leftmargin=*]
  \item Tentukan domain $f(x)=\sqrt{x-3}$.
  \item Tentukan domain $g(x)=\dfrac{1}{x+2}$.
  \item Tentukan asimtot tegak $f(x)=\dfrac{1}{x-1}$.
  \item Grafik $y=x^2$ digeser $3$ satuan ke atas. Tulis persamaan barunya.
  \item Tentukan range $f(x)=x^2+1$.
\end{enumerate}

\bagiansoal{B}{(Intermediate --- setara SNBT Sulit)}
\begin{enumerate}[leftmargin=*]
  \item Tentukan domain dan range $f(x)=\sqrt{4-x^2}$.
  \item Grafik $y=f(x)$ digeser $2$ ke kanan dan $1$ ke bawah. Tulis bentuk transformasinya.
  \item Tentukan asimtot datar $f(x)=\dfrac{2x+1}{x-3}$.
  \item Diberikan $f(x)=x^2$. Uraikan transformasi untuk memperoleh $y=-2f(x-1)$.
  \item Tentukan domain $f(x)=\dfrac{\sqrt{x-1}}{x-4}$.
  \item Fungsi $f(x)=\begin{cases}x+1,&x<0\\2x,&x\ge0\end{cases}$. Hitung $f(-2)$, $f(0)$, $f(3)$.
\end{enumerate}

\bagiansoal{C}{(Advanced --- setara Olimpiade SMA, gabungan antarbab)}
\begin{enumerate}[leftmargin=*]
  \item \textbf{(Grafik + Kuadrat).} Tentukan titik puncak dan sumbu simetri $y=2(x-3)^2+4$.
  \item \textbf{(Grafik + Eksponen).} Grafik $y=2^x$ menjadi $y=2^{x-1}+3$. Jelaskan transformasinya dan tentukan asimtot datarnya.
  \item \textbf{(Grafik + Rasional).} Tentukan domain serta asimtot $f(x)=\dfrac{x+2}{x^2-x-6}$.
  \item \textbf{(Grafik + Logaritma).} Tentukan domain $f(x)={}^{10}\!\log(x^2-4)$.
  \item \textbf{(Grafik + Simetri).} Tentukan apakah $f(x)=x^3-x$ genap, ganjil, atau bukan keduanya.
  \item \textbf{(Grafik + Komposisi).} Jika $f(x)=\sqrt{x}$ dan $g(x)=x^2-9$, tentukan domain $(f\circ g)(x)$.
\end{enumerate}

\bagiansoal{D}{(Expert --- setara tahun pertama universitas, sintesis \& pembuktian)}
\begin{enumerate}[leftmargin=*]
  \item \textbf{(Pembuktian).} Buktikan bahwa jumlah dua fungsi genap adalah genap dan jumlah dua fungsi ganjil adalah ganjil.
  \item \textbf{(Asimtot umum).} Untuk $f(x)=\dfrac{ax+b}{cx+d}$ dengan $c\neq0$, tentukan asimtot tegak dan datarnya.
  \item \textbf{(Bentuk puncak).} Dengan melengkapkan kuadrat, tunjukkan bahwa setiap $ax^2+bx+c$ dapat ditulis $a(x-h)^2+k$, dan nyatakan $h,k$ dalam $a,b,c$.
  \item \textbf{(Dekomposisi).} Buktikan bahwa setiap fungsi $f$ dapat ditulis sebagai jumlah fungsi genap dan fungsi ganjil: $f(x)=\tfrac12[f(x)+f(-x)]+\tfrac12[f(x)-f(-x)]$.
\end{enumerate}

\begin{challenge}
\begin{enumerate}[leftmargin=*,topsep=2pt]
  \item Sebuah fungsi kuadrat memiliki puncak di $(2,5)$ dan melalui $(0,1)$. Tentukan rumusnya.
  \item Sketsa $y=|x^2-4|$ dan jelaskan bagaimana nilai mutlak mencerminkan bagian grafik yang berada di bawah sumbu-$x$.
\end{enumerate}
\end{challenge}

% ============================================================
\chapter{Fungsi Rasional}
% ============================================================
\begin{tujuan}
Setelah mempelajari bab ini, peserta didik mampu:
\begin{itemize}[leftmargin=*,topsep=2pt]
  \item menentukan domain fungsi rasional dengan mengenali pembuat nol penyebut;
  \item menentukan asimtot tegak, datar, dan miring suatu fungsi rasional;
  \item menggambar sketsa grafik fungsi rasional menggunakan asimtot dan titik potong;
  \item menyelesaikan persamaan rasional dan memeriksa syarat keberlakuannya;
  \item menyelesaikan pertidaksamaan rasional dengan metode garis bilangan (tanda).
\end{itemize}
\textbf{Capaian Pembelajaran (Fase F).} Peserta didik dapat menganalisis fungsi rasional melalui domain, asimtot, dan grafiknya, serta menyelesaikan persamaan dan pertidaksamaan rasional.
\end{tujuan}

\begin{pemantik}
\begin{itemize}[leftmargin=*,topsep=2pt]
  \item Mengapa grafik $y=\dfrac{1}{x}$ "patah" menjadi dua bagian yang tak pernah bertemu?
  \item Apa beda antara garis yang \emph{tak boleh disentuh} grafik (asimtot) dan titik yang \emph{hilang} dari grafik (lubang)?
  \item Ketika menyelesaikan $\dfrac{1}{x-2}>3$, bolehkah kita langsung mengalikan kedua ruas dengan $x-2$? Apa risikonya?
\end{itemize}
\end{pemantik}

\begin{studikasus}{Konsentrasi Obat dalam Darah}
Setelah obat dissuntikkan, konsentrasinya (mg/L) sering dimodelkan $C(t)=\dfrac{20t}{t^2+4}$, dengan $t$ waktu (jam).
\begin{center}
\begin{tikzpicture}
\begin{axis}[axis lines=middle,width=8.2cm,height=4.6cm,
  xlabel={\scriptsize $t$ (jam)},ylabel={\scriptsize $C$},
  xmin=0,xmax=14,ymin=0,ymax=6,xtick=\empty,ytick=\empty,samples=80]
\addplot[very thick,biru,domain=0:14]{20*x/(x^2+4)};
\draw[dashed,oranye] (axis cs:0,0)--(axis cs:14,0);
\end{axis}
\end{tikzpicture}
\end{center}
\textbf{Amati dan duga.} (a) Mengapa konsentrasi mula-mula naik lalu turun? (b) Ke arah mana nilai $C$ menuju saat $t$ sangat besar? Garis $C=0$ yang didekati grafik tanpa pernah disentuh adalah \emph{asimtot datar}. Fungsi berbentuk hasil bagi dua polinomial seperti ini disebut \textbf{fungsi rasional}, dan ciri-ciri grafiknya ditentukan oleh asimtot dan titik potongnya.
\end{studikasus}

\section*{Peta Konsep Bab}
\addcontentsline{toc}{section}{Peta Konsep Bab}
\begin{center}
\begin{tikzpicture}[
  node distance=8mm,
  konsepbox/.style={draw=biru,fill=birumuda,rounded corners,align=center,inner sep=5pt,font=\small\bfseries},
  subbox/.style={draw=hijau,fill=hijaumuda,rounded corners,align=center,inner sep=4pt,font=\footnotesize},
  garis/.style={-Stealth,biru,thick}]
  \node[konsepbox] (akar) {FUNGSI RASIONAL};
  \node[subbox,below left=12mm and 16mm of akar] (dom) {Domain \&\\asimtot};
  \node[subbox,below=12mm of akar] (graf) {Grafik};
  \node[subbox,below right=12mm and 16mm of akar] (pp) {Persamaan \&\\pertidaksamaan};
  \node[subbox,below=10mm of dom] (as) {Tegak, datar,\\miring};
  \draw[garis] (akar)--(dom); \draw[garis] (akar)--(graf); \draw[garis] (akar)--(pp);
  \draw[garis] (dom)--(as);
\end{tikzpicture}\\
\small \textbf{Peta konsep.} Tentukan domain dan asimtot, gunakan untuk menggambar grafik, lalu selesaikan persamaan dan pertidaksamaannya.
\end{center}

\section{Domain dan Asimtot}
\begin{definisi}
\textbf{Fungsi rasional} adalah fungsi berbentuk $f(x)=\dfrac{P(x)}{Q(x)}$ dengan $P(x)$ dan $Q(x)$ polinomial dan $Q(x)\neq 0$.
\end{definisi}

\begin{konsep}{Domain}
Domain fungsi rasional adalah semua bilangan real \emph{kecuali} pembuat nol penyebut. Untuk $f(x)=\dfrac{P(x)}{Q(x)}$, syaratnya $Q(x)\neq 0$.
\end{konsep}

\begin{konsep}{Tiga Jenis Asimtot}
Untuk $f(x)=\dfrac{P(x)}{Q(x)}$ dalam bentuk paling sederhana (faktor sekutu sudah dicoret):
\begin{itemize}[leftmargin=*,topsep=2pt]
  \item \textbf{Asimtot tegak:} garis $x=a$ untuk tiap pembuat nol penyebut $a$ (yang tidak menjadi nol pembilang).
  \item \textbf{Asimtot datar:} bandingkan derajat $P$ dan $Q$.
  \begin{itemize}
    \item $\deg P<\deg Q \Rightarrow y=0$;
    \item $\deg P=\deg Q \Rightarrow y=\dfrac{\text{koefisien utama }P}{\text{koefisien utama }Q}$;
    \item $\deg P>\deg Q \Rightarrow$ tidak ada asimtot datar.
  \end{itemize}
  \item \textbf{Asimtot miring:} muncul tepat ketika $\deg P=\deg Q+1$. Dapatkan dengan pembagian polinomial; hasil bagi (linear) adalah persamaan asimtot miringnya.
\end{itemize}
\end{konsep}

\begin{center}
\begin{tikzpicture}
\begin{axis}[axis lines=middle,width=8.5cm,height=6cm,
  xlabel=$x$,ylabel=$y$,xmin=-4,xmax=4,ymin=-4.5,ymax=4.5,
  xtick=\empty,ytick=\empty,samples=80,clip=false]
% y = 1/x (two branches)
\addplot[very thick,biru,domain=-4:-0.2]{1/x};
\addplot[very thick,biru,domain=0.2:4]{1/x};
% Asymptote: x = 0 (vertical)
\draw[dashed,oranye,thick] (axis cs:0,-4.5)--(axis cs:0,4.5);
% Asymptote: y = 0 (horizontal)
\draw[dashed,hijau,thick] (axis cs:-4,0)--(axis cs:4,0);
% Labels for asymptotes
\node[oranye,font=\scriptsize,anchor=south west] at (axis cs:0.1,4.2)
  {asimtot tegak $x=0$};
\node[hijau,font=\scriptsize,anchor=north east] at (axis cs:3.8,-0.15)
  {asimtot datar $y=0$};
% "mendekati tapi tidak menyentuh" annotation
\node[biru,font=\footnotesize,align=center,fill=white,inner sep=2pt,draw=biru,rounded corners=2pt] at (axis cs:2.5,2.8)
  {mendekati\\tapi tak\\menyentuh};
\draw[-Stealth,biru,thin] (axis cs:2.5,2.15)--(axis cs:1.1,0.9);
\draw[-Stealth,biru,thin] (axis cs:2.5,2.15)--(axis cs:0.25,3.5);
\end{axis}
\end{tikzpicture}\\
\small Grafik $y=\tfrac{1}{x}$: dua asimtot (garis putus-putus) yang tak pernah disentuh grafik. Makin dekat $x$ ke $0$, makin besar $|y|$; makin besar $|x|$, makin mendekati $y=0$.
\end{center}

\begin{contoh}{Domain dan Asimtot Datar/Tegak}
Tentukan domain dan semua asimtot $f(x)=\dfrac{2x+1}{x-3}$.

\jawab Penyebut nol di $x=3$, jadi domain $\{x\mid x\neq 3\}$ dan asimtot tegak $x=3$. Karena $\deg P=\deg Q=1$, asimtot datar $y=\dfrac{2}{1}=2$.
\end{contoh}

\begin{contoh}{Asimtot Miring}
Tentukan asimtot miring $f(x)=\dfrac{x^2+1}{x-1}$.

\jawab Karena $\deg P=2=\deg Q+1$, lakukan pembagian: $x^2+1=(x-1)(x+1)+2$, sehingga $f(x)=x+1+\dfrac{2}{x-1}$. Saat $x\to\pm\infty$ suku $\dfrac{2}{x-1}\to 0$, jadi asimtot miringnya $y=x+1$.
\end{contoh}

\begin{kesalahan}
Jika pembilang dan penyebut punya faktor sekutu, pembuat nol itu menghasilkan \emph{lubang} (titik kosong), bukan asimtot tegak. Contoh $f(x)=\dfrac{x^2-4}{x-2}=x+2$ untuk $x\neq2$: grafiknya garis $y=x+2$ dengan lubang di $(2,4)$, tanpa asimtot tegak.
\end{kesalahan}

\section{Grafik Fungsi Rasional}
\begin{konsep}{Langkah Menyketsa Grafik}
\begin{enumerate}[leftmargin=*,topsep=2pt]
  \item Tentukan domain dan sederhanakan (catat lubang bila ada).
  \item Tentukan asimtot tegak, datar/miring.
  \item Cari titik potong sumbu: sumbu-$x$ ($P(x)=0$) dan sumbu-$y$ ($x=0$).
  \item Uji tanda di tiap selang yang dibatasi asimtot tegak dan titik nol, lalu sketsa tiap cabang.
\end{enumerate}
\end{konsep}

\begin{center}
\begin{tikzpicture}
\begin{axis}[axis lines=middle,width=8.5cm,height=5.4cm,
  xlabel=$x$,ylabel=$y$,xmin=-4,xmax=8,ymin=-4,ymax=8,
  xtick={3},ytick={2},samples=80]
\addplot[very thick,biru,domain=-4:2.7]{(2*x+1)/(x-3)};
\addplot[very thick,biru,domain=3.5:8]{(2*x+1)/(x-3)};
\draw[dashed,oranye] (axis cs:3,-4)--(axis cs:3,8);
\draw[dashed,oranye] (axis cs:-4,2)--(axis cs:8,2);
\end{axis}
\end{tikzpicture}\\
\small Grafik $f(x)=\dfrac{2x+1}{x-3}$ dengan asimtot tegak $x=3$ dan asimtot datar $y=2$.
\end{center}

\begin{dunianyata}
Fungsi rasional memodelkan banyak hal: konsentrasi obat dalam darah, biaya rata-rata per unit produksi $\overline{C}(x)=\dfrac{C(x)}{x}$, hambatan paralel pada rangkaian listrik, dan laju reaksi enzim (persamaan Michaelis--Menten). Asimtot datar sering bermakna "batas jenuh" yang tak terlampaui.
\end{dunianyata}

\section{Persamaan dan Pertidaksamaan Rasional}
\begin{konsep}{Persamaan Rasional}
Untuk menyelesaikan persamaan rasional: kalikan dengan KPK penyebut (atau samakan penyebut), selesaikan persamaan polinomial yang dihasilkan, lalu \textbf{buang akar yang membuat penyebut nol} (akar palsu).
\end{konsep}

\begin{contoh}{Persamaan Rasional}
Selesaikan $\dfrac{1}{x-1}+\dfrac{1}{x+1}=1$.

\jawab Samakan penyebut: $\dfrac{(x+1)+(x-1)}{(x-1)(x+1)}=\dfrac{2x}{x^2-1}=1$. Maka $2x=x^2-1\Rightarrow x^2-2x-1=0\Rightarrow x=1\pm\sqrt2$. Keduanya $\neq\pm1$, jadi HP $=\{1-\sqrt2,\,1+\sqrt2\}$.
\end{contoh}

\begin{konsep}{Pertidaksamaan Rasional}
\emph{Jangan} mengalikan silang dengan penyebut yang tandanya belum diketahui. Pindahkan semua ke satu ruas hingga berbentuk $\dfrac{P(x)}{Q(x)}\;\substack{<\\>}\;0$, tentukan pembuat nol pembilang dan penyebut, lalu gunakan \textbf{garis bilangan} untuk menguji tanda. Titik dari penyebut selalu \emph{terbuka}.
\end{konsep}

\begin{contoh}{Pertidaksamaan Rasional}
Selesaikan $\dfrac{x+3}{x-2}<1$.

\jawab Pindah ruas: $\dfrac{x+3}{x-2}-1<0\Rightarrow\dfrac{(x+3)-(x-2)}{x-2}=\dfrac{5}{x-2}<0$. Karena pembilang $5>0$, pecahan negatif hanya jika $x-2<0$, yaitu $x<2$. HP $=\{x\mid x<2\}$.
\end{contoh}

\begin{kesalahan}
Pada $\dfrac{x+3}{x-2}<1$, mengalikan langsung dengan $x-2$ keliru karena tandanya bisa positif atau negatif. Jika $x-2<0$, arah pertidaksamaan harus dibalik --- itulah mengapa metode garis bilangan lebih aman.
\end{kesalahan}

\section*{Ringkasan Bab}
\addcontentsline{toc}{section}{Ringkasan Bab}
\begin{itemize}[leftmargin=*,topsep=2pt]
  \item \textbf{Domain:} semua real kecuali pembuat nol penyebut.
  \item \textbf{Asimtot tegak:} pembuat nol penyebut (setelah disederhanakan); faktor sekutu menimbulkan \emph{lubang}.
  \item \textbf{Asimtot datar:} bandingkan derajat ($\deg P<\deg Q\Rightarrow y=0$; sama $\Rightarrow$ rasio koefisien utama; lebih besar $\Rightarrow$ tak ada).
  \item \textbf{Asimtot miring:} bila $\deg P=\deg Q+1$, gunakan pembagian polinomial.
  \item \textbf{Persamaan:} kalikan penyebut, cek akar palsu.
  \item \textbf{Pertidaksamaan:} satukan ruas, uji tanda dengan garis bilangan (titik penyebut terbuka).
\end{itemize}

\begin{refleksi}
\begin{itemize}[leftmargin=*,topsep=2pt]
  \item[\cek] Saya dapat menentukan domain fungsi rasional.
  \item[\cek] Saya dapat membedakan asimtot tegak, datar, dan miring.
  \item[\cek] Saya dapat menyketsa grafik fungsi rasional.
  \item[\cek] Saya dapat menyelesaikan persamaan rasional dan membuang akar palsu.
  \item[\cek] Saya dapat menyelesaikan pertidaksamaan rasional dengan garis bilangan.
\end{itemize}
\end{refleksi}

\section*{Soal Akhir Bab \thechapter}
\addcontentsline{toc}{section}{Soal Akhir Bab \thechapter}

\bagiansoal{A}{(Fundamental --- setara SNBT Dasar)}
\begin{enumerate}[leftmargin=*]
  \item Tentukan domain $f(x)=\dfrac{1}{x-3}$.
  \item Tentukan domain $f(x)=\dfrac{x+1}{x^2-4}$.
  \item Tentukan asimtot tegak $f(x)=\dfrac{5}{x+2}$.
  \item Tentukan asimtot datar $f(x)=\dfrac{3x+1}{x-4}$.
  \item Tentukan asimtot datar $f(x)=\dfrac{2x+1}{x^2-1}$.
  \item Selesaikan $\dfrac{2}{x}=4$.
  \item Selesaikan $\dfrac{x+1}{x-2}=0$.
  \item Tentukan semua asimtot tegak $f(x)=\dfrac{x+1}{x^2-x}$.
  \item Hitung $f(2)$ untuk $f(x)=\dfrac{x^2-1}{x+1}$.
  \item Tentukan titik potong sumbu-$x$ dari $f(x)=\dfrac{2x-6}{x+1}$.
\end{enumerate}

\bagiansoal{B}{(Intermediate --- setara SNBT Sulit)}
\begin{enumerate}[leftmargin=*]
  \item Tentukan asimtot tegak dan datar $f(x)=\dfrac{2x^2+3}{x^2-1}$.
  \item Tentukan asimtot miring $f(x)=\dfrac{x^2+1}{x-1}$.
  \item Selesaikan $\dfrac{1}{x-1}+\dfrac{1}{x+1}=1$.
  \item Selesaikan pertidaksamaan $\dfrac{x-2}{x+1}\ge 0$.
  \item Selesaikan pertidaksamaan $\dfrac{x+3}{x-2}<1$.
\end{enumerate}

\bagiansoal{C}{(Advanced --- setara Olimpiade SMA, gabungan antarbab)}
\begin{enumerate}[leftmargin=*]
  \item \textbf{(Rasional + Grafik).} Sederhanakan dan jelaskan grafik $f(x)=\dfrac{x^2-4}{x-2}$, termasuk ada/tidaknya lubang.
  \item \textbf{(Rasional + Limit).} Hitung $\displaystyle\lim_{x\to\infty}\dfrac{3x^2-x+1}{2x^2+5}$ dan kaitkan dengan asimtot datar.
  \item \textbf{(Rasional + Komposisi).} Diberikan $f(x)=\dfrac1x$. Tentukan $(f\circ f)(x)$ beserta domainnya.
  \item \textbf{(Rasional + Pertidaksamaan).} Selesaikan $\dfrac{x^2-1}{x^2-4}\le 0$.
  \item \textbf{(Rasional + Ketaksamaan).} Untuk $x>0$, tentukan nilai minimum $f(x)=x+\dfrac1x$ dan nilai $x$ yang mencapainya.
\end{enumerate}

\bagiansoal{D}{(Expert --- setara tahun pertama universitas, sintesis \& pembuktian)}
\begin{enumerate}[leftmargin=*]
  \item \textbf{(Pecahan Parsial).} Tuliskan $\dfrac{3x+5}{(x+1)(x+2)}$ dalam bentuk $\dfrac{A}{x+1}+\dfrac{B}{x+2}$.
  \item \textbf{(Asimtot Miring Umum).} Untuk $f(x)=\dfrac{ax^2+bx+c}{x+d}$ ($a\neq0$), buktikan asimtot miringnya $y=ax+(b-ad)$.
  \item \textbf{(Rasional + Turunan).} Dengan definisi turunan, buktikan bahwa turunan $f(x)=\dfrac1x$ adalah $f'(x)=-\dfrac{1}{x^2}$.
  \item \textbf{(Pembuktian Ketaksamaan).} Buktikan untuk setiap $x>0$ berlaku $x+\dfrac1x\ge 2$, dan tentukan kapan kesamaan tercapai.
  \item \textbf{(Rasional + Invers).} Tentukan $f^{-1}(x)$ untuk $f(x)=\dfrac{x+1}{x-1}$ dan tunjukkan bahwa $f$ adalah inversnya sendiri.
\end{enumerate}

\begin{challenge}
\begin{enumerate}[leftmargin=*,topsep=2pt]
  \item Selesaikan $\dfrac{x^2-3x+2}{x^2-5x+6}>0$.
  \item Tentukan nilai maksimum dan minimum $f(x)=\dfrac{x}{x^2+1}$ beserta range-nya.
  \item Buktikan bahwa $f(x)=\dfrac{x^2+x+1}{x^2-x+1}$ hanya bernilai pada selang $\left[\tfrac13,\,3\right]$.
\end{enumerate}
\end{challenge}

% ============================================================
\chapter{Nilai Mutlak dan Bentuk Irasional}
% ============================================================
\begin{tujuan}
Setelah mempelajari bab ini, peserta didik mampu:
\begin{itemize}[leftmargin=*,topsep=2pt]
  \item memaknai nilai mutlak sebagai jarak dan menyatakan definisinya secara piecewise;
  \item menggambar grafik fungsi nilai mutlak serta transformasinya;
  \item menyelesaikan persamaan, pertidaksamaan, dan sistem yang memuat nilai mutlak;
  \item melakukan operasi dan menyederhanakan bentuk akar serta merasionalkan penyebut;
  \item menyelesaikan persamaan dan pertidaksamaan irasional dengan memeriksa syaratnya.
\end{itemize}
\textbf{Capaian Pembelajaran (Fase F).} Peserta didik dapat menyelesaikan masalah yang melibatkan nilai mutlak dan bentuk akar, termasuk persamaan dan pertidaksamaannya.
\end{tujuan}

\begin{pemantik}
\begin{itemize}[leftmargin=*,topsep=2pt]
  \item Mengapa $\sqrt{x^2}=|x|$ dan bukan sekadar $x$?
  \item Sebuah mesin mengisi botol "500 mL" dengan toleransi $\pm5$ mL. Bagaimana menuliskan syarat ini dalam satu pertidaksamaan nilai mutlak?
  \item Saat menyelesaikan $\sqrt{2x+3}=x$, mengapa kita wajib memeriksa kembali jawabannya ke persamaan asli?
\end{itemize}
\end{pemantik}

\begin{studikasus}{Kendali Mutu dan Toleransi}
Sebuah pabrik memproduksi poros dengan diameter ideal $20$ mm. Poros diterima bila selisih diameternya dari nilai ideal tidak lebih dari $0{,}3$ mm. Jika $d$ diameter terukur, syarat "diterima" ditulis ringkas sebagai
\[
|d-20|\le 0{,}3 \quad\Longleftrightarrow\quad 19{,}7\le d\le 20{,}3 .
\]
\textbf{Amati dan duga.} Mengapa satu pertidaksamaan nilai mutlak setara dengan sebuah selang? Apa makna geometris dari $|d-20|$? Inilah inti nilai mutlak: ia mengukur \emph{jarak}, dan jarak selalu tak negatif.
\end{studikasus}

\section*{Peta Konsep Bab}
\addcontentsline{toc}{section}{Peta Konsep Bab}
\begin{center}
\begin{tikzpicture}[
  node distance=8mm,
  konsepbox/.style={draw=biru,fill=birumuda,rounded corners,align=center,inner sep=5pt,font=\small\bfseries},
  subbox/.style={draw=hijau,fill=hijaumuda,rounded corners,align=center,inner sep=4pt,font=\footnotesize},
  garis/.style={-Stealth,biru,thick}]
  \node[konsepbox] (akar) {NILAI MUTLAK \&\ IRASIONAL};
  \node[subbox,below left=12mm and 18mm of akar] (nm) {Nilai mutlak:\\definisi \&\ grafik};
  \node[subbox,below=12mm of akar] (pp) {Persamaan,\\pertidaksamaan,\\sistem};
  \node[subbox,below right=12mm and 18mm of akar] (ir) {Bentuk akar \&\\rasionalisasi};
  \node[subbox,below=10mm of ir] (psi) {Persamaan \&\\pertidaksamaan\\irasional};
  \draw[garis] (akar)--(nm); \draw[garis] (akar)--(pp); \draw[garis] (akar)--(ir);
  \draw[garis] (ir)--(psi);
\end{tikzpicture}\\
\small \textbf{Peta konsep.} Nilai mutlak (jarak) dan bentuk akar adalah dua tema yang sama-sama menuntut pemeriksaan syarat saat diselesaikan.
\end{center}

\section{Nilai Mutlak}
\begin{definisi}
\textbf{Nilai mutlak} dari bilangan real $x$ adalah jaraknya dari $0$ pada garis bilangan:
\[
|x|=\begin{cases} x, & x\ge 0\\ -x, & x<0 \end{cases}
\]
Akibatnya $|x|\ge0$ untuk setiap $x$, dan $\sqrt{x^2}=|x|$.
\end{definisi}

\begin{sifat}
Untuk $a\ge0$: $|x|=a\iff x=\pm a$;\quad $|x|<a\iff -a<x<a$;\quad $|x|>a\iff x<-a$ atau $x>a$. Selain itu $|xy|=|x|\,|y|$ dan $|a+b|\le|a|+|b|$ (ketaksamaan segitiga).
\end{sifat}

\begin{center}
\begin{tikzpicture}
\begin{axis}[axis lines=middle,width=8cm,height=4.8cm,
  xlabel=$x$,ylabel=$y$,xmin=-2,xmax=6,ymin=-0.5,ymax=4.5,
  xtick={2},ytick={1},samples=2]
\addplot[very thick,biru,domain=-2:2]{1+(2-x)};
\addplot[very thick,biru,domain=2:6]{1+(x-2)};
\fill[oranye](2,1)circle(2pt);
\end{axis}
\end{tikzpicture}\\
\small Grafik $y=|x-2|+1$: bentuk "V" dengan titik balik (verteks) $(2,1)$, diperoleh dari $y=|x|$ digeser $2$ ke kanan dan $1$ ke atas.
\end{center}

\begin{contoh}{Persamaan Nilai Mutlak}
Selesaikan $|x-3|=2$.

\jawab $x-3=2$ atau $x-3=-2$, sehingga $x=5$ atau $x=1$. HP $=\{1,5\}$.
\end{contoh}

\begin{contoh}{Pertidaksamaan Nilai Mutlak}
Selesaikan $|2x+1|\ge 5$.

\jawab $2x+1\ge5$ atau $2x+1\le-5$. Dari yang pertama $x\ge2$; dari yang kedua $x\le-3$. HP $=\{x\mid x\le-3\ \text{atau}\ x\ge2\}$.
\end{contoh}

\begin{konsep}{Sistem dan Bentuk $|f|=|g|$}
Persamaan $|f(x)|=|g(x)|$ setara dengan $f(x)=g(x)$ \emph{atau} $f(x)=-g(x)$ (atau kuadratkan kedua ruas). Untuk sistem, selesaikan tiap persamaan nilai mutlak lalu padukan dengan substitusi.
\end{konsep}

\begin{kesalahan}
$|x|<a$ saat $a\le0$ tidak punya penyelesaian (jarak tak mungkin negatif), sedangkan $|x|>a$ saat $a<0$ dipenuhi \emph{semua} $x$. Jangan otomatis menulis $-a<x<a$ tanpa memeriksa tanda $a$.
\end{kesalahan}

\section{Bentuk Akar (Irasional)}
\begin{konsep}{Operasi Bentuk Akar}
Untuk $a,b\ge0$: $\sqrt{a}\,\sqrt{b}=\sqrt{ab}$ dan $\dfrac{\sqrt a}{\sqrt b}=\sqrt{\dfrac ab}$ ($b>0$). Akar sejenis dapat dijumlahkan: $p\sqrt c+q\sqrt c=(p+q)\sqrt c$. Sederhanakan dengan mengeluarkan faktor kuadrat sempurna, mis. $\sqrt{50}=\sqrt{25\cdot2}=5\sqrt2$.
\end{konsep}

\begin{konsep}{Rasionalisasi Penyebut}
Hilangkan akar di penyebut dengan mengalikan bentuk sekawan:
\[
\frac{1}{\sqrt a}=\frac{\sqrt a}{a},\qquad
\frac{c}{\sqrt a-\sqrt b}=\frac{c(\sqrt a+\sqrt b)}{a-b}.
\]
Sekawan dari $\sqrt a-\sqrt b$ adalah $\sqrt a+\sqrt b$, memanfaatkan $(\sqrt a-\sqrt b)(\sqrt a+\sqrt b)=a-b$.
\end{konsep}

\begin{contoh}{Rasionalisasi}
Rasionalkan $\dfrac{3}{\sqrt5-\sqrt2}$.

\jawab Kalikan dengan sekawan: $\dfrac{3(\sqrt5+\sqrt2)}{(\sqrt5-\sqrt2)(\sqrt5+\sqrt2)}=\dfrac{3(\sqrt5+\sqrt2)}{5-2}=\sqrt5+\sqrt2.$
\end{contoh}

\begin{konsep}{Persamaan dan Pertidaksamaan Irasional}
Untuk persamaan $\sqrt{f(x)}=g(x)$: kuadratkan, namun \textbf{wajib} $g(x)\ge0$ dan $f(x)\ge0$; periksa kembali akar (buang akar palsu). Untuk pertidaksamaan, jaga syarat $f(x)\ge0$ dan perhatikan tanda saat mengkuadratkan. Misal $\sqrt{f(x)}<a$ ($a>0$) $\iff 0\le f(x)<a^2$.
\end{konsep}

\begin{contoh}{Persamaan Irasional}
Selesaikan $\sqrt{2x+3}=x$.

\jawab Syarat $x\ge0$ dan $2x+3\ge0$. Kuadratkan: $2x+3=x^2\Rightarrow x^2-2x-3=0\Rightarrow(x-3)(x+1)=0\Rightarrow x=3$ atau $x=-1$. Karena perlu $x\ge0$, hanya $x=3$ (cek: $\sqrt9=3$ \checkmark). HP $=\{3\}$.
\end{contoh}

\begin{kesalahan}
Mengkuadratkan dapat memunculkan \emph{akar palsu}. Pada $\sqrt{2x+3}=x$, nilai $x=-1$ memenuhi hasil kuadrat tetapi tidak persamaan asli (ruas kanan negatif). Selalu substitusikan kembali.
\end{kesalahan}

\section*{Ringkasan Bab}
\addcontentsline{toc}{section}{Ringkasan Bab}
\begin{itemize}[leftmargin=*,topsep=2pt]
  \item $|x|$ = jarak dari $0$; $\sqrt{x^2}=|x|$; $|x|\ge0$ selalu.
  \item $|x|=a\iff x=\pm a$; $|x|<a\iff -a<x<a$; $|x|>a\iff x<-a$ atau $x>a$ (untuk $a\ge0$).
  \item Grafik $y=|x-h|+k$: "V" berverteks $(h,k)$.
  \item Bentuk akar: keluarkan kuadrat sempurna; rasionalkan dengan sekawan.
  \item Persamaan/pertidaksamaan irasional: jaga syarat $\ge0$, kuadratkan, cek akar palsu.
\end{itemize}

\begin{refleksi}
\begin{itemize}[leftmargin=*,topsep=2pt]
  \item[\cek] Saya memahami nilai mutlak sebagai jarak dan bentuk piecewise-nya.
  \item[\cek] Saya dapat menyketsa grafik fungsi nilai mutlak.
  \item[\cek] Saya dapat menyelesaikan persamaan, pertidaksamaan, dan sistem nilai mutlak.
  \item[\cek] Saya dapat menyederhanakan dan merasionalkan bentuk akar.
  \item[\cek] Saya dapat menyelesaikan persamaan/pertidaksamaan irasional dan memeriksa akar palsu.
\end{itemize}
\end{refleksi}

\section*{Soal Akhir Bab \thechapter}
\addcontentsline{toc}{section}{Soal Akhir Bab \thechapter}

\bagiansoal{A}{(Fundamental --- setara SNBT Dasar)}
\begin{enumerate}[leftmargin=*]
  \item Hitung $|-7|$.
  \item Selesaikan $|x|=5$.
  \item Selesaikan $|x-3|=2$.
  \item Selesaikan $|2x|=8$.
  \item Sederhanakan $\sqrt{50}$.
  \item Sederhanakan $\sqrt{12}+\sqrt{27}$.
  \item Rasionalkan $\dfrac{1}{\sqrt3}$.
  \item Rasionalkan $\dfrac{2}{\sqrt5}$.
  \item Selesaikan $\sqrt{x}=4$.
  \item Hitung $|3-8|+|-2|$.
\end{enumerate}

\bagiansoal{B}{(Intermediate --- setara SNBT Sulit)}
\begin{enumerate}[leftmargin=*]
  \item Selesaikan $|x-2|=|2x+1|$.
  \item Selesaikan $|x-1|<3$.
  \item Selesaikan $|2x+1|\ge5$.
  \item Rasionalkan $\dfrac{3}{\sqrt5-\sqrt2}$.
  \item Selesaikan $\sqrt{2x+3}=x$.
\end{enumerate}

\bagiansoal{C}{(Advanced --- setara Olimpiade SMA, gabungan antarbab)}
\begin{enumerate}[leftmargin=*]
  \item \textbf{(Nilai mutlak + Grafik).} Jelaskan cara memperoleh grafik $y=|x-2|+1$ dari $y=|x|$ dan tentukan verteksnya.
  \item \textbf{(Nilai mutlak + Kuadrat).} Selesaikan $|x^2-4|=3$.
  \item \textbf{(Irasional + Pertidaksamaan).} Selesaikan $\sqrt{x-1}<2$.
  \item \textbf{(Nilai mutlak + Pertidaksamaan).} Selesaikan $|x|<|x-2|$.
  \item \textbf{(Irasional + Bilangan).} Sederhanakan $\sqrt{7+2\sqrt{10}}$.
\end{enumerate}

\bagiansoal{D}{(Expert --- setara tahun pertama universitas, sintesis \& pembuktian)}
\begin{enumerate}[leftmargin=*]
  \item \textbf{(Ketaksamaan Segitiga).} Buktikan $|a+b|\le|a|+|b|$ untuk setiap bilangan real $a,b$.
  \item \textbf{(Irasional + Limit).} Hitung $\displaystyle\lim_{x\to0}\dfrac{\sqrt{x+1}-1}{x}$ dengan rasionalisasi.
  \item \textbf{(Nilai mutlak + Turunan).} Tunjukkan bahwa $f(x)=|x|$ tidak mempunyai turunan di $x=0$.
  \item \textbf{(Pembuktian).} Buktikan bahwa $\sqrt2$ adalah bilangan irasional.
  \item \textbf{(Sistem Nilai Mutlak).} Selesaikan sistem $|x+y|=6$ dan $|x-y|=2$.
\end{enumerate}

\begin{challenge}
\begin{enumerate}[leftmargin=*,topsep=2pt]
  \item Selesaikan $|x-1|+|x-3|=4$.
  \item Selesaikan $\sqrt{x+4}-\sqrt{x-1}=1$.
  \item Tentukan nilai minimum $f(x)=|x-1|+|x-2|+|x-3|$ dan nilai $x$ yang mencapainya.
\end{enumerate}
\end{challenge}

% ============================================================
\chapter{Pemodelan Matematika}
% ============================================================
\begin{tujuan}
Setelah mempelajari bab ini, peserta didik mampu:
\begin{itemize}[leftmargin=*,topsep=2pt]
  \item mengenali situasi nyata yang cocok dimodelkan fungsi linear, kuadrat, eksponensial, atau logistik;
  \item menyusun model matematika dan menafsirkan parameternya;
  \item menggunakan model untuk prediksi dan optimasi sederhana;
  \item mengevaluasi kelayakan serta keterbatasan suatu model.
\end{itemize}
\textbf{Capaian Pembelajaran (Fase F+).} Peserta didik dapat memodelkan fenomena dunia nyata dengan fungsi yang sesuai dan mengevaluasinya secara kritis.
\end{tujuan}

\begin{pemantik}
\begin{itemize}[leftmargin=*,topsep=2pt]
  \item Tabungan bertambah tetap tiap bulan, sedangkan virus menyebar dengan melipat ganda --- model apa yang membedakannya?
  \item Mengapa pertumbuhan populasi tak bisa eksponensial selamanya?
  \item Bagaimana memilih harga jual yang memaksimalkan keuntungan?
\end{itemize}
\end{pemantik}

\begin{studikasus}{Menyebarnya Sebuah Tren}
Sebuah video viral ditonton $200$ kali pada hari ke-0, lalu jumlah penonton melipat sekitar $1{,}5\times$ tiap hari pada awalnya, sebelum melambat saat mendekati jenuh.
\begin{center}
\begin{tabular}{c|ccccc}
\toprule
Hari & 0 & 1 & 2 & 3 & 4 \\
Penonton (ribu) & 0{,}2 & 0{,}3 & 0{,}45 & 0{,}68 & 1{,}0 \\
\bottomrule
\end{tabular}
\end{center}
\textbf{Amati dan duga.} (a) Pola awal lebih cocok linear atau eksponensial? (b) Apa yang membatasi pertumbuhan ini pada akhirnya? (c) Fungsi apa yang menangkap "cepat di awal, jenuh di akhir"? Inilah perjalanan dari model \textbf{eksponensial} menuju model \textbf{logistik}.
\end{studikasus}

\section*{Peta Konsep Bab}
\addcontentsline{toc}{section}{Peta Konsep Bab}
\begin{center}
\begin{tikzpicture}[
  node distance=8mm,
  konsepbox/.style={draw=biru,fill=birumuda,rounded corners,align=center,inner sep=5pt,font=\small\bfseries},
  subbox/.style={draw=hijau,fill=hijaumuda,rounded corners,align=center,inner sep=4pt,font=\footnotesize},
  garis/.style={-Stealth,biru,thick}]
  \node[konsepbox] (akar) {PEMODELAN MATEMATIKA};
  \node[subbox,below left=12mm and 16mm of akar] (lin) {Linear \&\\kuadrat};
  \node[subbox,below=12mm of akar] (eks) {Eksponensial\\\&\ logaritma};
  \node[subbox,below right=12mm and 16mm of akar] (log) {Logistik\\(jenuh)};
  \node[subbox,below=10mm of eks] (eval) {Prediksi,\\optimasi, evaluasi};
  \draw[garis] (akar)--(lin); \draw[garis] (akar)--(eks); \draw[garis] (akar)--(log);
  \draw[garis] (eks)--(eval);
\end{tikzpicture}\\
\small \textbf{Peta konsep.} Pilih model (linear, kuadrat, eksponensial, logistik), taksir parameter, lalu pakai untuk prediksi/optimasi sambil mengevaluasi keterbatasannya.
\end{center}

\begin{center}
\begin{tikzpicture}[
  kotak/.style={draw=biru,fill=birumuda,rounded corners,align=center,font=\footnotesize,inner sep=5pt,minimum width=2.6cm,minimum height=0.72cm},
  panah/.style={-Stealth,biru,thick}]
\node[kotak] (A) at (0,0) {\textbf{Masalah}\\Nyata};
\node[kotak] (B) at (3.1,0) {Penyeder-\\hanaan};
\node[kotak] (C) at (6.2,0) {Model\\Matematika};
\node[kotak] (D) at (6.2,-1.9) {Perhitungan\\(Solusi)};
\node[kotak] (E) at (3.1,-1.9) {Interpretasi\\Hasil};
\node[kotak] (F) at (0,-1.9) {Validasi \&\\Revisi};
\draw[panah] (A)--(B);
\draw[panah] (B)--(C);
\draw[panah] (C)--(D);
\draw[panah] (D)--(E);
\draw[panah] (E)--(F);
\draw[panah] (F) -- ++(-0.65,0) -- ++(0,1.9) -- (A);
\end{tikzpicture}\\
\small \textbf{Siklus pemodelan matematika.} Dari masalah nyata ke model, hitung solusinya, tafsirkan dalam konteks asli, lalu validasi --- ulang jika model belum cukup akurat.
\end{center}

\section{Memilih dan Menyusun Model}
\begin{konsep}{Empat Model Dasar}
\begin{itemize}[leftmargin=*,topsep=2pt]
  \item \textbf{Linear} $y=a+bx$: laju perubahan tetap (selisih tetap).
  \item \textbf{Kuadrat} $y=ax^2+bx+c$: lintasan proyektil, optimasi luas/keuntungan.
  \item \textbf{Eksponensial} $y=A\,r^{t}$: pertumbuhan/peluruhan dengan rasio tetap.
  \item \textbf{Logistik} $y=\dfrac{K}{1+Ae^{-rt}}$: tumbuh cepat lalu jenuh di kapasitas $K$.
\end{itemize}
\end{konsep}

\begin{contoh}{Model Eksponensial}
Populasi bakteri menggandakan diri tiap jam, mulai dari $100$. Susun model dan tentukan kapan mencapai $1600$.

\jawab $P(t)=100\cdot2^{t}$. Saat $P=1600$: $2^t=16=2^4$, jadi $t=4$ jam.
\end{contoh}

\begin{contoh}{Model Kuadrat (Optimasi)}
Pendapatan $R(x)=x(100-x)$ dengan $x$ harga. Tentukan harga yang memaksimalkan pendapatan.

\jawab $R(x)=-x^2+100x$ berpuncak di $x=-\tfrac{100}{2(-1)}=50$. Maka harga $50$ memberi pendapatan maksimum $R(50)=2500$.
\end{contoh}

\section{Pertumbuhan, Peluruhan, dan Keterbatasan}
\begin{konsep}{Pertumbuhan dan Peluruhan}
\textbf{Pertumbuhan $p\%$ per periode:} $y=A(1+p)^t$. \textbf{Peluruhan:} $y=A(1-p)^t$. \textbf{Waktu paruh} $T$: $y=A\left(\tfrac12\right)^{t/T}$.
\end{konsep}

\begin{dunianyata}
Model eksponensial menjelaskan bunga majemuk, peluruhan radioaktif, dan fase awal wabah. Namun di dunia nyata, sumber daya terbatas; model \textbf{logistik} memperbaikinya dengan memperhitungkan \emph{kapasitas tampung} $K$, sehingga pertumbuhan melambat dan mendatar.
\end{dunianyata}

\begin{kesalahan}
Ekstrapolasi model eksponensial terlalu jauh menghasilkan ramalan tak masuk akal (populasi "tak hingga"). Selalu periksa apakah asumsi model masih berlaku pada rentang yang ditanyakan.
\end{kesalahan}

\section*{Ringkasan Bab}
\addcontentsline{toc}{section}{Ringkasan Bab}
\begin{itemize}[leftmargin=*,topsep=2pt]
  \item \textbf{Linear:} laju tetap; \textbf{kuadrat:} lintasan/optimasi; \textbf{eksponensial:} rasio tetap; \textbf{logistik:} jenuh.
  \item Pertumbuhan $A(1+p)^t$, peluruhan $A(1-p)^t$, paruh $A(\tfrac12)^{t/T}$.
  \item Model dipakai untuk prediksi dan optimasi, lalu \emph{dievaluasi} keterbatasannya.
\end{itemize}

\begin{refleksi}
\begin{itemize}[leftmargin=*,topsep=2pt]
  \item[\cek] Saya dapat memilih model yang cocok untuk suatu situasi.
  \item[\cek] Saya dapat menyusun model dan menafsirkan parameternya.
  \item[\cek] Saya dapat memakai model untuk prediksi dan optimasi.
  \item[\cek] Saya dapat menjelaskan keterbatasan sebuah model.
\end{itemize}
\end{refleksi}

\section*{Soal Akhir Bab \thechapter}
\addcontentsline{toc}{section}{Soal Akhir Bab \thechapter}

\bagiansoal{A}{(Fundamental --- setara SNBT Dasar)}
\begin{enumerate}[leftmargin=*]
  \item Model biaya $C(x)=5000+2000x$ (rupiah, $x$ unit). Hitung biaya untuk $10$ unit.
  \item Populasi bakteri awal $100$ menggandakan diri tiap jam. Tulis model $P(t)$.
  \item Harga mobil menyusut $10\%$ per tahun dari Rp$200$ juta. Tulis model dan harga setelah $1$ tahun.
  \item Tinggi proyektil $h(t)=20t-5t^2$. Tentukan tinggi saat $t=2$.
  \item Suhu kopi $T(t)=25+60(0{,}95)^t$. Berapa suhu awal ($t=0$)?
\end{enumerate}

\bagiansoal{B}{(Intermediate --- setara SNBT Sulit)}
\begin{enumerate}[leftmargin=*]
  \item Proyektil $h(t)=20t-5t^2$. Tentukan waktu mencapai tanah dan tinggi maksimumnya.
  \item Tabungan bunga majemuk $5\%$/tahun, awal Rp$1$ juta. Tulis model dan nilainya setelah $3$ tahun.
  \item Model $P(t)=100\cdot2^t$. Kapan populasi mencapai $1600$?
  \item Zat radioaktif berwaktu paruh $5$ tahun. Tulis model $N(t)$ dan fraksi tersisa setelah $15$ tahun.
  \item Sebuah model linear melalui $(0,3)$ dan $(4,11)$. Tentukan rumusnya.
  \item Pendapatan $R(x)=x(100-x)$. Tentukan harga $x$ yang memaksimalkan pendapatan.
\end{enumerate}

\bagiansoal{C}{(Advanced --- setara Olimpiade SMA, gabungan antarbab)}
\begin{enumerate}[leftmargin=*]
  \item \textbf{(Pemodelan + Eksponen/Log).} Populasi kota tumbuh $3\%$/tahun, kini $500.000$. Tentukan kapan mencapai $1.000.000$ (nyatakan dengan logaritma).
  \item \textbf{(Pemodelan + Kuadrat).} Tali sepanjang $40$ m dibuat persegi panjang. Tentukan ukuran yang memaksimalkan luas.
  \item \textbf{(Pemodelan + Logaritma).} Magnitudo gempa $M={}^{10}\!\log(I/I_0)$. Jika $I=1000\,I_0$, tentukan $M$.
  \item \textbf{(Pemodelan + Logistik).} Model $P(t)=\dfrac{1000}{1+9e^{-t}}$. Tentukan populasi awal dan batas saat $t\to\infty$.
  \item \textbf{(Pemodelan + Statistika).} Data pertumbuhan tampak lurus pada skala $\log$. Model apa yang cocok dan bagaimana menaksir parameternya?
  \item \textbf{(Pemodelan + Optimasi).} Keuntungan $P(x)=-2x^2+80x-200$. Tentukan $x$ yang memaksimalkan keuntungan dan keuntungan maksimumnya.
\end{enumerate}

\bagiansoal{D}{(Expert --- setara tahun pertama universitas, sintesis \& pembuktian)}
\begin{enumerate}[leftmargin=*]
  \item \textbf{(Pemodelan + Kalkulus).} Untuk $P(t)=P_0e^{kt}$, tunjukkan bahwa laju pertumbuhan sebanding dengan populasi, yaitu $\dfrac{dP}{dt}=kP$.
  \item \textbf{(Peluruhan).} Dari $N(t)=N_0e^{-\lambda t}$, turunkan rumus waktu paruh $T=\dfrac{\ln 2}{\lambda}$.
  \item \textbf{(Evaluasi model).} Jelaskan mengapa model eksponensial gagal jangka panjang untuk populasi nyata, dan bagaimana model logistik memperbaikinya.
  \item \textbf{(Pemodelan terbuka).} Rancang model sederhana penyebaran penyakit yang jumlah kasusnya menggandakan diri tiap $3$ hari pada awalnya. Tulis modelnya dan diskusikan keterbatasannya.
\end{enumerate}

\begin{challenge}
\begin{enumerate}[leftmargin=*,topsep=2pt]
  \item Dua cara menabung (ribuan rupiah): $A(t)=100+50t$ (linear) dan $B(t)=100\cdot(1{,}2)^t$ (eksponensial). Tentukan sekitar bulan ke berapa $B$ melampaui $A$.
  \item Pada model logistik $P(t)=\dfrac{K}{1+Ae^{-rt}}$, tunjukkan bahwa pertumbuhan tercepat (titik belok) terjadi saat $P=\dfrac{K}{2}$.
\end{enumerate}
\end{challenge}

% ============================================================
\chapter{Trigonometri Lanjutan}
% ============================================================
\begin{tujuan}
Setelah mempelajari bab ini, peserta didik mampu:
\begin{itemize}[leftmargin=*,topsep=2pt]
  \item menggunakan lingkaran satuan dan menggambar grafik fungsi trigonometri;
  \item menerapkan identitas dasar, rumus jumlah/selisih sudut, sudut rangkap, dan sudut paruh;
  \item menyelesaikan persamaan dan pertidaksamaan trigonometri;
  \item menggunakan aturan sinus dan cosinus pada segitiga;
  \item memodelkan gejala periodik (gelombang, pasang surut) dengan fungsi sinusoidal.
\end{itemize}
\textbf{Capaian Pembelajaran (Fase F+).} Peserta didik dapat memanipulasi identitas trigonometri dan menggunakannya untuk menyelesaikan persamaan serta memodelkan fenomena periodik.
\end{tujuan}

\begin{pemantik}
\begin{itemize}[leftmargin=*,topsep=2pt]
  \item Jika kamu tahu $\sin30^\circ$ dan $\sin45^\circ$, dapatkah kamu menemukan $\sin75^\circ$ tanpa kalkulator?
  \item Mengapa bunyi, cahaya, dan arus listrik semuanya dimodelkan dengan fungsi sinus?
  \item Persamaan $\sin x=\tfrac12$ punya berapa banyak penyelesaian?
\end{itemize}
\end{pemantik}

\begin{studikasus}{Pasang Surut Air Laut}
Ketinggian permukaan laut di sebuah dermaga naik-turun secara periodik. Model sederhananya berbentuk sinusoidal: $h(t)=A\sin\big(\omega t+\varphi\big)+c$, dengan $A$ amplitudo, $\omega$ menentukan periode, dan $c$ ketinggian rata-rata.
\begin{center}
\begin{tikzpicture}
\begin{axis}[axis lines=middle,width=10cm,height=4cm,xlabel={\small $t$ (jam)},ylabel={\small $h$ (m)},
  xmin=0,xmax=24,ymin=-0.5,ymax=4.5,xtick={0,6,12,18,24},ytick={2},samples=120]
\addplot[very thick,biru,domain=0:24]{2+1.5*sin(360/12*x)};
\end{axis}
\end{tikzpicture}\\
\small Ketinggian $h(t)=2+1{,}5\sin\!\big(\tfrac{360^\circ}{12}t\big)$: amplitudo $1{,}5$ m, periode $12$ jam, rata-rata $2$ m.
\end{center}
\textbf{Amati dan duga.} (a) Berapa selisih pasang tertinggi dan terendah? (b) Setiap berapa jam pola berulang? (c) Kapan air mencapai $2{,}75$ m? Menjawab (c) memerlukan \textbf{persamaan trigonometri}.
\end{studikasus}

\section*{Peta Konsep Bab}
\addcontentsline{toc}{section}{Peta Konsep Bab}
\begin{center}
\begin{tikzpicture}[
  node distance=8mm,
  konsepbox/.style={draw=biru,fill=birumuda,rounded corners,align=center,inner sep=5pt,font=\small\bfseries},
  subbox/.style={draw=hijau,fill=hijaumuda,rounded corners,align=center,inner sep=4pt,font=\footnotesize},
  garis/.style={-Stealth,biru,thick}]
  \node[konsepbox] (akar) {TRIGONOMETRI LANJUTAN};
  \node[subbox,below left=12mm and 18mm of akar] (id) {Identitas\\dasar};
  \node[subbox,below=12mm of akar] (sum) {Jumlah/selisih,\\rangkap, paruh};
  \node[subbox,below right=12mm and 16mm of akar] (pers) {Persamaan \&\\pertidaksamaan};
  \node[subbox,below=10mm of sum] (apl) {Aturan sinus/cosinus\\\&\ gelombang};
  \draw[garis] (akar)--(id); \draw[garis] (akar)--(sum); \draw[garis] (akar)--(pers);
  \draw[garis] (sum)--(apl);
\end{tikzpicture}\\
\small \textbf{Peta konsep.} Identitas dasar melahirkan rumus jumlah/selisih, sudut rangkap dan paruh, yang dipakai menyelesaikan persamaan dan memodelkan gelombang.
\end{center}

\begin{catatan}
\textbf{Pengingat Materi Kelas 10 (MAE Bab 4 --- Trigonometri Dasar).} Bab ini adalah kelanjutan langsung dari materi berikut. Tinjau ulang di modul MAE jika belum yakin:
\begin{itemize}[leftmargin=*,topsep=2pt]
  \item Perbandingan trigonometri pada segitiga siku-siku: $\sin\alpha=\tfrac{\text{depan}}{\text{miring}}$, $\cos\alpha=\tfrac{\text{samping}}{\text{miring}}$, $\tan\alpha=\tfrac{\text{depan}}{\text{samping}}$.
  \item Nilai sudut istimewa $0^\circ,30^\circ,45^\circ,60^\circ,90^\circ$ dan cara menurunkannya dari dua segitiga dasar.
  \item Tanda per kuadran: \textbf{A}ll--\textbf{S}in--\textbf{T}an--\textbf{C}os (ASTC).
  \item Identitas dasar: $\sin^2\alpha+\cos^2\alpha=1$.
  \item Rumus jumlah/selisih sudut, sudut rangkap, persamaan trigonometri sederhana (\emph{bagian Pengayaan} MAE Bab 4).
\end{itemize}
\end{catatan}

\section{Lingkaran Satuan dan Identitas Dasar}
\begin{center}
\begin{tikzpicture}[scale=1.4]
  \coordinate (O) at (0,0);
  \coordinate (A) at (1,0);
  \coordinate (P) at (0.6,0.8);
  \draw[->] (-1.3,0)--(1.5,0) node[right]{$x$};
  \draw[->] (0,-1.3)--(0,1.4) node[above]{$y$};
  \draw[biru,thick] (O) circle (1);
  \draw[oranye,very thick,-{Latex}] (O)--(P);
  \draw[dashed,gray] (0.6,0)--(P) (0,0.8)--(P);
  \node[below] at (0.6,0){\small $\cos\theta$};
  \node[left] at (0,0.8){\small $\sin\theta$};
  \pic[draw,angle radius=0.4cm,"$\theta$"]{angle=A--O--P};
\end{tikzpicture}
\qquad
\begin{tikzpicture}
\begin{axis}[axis lines=middle,width=6.5cm,height=4cm,
  xmin=0,xmax=360,ymin=-1.4,ymax=1.4,
  xtick={0,90,180,270,360},ytick={-1,1},xticklabel style={font=\tiny}]
\addplot[very thick,biru,domain=0:360,samples=100]{sin(x)};
\addplot[very thick,oranye,dashed,domain=0:360,samples=100]{cos(x)};
\end{axis}
\end{tikzpicture}\\
\small Kiri: definisi sin \& cos pada lingkaran satuan. Kanan: grafik $y=\sin x$ (biru) dan $y=\cos x$ (oranye).
\end{center}

\smallskip\noindent\textbf{Tanda fungsi di tiap kuadran.} Karena $\cos\theta$ adalah koordinat-$x$ dan $\sin\theta$ koordinat-$y$ pada lingkaran satuan, tandanya berganti tiap kuadran:
\begin{center}
\begin{tikzpicture}[scale=1.15,>=Stealth,line join=round]
\fill[birumuda,opacity=0.5] (0,0)--(2,0) arc(0:90:2)--cycle;
\fill[hijaumuda,opacity=0.5] (0,0)--(0,2) arc(90:180:2)--cycle;
\fill[oranyemuda,opacity=0.5] (0,0)--(-2,0) arc(180:270:2)--cycle;
\fill[ungumuda,opacity=0.5] (0,0)--(0,-2) arc(270:360:2)--cycle;
\draw[biru,thick] (0,0) circle (2);
\draw[->](-2.6,0)--(2.7,0) node[font=\scriptsize,right]{$x$};
\draw[->](0,-2.6)--(0,2.7) node[font=\scriptsize,above]{$y$};
\coordinate (Oc) at (0,0); \coordinate (Xa) at (2,0);
\draw[oranye,very thick,->](Oc)--(35:2) coordinate (Pp);
\pic[draw,angle radius=0.5cm,"$\theta$"]{angle=Xa--Oc--Pp};
\node[font=\scriptsize,align=center] at (1.25,1.25){\textbf{I}\\semua $+$};
\node[font=\scriptsize,align=center] at (-1.25,1.25){\textbf{II}\\$\sin+$};
\node[font=\scriptsize,align=center] at (-1.25,-1.25){\textbf{III}\\$\tan+$};
\node[font=\scriptsize,align=center] at (1.25,-1.25){\textbf{IV}\\$\cos+$};
\end{tikzpicture}\\
\small \textbf{ASTC.} Kuadran I: \emph{semua} positif; II: hanya $\sin$ (\&\ $\csc$); III: hanya $\tan$ (\&\ $\cot$); IV: hanya $\cos$ (\&\ $\sec$).
\end{center}

\begin{konsep}{Tanda per Kuadran dan Sudut Acuan}
Setiap sudut dapat dikaitkan dengan \textbf{sudut acuan} $\theta'$ (sudut lancip terhadap sumbu-$x$). Nilai fungsi sama dengan nilai pada $\theta'$, lalu diberi \emph{tanda} sesuai kuadran (ASTC):
\[
\text{KI: }\theta'=\theta,\quad \text{KII: }\theta'=180^\circ-\theta,\quad \text{KIII: }\theta'=\theta-180^\circ,\quad \text{KIV: }\theta'=360^\circ-\theta.
\]
\end{konsep}

\begin{konsep}{Identitas Dasar}
\[
\sin^2\theta + \cos^2\theta = 1, \qquad
1+\tan^2\theta=\sec^2\theta, \qquad
\tan\theta = \frac{\sin\theta}{\cos\theta}.
\]
\end{konsep}

\begin{contoh}{Membuktikan Identitas}
Buktikan $\dfrac{1-\cos^2\theta}{\sin\theta} = \sin\theta$.

\jawab $1-\cos^2\theta=\sin^2\theta$, sehingga $\dfrac{\sin^2\theta}{\sin\theta}=\sin\theta.$
\end{contoh}

\smallskip\noindent\textbf{Lingkaran satuan dan sudut-sudut istimewa.}
\begin{center}
\begin{tikzpicture}[scale=1.85,line join=round]
\draw[biru,thick] (0,0) circle (1);
\draw[->](-1.28,0)--(1.32,0) node[font=\scriptsize,right]{$x$};
\draw[->](0,-1.28)--(0,1.32) node[font=\scriptsize,above]{$y$};
\foreach \a in {0,30,45,60,90,120,135,150,180,210,225,240,270,300,315,330}{
  \fill[oranye] (\a:1) circle (0.5pt);
  \node[font=\tiny] at (\a:1.15){$\a^\circ$};
}
\fill (0,0) circle (0.6pt);
\end{tikzpicture}
\quad
\begin{tabular}[b]{c|ccccc}
\toprule
$\theta$ & $0^\circ$ & $30^\circ$ & $45^\circ$ & $60^\circ$ & $90^\circ$\\
\midrule
$\sin\theta$ & $0$ & $\tfrac12$ & $\tfrac{\sqrt2}{2}$ & $\tfrac{\sqrt3}{2}$ & $1$\\[2pt]
$\cos\theta$ & $1$ & $\tfrac{\sqrt3}{2}$ & $\tfrac{\sqrt2}{2}$ & $\tfrac12$ & $0$\\[2pt]
$\tan\theta$ & $0$ & $\tfrac{\sqrt3}{3}$ & $1$ & $\sqrt3$ & $-$\\
\bottomrule
\end{tabular}\\
\small Mnemonik: $\sin$ dari $0^\circ$ ke $90^\circ$ adalah $\tfrac{\sqrt0}{2},\tfrac{\sqrt1}{2},\tfrac{\sqrt2}{2},\tfrac{\sqrt3}{2},\tfrac{\sqrt4}{2}$; $\cos$ kebalikannya. Sudut di kuadran lain memakai sudut acuan $+$ tanda ASTC.
\end{center}

\smallskip\noindent\textbf{Sudut berelasi (reduksi).} Empat sudut dengan acuan sama $\theta$:
\begin{center}
\begin{tikzpicture}[scale=1.55,>=Stealth,line join=round]
\draw[biru,thick](0,0) circle (1);
\draw[->](-1.3,0)--(1.35,0) node[font=\scriptsize,right]{$x$};
\draw[->](0,-1.3)--(0,1.35) node[font=\scriptsize,above]{$y$};
\foreach \a in {40,140,220,320}{\draw[oranye,thick,->](0,0)--(\a:1); \fill[oranye](\a:1) circle (0.8pt);}
\node[font=\scriptsize,above right] at (40:1){$\theta$};
\node[font=\scriptsize,above left] at (140:1){$180^\circ{-}\theta$};
\node[font=\scriptsize,below left] at (220:1){$180^\circ{+}\theta$};
\node[font=\scriptsize,below right] at (320:1){$360^\circ{-}\theta$};
\end{tikzpicture}
\end{center}

\begin{konsep}{Rumus Sudut Berelasi}
\[
\begin{array}{lll}
\sin(180^\circ-\theta)=\sin\theta & \cos(180^\circ-\theta)=-\cos\theta & \tan(180^\circ-\theta)=-\tan\theta\\
\sin(180^\circ+\theta)=-\sin\theta & \cos(180^\circ+\theta)=-\cos\theta & \tan(180^\circ+\theta)=\tan\theta\\
\sin(360^\circ-\theta)=-\sin\theta & \cos(360^\circ-\theta)=\cos\theta & \tan(360^\circ-\theta)=-\tan\theta\\
\end{array}
\]
\textbf{Komplemen:} $\sin(90^\circ-\theta)=\cos\theta$ dan $\cos(90^\circ-\theta)=\sin\theta$.
\end{konsep}

\smallskip\noindent\textbf{Grafik tangen dan transformasi sinusoidal.}
\begin{center}
\begin{tikzpicture}
\begin{axis}[axis lines=middle,width=7cm,height=4.3cm,title={\scriptsize $y=\tan x$},
  xmin=0,xmax=360,ymin=-4,ymax=4,xtick={0,90,180,270,360},xticklabel style={font=\tiny},
  ytick={-1,1},samples=200,restrict y to domain=-4.5:4.5]
\addplot[very thick,biru,domain=0:88]{tan(x)};
\addplot[very thick,biru,domain=92:268]{tan(x)};
\addplot[very thick,biru,domain=272:360]{tan(x)};
\draw[dashed,oranye](axis cs:90,-4)--(axis cs:90,4);
\draw[dashed,oranye](axis cs:270,-4)--(axis cs:270,4);
\end{axis}
\end{tikzpicture}
\quad
\begin{tikzpicture}
\begin{axis}[axis lines=middle,width=7cm,height=4.3cm,title={\scriptsize $y=A\sin x+D$},
  xmin=0,xmax=360,ymin=-1.2,ymax=5,xtick={0,90,180,270,360},xticklabel style={font=\tiny},
  ytick={2},samples=120]
\addplot[gray,dashed,domain=0:360]{sin(x)};
\addplot[very thick,biru,domain=0:360]{2*sin(x)+2};
\draw[dashed,hijau!60!black](axis cs:0,2)--(axis cs:360,2);
\node[font=\tiny,hijau!60!black] at (axis cs:320,2.4){garis tengah $D$};
\node[font=\tiny,biru] at (axis cs:90,4.6){$A{=}2$};
\end{axis}
\end{tikzpicture}\\
\small Kiri: $y=\tan x$ berperiode $180^\circ$ dengan asimtot di $90^\circ,270^\circ$. Kanan: pada $y=A\sin\!\big(B(x-C)\big)+D$, $A$ amplitudo, periode $=\tfrac{360^\circ}{B}$, $C$ geseran fase, $D$ garis tengah.
\end{center}

\section{Rumus Jumlah, Selisih, Sudut Rangkap, dan Paruh}
\begin{konsep}{Rumus Penting}
\[
\sin(\alpha\pm\beta)=\sin\alpha\cos\beta\pm\cos\alpha\sin\beta,
\]
\[
\cos(\alpha\pm\beta)=\cos\alpha\cos\beta\mp\sin\alpha\sin\beta.
\]
\textbf{Sudut rangkap:} $\sin2\theta=2\sin\theta\cos\theta$, $\ \cos2\theta=\cos^2\theta-\sin^2\theta=2\cos^2\theta-1=1-2\sin^2\theta$.\\
\textbf{Sudut paruh:} $\sin^2\tfrac\theta2=\dfrac{1-\cos\theta}{2}$, $\ \cos^2\tfrac\theta2=\dfrac{1+\cos\theta}{2}$.
\end{konsep}

\begin{contoh}{Nilai Sudut Tak Baku}
Tentukan nilai $\sin75^\circ$.

\jawab $\sin75^\circ=\sin(45^\circ+30^\circ)=\sin45^\circ\cos30^\circ+\cos45^\circ\sin30^\circ
=\tfrac{\sqrt2}{2}\cdot\tfrac{\sqrt3}{2}+\tfrac{\sqrt2}{2}\cdot\tfrac12=\dfrac{\sqrt6+\sqrt2}{4}.$
\end{contoh}

\section{Persamaan dan Pertidaksamaan Trigonometri}
\begin{konsep}{Penyelesaian Umum}
Untuk $0^\circ\le x<360^\circ$ (lalu tambahkan periode):
\[
\sin x=\sin\alpha \Rightarrow x=\alpha \text{ atau } x=180^\circ-\alpha,
\]
\[
\cos x=\cos\alpha \Rightarrow x=\alpha \text{ atau } x=360^\circ-\alpha.
\]
Solusi umum diperoleh dengan menambahkan kelipatan $360^\circ$ (atau $180^\circ$ untuk $\tan$).
\end{konsep}

\begin{contoh}{Persamaan Trigonometri}
Selesaikan $2\sin x-1=0$ untuk $0^\circ\le x<360^\circ$.

\jawab $\sin x=\tfrac12\Rightarrow x=30^\circ$ atau $x=150^\circ$.
\end{contoh}

\begin{center}
\begin{tikzpicture}[scale=1.5,line join=round]
\draw[biru,thick](0,0) circle (1);
\draw[->](-1.3,0)--(1.35,0) node[font=\scriptsize,right]{$x$};
\draw[->](0,-1.3)--(0,1.35) node[font=\scriptsize,above]{$y$};
\draw[hijau!60!black,dashed](-1.3,0.5)--(1.3,0.5);
\node[hijau!60!black,font=\scriptsize,right] at (1.3,0.5){$y=\tfrac12$};
\draw[oranye,thick,->](0,0)--(30:1); \draw[oranye,thick,->](0,0)--(150:1);
\fill[oranye](30:1) circle (0.9pt) node[font=\scriptsize,above right]{$30^\circ$};
\fill[oranye](150:1) circle (0.9pt) node[font=\scriptsize,above left]{$150^\circ$};
\end{tikzpicture}\\
\small $\sin x=\tfrac12$ berarti koordinat-$y=\tfrac12$: garis mendatar memotong lingkaran di $30^\circ$ (KI) dan $150^\circ$ (KII) --- dua solusi pada $[0^\circ,360^\circ)$.
\end{center}

\section{Aturan Sinus dan Cosinus}
\begin{konsep}{Pada Segitiga}
\textbf{Aturan Sinus:} $\dfrac{a}{\sin A} = \dfrac{b}{\sin B} = \dfrac{c}{\sin C}$.\qquad
\textbf{Aturan Cosinus:} $c^2 = a^2 + b^2 - 2ab\cos C$.\\
\textbf{Luas segitiga:} $L=\tfrac12 ab\sin C$.
\end{konsep}

\begin{center}
\begin{tikzpicture}[scale=1.0,line join=round]
\coordinate (A) at (0,0);\coordinate (B) at (4,0);\coordinate (C) at (1.3,2.1);
\fill[birumuda,opacity=0.5] (A)--(B)--(C)--cycle;
\draw[biru,thick](A)--(B)--(C)--cycle;
\node[font=\scriptsize,below left] at (A){$A$};
\node[font=\scriptsize,below right] at (B){$B$};
\node[font=\scriptsize,above] at (C){$C$};
\node[oranye,font=\scriptsize,below] at ($(A)!0.5!(B)$){$c$};
\node[oranye,font=\scriptsize,above right] at ($(B)!0.5!(C)$){$a$};
\node[oranye,font=\scriptsize,above left] at ($(A)!0.5!(C)$){$b$};
\end{tikzpicture}\\
\small Konvensi: sisi $a$ \emph{di hadapan} sudut $A$, sisi $b$ di hadapan $B$, sisi $c$ di hadapan $C$. (Aturan sinus mengaitkan sisi dengan sudut di hadapannya.)
\end{center}

\begin{contoh}{Aturan Cosinus}
Pada segitiga, $a=7$, $b=5$, dan sudut $C=60^\circ$. Tentukan sisi $c$.

\jawab
$c^2 = 49 + 25 - 2(7)(5)\cos60^\circ = 74 - 70(\tfrac12) = 39$,
sehingga $c=\sqrt{39}\approx6{,}24$.
\end{contoh}

\begin{kesalahan}
$\sin(\alpha+\beta)\neq\sin\alpha+\sin\beta$. Fungsi trigonometri \emph{tidak} linear --- selalu pakai rumus jumlah sudut, bukan "membagikan" sin ke dalam kurung.
\end{kesalahan}

\begin{dunianyata}
Setiap gelombang --- suara, cahaya, sinyal radio, arus AC --- dimodelkan dengan $A\sin(\omega t+\varphi)$. Amplitudo $A$ menentukan kekuatan, $\omega$ menentukan frekuensi, dan $\varphi$ menentukan pergeseran fase. Analisis Fourier bahkan menyatakan \emph{semua} sinyal sebagai jumlah sinus.
\end{dunianyata}

\begin{konsep}{Bahasa Alternatif yang Sering Mengecoh di Soal}
\begin{itemize}[leftmargin=*,topsep=2pt]
  \item \textbf{Derajat vs radian}: $\pi$ rad $=180^\circ$. Pastikan mode kalkulator dan satuan sudut konsisten.
  \item $\sin^{-1}x$ (yaitu \emph{arcsin}, menghasilkan \emph{sudut}) $\neq (\sin x)^{-1}=\dfrac{1}{\sin x}=\csc x$. Namun $\sin^2 x$ berarti $(\sin x)^2$.
  \item \textbf{Invers trigonometri} (arcsin, arccos, arctan) memiliki \emph{rentang nilai utama} terbatas, jadi hanya memberi \emph{satu} sudut --- bukan semua solusi.
  \item \textbf{Periode}: $\sin$ dan $\cos$ berperiode $360^\circ$ ($2\pi$), tetapi $\tan$ berperiode $180^\circ$ ($\pi$).
  \item $\sin x=\sin\alpha$ memiliki \emph{tak hingga} solusi: tambahkan $k\cdot360^\circ$ (atau $k\cdot180^\circ$ untuk $\tan$).
  \item \textbf{Sudut elevasi} diukur ke \emph{atas} dari garis horizontal; \textbf{sudut depresi} ke \emph{bawah}.
  \item \textbf{Sudut acuan/berelasi} adalah sudut lancip terhadap sumbu-$x$ --- bukan sudut sebenarnya.
  \item \textbf{"Sudut rangkap"} ($2\theta$) berbeda dari "dua kali nilai" ($2\sin\theta$): $\sin2\theta\neq2\sin\theta$.
\end{itemize}
\end{konsep}

\section*{Ringkasan Bab}
\addcontentsline{toc}{section}{Ringkasan Bab}
\begin{itemize}[leftmargin=*,topsep=2pt]
  \item Identitas: $\sin^2+\cos^2=1$, $1+\tan^2=\sec^2$.
  \item Jumlah sudut: $\sin(\alpha\pm\beta)$, $\cos(\alpha\pm\beta)$; rangkap $\sin2\theta,\cos2\theta$; paruh.
  \item \textbf{Persamaan:} $\sin x=\sin\alpha\Rightarrow x=\alpha$ atau $180^\circ-\alpha$ ($+k\cdot360^\circ$).
  \item \textbf{Segitiga:} aturan sinus, cosinus, dan luas $\tfrac12ab\sin C$.
  \item Fungsi sinusoidal memodelkan gelombang.
\end{itemize}

\begin{refleksi}
\begin{itemize}[leftmargin=*,topsep=2pt]
  \item[\cek] Saya dapat membuktikan identitas trigonometri dasar.
  \item[\cek] Saya dapat memakai rumus jumlah/selisih, rangkap, dan paruh.
  \item[\cek] Saya dapat menyelesaikan persamaan dan pertidaksamaan trigonometri.
  \item[\cek] Saya dapat menerapkan aturan sinus/cosinus dan model gelombang.
\end{itemize}
\end{refleksi}

\section*{Soal Akhir Bab \thechapter}
\addcontentsline{toc}{section}{Soal Akhir Bab \thechapter}

\bagiansoal{A}{(Fundamental --- setara SNBT Dasar)}
\begin{enumerate}[leftmargin=*]
  \item Hitung $\sin30^\circ+\cos60^\circ$.
  \item Jika $\sin\theta=\tfrac35$ dan $\theta$ lancip, tentukan $\cos\theta$.
  \item Hitung $\sin2\theta$ jika $\sin\theta=\tfrac35$, $\cos\theta=\tfrac45$.
  \item Selesaikan $\sin x=\tfrac12$ untuk $0^\circ\le x<360^\circ$.
  \item Tentukan amplitudo dan periode $y=3\sin(2x)$.
\end{enumerate}

\bagiansoal{B}{(Intermediate --- setara SNBT Sulit)}
\begin{enumerate}[leftmargin=*]
  \item Buktikan $\dfrac{1-\cos^2\theta}{\sin\theta}=\sin\theta$.
  \item Tentukan nilai $\cos75^\circ$ dengan rumus selisih/jumlah sudut.
  \item Selesaikan $2\cos x+1=0$ untuk $0^\circ\le x<360^\circ$.
  \item Dengan aturan sinus, tentukan sisi $b$ jika $a=10$, $A=30^\circ$, $B=45^\circ$.
  \item Buktikan $\cos2\theta=1-2\sin^2\theta$.
  \item Sebuah segitiga memiliki $a=8$, $b=6$, $C=60^\circ$. Tentukan luasnya.
  \item Selesaikan $\sin2x=\sin x$ untuk $0^\circ\le x<360^\circ$.
\end{enumerate}

\bagiansoal{C}{(Advanced --- setara Olimpiade SMA, gabungan antarbab)}
\begin{enumerate}[leftmargin=*]
  \item \textbf{(Identitas + Aljabar).} Jika $\sin\theta+\cos\theta=\tfrac{7}{5}$, tentukan $\sin\theta\cos\theta$.
  \item \textbf{(Trigonometri + Persamaan Kuadrat).} Selesaikan $2\cos^2 x-3\cos x+1=0$ untuk $0^\circ\le x<360^\circ$.
  \item \textbf{(Trigonometri + Bilangan Kompleks).} Dengan De Moivre, buktikan $\sin3\theta=3\sin\theta-4\sin^3\theta$.
  \item \textbf{(Pertidaksamaan).} Tentukan himpunan penyelesaian $\sin x>\tfrac12$ untuk $0^\circ\le x<360^\circ$.
  \item \textbf{(Aturan Cosinus + Geometri).} Pada segitiga dengan sisi $5,7,8$, tentukan kosinus sudut terbesar.
  \item \textbf{(Pemodelan Gelombang).} Pasang surut dimodelkan $h(t)=2+1{,}5\sin\!\big(\tfrac{360^\circ}{12}t\big)$. Tentukan kapan pertama kali $h=2{,}75$ m.
\end{enumerate}

\bagiansoal{D}{(Expert --- setara tahun pertama universitas, sintesis \& pembuktian)}
\begin{enumerate}[leftmargin=*]
  \item \textbf{(Pembuktian).} Buktikan rumus $\cos(\alpha-\beta)=\cos\alpha\cos\beta+\sin\alpha\sin\beta$ menggunakan hasil kali titik dua vektor satuan.
  \item \textbf{(Penjumlahan menjadi perkalian).} Buktikan bahwa
  \[\sin A+\sin B=2\sin\tfrac{A+B}{2}\cos\tfrac{A-B}{2}.\]
  \item \textbf{(Trigonometri + Kalkulus).} Gunakan $\lim_{x\to0}\tfrac{\sin x}{x}=1$ untuk menghitung $\lim_{x\to0}\tfrac{1-\cos x}{x^2}$.
  \item \textbf{(Identitas + Deret).} Buktikan bahwa $\sin a+\sin(a+d)+\cdots+\sin(a+(n-1)d)=\dfrac{\sin\frac{nd}{2}}{\sin\frac{d}{2}}\sin\!\Big(a+\tfrac{(n-1)d}{2}\Big)$.
  \item \textbf{(Pembuktian).} Buktikan bahwa fungsi $f(x)=\sin x$ berperiode $2\pi$ dan tunjukkan bahwa $3\sin x+4\cos x$ dapat ditulis $R\sin(x+\varphi)$; tentukan $R$ dan $\varphi$.
\end{enumerate}

\begin{challenge}
\begin{enumerate}[leftmargin=*,topsep=2pt]
  \item Tanpa kalkulator, hitung $\cos20^\circ\cos40^\circ\cos80^\circ$.
  \item Buktikan bahwa $\tan1^\circ\tan2^\circ\tan3^\circ\cdots\tan89^\circ=1$.
\end{enumerate}
\end{challenge}

% ============================================================
\chapter{Logika Matematika}
% ============================================================
\begin{tujuan}
Setelah mempelajari bab ini, peserta didik mampu:
\begin{itemize}[leftmargin=*,topsep=2pt]
  \item membedakan proposisi dan bukan proposisi serta menyusun negasinya;
  \item menentukan nilai kebenaran konjungsi, disjungsi, implikasi, dan biimplikasi melalui tabel kebenaran;
  \item menyusun konvers, invers, dan kontraposisi serta mengenali ekuivalensi logika;
  \item menafsirkan dan menegasikan pernyataan berkuantor universal dan eksistensial;
  \item menyusun pembuktian langsung, dengan kontraposisi, dan dengan kontradiksi.
\end{itemize}
\textbf{Capaian Pembelajaran (Fase F).} Peserta didik dapat menggunakan logika matematika untuk menilai kebenaran pernyataan dan menyusun argumen serta pembuktian yang sahih.
\end{tujuan}

\begin{pemantik}
\begin{itemize}[leftmargin=*,topsep=2pt]
  \item Mengapa pernyataan "Jika $2>3$ maka bumi datar" justru bernilai \emph{benar} secara logika?
  \item Apa beda "Jika hujan maka jalan basah" dengan "Jika jalan basah maka hujan"? Apakah keduanya selalu sama benarnya?
  \item Bagaimana cara membuktikan sesuatu dengan mengandaikan kebalikannya benar?
\end{itemize}
\end{pemantik}

\begin{studikasus}{Logika di Balik Mesin dan Hukum}
Setiap baris kode "if \dots then \dots", setiap pasal hukum "barang siapa \dots maka \dots", dan setiap rangkaian saklar listrik adalah penerapan logika proposisi. Komputer memutuskan benar/salah ($1/0$) melalui gerbang logika AND, OR, NOT --- persis konjungsi, disjungsi, dan negasi.
\smallskip

\textbf{Amati dan duga.} Saat seorang programmer menulis \texttt{if (umur >= 17 \&\& punyaKTP)}, syarat apa yang membuat seluruh kondisi bernilai benar? Bagaimana menuliskan "kebalikan" dari kondisi itu? Bab ini memberi bahasa yang presisi untuk menjawabnya.
\end{studikasus}

\section*{Peta Konsep Bab}
\addcontentsline{toc}{section}{Peta Konsep Bab}
\begin{center}
\begin{tikzpicture}[
  node distance=8mm,
  konsepbox/.style={draw=biru,fill=birumuda,rounded corners,align=center,inner sep=5pt,font=\small\bfseries},
  subbox/.style={draw=hijau,fill=hijaumuda,rounded corners,align=center,inner sep=4pt,font=\footnotesize},
  garis/.style={-Stealth,biru,thick}]
  \node[konsepbox] (akar) {LOGIKA MATEMATIKA};
  \node[subbox,below left=12mm and 20mm of akar] (op) {Proposisi \&\\operator};
  \node[subbox,below=12mm of akar] (ki) {Konvers, invers,\\kontraposisi};
  \node[subbox,below right=12mm and 20mm of akar] (ku) {Kuantor};
  \node[subbox,below=10mm of ki] (bukti) {Metode\\pembuktian};
  \draw[garis] (akar)--(op); \draw[garis] (akar)--(ki); \draw[garis] (akar)--(ku);
  \draw[garis] (ki)--(bukti);
\end{tikzpicture}\\
\small \textbf{Peta konsep.} Dari proposisi dan tabel kebenaran menuju ekuivalensi, kuantor, hingga metode pembuktian formal.
\end{center}

\section{Proposisi dan Operator Logika}
\begin{definisi}
\textbf{Proposisi} (pernyataan) adalah kalimat yang bernilai \emph{benar} (B) atau \emph{salah} (S), tetapi tidak keduanya. Kalimat tanya, perintah, atau kalimat terbuka (memuat variabel bebas) bukan proposisi.
\end{definisi}

\begin{konsep}{Operator Logika dan Tabel Kebenaran}
Dari proposisi $p,q$ dibentuk:
\begin{itemize}[leftmargin=*,topsep=2pt]
  \item \textbf{Negasi} $\lnot p$ ("bukan $p$"): membalik nilai.
  \item \textbf{Konjungsi} $p\wedge q$ ("$p$ dan $q$"): benar hanya jika keduanya benar.
  \item \textbf{Disjungsi} $p\vee q$ ("$p$ atau $q$"): salah hanya jika keduanya salah.
  \item \textbf{Implikasi} $p\Rightarrow q$ ("jika $p$ maka $q$"): salah hanya jika $p$ benar tetapi $q$ salah.
  \item \textbf{Biimplikasi} $p\Leftrightarrow q$ ("$p$ jika dan hanya jika $q$"): benar jika $p$ dan $q$ bernilai sama.
\end{itemize}
\begin{center}
\begin{tabular}{c c | c c c c c}
\toprule
$p$ & $q$ & $\lnot p$ & $p\wedge q$ & $p\vee q$ & $p\Rightarrow q$ & $p\Leftrightarrow q$\\
\midrule
B & B & S & B & B & B & B\\
B & S & S & S & B & S & S\\
S & B & B & S & B & B & S\\
S & S & B & S & S & B & B\\
\bottomrule
\end{tabular}
\end{center}
\end{konsep}

\begin{contoh}{Nilai Kebenaran}
Tentukan nilai kebenaran "Jika $2>3$ maka $5>1$".

\jawab Anteseden $2>3$ bernilai S. Implikasi dengan anteseden salah selalu \textbf{benar} (B), tanpa memandang konsekuennya.
\end{contoh}

\begin{kesalahan}
Implikasi $p\Rightarrow q$ \emph{tidak} berarti $p$ menyebabkan $q$. Ia hanya salah pada satu kasus: $p$ benar tetapi $q$ salah. Inilah mengapa "jika premis salah" membuat implikasi otomatis benar (vacuously true).
\end{kesalahan}

\section{Konvers, Invers, Kontraposisi, dan Ekuivalensi}
\begin{konsep}{Tiga Bentuk Turunan Implikasi}
Dari implikasi $p\Rightarrow q$:
\[
\text{Konvers: } q\Rightarrow p,\qquad
\text{Invers: } \lnot p\Rightarrow\lnot q,\qquad
\text{Kontraposisi: } \lnot q\Rightarrow\lnot p.
\]
\textbf{Kunci:} $p\Rightarrow q$ ekuivalen dengan kontraposisinya $\lnot q\Rightarrow\lnot p$; sedangkan konvers ekuivalen dengan invers.
\end{konsep}

\begin{konsep}{Ekuivalensi Logika Penting}
\[
\begin{array}{ll}
p\Rightarrow q \equiv \lnot p\vee q & \text{(definisi implikasi)}\\
\lnot(p\wedge q)\equiv \lnot p\vee\lnot q & \text{(De Morgan)}\\
\lnot(p\vee q)\equiv \lnot p\wedge\lnot q & \text{(De Morgan)}\\
\lnot(p\Rightarrow q)\equiv p\wedge\lnot q & \text{(negasi implikasi)}\\
p\Leftrightarrow q \equiv (p\Rightarrow q)\wedge(q\Rightarrow p) &
\end{array}
\]
Dua pernyataan \textbf{ekuivalen} bila tabel kebenarannya identik. Pernyataan yang selalu benar disebut \emph{tautologi}.
\end{konsep}

\begin{contoh}{Kontraposisi}
Tulis kontraposisi dari "Jika hujan maka jalan basah".

\jawab $p$: hujan, $q$: jalan basah. Kontraposisi $\lnot q\Rightarrow\lnot p$: "Jika jalan tidak basah maka tidak hujan". Nilainya selalu sama dengan implikasi asal.
\end{contoh}

\section{Kuantor}
\begin{konsep}{Kuantor Universal dan Eksistensial}
\textbf{Universal} $\forall x,\,P(x)$: "untuk setiap $x$, berlaku $P(x)$". \textbf{Eksistensial} $\exists x,\,P(x)$: "ada (paling sedikit satu) $x$ yang memenuhi $P(x)$". Negasinya saling bertukar:
\[
\lnot\big(\forall x,\,P(x)\big)\equiv \exists x,\,\lnot P(x),\qquad
\lnot\big(\exists x,\,P(x)\big)\equiv \forall x,\,\lnot P(x).
\]
\end{konsep}

\begin{contoh}{Negasi Berkuantor}
Negasikan "Semua bilangan prima adalah ganjil".

\jawab Bentuk $\forall x,\,P(x)$. Negasinya $\exists x,\,\lnot P(x)$: "Ada bilangan prima yang tidak ganjil" (benar, yaitu $2$).
\end{contoh}

\section{Metode Pembuktian}
\begin{konsep}{Tiga Strategi Dasar}
\begin{itemize}[leftmargin=*,topsep=2pt]
  \item \textbf{Pembuktian langsung:} andaikan $p$ benar, turunkan $q$ melalui rangkaian implikasi.
  \item \textbf{Pembuktian kontraposisi:} buktikan $\lnot q\Rightarrow\lnot p$ (ekuivalen dengan $p\Rightarrow q$).
  \item \textbf{Pembuktian kontradiksi:} andaikan pernyataan salah (mis. andaikan $p\wedge\lnot q$), lalu turunkan kemustahilan.
\end{itemize}
\end{konsep}

\begin{contoh}{Pembuktian Kontraposisi}
Buktikan: jika $n^2$ ganjil maka $n$ ganjil.

\jawab Kontraposisi: "jika $n$ genap maka $n^2$ genap". Misal $n=2k$, maka $n^2=4k^2=2(2k^2)$ genap. Karena kontraposisi benar, pernyataan asal benar. $\blacksquare$
\end{contoh}

\begin{dunianyata}
Logika adalah fondasi pemrograman (percabangan, gerbang logika digital), basis data (kueri), kecerdasan buatan (sistem pakar, inferensi), dan penalaran hukum. Pembuktian formal menjamin algoritma kriptografi dan perangkat keras bekerja benar.
\end{dunianyata}

\section*{Ringkasan Bab}
\addcontentsline{toc}{section}{Ringkasan Bab}
\begin{itemize}[leftmargin=*,topsep=2pt]
  \item Proposisi bernilai B/S; operator: $\lnot,\wedge,\vee,\Rightarrow,\Leftrightarrow$ (lihat tabel kebenaran).
  \item $p\Rightarrow q\equiv\lnot p\vee q$; salah hanya saat $p$ benar, $q$ salah.
  \item $p\Rightarrow q$ ekuivalen kontraposisinya $\lnot q\Rightarrow\lnot p$; konvers $\equiv$ invers.
  \item De Morgan: $\lnot(p\wedge q)\equiv\lnot p\vee\lnot q$; $\lnot(p\vee q)\equiv\lnot p\wedge\lnot q$.
  \item Negasi kuantor: $\lnot\forall\to\exists\lnot$ dan $\lnot\exists\to\forall\lnot$.
  \item Pembuktian: langsung, kontraposisi, kontradiksi.
\end{itemize}

\begin{refleksi}
\begin{itemize}[leftmargin=*,topsep=2pt]
  \item[\cek] Saya dapat menyusun tabel kebenaran dan menentukan nilai pernyataan majemuk.
  \item[\cek] Saya dapat membentuk konvers, invers, dan kontraposisi.
  \item[\cek] Saya dapat menegasikan pernyataan majemuk dan berkuantor.
  \item[\cek] Saya dapat memilih dan menjalankan metode pembuktian yang tepat.
\end{itemize}
\end{refleksi}

\section*{Soal Akhir Bab \thechapter}
\addcontentsline{toc}{section}{Soal Akhir Bab \thechapter}

\bagiansoal{A}{(Fundamental --- setara SNBT Dasar)}
\begin{enumerate}[leftmargin=*]
  \item Tentukan negasi dari "Semua siswa rajin".
  \item Tentukan nilai kebenaran "$2+2=4$ dan $3>5$".
  \item Tentukan nilai kebenaran "$2>5$ atau $4=4$".
  \item Tentukan negasi dari "Hari ini hujan".
  \item Tulis kontraposisi dari "Jika hujan maka jalan basah".
  \item Tentukan nilai kebenaran "Jika $2>3$ maka $5>1$".
  \item Tentukan negasi dari $p\wedge q$.
  \item Tulis konvers dari "Jika $x>0$ maka $x^2>0$".
  \item Kapan biimplikasi $p\Leftrightarrow q$ bernilai benar?
  \item Tentukan negasi dari "$\exists x,\,P(x)$".
\end{enumerate}

\bagiansoal{B}{(Intermediate --- setara SNBT Sulit)}
\begin{enumerate}[leftmargin=*]
  \item Susun tabel kebenaran $p\Rightarrow q$ dan sebutkan kapan ia bernilai salah.
  \item Tentukan negasi dari "Jika $n$ genap maka $n^2$ genap".
  \item Tunjukkan dengan tabel kebenaran bahwa $p\Rightarrow q\equiv\lnot p\vee q$.
  \item Tentukan negasi dari "Untuk setiap bilangan real $x$, $x^2\ge0$".
  \item Diberikan $p$ benar, $q$ salah, $r$ benar. Tentukan nilai $(p\wedge\lnot q)\Rightarrow r$.
\end{enumerate}

\bagiansoal{C}{(Advanced --- setara Olimpiade SMA, gabungan antarbab)}
\begin{enumerate}[leftmargin=*]
  \item \textbf{(Logika + Bilangan).} Buktikan kontraposisi "jika $n^2$ genap maka $n$ genap", yaitu "jika $n$ ganjil maka $n^2$ ganjil".
  \item \textbf{(Logika + Tabel).} Tunjukkan $[(p\Rightarrow q)\wedge p]\Rightarrow q$ adalah tautologi (modus ponens).
  \item \textbf{(Logika + Kuantor).} Negasikan "$\forall x\,\exists y,\ x+y=0$".
  \item \textbf{(Logika + Bilangan).} Tentukan nilai kebenaran "$\forall x\in\mathbb{R}\,\exists y\in\mathbb{R},\ y^2=x$".
  \item \textbf{(Logika + Ekuivalensi).} Sederhanakan $\lnot(p\Rightarrow q)$ dan $\lnot(p\Leftrightarrow q)$.
\end{enumerate}

\bagiansoal{D}{(Expert --- setara tahun pertama universitas, sintesis \& pembuktian)}
\begin{enumerate}[leftmargin=*]
  \item \textbf{(Kontradiksi).} Buktikan bahwa banyaknya bilangan prima tak hingga (Euclid).
  \item \textbf{(Tautologi).} Buktikan $[(p\Rightarrow q)\wedge(q\Rightarrow r)]\Rightarrow(p\Rightarrow r)$ tautologi (silogisme hipotetis).
  \item \textbf{(Pembuktian Langsung).} Buktikan bahwa jumlah dua bilangan genap selalu genap.
  \item \textbf{(Kontraposisi).} Buktikan: jika $3n+2$ ganjil maka $n$ ganjil.
  \item \textbf{(Kuantor).} Buktikan benar bahwa "$\forall x\in\mathbb{R},\ x^2-2x+2>0$".
\end{enumerate}

\begin{challenge}
\begin{enumerate}[leftmargin=*,topsep=2pt]
  \item Di sebuah pulau, ksatria selalu berkata benar dan penipu selalu berbohong. Penduduk $A$ berkata: "Saya penipu". Tunjukkan dengan logika bahwa tidak ada penduduk yang dapat mengucapkan kalimat itu.
  \item Buktikan dengan kontradiksi: jika $3n+2$ genap maka $n$ genap.
  \item Buktikan bahwa $\sqrt2+\sqrt3$ adalah bilangan irasional.
\end{enumerate}
\end{challenge}

% ============================================================
\chapter{Induksi Matematika}
% ============================================================
\begin{tujuan}
Setelah mempelajari bab ini, peserta didik mampu:
\begin{itemize}[leftmargin=*,topsep=2pt]
  \item menjelaskan prinsip induksi matematika (basis dan langkah induksi);
  \item membuktikan rumus jumlah deret dengan induksi;
  \item membuktikan sifat keterbagian dengan induksi;
  \item membuktikan ketaksamaan dengan induksi;
  \item menerapkan induksi pada barisan dan pangkat matriks.
\end{itemize}
\textbf{Capaian Pembelajaran (Fase F).} Peserta didik dapat menyusun pembuktian pernyataan tentang bilangan asli menggunakan prinsip induksi matematika.
\end{tujuan}

\begin{pemantik}
\begin{itemize}[leftmargin=*,topsep=2pt]
  \item Jika kartu domino disusun berderet, cukup syarat apa agar \emph{semua} pasti roboh ketika yang pertama dijatuhkan?
  \item Mengapa membuktikan "berlaku untuk $1$" dan "jika berlaku untuk $k$ maka untuk $k+1$" sudah cukup untuk semua bilangan asli?
  \item Apa bedanya mengecek $1000$ kasus dengan membuktikan satu langkah induksi?
\end{itemize}
\end{pemantik}

\begin{studikasus}{Efek Domino}
Bayangkan barisan domino tak berujung. Anda ingin memastikan \emph{semua} domino roboh. Tidak mungkin mendorong satu per satu selamanya. Cukup jamin dua hal: (1) domino \textbf{pertama} roboh; (2) \textbf{setiap} domino yang roboh pasti menjatuhkan domino berikutnya.
\smallskip

\textbf{Amati dan duga.} Jika kedua syarat itu dipenuhi, mengapa kita yakin domino ke-$1000$ --- bahkan ke-sejuta --- pasti roboh tanpa harus mengeceknya? Inilah jantung \emph{induksi matematika}: \textbf{basis} (domino pertama) dan \textbf{langkah induksi} (yang ke-$k$ menjatuhkan yang ke-$k+1$).
\end{studikasus}

\section*{Peta Konsep Bab}
\addcontentsline{toc}{section}{Peta Konsep Bab}
\begin{center}
\begin{tikzpicture}[
  node distance=8mm,
  konsepbox/.style={draw=biru,fill=birumuda,rounded corners,align=center,inner sep=5pt,font=\small\bfseries},
  subbox/.style={draw=hijau,fill=hijaumuda,rounded corners,align=center,inner sep=4pt,font=\footnotesize},
  garis/.style={-Stealth,biru,thick}]
  \node[konsepbox] (akar) {INDUKSI MATEMATIKA};
  \node[subbox,below left=12mm and 20mm of akar] (pr) {Prinsip:\\basis \&\ langkah};
  \node[subbox,below=12mm of akar] (jml) {Jumlah \&\\barisan};
  \node[subbox,below right=12mm and 20mm of akar] (kb) {Keterbagian \&\\ketaksamaan};
  \node[subbox,below=10mm of jml] (mat) {Pangkat\\matriks};
  \draw[garis] (akar)--(pr); \draw[garis] (akar)--(jml); \draw[garis] (akar)--(kb);
  \draw[garis] (jml)--(mat);
\end{tikzpicture}\\
\small \textbf{Peta konsep.} Satu prinsip (basis + langkah) diterapkan pada jumlah, keterbagian, ketaksamaan, hingga matriks.
\end{center}

\section{Prinsip Induksi Matematika}
\begin{konsep}{Prinsip Induksi}
Misal $P(n)$ pernyataan untuk setiap bilangan asli $n\ge n_0$. Jika:
\begin{enumerate}[leftmargin=*,topsep=2pt,label=(\arabic*)]
  \item \textbf{Basis induksi:} $P(n_0)$ benar; dan
  \item \textbf{Langkah induksi:} untuk sembarang $k\ge n_0$, \emph{andaikan} $P(k)$ benar (hipotesis induksi), lalu tunjukkan $P(k+1)$ benar,
\end{enumerate}
maka $P(n)$ benar untuk semua $n\ge n_0$.
\end{konsep}

\begin{contoh}{Pembuktian Jumlah}
Buktikan $1+2+3+\cdots+n=\dfrac{n(n+1)}{2}$ untuk setiap $n\ge1$.

\jawab \textbf{Basis} $n=1$: ruas kiri $=1$, ruas kanan $=\tfrac{1\cdot2}{2}=1$. Benar.

\textbf{Langkah.} Andaikan benar untuk $n=k$: $1+2+\cdots+k=\dfrac{k(k+1)}{2}$. Maka
\[
1+2+\cdots+k+(k+1)=\frac{k(k+1)}{2}+(k+1)=(k+1)\!\left(\frac{k}{2}+1\right)=\frac{(k+1)(k+2)}{2},
\]
yang persis bentuk rumus untuk $n=k+1$. Jadi pernyataan benar untuk semua $n\ge1$. $\blacksquare$
\end{contoh}

\begin{kesalahan}
Langkah induksi \emph{harus} memakai hipotesis $P(k)$ untuk menurunkan $P(k+1)$. Membuktikan $P(k+1)$ secara terpisah (tanpa menggunakan $P(k)$) bukan induksi. Sebaliknya, basis tidak boleh dilupakan --- tanpa "domino pertama", rantai tak pernah mulai.
\end{kesalahan}

\section{Pembuktian Keterbagian dan Ketaksamaan}
\begin{contoh}{Keterbagian}
Buktikan $3\mid(n^3-n)$ untuk setiap $n\ge1$.

\jawab \textbf{Basis} $n=1$: $1-1=0$ habis dibagi $3$. \textbf{Langkah.} Andaikan $3\mid(k^3-k)$. Maka
\[
(k+1)^3-(k+1)=k^3+3k^2+3k+1-k-1=(k^3-k)+3(k^2+k).
\]
Suku pertama habis dibagi $3$ (hipotesis), suku kedua jelas kelipatan $3$. Jadi jumlahnya habis dibagi $3$. $\blacksquare$
\end{contoh}

\begin{contoh}{Ketaksamaan}
Buktikan $2^n>n^2$ untuk setiap $n\ge5$.

\jawab \textbf{Basis} $n=5$: $2^5=32>25=5^2$. \textbf{Langkah.} Andaikan $2^k>k^2$ untuk suatu $k\ge5$. Maka $2^{k+1}=2\cdot2^k>2k^2$. Cukup tunjukkan $2k^2\ge(k+1)^2$, yaitu $2k^2\ge k^2+2k+1\iff k^2-2k-1\ge0$, yang benar untuk $k\ge5$ (sebab $k^2-2k-1=(k-1)^2-2>0$). Maka $2^{k+1}>(k+1)^2$. $\blacksquare$
\end{contoh}

\section{Induksi pada Barisan dan Matriks}
\begin{contoh}{Pangkat Matriks}
Buktikan $\begin{pmatrix}1&1\\0&1\end{pmatrix}^{\!n}=\begin{pmatrix}1&n\\0&1\end{pmatrix}$ untuk $n\ge1$.

\jawab \textbf{Basis} $n=1$: jelas benar. \textbf{Langkah.} Andaikan benar untuk $n=k$. Maka
\[
\begin{pmatrix}1&1\\0&1\end{pmatrix}^{\!k+1}
=\begin{pmatrix}1&k\\0&1\end{pmatrix}\begin{pmatrix}1&1\\0&1\end{pmatrix}
=\begin{pmatrix}1&k+1\\0&1\end{pmatrix}.
\]
Sesuai bentuk untuk $n=k+1$. $\blacksquare$
\end{contoh}

\begin{dunianyata}
Induksi adalah tulang punggung pembuktian \emph{kebenaran algoritma} (loop invariant) dan analisis kompleksitas dalam ilmu komputer, serta rekursi pada struktur data. Banyak rumus deret keuangan dan model populasi diskret divalidasi dengan induksi.
\end{dunianyata}

\section*{Ringkasan Bab}
\addcontentsline{toc}{section}{Ringkasan Bab}
\begin{itemize}[leftmargin=*,topsep=2pt]
  \item Dua langkah wajib: \textbf{basis} $P(n_0)$ dan \textbf{langkah} $P(k)\Rightarrow P(k+1)$.
  \item Untuk \textbf{jumlah}: tambahkan suku ke-$(k+1)$ ke hipotesis, lalu faktorkan.
  \item Untuk \textbf{keterbagian}: pisahkan bentuk $P(k+1)$ menjadi "(bagian hipotesis) $+$ (kelipatan pembagi)".
  \item Untuk \textbf{ketaksamaan}: gunakan hipotesis lalu perkuat dengan estimasi tambahan.
  \item Induksi juga berlaku untuk barisan rekursif dan pangkat matriks.
\end{itemize}

\begin{refleksi}
\begin{itemize}[leftmargin=*,topsep=2pt]
  \item[\cek] Saya dapat menuliskan basis dan hipotesis induksi dengan benar.
  \item[\cek] Saya dapat membuktikan rumus jumlah deret dengan induksi.
  \item[\cek] Saya dapat membuktikan keterbagian dan ketaksamaan dengan induksi.
  \item[\cek] Saya dapat menerapkan induksi pada matriks/barisan rekursif.
\end{itemize}
\end{refleksi}

\section*{Soal Akhir Bab \thechapter}
\addcontentsline{toc}{section}{Soal Akhir Bab \thechapter}

\bagiansoal{A}{(Fundamental --- setara SNBT Dasar)}
\begin{enumerate}[leftmargin=*]
  \item Tentukan basis induksi (cek $n=1$) untuk $1+2+\cdots+n=\dfrac{n(n+1)}{2}$.
  \item Tuliskan hipotesis induksi $P(k)$ untuk pernyataan $1+3+5+\cdots+(2n-1)=n^2$.
  \item Hitung $1+2+3+\cdots+10$.
  \item Hitung $1+3+5+\cdots+(2\cdot5-1)$.
  \item Periksa basis $n=1$ untuk pernyataan $2^n>n$.
  \item Pada pembuktian $1+2+\cdots+n$, suku apa yang ditambahkan saat beralih dari $n=k$ ke $n=k+1$?
  \item Periksa apakah $3\mid(n^3-n)$ untuk $n=2$.
  \item Tuliskan bentuk $P(k+1)$ dari $P(n):\ n^2$.
  \item Hitung $2+4+6+8$ dan bandingkan dengan $n(n+1)$ untuk $n=4$.
  \item Berapa nilai $n_0$ (basis terkecil) untuk membuktikan $2^n>n^2$?
\end{enumerate}

\bagiansoal{B}{(Intermediate --- setara SNBT Sulit)}
\begin{enumerate}[leftmargin=*]
  \item Buktikan dengan induksi: $1+2+\cdots+n=\dfrac{n(n+1)}{2}$.
  \item Buktikan dengan induksi: $1+3+5+\cdots+(2n-1)=n^2$.
  \item Buktikan dengan induksi: $1^2+2^2+\cdots+n^2=\dfrac{n(n+1)(2n+1)}{6}$.
  \item Buktikan dengan induksi: $3\mid(n^3-n)$ untuk semua $n\ge1$.
  \item Buktikan dengan induksi: $1+2+2^2+\cdots+2^{n-1}=2^n-1$.
\end{enumerate}

\bagiansoal{C}{(Advanced --- setara Olimpiade SMA, gabungan antarbab)}
\begin{enumerate}[leftmargin=*]
  \item \textbf{(Induksi + Deret Geometri).} Buktikan $a+ar+\cdots+ar^{n-1}=\dfrac{a(r^n-1)}{r-1}$ untuk $r\neq1$.
  \item \textbf{(Induksi + Ketaksamaan).} Buktikan $2^n>n^2$ untuk $n\ge5$.
  \item \textbf{(Induksi + Keterbagian).} Buktikan $6\mid(n^3+5n)$ untuk semua $n\ge1$.
  \item \textbf{(Induksi + Ketaksamaan Bernoulli).} Buktikan $(1+x)^n\ge1+nx$ untuk $x\ge-1$ dan $n\ge1$.
  \item \textbf{(Induksi + Matriks).} Buktikan $\begin{pmatrix}1&1\\0&1\end{pmatrix}^{\!n}=\begin{pmatrix}1&n\\0&1\end{pmatrix}$.
\end{enumerate}

\bagiansoal{D}{(Expert --- setara tahun pertama universitas, sintesis \& pembuktian)}
\begin{enumerate}[leftmargin=*]
  \item \textbf{(Induksi Kuat).} Buktikan setiap bilangan asli $n\ge2$ dapat dinyatakan sebagai hasil kali bilangan-bilangan prima.
  \item \textbf{(Induksi + Matriks Fibonacci).} Dengan $F_1=F_2=1$, buktikan $\begin{pmatrix}1&1\\1&0\end{pmatrix}^{\!n}=\begin{pmatrix}F_{n+1}&F_n\\F_n&F_{n-1}\end{pmatrix}$.
  \item \textbf{(Induksi + Ketaksamaan Harmonik).} Buktikan $1+\dfrac12+\dfrac13+\cdots+\dfrac{1}{2^n}\ge1+\dfrac{n}{2}$.
  \item \textbf{(Induksi + Kombinatorik).} Buktikan bahwa himpunan dengan $n$ anggota memiliki tepat $2^n$ himpunan bagian.
  \item \textbf{(Induksi + Bilangan Kompleks).} Buktikan teorema De Moivre $(\cos\theta+i\sin\theta)^n=\cos n\theta+i\sin n\theta$.
\end{enumerate}

\begin{challenge}
\begin{enumerate}[leftmargin=*,topsep=2pt]
  \item Tunjukkan letak kesalahan pada "bukti" bahwa semua kuda berwarna sama: basis $n=1$ jelas benar; pada langkah, dari sembarang $k$ kuda berwarna sama disimpulkan $k+1$ kuda berwarna sama. Di mana argumen ini runtuh?
  \item Buktikan dengan induksi bahwa $n!>2^n$ untuk setiap $n\ge4$.
  \item Buktikan dengan induksi bahwa $F_1+F_2+\cdots+F_n=F_{n+2}-1$ (jumlah bilangan Fibonacci).
\end{enumerate}
\end{challenge}

% ============================================================
\part*{\centering FASE F --- KELAS 12}
\addcontentsline{toc}{part}{FASE F --- KELAS 12}

% ============================================================
\chapter{Aturan Pencacahan}
% ============================================================
\begin{tujuan}
Setelah mempelajari bab ini, peserta didik mampu:
\begin{itemize}[leftmargin=*,topsep=2pt]
  \item menerapkan aturan penjumlahan dan aturan perkalian dalam pencacahan;
  \item menghitung dan memaknai faktorial;
  \item membedakan dan menghitung permutasi (termasuk siklis dan unsur sama) serta kombinasi;
  \item menggunakan kombinasi untuk menyusun ekspansi Binomial Newton;
  \item memilih teknik pencacahan yang tepat untuk suatu masalah.
\end{itemize}
\textbf{Capaian Pembelajaran (Fase F).} Peserta didik dapat menyelesaikan masalah pencacahan menggunakan aturan penjumlahan, perkalian, permutasi, dan kombinasi.
\end{tujuan}

\begin{pemantik}
\begin{itemize}[leftmargin=*,topsep=2pt]
  \item Kapan kita "menjumlahkan" pilihan dan kapan "mengalikan"-nya?
  \item Mengapa urutan penting pada "juara 1, 2, 3" tetapi tidak penting pada "memilih 3 anggota tim"?
  \item Berapa banyak kata sandi 4 digit yang mungkin? Mengapa angka itu terasa "aman" padahal kita hanya memakai 10 angka?
\end{itemize}
\end{pemantik}

\begin{studikasus}{Menu, PIN, dan Plat Nomor}
Sebuah kantin menawarkan $3$ makanan dan $4$ minuman. Jika satu paket terdiri atas satu makanan \emph{dan} satu minuman, banyak paket $=3\times4=12$ (aturan perkalian). Tetapi jika seseorang hanya ingin membeli satu item --- makanan \emph{atau} minuman --- banyak pilihan $=3+4=7$ (aturan penjumlahan).
\smallskip

\textbf{Amati dan duga.} Kata "dan" mengarah ke perkalian, "atau" ke penjumlahan. Bagaimana dengan PIN ATM $6$ digit? Berapa banyak kemungkinannya, dan mengapa urutan digit membuat jumlahnya meledak? Bab ini membangun alat untuk mencacah kemungkinan secara sistematis.
\end{studikasus}

\section*{Peta Konsep Bab}
\addcontentsline{toc}{section}{Peta Konsep Bab}
\begin{center}
\begin{tikzpicture}[
  node distance=8mm,
  konsepbox/.style={draw=biru,fill=birumuda,rounded corners,align=center,inner sep=5pt,font=\small\bfseries},
  subbox/.style={draw=hijau,fill=hijaumuda,rounded corners,align=center,inner sep=4pt,font=\footnotesize},
  garis/.style={-Stealth,biru,thick}]
  \node[konsepbox] (akar) {ATURAN PENCACAHAN};
  \node[subbox,below left=12mm and 20mm of akar] (ap) {Penjumlahan \&\\perkalian};
  \node[subbox,below=12mm of akar] (perm) {Faktorial \&\\permutasi};
  \node[subbox,below right=12mm and 20mm of akar] (komb) {Kombinasi};
  \node[subbox,below=10mm of komb] (bin) {Binomial\\Newton};
  \draw[garis] (akar)--(ap); \draw[garis] (akar)--(perm); \draw[garis] (akar)--(komb);
  \draw[garis] (komb)--(bin);
\end{tikzpicture}\\
\small \textbf{Peta konsep.} Dua aturan dasar melahirkan permutasi (urutan penting) dan kombinasi (urutan tak penting), yang menuntun ke Binomial Newton.
\end{center}

\section{Aturan Penjumlahan dan Perkalian}
\begin{konsep}{Dua Aturan Dasar}
\textbf{Aturan penjumlahan:} jika kejadian $A$ dapat terjadi $m$ cara dan kejadian $B$ (yang \emph{saling lepas} dengan $A$) dapat terjadi $n$ cara, maka "$A$ atau $B$" terjadi $m+n$ cara. \textbf{Aturan perkalian:} jika suatu prosedur terdiri atas tahap berurutan dengan $m$ lalu $n$ cara, maka seluruh prosedur dapat dilakukan $m\times n$ cara.
\end{konsep}

\begin{contoh}{Aturan Perkalian}
Berapa banyak PIN $4$ digit (angka $0$--$9$, boleh berulang)?

\jawab Tiap posisi $10$ pilihan, empat posisi: $10\times10\times10\times10=10^4=10.000$.
\end{contoh}

\section{Faktorial dan Permutasi}
\begin{definisi}
\textbf{Faktorial} $n!=n\cdot(n-1)\cdots2\cdot1$ untuk $n\ge1$, dan $0!=1$. \textbf{Permutasi} $r$ unsur dari $n$ unsur (urutan diperhatikan):
\[
P(n,r)={}^{n}P_{r}=\frac{n!}{(n-r)!}.
\]
\end{definisi}

\begin{konsep}{Permutasi Khusus}
\textbf{Unsur sama:} banyak susunan $n$ objek dengan kelompok identik berukuran $n_1,n_2,\dots,n_k$ adalah $\dfrac{n!}{n_1!\,n_2!\cdots n_k!}$. \textbf{Siklis (melingkar):} banyak susunan $n$ objek berbeda dalam lingkaran adalah $(n-1)!$ (karena rotasi dianggap sama).
\end{konsep}

\begin{contoh}{Permutasi Unsur Sama}
Berapa banyak susunan huruf berbeda dari kata "BUKU"?

\jawab Huruf: B, U, K, U --- huruf U muncul $2$ kali. Banyak susunan $=\dfrac{4!}{2!}=\dfrac{24}{2}=12$.
\end{contoh}

\begin{kesalahan}
Jangan memakai permutasi ($P$) saat urutan tidak penting --- itu akan menghitung berlebih. "Memilih 3 pemenang berbeda peringkat" memakai $P$; "memilih 3 anggota tim" memakai kombinasi $C$.
\end{kesalahan}

\section{Kombinasi dan Binomial Newton}
\begin{konsep}{Kombinasi}
\textbf{Kombinasi} $r$ unsur dari $n$ unsur (urutan \emph{tidak} diperhatikan):
\[
C(n,r)={}^{n}C_{r}=\binom{n}{r}=\frac{n!}{r!\,(n-r)!}.
\]
Hubungannya dengan permutasi: $P(n,r)=r!\cdot C(n,r)$. Sifat: $\binom{n}{r}=\binom{n}{n-r}$.
\end{konsep}

\begin{konsep}{Binomial Newton (Pengantar)}
\[
(a+b)^n=\sum_{k=0}^{n}\binom{n}{k}a^{n-k}b^{k}.
\]
Suku ke-$(k+1)$ adalah $\binom{n}{k}a^{n-k}b^{k}$. Koefisien $\binom nk$ membentuk \emph{segitiga Pascal}, dengan $\binom nk=\binom{n-1}{k-1}+\binom{n-1}{k}$.
\end{konsep}

\begin{center}
\small
$\begin{array}{ccccccccc}
 & & & & 1 & & & &\\
 & & & 1 & & 1 & & &\\
 & & 1 & & 2 & & 1 & &\\
 & 1 & & 3 & & 3 & & 1 &\\
1 & & 4 & & 6 & & 4 & & 1
\end{array}$\\
\small Segitiga Pascal: baris ke-$n$ memuat koefisien $(a+b)^n$.
\end{center}

\begin{contoh}{Suku pada Ekspansi Binomial}
Tentukan koefisien $x^2$ pada $(1+x)^5$.

\jawab Suku memuat $x^2$ adalah $\binom{5}{2}(1)^{3}x^{2}=10x^2$. Koefisiennya $10$.
\end{contoh}

\begin{dunianyata}
Pencacahan mendasari kriptografi (ruang kunci sandi), pengkodean (kode pos, plat nomor), genetika (kombinasi gen), dan tentu saja peluang. Estimasi "berapa banyak kemungkinan" menentukan apakah sebuah brute-force layak atau mustahil.
\end{dunianyata}

\section*{Ringkasan Bab}
\addcontentsline{toc}{section}{Ringkasan Bab}
\begin{itemize}[leftmargin=*,topsep=2pt]
  \item "Atau" (saling lepas) $\to$ jumlah; "dan"/bertahap $\to$ kali.
  \item $n!=n(n-1)\cdots1$, $0!=1$.
  \item Permutasi (urutan penting): $P(n,r)=\dfrac{n!}{(n-r)!}$; siklis $(n-1)!$; unsur sama $\dfrac{n!}{n_1!\cdots n_k!}$.
  \item Kombinasi (urutan tak penting): $\binom nr=\dfrac{n!}{r!(n-r)!}$, $\binom nr=\binom{n}{n-r}$.
  \item Binomial: $(a+b)^n=\sum_{k=0}^n\binom nk a^{n-k}b^k$.
\end{itemize}

\begin{refleksi}
\begin{itemize}[leftmargin=*,topsep=2pt]
  \item[\cek] Saya dapat memilih antara aturan penjumlahan dan perkalian.
  \item[\cek] Saya dapat menghitung permutasi, termasuk siklis dan unsur sama.
  \item[\cek] Saya dapat menghitung kombinasi dan membedakannya dari permutasi.
  \item[\cek] Saya dapat menuliskan ekspansi dan suku Binomial Newton.
\end{itemize}
\end{refleksi}

\section*{Soal Akhir Bab \thechapter}
\addcontentsline{toc}{section}{Soal Akhir Bab \thechapter}

\bagiansoal{A}{(Fundamental --- setara SNBT Dasar)}
\begin{enumerate}[leftmargin=*]
  \item Hitung $5!$.
  \item Hitung $0!$.
  \item Hitung $P(5,2)$.
  \item Hitung $C(5,2)$.
  \item Seseorang memilih satu item dari $3$ baju \emph{atau} $2$ celana. Berapa pilihannya?
  \item Berapa pasang (baju, celana) dari $3$ baju dan $2$ celana?
  \item Hitung $\dfrac{6!}{4!}$.
  \item Hitung $\binom{6}{0}+\binom{6}{6}$.
  \item Berapa banyak susunan huruf berbeda dari kata "PQRS"?
  \item Hitung $\binom{7}{5}$.
\end{enumerate}

\bagiansoal{B}{(Intermediate --- setara SNBT Sulit)}
\begin{enumerate}[leftmargin=*]
  \item Berapa banyak bilangan $3$ angka berbeda yang dapat dibentuk dari $\{1,2,3,4,5\}$?
  \item Berapa banyak cara $6$ orang duduk mengelilingi meja bundar?
  \item Berapa banyak susunan huruf berbeda dari kata "MATEMATIKA"?
  \item Dari $10$ orang dipilih $3$ untuk panitia (tanpa jabatan). Berapa cara?
  \item Tentukan koefisien $x^2$ pada $(1+x)^5$.
\end{enumerate}

\bagiansoal{C}{(Advanced --- setara Olimpiade SMA, gabungan antarbab)}
\begin{enumerate}[leftmargin=*]
  \item \textbf{(Pencacahan + Peluang).} Berapa banyak cara memilih $5$ kartu dari $52$ kartu?
  \item \textbf{(Pencacahan + Bilangan).} Berapa banyak bilangan genap $3$ angka berbeda dari $\{1,2,3,4,5\}$?
  \item \textbf{(Binomial).} Tentukan suku ke-$4$ dari ekspansi $(2x+3)^5$.
  \item \textbf{(Permutasi + Restriksi).} Berapa banyak susunan huruf "DUTA" dengan kedua vokal selalu berdampingan?
  \item \textbf{(Kombinasi + Pembuktian).} Buktikan secara kombinatorial bahwa $\binom nk=\binom{n}{n-k}$.
\end{enumerate}

\bagiansoal{D}{(Expert --- setara tahun pertama universitas, sintesis \& pembuktian)}
\begin{enumerate}[leftmargin=*]
  \item \textbf{(Teorema Binomial).} Buktikan $\displaystyle\sum_{k=0}^{n}\binom nk=2^n$.
  \item \textbf{(Identitas Pascal).} Buktikan $\binom nk=\binom{n-1}{k-1}+\binom{n-1}{k}$.
  \item \textbf{(Bintang dan Batang).} Berapa banyak solusi bilangan bulat tak negatif dari $x_1+x_2+x_3=10$?
  \item \textbf{(Inklusi--Eksklusi).} Berapa banyak bilangan dari $1$ sampai $100$ yang habis dibagi $2$ atau $3$?
  \item \textbf{(Binomial + Identitas).} Buktikan $\displaystyle\sum_{k=0}^{n}k\binom nk=n\,2^{n-1}$.
\end{enumerate}

\begin{challenge}
\begin{enumerate}[leftmargin=*,topsep=2pt]
  \item Berapa banyak diagonal pada segi-$n$ ($n\ge3$)?
  \item Pada grid persegi, berapa banyak jalur terpendek dari $(0,0)$ ke $(m,n)$ jika hanya boleh bergerak ke kanan atau ke atas?
  \item Buktikan identitas $\displaystyle\sum_{k=0}^{n}\binom nk^{2}=\binom{2n}{n}$.
\end{enumerate}
\end{challenge}

% ============================================================
\chapter{Peluang}
% ============================================================
\begin{tujuan}
Setelah mempelajari bab ini, peserta didik mampu:
\begin{itemize}[leftmargin=*,topsep=2pt]
  \item menentukan ruang sampel dan kejadian suatu percobaan;
  \item menghitung peluang kejadian, peluang komplemen, dan peluang gabungan;
  \item menghitung peluang bersyarat dan memeriksa kebebasan (independensi) kejadian;
  \item menerapkan teorema Bayes pada masalah sederhana.
\end{itemize}
\textbf{Capaian Pembelajaran (Fase F).} Peserta didik dapat memodelkan ketidakpastian melalui peluang, termasuk peluang bersyarat dan teorema Bayes.
\end{tujuan}

\begin{pemantik}
\begin{itemize}[leftmargin=*,topsep=2pt]
  \item Jika hasil tes penyakit "positif", seberapa yakin Anda benar-benar sakit? Mengapa jawabannya bisa jauh lebih kecil dari yang diduga?
  \item Apa beda "peluang hujan" dengan "peluang hujan \emph{jika} langit mendung"?
  \item Mengapa pada kuis berhadiah, berpindah pilihan setelah satu opsi salah dibuka justru menguntungkan?
\end{itemize}
\end{pemantik}

\begin{studikasus}{Tes Medis dan Intuisi yang Menyesatkan}
Suatu penyakit menjangkiti $1\%$ populasi. Sebuah tes mendeteksi $99\%$ orang sakit dengan benar, tetapi juga keliru memberi hasil positif pada $5\%$ orang sehat. Seseorang dites dan hasilnya \textbf{positif}.
\smallskip

\textbf{Amati dan duga.} Banyak orang menebak peluang ia benar-benar sakit sekitar $99\%$. Padahal karena orang sehat jauh lebih banyak, "positif palsu" menumpuk. Jawaban sebenarnya hanya sekitar $17\%$! Untuk memahami mengapa, kita perlu \emph{peluang bersyarat} dan \emph{teorema Bayes} --- inti bab ini.
\end{studikasus}

\section*{Peta Konsep Bab}
\addcontentsline{toc}{section}{Peta Konsep Bab}
\begin{center}
\begin{tikzpicture}[
  node distance=8mm,
  konsepbox/.style={draw=biru,fill=birumuda,rounded corners,align=center,inner sep=5pt,font=\small\bfseries},
  subbox/.style={draw=hijau,fill=hijaumuda,rounded corners,align=center,inner sep=4pt,font=\footnotesize},
  garis/.style={-Stealth,biru,thick}]
  \node[konsepbox] (akar) {PELUANG};
  \node[subbox,below left=12mm and 18mm of akar] (rs) {Ruang sampel\\\&\ kejadian};
  \node[subbox,below=12mm of akar] (gab) {Komplemen \&\\gabungan};
  \node[subbox,below right=12mm and 18mm of akar] (br) {Bersyarat \&\\independensi};
  \node[subbox,below=10mm of br] (bayes) {Teorema\\Bayes};
  \draw[garis] (akar)--(rs); \draw[garis] (akar)--(gab); \draw[garis] (akar)--(br);
  \draw[garis] (br)--(bayes);
\end{tikzpicture}\\
\small \textbf{Peta konsep.} Dari ruang sampel ke peluang dasar, lalu komplemen/gabungan, hingga peluang bersyarat dan Bayes.
\end{center}

\begin{catatan}
\textbf{Pengingat Materi Kelas 10 (MAE Bab 8 --- Peluang).} Bab ini memperdalam teori peluang yang diperkenalkan di Kelas 10. Tinjau ulang konsep-konsep ini:
\begin{itemize}[leftmargin=*,topsep=2pt]
  \item Percobaan acak, ruang sampel $S$, dan kejadian $A \subseteq S$.
  \item Peluang klasik: $P(A)=\dfrac{n(A)}{n(S)}$ untuk ruang sampel seragam.
  \item Frekuensi relatif sebagai pendekatan peluang empiris.
  \item Sifat dasar: $0\le P(A)\le1$, $P(S)=1$, $P(\emptyset)=0$.
  \item Aturan penjumlahan untuk kejadian saling lepas: $P(A\cup B)=P(A)+P(B)$.
\end{itemize}
\end{catatan}

\section{Ruang Sampel, Kejadian, dan Peluang Dasar}
\begin{definisi}
\textbf{Ruang sampel} $S$ adalah himpunan semua hasil yang mungkin; \textbf{kejadian} $A$ adalah himpunan bagian dari $S$. Untuk ruang sampel berhingga dengan hasil berpeluang sama:
\[
P(A)=\frac{n(A)}{n(S)},\qquad 0\le P(A)\le1.
\]
\end{definisi}

\begin{contoh}{Peluang Dasar}
Sebuah dadu dilempar. Tentukan peluang muncul mata genap.

\jawab $S=\{1,2,3,4,5,6\}$, kejadian $A=\{2,4,6\}$. $P(A)=\dfrac{3}{6}=\dfrac12$.
\end{contoh}

\section{Komplemen dan Peluang Gabungan}
\begin{konsep}{Komplemen dan Gabungan}
\textbf{Komplemen:} $P(A^c)=1-P(A)$ (berguna untuk "minimal satu"). \textbf{Gabungan:}
\[
P(A\cup B)=P(A)+P(B)-P(A\cap B).
\]
Jika $A$ dan $B$ \emph{saling lepas} ($A\cap B=\varnothing$), maka $P(A\cup B)=P(A)+P(B)$.
\end{konsep}

\begin{contoh}{Peluang Gabungan}
Dari satu set $52$ kartu, tentukan peluang terambil kartu King atau kartu hati.

\jawab $P(\text{King})=\tfrac{4}{52}$, $P(\text{hati})=\tfrac{13}{52}$, $P(\text{King hati})=\tfrac{1}{52}$. Maka $P=\tfrac{4+13-1}{52}=\tfrac{16}{52}=\tfrac{4}{13}$.
\end{contoh}

\section{Peluang Bersyarat, Independensi, dan Teorema Bayes}
\begin{konsep}{Peluang Bersyarat dan Independensi}
\textbf{Peluang bersyarat} $A$ jika $B$ diketahui terjadi:
\[
P(A\mid B)=\frac{P(A\cap B)}{P(B)},\quad P(B)>0.
\]
Akibatnya $P(A\cap B)=P(B)\,P(A\mid B)$. Kejadian $A,B$ \textbf{independen} bila $P(A\cap B)=P(A)\,P(B)$, ekuivalen dengan $P(A\mid B)=P(A)$.
\end{konsep}

\begin{center}
\begin{tikzpicture}[
  akar/.style={draw=biru,fill=birumuda,rounded corners,align=center,font=\scriptsize\bfseries,inner sep=5pt},
  cabangL/.style={draw=hijau,fill=hijaumuda,rounded corners,align=center,font=\scriptsize,inner sep=4pt},
  cabangP/.style={draw=oranye,fill=oranyemuda,rounded corners,align=center,font=\scriptsize,inner sep=4pt},
  daun/.style={draw=abu,rounded corners,align=center,font=\scriptsize,inner sep=3pt},
  panah/.style={-Stealth,biru,thick}]
% Root
\node[akar] (R) at (4,3.2) {Populasi};
% Level 1
\node[cabangL] (L) at (1.2,1.6) {Laki-laki};
\node[cabangP] (P) at (6.8,1.6) {Perempuan};
% Level 2 left (Laki-laki)
\node[daun] (LL) at (0,0) {Lulus};
\node[daun] (LT) at (2.4,0) {Tidak Lulus};
% Level 2 right (Perempuan)
\node[daun] (PL) at (5.4,0) {Lulus};
\node[daun] (PT) at (8,0) {Tidak Lulus};
% Edges
\draw[panah] (R)--(L) node[midway,above left,font=\tiny]{$P(\mathrm{lk})$};
\draw[panah] (R)--(P) node[midway,above right,font=\tiny]{$P(\mathrm{pr})$};
\draw[panah] (L)--(LL) node[midway,left,font=\tiny]{$P(L|\mathrm{lk})$};
\draw[panah] (L)--(LT) node[midway,right,font=\tiny]{$P(\bar{L}|\mathrm{lk})$};
\draw[panah] (P)--(PL) node[midway,left,font=\tiny]{$P(L|\mathrm{pr})$};
\draw[panah] (P)--(PT) node[midway,right,font=\tiny]{$P(\bar{L}|\mathrm{pr})$};
% Probability product annotation
\node[font=\tiny,biru,align=left] at (4,-0.65)
  {Peluang suatu jalur $=$ \textbf{perkalian label sepanjang cabang}};
\node[font=\tiny,biru] at (4,-1.1)
  {Contoh: $P(\mathrm{lk}\cap L)=P(\mathrm{lk})\cdot P(L|\mathrm{lk})$};
\end{tikzpicture}\\
\small \textbf{Diagram pohon peluang bersyarat.} Cabang pertama = peluang prior; cabang kedua = peluang bersyarat; ujung pohon = peluang gabungan (dikumpulkan untuk hukum peluang total).
\end{center}

\begin{konsep}{Teorema Bayes}
Jika $P(A),P(B)>0$:
\[
P(A\mid B)=\frac{P(B\mid A)\,P(A)}{P(B)},\qquad
P(B)=P(B\mid A)P(A)+P(B\mid A^c)P(A^c).
\]
Bentuk kedua (hukum peluang total) sering dipakai untuk menghitung penyebut.
\end{konsep}

\begin{contoh}{Teorema Bayes}
Mesin A memproduksi $60\%$ barang (cacat $2\%$), mesin B $40\%$ (cacat $5\%$). Sebuah barang cacat diambil. Berapa peluang ia dari mesin A?

\jawab $P(\text{cacat})=0{,}6\cdot0{,}02+0{,}4\cdot0{,}05=0{,}012+0{,}020=0{,}032$. Maka
$P(A\mid\text{cacat})=\dfrac{0{,}012}{0{,}032}=0{,}375.$
\end{contoh}

\begin{kesalahan}
$P(A\mid B)$ dan $P(B\mid A)$ \emph{tidak} sama (kekeliruan jaksa). "Peluang positif jika sakit" berbeda jauh dari "peluang sakit jika positif" --- selisihnya bergantung pada prevalensi.
\end{kesalahan}

\begin{dunianyata}
Peluang bersyarat dan Bayes adalah mesin di balik filter spam, diagnosis medis, machine learning (klasifikasi naive Bayes), sistem rekomendasi, dan asuransi. Memahami "positif palsu" penting agar tak salah menafsir hasil tes.
\end{dunianyata}

\section*{Ringkasan Bab}
\addcontentsline{toc}{section}{Ringkasan Bab}
\begin{itemize}[leftmargin=*,topsep=2pt]
  \item $P(A)=\dfrac{n(A)}{n(S)}$, $0\le P(A)\le1$.
  \item $P(A^c)=1-P(A)$; $P(A\cup B)=P(A)+P(B)-P(A\cap B)$.
  \item Bersyarat: $P(A\mid B)=\dfrac{P(A\cap B)}{P(B)}$; independen jika $P(A\cap B)=P(A)P(B)$.
  \item Bayes: $P(A\mid B)=\dfrac{P(B\mid A)P(A)}{P(B)}$ dengan $P(B)$ dari hukum peluang total.
\end{itemize}

\begin{refleksi}
\begin{itemize}[leftmargin=*,topsep=2pt]
  \item[\cek] Saya dapat menentukan ruang sampel dan menghitung peluang dasar.
  \item[\cek] Saya dapat menggunakan komplemen dan rumus gabungan.
  \item[\cek] Saya dapat menghitung peluang bersyarat dan menguji independensi.
  \item[\cek] Saya dapat menerapkan teorema Bayes pada masalah nyata.
\end{itemize}
\end{refleksi}

\section*{Soal Akhir Bab \thechapter}
\addcontentsline{toc}{section}{Soal Akhir Bab \thechapter}

\bagiansoal{A}{(Fundamental --- setara SNBT Dasar)}
\begin{enumerate}[leftmargin=*]
  \item Tuliskan ruang sampel pelemparan satu dadu.
  \item Tentukan peluang muncul mata genap pada satu dadu.
  \item Tentukan peluang muncul mata lebih dari $4$ pada satu dadu.
  \item Sebuah koin dilempar. Tentukan peluang muncul gambar.
  \item Dari $52$ kartu, tentukan peluang terambil kartu hati.
  \item Jika $P(A)=0{,}3$, tentukan $P(A^c)$.
  \item Dua koin dilempar. Ada berapa anggota ruang sampelnya?
  \item Dua koin dilempar. Tentukan peluang muncul dua gambar.
  \item Jika $P(A)=0{,}5$, $P(B)=0{,}3$, dan $A,B$ saling lepas, tentukan $P(A\cup B)$.
  \item Kantong berisi $3$ bola merah dan $2$ biru. Tentukan peluang terambil bola merah.
\end{enumerate}

\bagiansoal{B}{(Intermediate --- setara SNBT Sulit)}
\begin{enumerate}[leftmargin=*]
  \item Dua dadu dilempar. Tentukan peluang jumlah mata $=7$.
  \item Diketahui $P(A)=0{,}6$, $P(B)=0{,}5$, $P(A\cap B)=0{,}2$. Tentukan $P(A\cup B)$.
  \item Diketahui $P(A\cap B)=0{,}12$ dan $P(B)=0{,}3$. Tentukan $P(A\mid B)$.
  \item Dari $52$ kartu, tentukan peluang terambil King atau kartu hati.
  \item Kotak berisi $5$ bola merah dan $3$ putih. Diambil $2$ bola tanpa pengembalian. Tentukan peluang keduanya merah.
\end{enumerate}

\bagiansoal{C}{(Advanced --- setara Olimpiade SMA, gabungan antarbab)}
\begin{enumerate}[leftmargin=*]
  \item \textbf{(Peluang + Pencacahan).} Dari $10$ orang dipilih $3$ secara acak. Tentukan peluang Andi (salah seorang) terpilih.
  \item \textbf{(Peluang + Independensi).} $A,B$ independen dengan $P(A)=0{,}4$, $P(B)=0{,}5$. Tentukan $P(A\cap B)$ dan $P(A\cup B)$.
  \item \textbf{(Peluang Bersyarat).} Sebuah dadu dilempar. Diketahui hasilnya ganjil. Tentukan peluang hasilnya $3$.
  \item \textbf{(Peluang + Komplemen).} Tiga koin dilempar. Tentukan peluang muncul minimal satu gambar.
  \item \textbf{(Bayes).} Mesin A ($60\%$, cacat $2\%$) dan mesin B ($40\%$, cacat $5\%$). Diambil barang cacat. Tentukan peluang berasal dari mesin A.
\end{enumerate}

\bagiansoal{D}{(Expert --- setara tahun pertama universitas, sintesis \& pembuktian)}
\begin{enumerate}[leftmargin=*]
  \item \textbf{(Bayes --- Tes Medis).} Penyakit menjangkiti $1\%$ populasi; tes benar $99\%$ pada yang sakit dan positif palsu $5\%$ pada yang sehat. Hitung peluang seseorang benar sakit jika hasil tesnya positif.
  \item \textbf{(Pembuktian).} Dari aksioma peluang, buktikan $P(A^c)=1-P(A)$.
  \item \textbf{(Independensi).} Buktikan bahwa jika $A$ dan $B$ independen, maka $A$ dan $B^c$ juga independen.
  \item \textbf{(Nilai Harapan).} Sebuah permainan membayar (dalam ribu rupiah) sebesar mata dadu yang muncul. Hitung nilai harapan pembayaran.
  \item \textbf{(Hukum Peluang Total).} Kotak I berisi $2$ merah, $3$ putih; Kotak II berisi $4$ merah, $1$ putih. Sebuah kotak dipilih acak (peluang sama), lalu diambil satu bola. Tentukan peluang terambil bola merah.
\end{enumerate}

\begin{challenge}
\begin{enumerate}[leftmargin=*,topsep=2pt]
  \item \textbf{(Monty Hall).} Pada kuis tiga pintu (satu berhadiah), Anda memilih satu pintu; pembawa acara membuka satu pintu lain berisi kambing. Apakah sebaiknya berpindah? Tunjukkan peluang menangnya.
  \item \textbf{(Masalah Ulang Tahun).} Susun ekspresi peluang bahwa di antara $23$ orang terdapat (sedikitnya) dua yang berulang tahun pada tanggal sama, dan jelaskan mengapa nilainya melebihi $\tfrac12$.
  \item \textbf{(Tepat Dua).} Tiga penembak mengenai sasaran secara independen dengan peluang $0{,}6$, $0{,}7$, dan $0{,}8$. Tentukan peluang tepat dua di antaranya mengenai sasaran.
\end{enumerate}
\end{challenge}

% ============================================================
\chapter{Irisan Kerucut}
% ============================================================
\begin{tujuan}
Setelah mempelajari bab ini, peserta didik mampu:
\begin{itemize}[leftmargin=*,topsep=2pt]
  \item menjelaskan definisi geometris parabola, elips, dan hiperbola (fokus, direktriks, eksentrisitas);
  \item menuliskan persamaan baku tiap irisan kerucut serta menggambar grafiknya;
  \item menentukan unsur-unsur: puncak, fokus, sumbu, direktriks, asimtot, dan eksentrisitas;
  \item menangani bentuk yang bergeser (pusat bukan di $O$) lewat melengkapkan kuadrat;
  \item menerapkan irisan kerucut pada lintasan orbit dan teknologi (antena, satelit).
\end{itemize}
\textbf{Capaian Pembelajaran (Fase F+).} Peserta didik dapat menganalisis irisan kerucut secara aljabar dan geometris serta menggunakannya untuk memodelkan lintasan dan bentuk nyata.
\end{tujuan}

\begin{pemantik}
\begin{itemize}[leftmargin=*,topsep=2pt]
  \item Mengapa orbit planet berbentuk elips, dan apa peran Matahari di dalamnya?
  \item Apa yang sama dari parabola, elips, dan hiperbola sehingga ketiganya disebut "irisan kerucut"?
  \item Mengapa antena parabola memantulkan semua sinyal ke satu titik?
\end{itemize}
\end{pemantik}

\begin{studikasus}{Orbit dan Antena}
Ketiga irisan kerucut muncul saat sebuah bidang memotong kerucut pada sudut berbeda. Mereka punya satu definisi pemersatu: \emph{himpunan titik dengan rasio jarak ke fokus dan ke direktriks tetap}, yaitu \textbf{eksentrisitas} $e$.
\begin{center}
\begin{tabular}{c|c|c}
\toprule
Irisan kerucut & Eksentrisitas $e$ & Contoh nyata \\
\midrule
Lingkaran & $e=0$ & roda \\
Elips & $0<e<1$ & orbit planet \\
Parabola & $e=1$ & antena, lintasan proyektil \\
Hiperbola & $e>1$ & orbit komet lepas, navigasi \\
\bottomrule
\end{tabular}
\end{center}
\textbf{Amati dan duga.} Bagaimana satu bilangan $e$ menentukan "bentuk" sebuah kurva, dari lingkaran sempurna sampai hiperbola terbuka? Inilah benang merah seluruh bab.
\end{studikasus}

\smallskip\noindent\textbf{Mengapa disebut "irisan kerucut".} Keempatnya muncul saat sebuah \emph{bidang} mengiris kerucut ganda pada kemiringan berbeda:
\begin{center}
\begin{tabular}{cccc}
% Lingkaran: potong mendatar
\begin{tikzpicture}[scale=0.5,line join=round]
\draw[gray] (0,0)--(1,2.5) (0,0)--(-1,2.5) (0,0)--(1,-2.5) (0,0)--(-1,-2.5);
\draw[gray] (0,2.5) ellipse (1 and 0.22); \draw[gray] (0,-2.5) ellipse (1 and 0.22);
\draw[merah,very thick] (-0.72,1.8)--(0.72,1.8);
\end{tikzpicture}
&
% Elips: potong miring (satu nappe)
\begin{tikzpicture}[scale=0.5,line join=round]
\draw[gray] (0,0)--(1,2.5) (0,0)--(-1,2.5) (0,0)--(1,-2.5) (0,0)--(-1,-2.5);
\draw[gray] (0,2.5) ellipse (1 and 0.22); \draw[gray] (0,-2.5) ellipse (1 and 0.22);
\draw[merah,very thick] (-0.55,1.2)--(0.8,2.3);
\end{tikzpicture}
&
% Parabola: sejajar garis pelukis
\begin{tikzpicture}[scale=0.5,line join=round]
\draw[gray] (0,0)--(1,2.5) (0,0)--(-1,2.5) (0,0)--(1,-2.5) (0,0)--(-1,-2.5);
\draw[gray] (0,2.5) ellipse (1 and 0.22); \draw[gray] (0,-2.5) ellipse (1 and 0.22);
\draw[merah,very thick] (-0.45,0.3)--(0.42,2.5);
\end{tikzpicture}
&
% Hiperbola: potong kedua nappe
\begin{tikzpicture}[scale=0.5,line join=round]
\draw[gray] (0,0)--(1,2.5) (0,0)--(-1,2.5) (0,0)--(1,-2.5) (0,0)--(-1,-2.5);
\draw[gray] (0,2.5) ellipse (1 and 0.22); \draw[gray] (0,-2.5) ellipse (1 and 0.22);
\draw[merah,very thick] (0.5,-2.5)--(0.5,2.5);
\end{tikzpicture}
\\[-1pt]
\footnotesize \textbf{Lingkaran} & \footnotesize \textbf{Elips} & \footnotesize \textbf{Parabola} & \footnotesize \textbf{Hiperbola}\\
\footnotesize iris mendatar & \footnotesize iris miring & \footnotesize sejajar pelukis & \footnotesize iris kedua kerucut\\
\end{tabular}
\end{center}

\section*{Peta Konsep Bab}
\addcontentsline{toc}{section}{Peta Konsep Bab}
\begin{center}
\begin{tikzpicture}[
  node distance=8mm,
  konsepbox/.style={draw=biru,fill=birumuda,rounded corners,align=center,inner sep=5pt,font=\small\bfseries},
  subbox/.style={draw=hijau,fill=hijaumuda,rounded corners,align=center,inner sep=4pt,font=\footnotesize},
  garis/.style={-Stealth,biru,thick}]
  \node[konsepbox] (akar) {IRISAN KERUCUT};
  \node[subbox,below left=12mm and 14mm of akar] (par) {Parabola\\$e=1$};
  \node[subbox,below=12mm of akar] (eli) {Elips\\$0<e<1$};
  \node[subbox,below right=12mm and 14mm of akar] (hip) {Hiperbola\\$e>1$};
  \node[subbox,below=10mm of eli] (uns) {Fokus, direktriks,\\eksentrisitas};
  \draw[garis] (akar)--(par); \draw[garis] (akar)--(eli); \draw[garis] (akar)--(hip);
  \draw[garis] (eli)--(uns);
\end{tikzpicture}\\
\small \textbf{Peta konsep.} Parabola, elips, dan hiperbola dibedakan oleh eksentrisitas, namun berbagi unsur yang sama: fokus dan direktriks.
\end{center}

\section{Bentuk Baku dan Unsur-Unsurnya}
\begin{konsep}{Bentuk Baku (berpusat di $O$)}
\[
\textbf{Parabola: } y^2 = 4px, \qquad
\textbf{Elips: } \frac{x^2}{a^2}+\frac{y^2}{b^2}=1, \qquad
\textbf{Hiperbola: } \frac{x^2}{a^2}-\frac{y^2}{b^2}=1.
\]
\textbf{Parabola} $y^2=4px$: fokus $(p,0)$, direktriks $x=-p$. \textbf{Elips} ($a>b$): fokus $(\pm c,0)$ dengan $c^2=a^2-b^2$, $e=\tfrac ca$. \textbf{Hiperbola:} fokus $(\pm c,0)$ dengan $c^2=a^2+b^2$, asimtot $y=\pm\tfrac ba x$, $e=\tfrac ca$.
\end{konsep}

\begin{center}
\begin{tikzpicture}
\begin{axis}[name=p,axis lines=middle,width=4.7cm,height=4.2cm,
  title=\small Parabola,xtick=\empty,ytick=\empty,xmin=-0.5,xmax=4,ymin=-3,ymax=3]
\addplot[very thick,biru,domain=-2.8:2.8,samples=50]({x*x/2},{x});
\end{axis}
\begin{axis}[name=e,axis lines=middle,width=4.7cm,height=4.2cm,
  at={($(p.east)+(1cm,0)$)},anchor=west,
  title=\small Elips,xtick=\empty,ytick=\empty,xmin=-3,xmax=3,ymin=-2.5,ymax=2.5]
\addplot[very thick,hijau,domain=0:360,samples=80]({2.2*cos(x)},{1.6*sin(x)});
\end{axis}
\begin{axis}[name=h,axis lines=middle,width=4.7cm,height=4.2cm,
  at={($(e.east)+(1cm,0)$)},anchor=west,
  title=\small Hiperbola,xtick=\empty,ytick=\empty,xmin=-4,xmax=4,ymin=-3,ymax=3]
\addplot[very thick,oranye,domain=-1.6:1.6,samples=40]({1.2*cosh(x)},{1.6*sinh(x)});
\addplot[very thick,oranye,domain=-1.6:1.6,samples=40]({-1.2*cosh(x)},{1.6*sinh(x)});
\end{axis}
\end{tikzpicture}
\end{center}

\smallskip\noindent\textbf{Anatomi parabola} $y^2=4px$ --- definisi: jarak $P$ ke \emph{fokus} $=$ jarak $P$ ke \emph{direktriks}.
\begin{center}
\begin{tikzpicture}[scale=0.95,>=Stealth,line join=round]
\draw[very thick,biru,smooth,samples=60,domain=-2.7:2.7] plot ({(\x*\x)/4},{\x});
\draw[gray,->](-1.6,0)--(3,0) node[font=\scriptsize,right]{sumbu};
\draw[oranye,very thick](-1,-2.7)--(-1,2.7);
\node[oranye,font=\scriptsize,align=center] at (-1.0,-3.05){direktriks $x=-p$};
\fill (1,0) circle (1.8pt) node[font=\scriptsize,below right]{$F(p,0)$};
\fill (0,0) circle (1.4pt) node[font=\scriptsize,below left]{puncak};
\coordinate (P) at (1.5625,2.5);
\fill[merah] (P) circle (1.5pt) node[font=\scriptsize,right]{$P$};
\draw[hijau!60!black,thick] (P)--(1,0);
\draw[hijau!60!black,thick,dashed] (P)--(-1,2.5);
\fill[merah] (-1,2.5) circle (1pt);
\draw[ungu,very thick] (1,-2)--(1,2);
\node[ungu,font=\scriptsize,right] at (1,-1.55){latus rectum $=4p$};
\end{tikzpicture}\\
\small Jarak hijau (ke fokus) $=$ jarak hijau putus-putus (ke direktriks). Tali busur tegak melalui fokus (\emph{latus rectum}) panjangnya $4p$.
\end{center}

\smallskip\noindent\textbf{Anatomi elips} $\dfrac{x^2}{a^2}+\dfrac{y^2}{b^2}=1$ --- definisi: $PF_1+PF_2=2a$.
\begin{center}
\begin{tikzpicture}[scale=0.85,>=Stealth,line join=round]
\draw[very thick,hijau,smooth,samples=90,domain=0:360] plot ({3*cos(\x)},{2*sin(\x)});
\draw[gray,->](-3.7,0)--(3.7,0); \draw[gray,->](0,-2.7)--(0,2.7);
\fill (2.236,0) circle (1.6pt) node[font=\scriptsize,below right]{$F_2$};
\fill (-2.236,0) circle (1.6pt) node[font=\scriptsize,below left]{$F_1$};
\fill (0,0) circle (1.2pt) node[font=\scriptsize,below left]{$O$};
\fill (3,0) circle (1pt) node[font=\scriptsize,below]{$(a,0)$};
\fill (0,2) circle (1pt) node[font=\scriptsize,above right]{$(0,b)$};
\coordinate (P) at ({3*cos(115)},{2*sin(115)});
\fill[merah] (P) circle (1.4pt) node[font=\scriptsize,above]{$P$};
\draw[oranye,thick] (-2.236,0)--(P)--(2.236,0);
\node[oranye,font=\scriptsize] at (-1.9,1.75){$PF_1{+}PF_2{=}2a$};
\draw[gray,dashed] (0,2)--(2.236,0);
\node[font=\scriptsize] at (1.4,1.2){$a$};
\node[font=\scriptsize,left] at (0,1){$b$};
\node[font=\scriptsize,below] at (1.12,0){$c$};
\end{tikzpicture}\\
\small Jumlah jarak ke dua fokus tetap $=2a$. Segitiga siku-siku memberi $a^2=b^2+c^2$, sehingga $c=\sqrt{a^2-b^2}$ dan eksentrisitas $e=\tfrac ca$.
\end{center}

\smallskip\noindent\textbf{Anatomi hiperbola} $\dfrac{x^2}{a^2}-\dfrac{y^2}{b^2}=1$ --- definisi: $|PF_1-PF_2|=2a$.
\begin{center}
\begin{tikzpicture}[scale=0.8,>=Stealth,line join=round]
\draw[gray,dashed](-3.3,-2.475)--(3.3,2.475);
\draw[gray,dashed](-3.3,2.475)--(3.3,-2.475);
\draw[very thick,oranye,smooth,samples=60,domain=-58:58] plot ({2/cos(\x)},{1.5*tan(\x)});
\draw[very thick,oranye,smooth,samples=60,domain=-58:58] plot ({-2/cos(\x)},{1.5*tan(\x)});
\draw[gray,->](-4,0)--(4,0); \draw[gray,->](0,-3)--(0,3);
\fill (2.5,0) circle (1.6pt) node[font=\scriptsize,below right]{$F_2$};
\fill (-2.5,0) circle (1.6pt) node[font=\scriptsize,below left]{$F_1$};
\fill (2,0) circle (1pt) node[font=\scriptsize,above left]{$(a,0)$};
\node[gray,font=\scriptsize,fill=white,inner sep=1pt,anchor=south west] at (2.0,1.65){asimtot $y{=}\pm\tfrac{b}{a}x$};
\coordinate (P) at (2.828,1.5);
\fill[merah] (P) circle (1.4pt) node[font=\scriptsize,right]{$P$};
\draw[hijau!60!black,thick] (-2.5,0)--(P)--(2.5,0);
\node[hijau!60!black,font=\scriptsize,fill=white,inner sep=1pt] at (-0.4,1.1){$|PF_1{-}PF_2|{=}2a$};
\end{tikzpicture}\\
\small Selisih jarak ke dua fokus tetap $=2a$. Kurva mendekati \emph{asimtot} $y=\pm\tfrac ba x$; di sini $c=\sqrt{a^2+b^2}$.
\end{center}

\begin{contoh}{Unsur Elips}
Tentukan puncak, fokus, dan eksentrisitas dari $\dfrac{x^2}{25}+\dfrac{y^2}{9}=1$.

\jawab $a^2=25\Rightarrow a=5$, $b^2=9\Rightarrow b=3$. Puncak $(\pm5,0)$ dan $(0,\pm3)$.
$c=\sqrt{25-9}=4$, fokus $(\pm4,0)$, eksentrisitas $e=\tfrac45$.
\end{contoh}

\begin{contoh}{Fokus Parabola}
Tentukan fokus dan direktriks parabola $y^2=12x$.

\jawab $4p=12\Rightarrow p=3$. Fokus $(3,0)$, direktriks $x=-3$.
\end{contoh}

\section{Bentuk Bergeser (Pusat Bukan di $O$)}
\begin{konsep}{Translasi Irisan Kerucut}
Jika pusat/puncak berada di $(h,k)$, ganti $x\to x-h$ dan $y\to y-k$. Misalnya elips berpusat $(h,k)$:
\[
\frac{(x-h)^2}{a^2}+\frac{(y-k)^2}{b^2}=1.
\]
Bentuk umum diolah ke bentuk baku dengan \emph{melengkapkan kuadrat}.
\end{konsep}

\begin{contoh}{Melengkapkan Kuadrat}
Ubah $x^2+4y^2-4x+8y+4=0$ ke bentuk baku elips.

\jawab Kelompokkan: $(x^2-4x)+4(y^2+2y)+4=0\Rightarrow(x-2)^2-4+4[(y+1)^2-1]+4=0$
$\Rightarrow(x-2)^2+4(y+1)^2=4\Rightarrow\dfrac{(x-2)^2}{4}+(y+1)^2=1.$ Pusat $(2,-1)$.
\end{contoh}

\begin{kesalahan}
Pada hiperbola, $c^2=a^2+b^2$ (dijumlah), sedangkan pada elips $c^2=a^2-b^2$ (dikurang). Tertukar tanda ini membuat letak fokus salah total.
\end{kesalahan}

\smallskip\noindent\textbf{Sifat pemantul parabola.} Semua sinar yang datang sejajar sumbu dipantulkan menuju \emph{satu titik}: fokus. Inilah prinsip antena parabola, teleskop, dan lampu sorot.
\begin{center}
\begin{tikzpicture}[scale=0.95,>=Stealth,line join=round]
\draw[very thick,biru,smooth,samples=60,domain=-2.7:2.7] plot ({(\x*\x)/4},{\x});
\draw[gray,->](-0.4,0)--(3.3,0) node[font=\scriptsize,right]{sumbu};
\fill (1,0) circle (1.8pt) node[font=\scriptsize,below right]{$F$};
\foreach \y in {0.9,1.7,2.4}{
  \pgfmathsetmacro\xx{(\y*\y)/4}
  \draw[oranye,thick,->] (3.15,\y)--(\xx,\y);
  \draw[hijau!60!black,thick,->] (\xx,\y)--(1,0);
}
\end{tikzpicture}\\
\small Sinar sejajar (oranye) memantul pada parabola dan berkumpul di fokus $F$ (hijau).
\end{center}

\begin{dunianyata}
Hukum Kepler menyatakan orbit planet berupa \textbf{elips} dengan Matahari di salah satu fokus. Antena dan reflektor lampu sorot berbentuk \textbf{parabola} karena memusatkan sinar sejajar ke fokus. Sistem navigasi lama (LORAN) memanfaatkan \textbf{hiperbola} sebagai tempat kedudukan selisih jarak tetap.
\end{dunianyata}

\section{Geometri Analitik Lanjutan: Sumbu Radikal dan Berkas Lingkaran}
\begin{konsep}{Kuasa Titik dan Sumbu Radikal}
Untuk lingkaran $L:\ x^2+y^2+Dx+Ey+F=0$, \textbf{kuasa} titik $P(x_0,y_0)$ adalah $x_0^2+y_0^2+Dx_0+Ey_0+F$. \textbf{Sumbu radikal} dua lingkaran $L_1=0$ dan $L_2=0$ (keduanya berkoefisien $x^2,y^2$ sama, yaitu $1$) adalah tempat kedudukan titik berkuasa sama, diperoleh dengan \emph{mengurangkan} kedua persamaan:
\[
L_1-L_2=0 \quad(\text{persamaan linear}).
\]
Sumbu radikal tegak lurus garis hubung pusat; bila kedua lingkaran berpotongan, sumbu radikal adalah \textbf{tali busur persekutuan} (garis melalui kedua titik potong).
\end{konsep}

\begin{konsep}{Berkas (Pencil) Lingkaran}
Semua lingkaran yang melalui titik potong $L_1=0$ dan $L_2=0$ dapat ditulis sebagai \textbf{berkas}
\[
L_1+\lambda L_2=0,\qquad \lambda\in\mathbb{R},
\]
dengan $\lambda=-1$ memberi sumbu radikalnya. \textbf{Irisan dua kerucut} (mis. garis$\cap$lingkaran, atau lingkaran$\cap$elips) dicari dengan menyelesaikan sistem persamaannya; secara umum dua kerucut berpotongan di paling banyak \emph{empat} titik.
\end{konsep}

\begin{contoh}{Sumbu Radikal dan Irisan Dua Lingkaran}
Tentukan sumbu radikal dan titik potong $L_1:\ x^2+y^2=25$ dan $L_2:\ x^2+y^2-12x+11=0$.

\jawab Kurangkan: $L_1-L_2:\ (x^2+y^2-25)-(x^2+y^2-12x+11)=12x-36=0\Rightarrow x=3$ (sumbu radikal / tali busur persekutuan). Substitusi $x=3$ ke $L_1$: $9+y^2=25\Rightarrow y=\pm4$. Titik potong $(3,4)$ dan $(3,-4)$.
\end{contoh}

\begin{kesalahan}
Sebelum mengurangkan dua lingkaran untuk mendapat sumbu radikal, pastikan koefisien $x^2$ dan $y^2$ \emph{sama} (normalkan menjadi $1$). Jika tidak, hasil pengurangan bukan persamaan garis yang benar.
\end{kesalahan}

\begin{latihan}
(1) Tentukan sumbu radikal $x^2+y^2-4=0$ dan $x^2+y^2-2x-3=0$.\quad (2) Tentukan titik potong $x^2+y^2=20$ dan $x^2+y^2-6x-2y=0$.\quad (3) Tuliskan berkas lingkaran yang melalui irisan $x^2+y^2-4=0$ dan $x^2+y^2-2x=0$.
\end{latihan}

\begin{konsep}{Bahasa Alternatif yang Sering Mengecoh di Soal}
\begin{itemize}[leftmargin=*,topsep=2pt]
  \item \textbf{Puncak/verteks} (titik pada kurva) berbeda dari \textbf{pusat} (titik tengah elips/hiperbola) dan \textbf{fokus}. Parabola tak punya pusat.
  \item \textbf{Sumbu mayor} $=2a$, \textbf{sumbu minor} $=2b$ (panjang \emph{penuh}); $a$ dan $b$ adalah \emph{setengah}-sumbu. Jangan tertukar.
  \item \textbf{$c^2=a^2-b^2$ (elips)} vs \textbf{$c^2=a^2+b^2$ (hiperbola)}. Pada elips $a>b$ selalu, dan $a$ ada di bawah suku yang lebih besar.
  \item \textbf{Eksentrisitas} $e=\tfrac ca$: lingkaran $0$, elips $0<e<1$, parabola $1$, hiperbola $>1$.
  \item \textbf{$4p$} pada $y^2=4px$ adalah \emph{empat kali} jarak fokus--puncak; fokus di $(p,0)$, bukan $(4p,0)$.
  \item \textbf{Direktriks} = garis (bukan titik); \textbf{asimtot} hanya milik hiperbola (bukan elips/parabola).
  \item \textbf{"Latus rectum"} = tali busur melalui fokus tegak lurus sumbu (elips/hiperbola: $\tfrac{2b^2}{a}$; parabola: $4p$).
  \item \textbf{Sumbu mayor pada sumbu-$y$}: jika penyebut terbesar ada di bawah $y^2$, maka elips "berdiri" --- fokus pada sumbu-$y$.
\end{itemize}
\end{konsep}

\section*{Ringkasan Bab}
\addcontentsline{toc}{section}{Ringkasan Bab}
\begin{itemize}[leftmargin=*,topsep=2pt]
  \item Definisi pemersatu lewat \textbf{eksentrisitas} $e$: lingkaran $0$, elips $0$--$1$, parabola $1$, hiperbola $>1$.
  \item \textbf{Parabola} $y^2=4px$: fokus $(p,0)$, direktriks $x=-p$.
  \item \textbf{Elips:} $c^2=a^2-b^2$, fokus $(\pm c,0)$, $e=\tfrac ca$.
  \item \textbf{Hiperbola:} $c^2=a^2+b^2$, asimtot $y=\pm\tfrac ba x$.
  \item Bentuk bergeser ditangani dengan melengkapkan kuadrat.
  \item \textbf{Sumbu radikal} $=L_1-L_2$ (tali busur persekutuan); \textbf{berkas lingkaran} $L_1+\lambda L_2=0$; dua kerucut berpotongan $\le4$ titik.
\end{itemize}

\begin{refleksi}
\begin{itemize}[leftmargin=*,topsep=2pt]
  \item[\cek] Saya dapat menjelaskan fokus, direktriks, dan eksentrisitas.
  \item[\cek] Saya dapat menuliskan dan menggambar persamaan baku tiap irisan kerucut.
  \item[\cek] Saya dapat menentukan unsur-unsurnya.
  \item[\cek] Saya dapat mengolah bentuk umum ke bentuk baku.
\end{itemize}
\end{refleksi}

\section*{Soal Akhir Bab \thechapter}
\addcontentsline{toc}{section}{Soal Akhir Bab \thechapter}

\bagiansoal{A}{(Fundamental --- setara SNBT Dasar)}
\begin{enumerate}[leftmargin=*]
  \item Tentukan fokus parabola $y^2=12x$.
  \item Tentukan puncak elips $\dfrac{x^2}{25}+\dfrac{y^2}{9}=1$.
  \item Tentukan asimtot hiperbola $\dfrac{x^2}{16}-\dfrac{y^2}{9}=1$.
  \item Tentukan persamaan parabola berpuncak $O$, fokus $(2,0)$.
  \item Sebutkan jenis irisan kerucut dengan eksentrisitas $e=1$.
\end{enumerate}

\bagiansoal{B}{(Intermediate --- setara SNBT Sulit)}
\begin{enumerate}[leftmargin=*]
  \item Tentukan fokus dan eksentrisitas elips $\dfrac{x^2}{169}+\dfrac{y^2}{144}=1$.
  \item Tentukan fokus hiperbola $\dfrac{x^2}{9}-\dfrac{y^2}{16}=1$.
  \item Tentukan persamaan elips berpusat $O$, puncak $(\pm5,0)$, fokus $(\pm3,0)$.
  \item Ubah $x^2+4y^2-4x+8y+4=0$ ke bentuk baku dan tentukan pusatnya.
  \item Tentukan puncak dan fokus parabola $(y-1)^2=8(x-2)$.
  \item Tentukan persamaan direktriks parabola $y^2=-16x$.
\end{enumerate}

\bagiansoal{C}{(Advanced --- setara Olimpiade SMA, gabungan antarbab)}
\begin{enumerate}[leftmargin=*]
  \item \textbf{(Elips + Definisi).} Buktikan bahwa jumlah jarak titik pada elips $\tfrac{x^2}{25}+\tfrac{y^2}{9}=1$ ke kedua fokus selalu $10$.
  \item \textbf{(Irisan Kerucut + Garis).} Tentukan titik potong garis $y=x+1$ dengan parabola $y^2=4x$.
  \item \textbf{(Parabola + Garis Singgung).} Tentukan persamaan garis singgung parabola $y^2=8x$ di titik $(2,4)$.
  \item \textbf{(Eksentrisitas).} Sebuah elips memiliki sumbu mayor $10$ dan eksentrisitas $0{,}6$. Tentukan persamaannya (pusat $O$, sumbu mayor pada sumbu-$x$).
  \item \textbf{(Hiperbola + Asimtot).} Tentukan persamaan hiperbola berpusat $O$ dengan asimtot $y=\pm\tfrac34 x$ dan melalui $(4,0)$.
  \item \textbf{(Irisan Kerucut + Astronomi).} Orbit sebuah planet berbentuk elips dengan jarak terdekat (perihelion) $a-c$ dan terjauh (aphelion) $a+c$ dari Matahari (di salah satu fokus). Jika perihelion $2$ SA dan aphelion $8$ SA, tentukan $a$, $c$, dan eksentrisitasnya.
\end{enumerate}

\bagiansoal{D}{(Expert --- setara tahun pertama universitas, sintesis \& pembuktian)}
\begin{enumerate}[leftmargin=*]
  \item \textbf{(Pembuktian).} Turunkan persamaan elips $\tfrac{x^2}{a^2}+\tfrac{y^2}{b^2}=1$ dari definisi "jumlah jarak ke dua fokus $(\pm c,0)$ sama dengan $2a$".
  \item \textbf{(Sifat Reflektif Parabola).} Jelaskan dan buktikan secara garis besar mengapa sinar sejajar sumbu parabola dipantulkan menuju fokus.
  \item \textbf{(Irisan Kerucut + Bilangan Kompleks/Trigonometri).} Tunjukkan bahwa parametrisasi $(a\cos t,\ b\sin t)$ memenuhi persamaan elips untuk setiap $t$, lalu tentukan luas elips memakai parametrisasi ini (boleh menyebut hasil integral).
  \item \textbf{(Eksentrisitas + Definisi fokus--direktriks).} Buktikan bahwa himpunan titik dengan rasio jarak ke fokus dan ke direktriks sama dengan $e$ menghasilkan elips bila $0<e<1$.
  \item \textbf{(Diskriminan Irisan Kerucut).} Untuk $Ax^2+Bxy+Cy^2+Dx+Ey+F=0$, jelaskan bagaimana tanda $B^2-4AC$ menentukan jenis irisan kerucut.
\end{enumerate}

\begin{challenge}
\begin{enumerate}[leftmargin=*,topsep=2pt]
  \item Sebuah komet bergerak pada lintasan parabola dengan Matahari di fokus. Pada jarak terdekat $r_0$ dari Matahari, tentukan persamaan lintasannya (pusat di puncak) dan jelaskan mengapa komet tak pernah kembali.
  \item Tentukan panjang \emph{latus rectum} (tali busur melalui fokus tegak lurus sumbu) elips $\tfrac{x^2}{a^2}+\tfrac{y^2}{b^2}=1$, dan buktikan rumusnya adalah $\tfrac{2b^2}{a}$.
\end{enumerate}
\end{challenge}

% ============================================================
\chapter{Limit Fungsi}
% ============================================================
\begin{tujuan}
Setelah mempelajari bab ini, peserta didik mampu:
\begin{itemize}[leftmargin=*,topsep=2pt]
  \item memahami konsep limit secara intuitif beserta limit kiri dan kanan;
  \item menghitung limit di suatu titik, limit tak hingga, dan limit di tak hingga;
  \item mengenali bentuk tak tentu dan memilih teknik penyelesaian yang tepat;
  \item menggunakan faktorisasi, perkalian sekawan, substitusi, dan pembagian suku tertinggi;
  \item menghubungkan limit dengan kontinuitas fungsi;
  \item menggunakan aturan L'Hôpital untuk bentuk tak tentu $\tfrac00$ dan $\tfrac\infty\infty$.
\end{itemize}
\textbf{Capaian Pembelajaran (Fase F+).} Peserta didik dapat memahami dan menghitung limit sebagai fondasi kalkulus.
\end{tujuan}

\begin{pemantik}
\begin{itemize}[leftmargin=*,topsep=2pt]
  \item Apa nilai $\dfrac{x^2-1}{x-1}$ saat $x$ \emph{sangat dekat} ke $1$, padahal di $x=1$ ia berbentuk $\tfrac00$?
  \item Bagaimana sebuah fungsi bisa "menuju" suatu nilai tanpa pernah benar-benar mencapainya?
  \item Apa yang terjadi pada $\dfrac{1}{x}$ ketika $x$ membesar tanpa batas?
\end{itemize}
\end{pemantik}

\begin{studikasus}{Kecepatan Sesaat}
Bagaimana mengukur kecepatan sebuah mobil \emph{pada satu saat persis}, bukan rata-rata? Kita hitung kecepatan rata-rata pada selang waktu yang makin pendek --- $1$ detik, $0{,}1$ detik, $0{,}01$ detik --- dan melihat ke arah mana angkanya \emph{menuju}.
\begin{center}
\begin{tabular}{c|ccccc}
\toprule
Selang $h$ (s) & 1 & 0{,}5 & 0{,}1 & 0{,}01 & $\to0$ \\
$\frac{s(2+h)-s(2)}{h}$ & 5 & 4{,}5 & 4{,}1 & 4{,}01 & $\to4$ \\
\bottomrule
\end{tabular}
\end{center}
\textbf{Amati dan duga.} Walau di $h=0$ rumusnya berbentuk $\tfrac00$, angkanya jelas menuju $4$. Gagasan "nilai yang dituju" inilah \textbf{limit} --- jantung dari kalkulus.
\end{studikasus}

\section*{Peta Konsep Bab}
\addcontentsline{toc}{section}{Peta Konsep Bab}
\begin{center}
\begin{tikzpicture}[
  node distance=8mm,
  konsepbox/.style={draw=biru,fill=birumuda,rounded corners,align=center,inner sep=5pt,font=\small\bfseries},
  subbox/.style={draw=hijau,fill=hijaumuda,rounded corners,align=center,inner sep=4pt,font=\footnotesize},
  garis/.style={-Stealth,biru,thick}]
  \node[konsepbox] (akar) {LIMIT FUNGSI};
  \node[subbox,below left=12mm and 16mm of akar] (sat) {Limit kiri,\\kanan};
  \node[subbox,below=12mm of akar] (tt) {Bentuk\\tak tentu};
  \node[subbox,below right=12mm and 16mm of akar] (inf) {Limit di\\tak hingga};
  \node[subbox,below=10mm of tt] (tek) {Faktorisasi, sekawan,\\suku tertinggi};
  \draw[garis] (akar)--(sat); \draw[garis] (akar)--(tt); \draw[garis] (akar)--(inf);
  \draw[garis] (tt)--(tek);
\end{tikzpicture}\\
\small \textbf{Peta konsep.} Dari konsep limit (termasuk kiri/kanan dan tak hingga) muncul bentuk tak tentu yang diselesaikan dengan beragam teknik.
\end{center}

\begin{catatan}
\textbf{Pengingat Materi Kelas 10 (MAE Bab 1 \&\ Bab 6).} Limit, turunan, dan integral banyak bekerja dengan fungsi-fungsi berikut yang sudah dipelajari di Kelas 10. Tinjau ulang jika perlu:
\begin{itemize}[leftmargin=*,topsep=2pt]
  \item \textbf{Eksponen \&\ Logaritma (MAE Bab 1):} sifat $a^m\cdot a^n=a^{m+n}$, $\log_a(xy)=\log_a x+\log_a y$, grafik $f(x)=a^x$ dan $f(x)=\log_a x$.
  \item \textbf{Fungsi Kuadrat (MAE Bab 6):} faktorisasi $ax^2+bx+c$, akar-akar, sumbu simetri $x=-\tfrac{b}{2a}$, dan sketsa parabola.
  \item Manipulasi aljabar: pemfaktoran selisih kuadrat/kuadrat sempurna, penyederhanaan pecahan, merasionalkan penyebut bentuk akar.
\end{itemize}
\end{catatan}

\section{Konsep Limit dan Limit Sepihak}
\begin{konsep}{Pengertian Limit}
$\displaystyle\lim_{x\to a}f(x)=L$ berarti nilai $f(x)$ menuju $L$ ketika $x$ mendekati $a$ (dari kedua sisi). Limit \textbf{ada} jika limit kiri dan kanan sama:
\[
\lim_{x\to a^-}f(x)=\lim_{x\to a^+}f(x)=L.
\]
\end{konsep}

\begin{center}
\begin{tikzpicture}
\begin{axis}[axis lines=middle,width=9cm,height=5.5cm,
  xlabel=$x$,ylabel=$y$,xmin=-0.3,xmax=4.5,ymin=-0.3,ymax=5.5,
  xtick={2},xticklabels={$a$},ytick={4},yticklabels={$L$},
  samples=60,clip=false]
% f(x) = x+2 with a hole at x=2 (limit exists but f(2) not defined)
\addplot[very thick,biru,domain=0.1:1.88]{x+2};
\addplot[very thick,biru,domain=2.12:4.2]{x+2};
\draw[biru,fill=white,thick] (2,4) circle(2.5pt);
% Dashed reference lines
\draw[dashed,gray,thin] (axis cs:0,4)--(axis cs:2,4);
\draw[dashed,gray,thin] (axis cs:2,0)--(axis cs:2,4);
% Left approach arrow (x → 2⁻)
\draw[-Stealth,oranye,very thick] (axis cs:0.6,2.6)--(axis cs:1.75,3.75)
  node[midway,above left,font=\footnotesize]{$x\!\to\!a^-$};
% Right approach arrow (x → 2⁺)
\draw[-Stealth,hijau,very thick] (axis cs:3.6,5.6)--(axis cs:2.25,4.25)
  node[midway,above right,font=\footnotesize]{$x\!\to\!a^+$};
% Label at hole
\node[biru,font=\tiny,below right] at (axis cs:2.05,3.9){lubang ($f(a)$ tak tentu)};
\end{axis}
\end{tikzpicture}\\
\small \textbf{Limit dari kiri dan kanan.} Kedua panah menuju titik yang sama ($L$), sehingga limit ada --- meski $f(a)$ sendiri tidak terdefinisi. Limit hanya peduli pada perilaku \emph{sekitar} $a$, bukan \emph{di} $a$.
\end{center}

\begin{contoh}{Substitusi Langsung}
Hitung $\displaystyle\lim_{x\to2}(3x+1)$.

\jawab Fungsi kontinu, substitusi langsung: $3(2)+1=7.$
\end{contoh}

\section{Bentuk Tak Tentu dan Tekniknya}
\begin{konsep}{Bentuk Tak Tentu}
Bentuk seperti $\dfrac00$, $\dfrac{\infty}{\infty}$, $\infty-\infty$, $0\cdot\infty$ tidak langsung bernilai; perlu dimanipulasi. Teknik utama:
\begin{itemize}[leftmargin=*,topsep=2pt]
  \item \textbf{Faktorisasi} lalu coret faktor penyebab nol;
  \item \textbf{Perkalian sekawan} untuk bentuk akar;
  \item \textbf{Bagi dengan pangkat tertinggi} untuk limit di tak hingga.
\end{itemize}
\end{konsep}

\begin{catatan}
\textbf{Catatan.} Selain ketiga teknik di atas, ada satu cara ampuh lain untuk bentuk $\tfrac00$ dan $\tfrac\infty\infty$, yaitu \emph{aturan L'Hôpital}, yang dibahas pada bagian akhir bab ini.
\end{catatan}

\begin{contoh}{Faktorisasi ($\tfrac00$)}
Hitung $\displaystyle\lim_{x\to2}\frac{x^2-4}{x-2}$.

\jawab $\dfrac{(x-2)(x+2)}{x-2}=x+2\to4.$
\end{contoh}

\begin{contoh}{Perkalian Sekawan}
Hitung $\displaystyle\lim_{x\to0}\frac{\sqrt{x+4}-2}{x}$.

\jawab Kalikan sekawan: $\dfrac{(\sqrt{x+4}-2)(\sqrt{x+4}+2)}{x(\sqrt{x+4}+2)}=\dfrac{x}{x(\sqrt{x+4}+2)}=\dfrac{1}{\sqrt{x+4}+2}\to\dfrac14.$
\end{contoh}

\begin{contoh}{Limit di Tak Hingga}
Hitung $\displaystyle\lim_{x\to\infty}\frac{3x^2+2x}{x^2-1}$.

\jawab Bagi pembilang dan penyebut dengan $x^2$: $\dfrac{3+2/x}{1-1/x^2}\to\dfrac{3}{1}=3.$
\end{contoh}

\begin{kesalahan}
$\tfrac00$ \textbf{bukan} berarti limitnya $0$ atau tak ada --- ia "tak tentu", artinya bisa bernilai berapa pun tergantung fungsinya. Selalu manipulasi dulu (faktorkan/sekawan) sebelum menyimpulkan.
\end{kesalahan}

\begin{dunianyata}
Limit mendasari definisi \emph{kecepatan sesaat}, \emph{kemiringan garis singgung}, dan \emph{laju perubahan}. Tanpa limit, konsep turunan dan integral pada bab-bab berikut tidak terdefinisi.
\end{dunianyata}

\section{Aturan L'Hôpital}
\begin{konsep}{Aturan L'Hôpital}
Untuk bentuk tak tentu $\dfrac00$ atau $\dfrac{\infty}{\infty}$, jika $f$ dan $g$ dapat diturunkan, maka
\[
\lim_{x\to a}\frac{f(x)}{g(x)}=\lim_{x\to a}\frac{f'(x)}{g'(x)},
\]
selama limit ruas kanan ada; aturan boleh \emph{diulang} bila hasilnya masih tak tentu. Turunan yang diperlukan di sini sederhana:
\[
\tfrac{d}{dx}x^n=nx^{n-1},\quad (\sin x)'=\cos x,\quad (\cos x)'=-\sin x,\quad (e^x)'=e^x,\quad (\ln x)'=\tfrac1x,
\]
(dibahas lengkap pada bab Turunan Fungsi).
\end{konsep}

\begin{contoh}{L'Hôpital untuk $\tfrac00$}
Hitung $\displaystyle\lim_{x\to2}\frac{x^2-4}{x-2}$ dengan L'Hôpital.

\jawab Substitusi memberi $\tfrac00$. Turunkan pembilang dan penyebut \emph{terpisah}:
$\displaystyle\lim_{x\to2}\frac{2x}{1}=4$ (sama dengan hasil faktorisasi).
\end{contoh}

\begin{contoh}{L'Hôpital Berulang}
Hitung $\displaystyle\lim_{x\to0}\frac{1-\cos x}{x^2}$.

\jawab $\tfrac00$ $\Rightarrow$ $\displaystyle\lim_{x\to0}\frac{\sin x}{2x}$ (masih $\tfrac00$) $\Rightarrow$ $\displaystyle\lim_{x\to0}\frac{\cos x}{2}=\frac12.$
\end{contoh}

\begin{kesalahan}
L'Hôpital \textbf{hanya} berlaku untuk bentuk tak tentu $\tfrac00$ atau $\tfrac\infty\infty$. Jangan terapkan pada bentuk tentu (mis. $\tfrac35$) --- hasilnya akan salah. Turunkan pembilang dan penyebut \emph{terpisah}; ini \textbf{bukan} aturan hasil bagi $\left(\tfrac fg\right)'$.
\end{kesalahan}

\section{Limit Fungsi Trigonometri}
\begin{konsep}{Limit Istimewa Trigonometri}
Untuk $x\to0$ (sudut dalam radian):
\[
\lim_{x\to0}\frac{\sin x}{x}=1,\quad
\lim_{x\to0}\frac{\tan x}{x}=1,\quad
\lim_{x\to0}\frac{\sin ax}{bx}=\frac ab,\quad
\lim_{x\to0}\frac{1-\cos x}{x^2}=\frac12.
\]
Semua berlaku hanya untuk bentuk $\tfrac00$ dengan argumen menuju $0$.
\end{konsep}

\begin{contoh}{Limit Trigonometri}
Hitung $\displaystyle\lim_{x\to0}\frac{\sin 5x}{\sin 3x}$ dan $\displaystyle\lim_{x\to0}\frac{1-\cos 2x}{x^2}$.

\jawab (i) $\dfrac{\sin5x}{\sin3x}=\dfrac{\sin5x}{5x}\cdot\dfrac{3x}{\sin3x}\cdot\dfrac{5x}{3x}\to1\cdot1\cdot\dfrac53=\dfrac53$. \\
(ii) $1-\cos2x=2\sin^2x$, jadi $\dfrac{2\sin^2x}{x^2}=2\left(\dfrac{\sin x}{x}\right)^2\to2$.
\end{contoh}

\begin{latihan}
Hitung: (1) $\displaystyle\lim_{x\to0}\frac{\sin 3x}{x}$;\quad (2) $\displaystyle\lim_{x\to0}\frac{\tan x}{2x}$;\quad (3) $\displaystyle\lim_{x\to0}\frac{1-\cos x}{x\sin x}$.
\end{latihan}

\section{Limit Eksponensial dan Logaritma}
\begin{konsep}{Limit Istimewa $e$, Eksponen, dan Logaritma}
\[
\lim_{x\to\infty}\!\left(1+\frac1x\right)^{x}=e,\quad
\lim_{x\to0}(1+x)^{1/x}=e,\quad
\lim_{x\to0}\frac{e^x-1}{x}=1,\quad
\lim_{x\to0}\frac{\ln(1+x)}{x}=1,\quad
\lim_{x\to0}\frac{a^x-1}{x}=\ln a.
\]
\end{konsep}

\begin{contoh}{Limit Eksponen dan Logaritma}
Hitung $\displaystyle\lim_{x\to0}\frac{e^{3x}-1}{x}$ dan $\displaystyle\lim_{x\to\infty}\left(1+\frac2x\right)^{x}$.

\jawab (i) $\dfrac{e^{3x}-1}{x}=3\cdot\dfrac{e^{3x}-1}{3x}\to3\cdot1=3$. \\
(ii) Tulis $\left(1+\tfrac2x\right)^x=\left[\left(1+\tfrac2x\right)^{x/2}\right]^{2}\to e^{2}$.
\end{contoh}

\begin{latihan}
Hitung: (1) $\displaystyle\lim_{x\to0}\frac{e^{2x}-1}{\sin x}$;\quad (2) $\displaystyle\lim_{x\to\infty}\left(1+\frac3x\right)^{x}$;\quad (3) $\displaystyle\lim_{x\to0}\frac{\ln(1+5x)}{x}$.
\end{latihan}

\section{Asimtot dari Limit}
\begin{konsep}{Asimtot lewat Limit}
\begin{itemize}[leftmargin=*,topsep=2pt]
  \item \textbf{Asimtot tegak} $x=a$ jika $\displaystyle\lim_{x\to a^\pm}f(x)=\pm\infty$.
  \item \textbf{Asimtot datar} $y=L$ jika $\displaystyle\lim_{x\to\pm\infty}f(x)=L$.
  \item \textbf{Asimtot miring} $y=mx+c$ dengan $m=\displaystyle\lim_{x\to\pm\infty}\frac{f(x)}{x}$ dan $c=\displaystyle\lim_{x\to\pm\infty}\big(f(x)-mx\big)$.
\end{itemize}
\end{konsep}

\begin{contoh}{Asimtot Miring}
Tentukan asimtot $f(x)=\dfrac{2x^2+1}{x-1}$.

\jawab Asimtot tegak $x=1$ (penyebut nol). Untuk miring, $2x^2+1=(x-1)(2x+2)+3$, jadi $f(x)=2x+2+\dfrac{3}{x-1}$; saat $x\to\pm\infty$ sukunya $\to0$, sehingga asimtot miring $y=2x+2$.
\end{contoh}

\begin{latihan}
(1) Tentukan asimtot datar $f(x)=\dfrac{3x-1}{x+2}$.\quad (2) Tentukan asimtot tegak $f(x)=\dfrac{1}{x^2-4}$.\quad (3) Tentukan asimtot miring $f(x)=\dfrac{x^2+x+1}{x}$.
\end{latihan}

\section*{Ringkasan Bab}
\addcontentsline{toc}{section}{Ringkasan Bab}
\begin{itemize}[leftmargin=*,topsep=2pt]
  \item Limit ada jika limit kiri $=$ limit kanan.
  \item Bentuk tak tentu ($\tfrac00$, $\tfrac\infty\infty$, $\infty-\infty$) memerlukan manipulasi.
  \item Teknik: faktorisasi, perkalian sekawan, bagi pangkat tertinggi.
  \item Limit $\dfrac{\sin x}{x}\to1$ saat $x\to0$; juga $\dfrac{\tan x}{x}\to1$, $\dfrac{1-\cos x}{x^2}\to\tfrac12$.
  \item Eksponen/logaritma: $\left(1+\tfrac1x\right)^x\to e$, $\dfrac{e^x-1}{x}\to1$, $\dfrac{\ln(1+x)}{x}\to1$.
  \item \textbf{Asimtot:} tegak ($\lim=\pm\infty$), datar ($\lim_{x\to\infty}=L$), miring ($m=\lim\tfrac fx$, $c=\lim(f-mx)$).
  \item \textbf{L'Hôpital:} untuk $\tfrac00$ atau $\tfrac\infty\infty$, $\lim\dfrac fg=\lim\dfrac{f'}{g'}$.
  \item Limit adalah fondasi turunan dan integral.
\end{itemize}

\begin{refleksi}
\begin{itemize}[leftmargin=*,topsep=2pt]
  \item[\cek] Saya memahami limit dan limit sepihak.
  \item[\cek] Saya dapat mengenali bentuk tak tentu.
  \item[\cek] Saya dapat memilih teknik (faktorisasi/sekawan/pangkat tertinggi).
  \item[\cek] Saya dapat memakai aturan L'Hôpital untuk bentuk $\tfrac00$ dan $\tfrac\infty\infty$.
  \item[\cek] Saya dapat menghubungkan limit dengan kontinuitas.
\end{itemize}
\end{refleksi}

\section*{Soal Akhir Bab \thechapter}
\addcontentsline{toc}{section}{Soal Akhir Bab \thechapter}

\bagiansoal{A}{(Fundamental --- setara SNBT Dasar)}
\begin{enumerate}[leftmargin=*]
  \item Hitung $\displaystyle\lim_{x\to2}(3x+1)$.
  \item Hitung $\displaystyle\lim_{x\to3}\frac{x^2-9}{x-3}$.
  \item Hitung $\displaystyle\lim_{x\to\infty}\frac{2x+1}{x-3}$.
  \item Hitung $\displaystyle\lim_{x\to1}(x^2+2x)$.
  \item Tuliskan nilai $\displaystyle\lim_{x\to0}\frac{\sin x}{x}$.
\end{enumerate}

\bagiansoal{B}{(Intermediate --- setara SNBT Sulit)}
\begin{enumerate}[leftmargin=*]
  \item Hitung $\displaystyle\lim_{x\to2}\frac{x^2-4}{x-2}$.
  \item Hitung $\displaystyle\lim_{x\to\infty}\frac{3x^2+2x}{x^2-1}$.
  \item Hitung $\displaystyle\lim_{x\to0}\frac{\sqrt{x+4}-2}{x}$.
  \item Hitung $\displaystyle\lim_{x\to\infty}\left(\sqrt{x^2+x}-x\right)$.
  \item Selidiki $\displaystyle\lim_{x\to3}\frac{x-3}{|x-3|}$ dengan limit kiri dan kanan.
  \item Hitung $\displaystyle\lim_{x\to0}\frac{\sin 3x}{2x}$.
  \item \textbf{(L'Hôpital).} Hitung $\displaystyle\lim_{x\to3}\frac{x^2-9}{x-3}$ dengan aturan L'Hôpital.
  \item \textbf{(L'Hôpital).} Hitung $\displaystyle\lim_{x\to0}\frac{\sin x}{x}$ dengan aturan L'Hôpital.
\end{enumerate}

\bagiansoal{C}{(Advanced --- setara Olimpiade SMA, gabungan antarbab)}
\begin{enumerate}[leftmargin=*]
  \item \textbf{(Limit + Faktorisasi).} Hitung $\displaystyle\lim_{x\to1}\frac{x^3-1}{x-1}$.
  \item \textbf{(Limit + Trigonometri).} Hitung $\displaystyle\lim_{x\to0}\frac{1-\cos x}{x^2}$.
  \item \textbf{(Limit + Akar).} Hitung $\displaystyle\lim_{x\to4}\frac{x-4}{\sqrt{x}-2}$.
  \item \textbf{(Limit Tak Hingga).} Hitung $\displaystyle\lim_{x\to\infty}\left(\sqrt{4x^2+x}-2x\right)$.
  \item \textbf{(Limit + Definisi Turunan).} Hitung $\displaystyle\lim_{h\to0}\frac{(2+h)^2-4}{h}$ dan jelaskan maknanya.
  \item \textbf{(Limit + Kontinuitas).} Tentukan $a$ agar $f(x)=\begin{cases}\frac{x^2-1}{x-1},&x\neq1\\ a,&x=1\end{cases}$ kontinu di $x=1$.
  \item \textbf{(L'Hôpital Berulang).} Hitung $\displaystyle\lim_{x\to0}\frac{x-\sin x}{x^3}$ dengan L'Hôpital.
  \item \textbf{(L'Hôpital, bentuk $\tfrac\infty\infty$).} Hitung $\displaystyle\lim_{x\to\infty}\frac{\ln x}{x}$ dengan L'Hôpital.
\end{enumerate}

\bagiansoal{D}{(Expert --- setara tahun pertama universitas, sintesis \& pembuktian)}
\begin{enumerate}[leftmargin=*]
  \item \textbf{(Teorema Apit).} Gunakan teorema apit untuk membuktikan $\displaystyle\lim_{x\to0}x^2\sin\frac1x=0$.
  \item \textbf{(Bilangan $e$).} Jelaskan mengapa $\displaystyle\lim_{x\to\infty}\left(1+\frac1x\right)^x=e$ dan tentukan $\displaystyle\lim_{x\to\infty}\left(1+\frac2x\right)^x$.
  \item \textbf{(Limit Tak Tentu).} Hitung $\displaystyle\lim_{x\to0}\frac{\tan x-\sin x}{x^3}$.
  \item \textbf{(Pembuktian Geometris).} Buktikan $\displaystyle\lim_{x\to0}\frac{\sin x}{x}=1$ dengan argumen luas pada lingkaran satuan.
  \item \textbf{(L'Hôpital, bentuk $0\cdot\infty$).} Hitung $\displaystyle\lim_{x\to0^+}x\ln x$ dengan mengubahnya ke bentuk $\tfrac\infty\infty$ lalu menerapkan L'Hôpital.
\end{enumerate}

\begin{challenge}
\begin{enumerate}[leftmargin=*,topsep=2pt]
  \item Hitung $\displaystyle\lim_{x\to0}\frac{\sqrt{1+x}-\sqrt{1-x}}{x}$.
  \item Tentukan $a$ dan $b$ agar $\displaystyle\lim_{x\to\infty}\left(\frac{x^2+1}{x+1}-ax-b\right)=0$.
\end{enumerate}
\end{challenge}

% ============================================================
\chapter{Turunan Fungsi}
% ============================================================
\begin{tujuan}
Setelah mempelajari bab ini, peserta didik mampu:
\begin{itemize}[leftmargin=*,topsep=2pt]
  \item memahami turunan sebagai limit dan sebagai gradien garis singgung;
  \item menggunakan aturan turunan: pangkat, hasil kali, hasil bagi, dan rantai;
  \item menentukan turunan fungsi trigonometri, eksponen, dan logaritma;
  \item menerapkan turunan pada garis singgung, kecepatan sesaat, dan laju perubahan;
  \item menyelesaikan masalah maksimum--minimum (optimasi).
\end{itemize}
\textbf{Capaian Pembelajaran (Fase F+).} Peserta didik dapat menghitung dan menerapkan turunan untuk menganalisis perubahan serta mengoptimalkan besaran.
\end{tujuan}

\begin{pemantik}
\begin{itemize}[leftmargin=*,topsep=2pt]
  \item Bagaimana menemukan kemiringan kurva yang melengkung di satu titik?
  \item Apa hubungan antara "laju perubahan" dan kemiringan garis singgung?
  \item Bagaimana kalkulus membantu memilih ukuran kemasan yang menghemat bahan?
\end{itemize}
\end{pemantik}

\begin{studikasus}{Mendesain Kaleng Hemat Bahan}
Sebuah pabrik ingin membuat kaleng bervolume tetap dengan bahan sesedikit mungkin. Luas bahan bergantung pada jari-jari; saat jari-jari membesar, luas mula-mula turun lalu naik kembali. Titik "paling hemat" adalah saat \emph{laju perubahan luas nol} --- yaitu turunannya nol.
\begin{center}
\begin{tikzpicture}
\begin{axis}[axis lines=middle,width=8.5cm,height=4.8cm,xlabel={\small jari-jari $r$},ylabel={\small luas bahan},
  xmin=0,xmax=2.8,ymin=0,ymax=9,xtick=\empty,ytick=\empty,samples=100,clip=false]
\addplot[very thick,biru,domain=0.55:2.5]{x^2+4/x};
\draw[dashed,gray] (axis cs:1.26,0)--(axis cs:1.26,4.76);
\draw[oranye,thick] (axis cs:0.66,4.76)--(axis cs:1.86,4.76);
\addplot[only marks,mark=*,oranye] coordinates {(1.26,4.76)};
\node[oranye,font=\scriptsize,anchor=south] at (axis cs:1.26,5.2){minimum (turunan $=0$)};
\node[font=\scriptsize,anchor=north] at (axis cs:1.26,-0.1){$r^*$};
\end{axis}
\end{tikzpicture}\\
\small Luas bahan terhadap jari-jari: titik terendah (turunan nol) memberi desain paling hemat.
\end{center}
\textbf{Amati dan duga.} Mengapa "titik terbaik" selalu berada di tempat garis singgung mendatar? Inilah inti optimasi dengan turunan.
\end{studikasus}

\section*{Peta Konsep Bab}
\addcontentsline{toc}{section}{Peta Konsep Bab}
\begin{center}
\begin{tikzpicture}[
  node distance=8mm,
  konsepbox/.style={draw=biru,fill=birumuda,rounded corners,align=center,inner sep=5pt,font=\small\bfseries},
  subbox/.style={draw=hijau,fill=hijaumuda,rounded corners,align=center,inner sep=4pt,font=\footnotesize},
  garis/.style={-Stealth,biru,thick}]
  \node[konsepbox] (akar) {TURUNAN FUNGSI};
  \node[subbox,below left=12mm and 16mm of akar] (def) {Definisi \&\\aturan pangkat};
  \node[subbox,below=12mm of akar] (rule) {Hasil kali, bagi,\\rantai};
  \node[subbox,below right=12mm and 16mm of akar] (fungsi) {Trig, ekspon,\\logaritma};
  \node[subbox,below=10mm of rule] (apl) {Garis singgung,\\optimasi};
  \draw[garis] (akar)--(def); \draw[garis] (akar)--(rule); \draw[garis] (akar)--(fungsi);
  \draw[garis] (rule)--(apl);
\end{tikzpicture}\\
\small \textbf{Peta konsep.} Dari definisi limit tumbuh aturan-aturan turunan, yang diterapkan pada garis singgung, kinematika, dan optimasi.
\end{center}

\section{Definisi dan Aturan Turunan}
\begin{konsep}{Definisi dan Aturan Dasar}
\[
f'(x) = \lim_{h\to0}\frac{f(x+h)-f(x)}{h}.
\]
\textbf{Aturan pangkat:} $\dfrac{d}{dx}x^n=nx^{n-1}$. \textbf{Linear:} $(af+bg)'=af'+bg'$. Turunan adalah \emph{gradien garis singgung}.
\end{konsep}

\begin{center}
\begin{tikzpicture}
\begin{axis}[axis lines=middle,width=7.5cm,height=5cm,
  xlabel=$x$,ylabel=$y$,xmin=-0.5,xmax=3,ymin=-1,ymax=5,xtick=\empty,ytick=\empty]
\addplot[very thick,biru,domain=-0.3:2.6,samples=50]{x^2};
\addplot[very thick,oranye,domain=0.3:2.5,samples=2]{2*1.5*x-1.5^2};
\fill (1.5,2.25) circle (2pt) node[above left]{\small singgung di $x=1{,}5$};
\end{axis}
\end{tikzpicture}\\
\small Garis singgung kurva $y=x^2$; gradiennya $=f'(1{,}5)=2(1{,}5)=3$.
\end{center}

\begin{center}
\begin{tikzpicture}
\begin{axis}[axis lines=middle,width=9cm,height=5.8cm,
  xlabel=$x$,ylabel=$y$,xmin=-0.3,xmax=3.8,ymin=-0.5,ymax=6.5,
  xtick={1.5},xticklabels={$a$},ytick=\empty,samples=80,clip=false]
% Parabola y=x²
\addplot[very thick,biru,domain=0:3]{x^2};
% Point of tangency a=1.5
\addplot[only marks,mark=*,oranye,mark size=3pt] coordinates {(1.5,2.25)};
% Tangent at x=1.5: y=3x-2.25
\addplot[oranye,very thick,domain=0.4:2.6]{3*x-2.25};
% Secant 1: a=1.5 to x=3.0, slope=(9-2.25)/(3-1.5)=4.5
\addplot[gray!40,thick,dashed,domain=0.3:2.6]{4.5*x-4.5};
% Secant 2: a=1.5 to x=2.5, slope=(6.25-2.25)/1=4.0
\addplot[gray!65,thick,dashed,domain=0.4:2.6]{4*x-3.75};
% Secant 3: a=1.5 to x=2.0, slope=(4-2.25)/0.5=3.5
\addplot[gray,thick,dashed,domain=0.6:2.6]{3.5*x-3};
% Labels - tangent
\node[oranye,font=\small,right,fill=white,inner sep=1pt] at (axis cs:2.55,5.6){tangent ($h\!\to\!0$)};
% Labels - secants (positioned away from the lines, with fill=white background)
\node[black!60,font=\small,right,fill=white,inner sep=1pt] at (axis cs:0.2,5.5){sekans ($h$ besar)};
\draw[-Stealth,black!60,thin] (axis cs:0.85,5.3)--(axis cs:1.15,4.7);
% Arrow showing Q approaching a
\draw[-Stealth,biru,thick] (axis cs:3.2,3.8)--(axis cs:2.15,3.1);
\node[biru,font=\small,fill=white,inner sep=1pt] at (axis cs:3.35,4.1){$Q\!\to\!a$};
\end{axis}
\end{tikzpicture}\\
\small \textbf{Sekans $\to$ tangent.} Saat titik $Q$ mendekati titik singgung $a$, garis sekans (putus-putus, makin gelap) mendekati garis singgung (oranye). Gradiennya menuju $f'(a)=\lim_{h\to0}\tfrac{f(a+h)-f(a)}{h}$.
\end{center}

\begin{konsep}{Aturan Hasil Kali, Hasil Bagi, dan Rantai}
\[
(fg)'=f'g+fg', \qquad \left(\frac fg\right)'=\frac{f'g-fg'}{g^2}, \qquad \big(f(g(x))\big)'=f'(g(x))\,g'(x).
\]
\textbf{Turunan khusus:} $(\sin x)'=\cos x$, $(\cos x)'=-\sin x$, $(e^x)'=e^x$, $(\ln x)'=\tfrac1x$.
\end{konsep}

\begin{contoh}{Aturan Rantai}
Tentukan turunan $f(x)=(2x+1)^5$.

\jawab $f'(x)=5(2x+1)^4\cdot2=10(2x+1)^4.$
\end{contoh}

\section{Penerapan: Garis Singgung dan Optimasi}
\begin{contoh}{Optimasi Sederhana}
Tentukan nilai maksimum $f(x)=-x^2+4x+1$.

\jawab $f'(x)=-2x+4=0\Rightarrow x=2$. Karena parabola terbuka ke bawah, ini maksimum: $f(2)=5.$
\end{contoh}

\begin{kesalahan}
Turunan hasil kali \textbf{bukan} $f'g'$. Gunakan $(fg)'=f'g+fg'$. Demikian pula jangan lupa faktor "turunan dalam" $g'(x)$ pada aturan rantai.
\end{kesalahan}

\begin{dunianyata}
Turunan adalah laju perubahan: kecepatan ($\tfrac{ds}{dt}$), percepatan, laju reaksi kimia, laju pertumbuhan, hingga "kemiringan" biaya marginal dalam ekonomi. Optimasi dengan turunan dipakai merancang struktur, menekan biaya, dan melatih model AI.
\end{dunianyata}

\section{Turunan Implisit}
\begin{konsep}{Penurunan Implisit}
Jika $y$ terkait $x$ secara \emph{implisit} (tidak dapat dipisah menjadi $y=f(x)$), turunkan kedua ruas terhadap $x$ dengan memperlakukan $y$ sebagai fungsi $x$ (gunakan aturan rantai: setiap kali menurunkan suku berisi $y$, kalikan $\dfrac{dy}{dx}$), lalu selesaikan untuk $\dfrac{dy}{dx}$.
\end{konsep}

\begin{contoh}{Turunan Implisit}
Tentukan $\dfrac{dy}{dx}$ dari $x^2+y^2=25$, lalu gradien di $(3,4)$.

\jawab Turunkan: $2x+2y\,y'=0\Rightarrow y'=-\dfrac xy$. Di $(3,4)$: $y'=-\dfrac34$.
\end{contoh}

\begin{contoh}{Implisit dengan Suku Campuran}
Tentukan $y'$ dari $x^2+xy+y^2=7$.

\jawab $2x+(y+x y')+2y\,y'=0\Rightarrow y'(x+2y)=-(2x+y)\Rightarrow y'=-\dfrac{2x+y}{x+2y}$.
\end{contoh}

\begin{kesalahan}
Saat menurunkan suku berisi $y$, jangan lupa faktor $y'$ (aturan rantai). Misalnya $\dfrac{d}{dx}\,y^2=2y\,y'$, \emph{bukan} $2y$.
\end{kesalahan}

\begin{latihan}
(1) Tentukan $y'$ dari $x^3+y^3=6xy$.\quad (2) Tentukan $y'$ dari $\sin y+x^2=y$.\quad (3) Tentukan persamaan garis singgung $x^2+y^2=25$ di $(3,4)$.
\end{latihan}

\section{Kemonotonan, Kecekungan, dan Titik Belok}
\begin{konsep}{Turunan Pertama dan Kedua}
\textbf{Naik--turun:} $f'(x)>0\Rightarrow f$ naik; $f'(x)<0\Rightarrow f$ turun. \textbf{Kecekungan:} $f''(x)>0\Rightarrow$ cekung ke atas; $f''(x)<0\Rightarrow$ cekung ke bawah. \textbf{Titik belok:} tempat $f''$ berubah tanda. \textbf{Uji turunan kedua:} bila $f'(c)=0$, maka $f''(c)<0$ memberi maksimum lokal dan $f''(c)>0$ memberi minimum lokal.
\end{konsep}

\begin{contoh}{Analisis Kurva}
Selidiki kemonotonan, jenis titik stasioner, dan titik belok $f(x)=x^3-3x^2$.

\jawab $f'(x)=3x^2-6x=3x(x-2)$: $f$ naik untuk $x<0$ dan $x>2$, turun untuk $0<x<2$. $f''(x)=6x-6$, nol di $x=1$ (titik belok $(1,-2)$). Karena $f''(0)=-6<0$, titik $(0,0)$ maksimum lokal; karena $f''(2)=6>0$, titik $(2,-4)$ minimum lokal.
\end{contoh}

\begin{latihan}
(1) Tentukan interval naik/turun dan titik belok $f(x)=x^3-3x$.\quad (2) Tentukan selang kecekungan $f(x)=x^4-2x^2$.\quad (3) Gunakan uji turunan kedua untuk menggolongkan titik stasioner $f(x)=x^3-12x$.
\end{latihan}

\section*{Ringkasan Bab}
\addcontentsline{toc}{section}{Ringkasan Bab}
\begin{itemize}[leftmargin=*,topsep=2pt]
  \item $f'(x)=\lim_{h\to0}\tfrac{f(x+h)-f(x)}{h}$; aturan pangkat $\tfrac{d}{dx}x^n=nx^{n-1}$.
  \item Hasil kali, hasil bagi, dan rantai; turunan trig/eksponen/logaritma.
  \item \textbf{Turunan implisit:} turunkan kedua ruas, suku berisi $y$ mendapat faktor $y'$.
  \item $f'>0$ naik, $f'<0$ turun; $f''>0$ cekung atas, $f''<0$ cekung bawah; \textbf{titik belok} saat $f''$ berganti tanda.
  \item Garis singgung di $x=a$ bergradien $f'(a)$.
  \item Titik stasioner ($f'=0$) menandai maksimum/minimum lokal (uji $f''$).
\end{itemize}

\begin{refleksi}
\begin{itemize}[leftmargin=*,topsep=2pt]
  \item[\cek] Saya memahami turunan sebagai limit dan gradien.
  \item[\cek] Saya dapat memakai aturan pangkat, hasil kali, bagi, dan rantai.
  \item[\cek] Saya dapat menurunkan fungsi trig, eksponen, dan logaritma.
  \item[\cek] Saya dapat menyelesaikan masalah optimasi.
\end{itemize}
\end{refleksi}

\section*{Soal Akhir Bab \thechapter}
\addcontentsline{toc}{section}{Soal Akhir Bab \thechapter}

\bagiansoal{A}{(Fundamental --- setara SNBT Dasar)}
\begin{enumerate}[leftmargin=*]
  \item Tentukan $f'(x)$ dari $f(x)=3x^2-4x+1$.
  \item Tentukan $f'(x)$ dari $f(x)=x^3-2x^2+5$.
  \item Tentukan turunan $f(x)=5x^4$.
  \item Tentukan gradien garis singgung $y=x^2-3x$ di $x=2$.
  \item Tentukan turunan $f(x)=\sqrt{x}$.
\end{enumerate}

\bagiansoal{B}{(Intermediate --- setara SNBT Sulit)}
\begin{enumerate}[leftmargin=*]
  \item Tentukan turunan $f(x)=(x^2+1)(x-3)$ dengan aturan hasil kali.
  \item Tentukan turunan $f(x)=\dfrac{x}{x+1}$ dengan aturan hasil bagi.
  \item Tentukan turunan $f(x)=(2x+1)^5$.
  \item Tentukan turunan $f(x)=\sin(3x)$.
  \item Tentukan turunan $f(x)=e^{2x}$.
  \item Tentukan turunan $f(x)=\ln(x^2+1)$.
  \item Tentukan titik stasioner $f(x)=x^3-3x$.
\end{enumerate}

\bagiansoal{C}{(Advanced --- setara Olimpiade SMA, gabungan antarbab)}
\begin{enumerate}[leftmargin=*]
  \item \textbf{(Turunan + Optimasi).} Sebuah kawat sepanjang $20$ cm dibentuk persegi panjang. Dengan turunan, tentukan ukuran yang memaksimalkan luas.
  \item \textbf{(Turunan + Garis Singgung).} Tentukan persamaan garis singgung $y=x^2+1$ di titik berabsis $x=1$.
  \item \textbf{(Turunan + Kinematika).} Posisi $s(t)=t^3-6t^2+9t$. Tentukan kecepatan dan saat benda berhenti sesaat.
  \item \textbf{(Turunan + Trigonometri).} Tentukan turunan $f(x)=\sin^2 x$.
  \item \textbf{(Turunan + Maks-Min).} Tentukan nilai maksimum $f(x)=-x^2+4x+1$.
  \item \textbf{(Aturan Rantai).} Tentukan turunan $f(x)=\sqrt{x^2+1}$.
\end{enumerate}

\bagiansoal{D}{(Expert --- setara tahun pertama universitas, sintesis \& pembuktian)}
\begin{enumerate}[leftmargin=*]
  \item \textbf{(Definisi).} Gunakan definisi limit untuk membuktikan bahwa turunan $f(x)=x^2$ adalah $2x$.
  \item \textbf{(Pembuktian Aturan Hasil Kali).} Buktikan $(fg)'=f'g+fg'$ dari definisi turunan.
  \item \textbf{(Optimasi Terapan).} Sebuah silinder bervolume tetap $V$ akan dibuat dengan luas permukaan minimum. Tunjukkan bahwa solusinya memenuhi tinggi $=2\times$ jari-jari.
  \item \textbf{(Turunan Trigonometri).} Buktikan $(\sin x)'=\cos x$ dari definisi limit (gunakan limit istimewa $\tfrac{\sin h}{h}\to1$ dan $\tfrac{1-\cos h}{h}\to0$).
  \item \textbf{(Teorema Nilai Rata-Rata).} Jelaskan mengapa jika $f'(x)=0$ pada seluruh interval, maka $f$ konstan di interval itu.
\end{enumerate}

\begin{challenge}
\begin{enumerate}[leftmargin=*,topsep=2pt]
  \item Tentukan turunan ke-$n$ dari $f(x)=\dfrac1x$.
  \item Untuk $x>0$, tentukan nilai minimum $f(x)=x+\dfrac1x$ dengan turunan, lalu cocokkan dengan ketaksamaan AM--GM.
\end{enumerate}
\end{challenge}

% ============================================================
\chapter{Integral}
% ============================================================
\begin{tujuan}
Setelah mempelajari bab ini, peserta didik mampu:
\begin{itemize}[leftmargin=*,topsep=2pt]
  \item menghitung integral tak tentu sebagai antiturunan;
  \item menghitung integral tentu dan menggunakan sifat-sifatnya;
  \item memakai substitusi sederhana;
  \item menghitung luas daerah dan volume benda putar;
  \item menjelaskan Teorema Dasar Kalkulus (hubungan turunan dan integral).
\end{itemize}
\textbf{Capaian Pembelajaran (Fase F+).} Peserta didik dapat menghitung integral dan menggunakannya untuk mengukur luas, volume, dan akumulasi.
\end{tujuan}

\begin{pemantik}
\begin{itemize}[leftmargin=*,topsep=2pt]
  \item Jika turunan "memecah" fungsi menjadi lajunya, adakah operasi yang "menyatukannya kembali"?
  \item Bagaimana menghitung luas daerah berbatas kurva lengkung?
  \item Apa hubungan antara luas di bawah kurva kecepatan dan jarak tempuh?
\end{itemize}
\end{pemantik}

\begin{studikasus}{Jarak dari Grafik Kecepatan}
Sebuah mobil melaju dengan kecepatan berubah-ubah. Jarak tempuhnya sama dengan \emph{luas daerah di bawah kurva kecepatan--waktu}. Untuk kecepatan yang berubah mulus, luas itu dihitung dengan \textbf{integral}.
\begin{center}
\begin{tikzpicture}
\begin{axis}[axis lines=middle,width=8cm,height=4.2cm,xlabel={\small $t$},ylabel={\small $v$},
  xmin=-0.3,xmax=3,ymin=-0.3,ymax=5,xtick={1,2},ytick=\empty,samples=50]
\addplot[very thick,biru,domain=-0.2:2.6]{x^2};
\addplot[draw=none,fill=birumuda,domain=1:2,samples=40] {x^2} \closedcycle;
\node at (1.5,1) {\scriptsize jarak $=\int_1^2 v\,dt$};
\end{axis}
\end{tikzpicture}\\
\small Luas daerah di bawah kurva kecepatan memberi jarak tempuh.
\end{center}
\textbf{Amati dan duga.} Bagaimana "menjumlahkan" luas potongan-potongan tipis tak hingga banyaknya menjadi satu angka pasti? Itulah yang dilakukan integral tentu.
\end{studikasus}

\section*{Peta Konsep Bab}
\addcontentsline{toc}{section}{Peta Konsep Bab}
\begin{center}
\begin{tikzpicture}[
  node distance=8mm,
  konsepbox/.style={draw=biru,fill=birumuda,rounded corners,align=center,inner sep=5pt,font=\small\bfseries},
  subbox/.style={draw=hijau,fill=hijaumuda,rounded corners,align=center,inner sep=4pt,font=\footnotesize},
  garis/.style={-Stealth,biru,thick}]
  \node[konsepbox] (akar) {INTEGRAL};
  \node[subbox,below left=12mm and 14mm of akar] (tt) {Tak tentu\\(antiturunan)};
  \node[subbox,below=12mm of akar] (tentu) {Tentu \&\\sifat};
  \node[subbox,below right=12mm and 14mm of akar] (sub) {Substitusi};
  \node[subbox,below=10mm of tentu] (apl) {Luas, volume,\\TDK};
  \draw[garis] (akar)--(tt); \draw[garis] (akar)--(tentu); \draw[garis] (akar)--(sub);
  \draw[garis] (tentu)--(apl);
\end{tikzpicture}\\
\small \textbf{Peta konsep.} Integral tak tentu (antiturunan) dan integral tentu (luas) dihubungkan oleh Teorema Dasar Kalkulus.
\end{center}

\section{Integral Tak Tentu dan Tentu}
\begin{konsep}{Rumus Dasar}
\[
\int x^n\,dx = \frac{x^{n+1}}{n+1}+C\ (n\neq-1), \qquad
\int_a^b f(x)\,dx = F(b)-F(a),
\]
dengan $F$ antiturunan dari $f$ ($F'=f$). Integral tentu mengukur luas (bertanda) di bawah kurva.
\end{konsep}

\begin{center}
\begin{tikzpicture}
\begin{axis}[axis lines=middle,width=7.5cm,height=4.8cm,
  xlabel=$x$,ylabel=$y$,xmin=-0.5,xmax=3,ymin=-0.5,ymax=5,xtick={1,2},ytick=\empty]
\addplot[very thick,biru,domain=-0.3:2.6,samples=50]{x^2};
\addplot[draw=none,fill=birumuda,domain=1:2,samples=40] {x^2} \closedcycle;
\node at (1.5,1) {\small $\int_1^2 x^2 dx$};
\end{axis}
\end{tikzpicture}
\end{center}

\begin{contoh}{Integral Tentu}
Hitung $\displaystyle\int_1^2 x^2\,dx$.

\jawab
$\left[\dfrac{x^3}{3}\right]_1^2 = \dfrac{8}{3}-\dfrac{1}{3} = \dfrac{7}{3}.$
\end{contoh}

\section{Substitusi, Luas, dan Teorema Dasar}
\begin{konsep}{Substitusi dan Teorema Dasar Kalkulus}
\textbf{Substitusi:} jika $u=g(x)$ maka $\displaystyle\int f(g(x))g'(x)\,dx=\int f(u)\,du$.\\
\textbf{Teorema Dasar Kalkulus:} $\dfrac{d}{dx}\displaystyle\int_a^x f(t)\,dt=f(x)$ --- integral dan turunan saling membatalkan.
\end{konsep}

\begin{contoh}{Substitusi}
Hitung $\displaystyle\int 2x(x^2+1)^3\,dx$.

\jawab Misal $u=x^2+1$, $du=2x\,dx$: $\displaystyle\int u^3\,du=\dfrac{u^4}{4}+C=\dfrac{(x^2+1)^4}{4}+C.$
\end{contoh}

\begin{contoh}{Luas Daerah}
Hitung luas daerah antara $y=x$ dan $y=x^2$ untuk $0\le x\le1$.

\jawab $\displaystyle\int_0^1(x-x^2)\,dx=\left[\dfrac{x^2}{2}-\dfrac{x^3}{3}\right]_0^1=\dfrac12-\dfrac13=\dfrac16.$
\end{contoh}

\begin{kesalahan}
Pada integral tak tentu, \textbf{jangan lupa $+C$}. Pada integral tentu, perhatikan bahwa luas di bawah sumbu-$x$ bernilai negatif --- untuk luas geometris murni, ambil nilai mutlaknya.
\end{kesalahan}

\begin{dunianyata}
Integral menghitung akumulasi: jarak dari kecepatan, kerja dari gaya, total muatan dari arus, dan luas/volume bentuk tak beraturan. Dalam statistika, luas di bawah kurva peluang menyatakan probabilitas.
\end{dunianyata}

\section{Teknik Lanjutan: Parsial, Trigonometri, dan Rasional}
\begin{konsep}{Integral Parsial}
\[
\int u\,dv = uv-\int v\,du.
\]
Pilih $u$ yang turunannya menyederhana (urutan prioritas \emph{LIATE}: Logaritma, Invers trig, Aljabar, Trigonometri, Eksponen).
\end{konsep}

\begin{contoh}{Integral Parsial}
Hitung $\displaystyle\int x\,e^{x}\,dx$.

\jawab Pilih $u=x$ ($du=dx$), $dv=e^x dx$ ($v=e^x$): $\displaystyle\int x e^x dx=xe^x-\int e^x dx=xe^x-e^x+C=e^x(x-1)+C.$
\end{contoh}

\begin{konsep}{Integral Trigonometri dan Rasional}
\textbf{Trigonometri dasar:} $\int\sin x\,dx=-\cos x+C$, $\int\cos x\,dx=\sin x+C$, $\int\sec^2x\,dx=\tan x+C$. Untuk pangkat genap gunakan identitas setengah sudut, mis. $\sin^2x=\tfrac{1-\cos2x}{2}$.\\
\textbf{Rasional:} uraikan menjadi \emph{pecahan parsial} lalu integralkan tiap suku; sering muncul $\int\tfrac{dx}{x-a}=\ln|x-a|+C$.
\end{konsep}

\begin{contoh}{Integral Trigonometri dan Rasional}
Hitung $\displaystyle\int\sin^2x\,dx$ dan $\displaystyle\int\frac{dx}{x^2-1}$.

\jawab (i) $\displaystyle\int\frac{1-\cos2x}{2}\,dx=\frac x2-\frac{\sin2x}{4}+C$.\\
(ii) $\dfrac{1}{x^2-1}=\dfrac12\!\left(\dfrac{1}{x-1}-\dfrac{1}{x+1}\right)$, jadi $\displaystyle\int\frac{dx}{x^2-1}=\frac12\ln\left|\frac{x-1}{x+1}\right|+C$.
\end{contoh}

\begin{latihan}
Hitung: (1) $\displaystyle\int x\,e^{2x}\,dx$;\quad (2) $\displaystyle\int\cos^2x\,dx$;\quad (3) $\displaystyle\int\frac{dx}{x^2-4}$.
\end{latihan}

\section{Volume Benda Putar dan Panjang Kurva}
\begin{konsep}{Rumus Volume dan Panjang}
\textbf{Cakram (terhadap sumbu-$x$):} $V=\pi\displaystyle\int_a^b [f(x)]^2\,dx$. \textbf{Cincin:} $V=\pi\displaystyle\int_a^b\big([R(x)]^2-[r(x)]^2\big)\,dx$. \textbf{Kulit tabung (terhadap sumbu-$y$):} $V=2\pi\displaystyle\int_a^b x\,f(x)\,dx$. \textbf{Panjang kurva:} $L=\displaystyle\int_a^b\sqrt{1+[f'(x)]^2}\,dx$.
\end{konsep}

\begin{contoh}{Volume dan Panjang Kurva}
(i) Daerah di bawah $y=\sqrt x$, $0\le x\le4$, diputar terhadap sumbu-$x$. (ii) Panjang kurva $y=\tfrac23 x^{3/2}$ dari $x=0$ sampai $x=3$.

\jawab (i) $V=\pi\displaystyle\int_0^4 (\sqrt x)^2\,dx=\pi\int_0^4 x\,dx=\pi\left[\tfrac{x^2}{2}\right]_0^4=8\pi$.\\
(ii) $y'=x^{1/2}$, $1+(y')^2=1+x$; $L=\displaystyle\int_0^3\sqrt{1+x}\,dx=\left[\tfrac23(1+x)^{3/2}\right]_0^3=\tfrac23(8-1)=\tfrac{14}{3}$.
\end{contoh}

\begin{latihan}
(1) Volume putar $y=x$, $0\le x\le1$, terhadap sumbu-$x$.\quad (2) Volume putar $y=x^2$, $0\le x\le1$, terhadap sumbu-$x$.\quad (3) Panjang kurva $y=\tfrac23 x^{3/2}$ dari $x=0$ sampai $x=1$.
\end{latihan}

\section*{Ringkasan Bab}
\addcontentsline{toc}{section}{Ringkasan Bab}
\begin{itemize}[leftmargin=*,topsep=2pt]
  \item $\int x^n dx=\tfrac{x^{n+1}}{n+1}+C$; jangan lupa konstanta $C$.
  \item Integral tentu $\int_a^b f=F(b)-F(a)$ mengukur luas bertanda.
  \item \textbf{Substitusi:} $\int f(g)g'\,dx=\int f(u)\,du$; \textbf{parsial:} $\int u\,dv=uv-\int v\,du$.
  \item \textbf{Trigonometri} (identitas setengah sudut) dan \textbf{rasional} (pecahan parsial).
  \item \textbf{Luas} antar kurva $=\int(\text{atas}-\text{bawah})$; \textbf{volume putar} $=\pi\int y^2\,dx$; \textbf{panjang kurva} $=\int\sqrt{1+(y')^2}\,dx$.
  \item \textbf{TDK:} turunan dan integral saling membalik.
\end{itemize}

\begin{refleksi}
\begin{itemize}[leftmargin=*,topsep=2pt]
  \item[\cek] Saya dapat menghitung integral tak tentu dan tentu.
  \item[\cek] Saya dapat memakai substitusi sederhana.
  \item[\cek] Saya dapat menghitung luas daerah dan volume putar.
  \item[\cek] Saya dapat menjelaskan Teorema Dasar Kalkulus.
\end{itemize}
\end{refleksi}

\section*{Soal Akhir Bab \thechapter}
\addcontentsline{toc}{section}{Soal Akhir Bab \thechapter}

\bagiansoal{A}{(Fundamental --- setara SNBT Dasar)}
\begin{enumerate}[leftmargin=*]
  \item Hitung $\displaystyle\int x^2\,dx$.
  \item Hitung $\displaystyle\int (2x+1)\,dx$.
  \item Hitung $\displaystyle\int_0^2 (2x+1)\,dx$.
  \item Hitung $\displaystyle\int_1^2 x^2\,dx$.
  \item Hitung $\displaystyle\int 3x^2\,dx$.
\end{enumerate}

\bagiansoal{B}{(Intermediate --- setara SNBT Sulit)}
\begin{enumerate}[leftmargin=*]
  \item Hitung $\displaystyle\int (x^3-2x^2+5)\,dx$.
  \item Hitung $\displaystyle\int_0^1 (3x^2+2x)\,dx$.
  \item Hitung $\displaystyle\int \frac{x^2+1}{\sqrt{x}}\,dx$.
  \item Hitung $\displaystyle\int_0^{\pi} \sin x\,dx$.
  \item Hitung $\displaystyle\int 2x(x^2+1)^3\,dx$ dengan substitusi.
  \item Hitung luas daerah di bawah $y=x^2$ dari $x=0$ sampai $x=3$.
\end{enumerate}

\bagiansoal{C}{(Advanced --- setara Olimpiade SMA, gabungan antarbab)}
\begin{enumerate}[leftmargin=*]
  \item \textbf{(Integral + Luas Antar Kurva).} Hitung luas daerah antara $y=x^2$ dan $y=x$ untuk $0\le x\le1$.
  \item \textbf{(Substitusi + Eksponen).} Hitung $\displaystyle\int x\,e^{x^2}\,dx$.
  \item \textbf{(Teorema Dasar).} Hitung $\dfrac{d}{dx}\displaystyle\int_0^x \sin(t^2)\,dt$.
  \item \textbf{(Volume Putar).} Daerah di bawah $y=\sqrt{x}$ dari $x=0$ sampai $x=4$ diputar mengelilingi sumbu-$x$. Tentukan volumenya.
  \item \textbf{(Integral + Simetri).} Hitung $\displaystyle\int_{-2}^{2} x^3\,dx$ dan jelaskan dengan kesimetrian.
  \item \textbf{(Integral + Kinematika).} Kecepatan $v(t)=3t^2-6t$. Jika $s(0)=2$, tentukan fungsi posisi $s(t)$.
\end{enumerate}

\bagiansoal{D}{(Expert --- setara tahun pertama universitas, sintesis \& pembuktian)}
\begin{enumerate}[leftmargin=*]
  \item \textbf{(Teorema Dasar Kalkulus).} Nyatakan kedua bagian Teorema Dasar Kalkulus dan jelaskan mengapa integral merupakan kebalikan turunan.
  \item \textbf{(Substitusi Trigonometri).} Hitung $\displaystyle\int_0^{\pi/2} \sin x\cos x\,dx$.
  \item \textbf{(Sifat Integral).} Buktikan $\displaystyle\int_a^b f=-\int_b^a f$ dan $\displaystyle\int_a^a f=0$ dari definisi $F(b)-F(a)$.
  \item \textbf{(Volume Bola).} Turunkan rumus volume bola berjari-jari $R$, yaitu $\tfrac43\pi R^3$, dengan integral putar.
\end{enumerate}

\begin{challenge}
\begin{enumerate}[leftmargin=*,topsep=2pt]
  \item Hitung $\displaystyle\int_0^1 \frac{x}{1+x^2}\,dx$.
  \item Buktikan $\displaystyle\int_0^a f(x)\,dx=\int_0^a f(a-x)\,dx$, lalu gunakan untuk menghitung $\displaystyle\int_0^{\pi/2}\frac{\sin x}{\sin x+\cos x}\,dx$.
\end{enumerate}
\end{challenge}

% ============================================================
\chapter{Distribusi Binomial}
% ============================================================
\begin{tujuan}
Setelah mempelajari bab ini, peserta didik mampu:
\begin{itemize}[leftmargin=*,topsep=2pt]
  \item mengenali ciri percobaan Bernoulli dan percobaan binomial;
  \item menghitung peluang melalui rumus distribusi binomial;
  \item menentukan nilai harapan dan variansi distribusi binomial;
  \item menerapkan distribusi binomial pada masalah nyata.
\end{itemize}
\textbf{Capaian Pembelajaran (Fase F).} Peserta didik dapat memodelkan banyaknya keberhasilan dalam percobaan berulang menggunakan distribusi binomial beserta nilai harapan dan variansinya.
\end{tujuan}

\begin{pemantik}
\begin{itemize}[leftmargin=*,topsep=2pt]
  \item Jika menebak acak $10$ soal pilihan ganda (4 opsi), berapa soal yang "diharapkan" benar?
  \item Mengapa peluang "tepat $5$ gambar dari $10$ lemparan" tidak sama dengan $\tfrac12$, padahal koinnya adil?
  \item Kapan suatu percobaan boleh dimodelkan sebagai binomial, dan kapan tidak?
\end{itemize}
\end{pemantik}

\begin{studikasus}{Kendali Mutu Produksi}
Sebuah pabrik tahu bahwa $5\%$ produknya cacat. Setiap hari, petugas mengambil sampel $20$ produk dan mencatat banyaknya yang cacat.
\smallskip

\textbf{Amati dan duga.} Banyaknya produk cacat dalam sampel berubah-ubah tiap hari --- inilah \emph{variabel acak}. Tiap produk hanya punya dua kemungkinan (cacat/tidak), peluang cacat tetap $5\%$, dan antarpemeriksaan saling bebas. Pola tepat seperti ini --- $n$ percobaan bebas dengan dua hasil dan peluang tetap --- dimodelkan oleh \textbf{distribusi binomial}, yang memungkinkan kita menghitung "berapa peluang menemukan tepat $2$ cacat" atau "rata-rata berapa cacat per hari".
\end{studikasus}

\section*{Peta Konsep Bab}
\addcontentsline{toc}{section}{Peta Konsep Bab}
\begin{center}
\begin{tikzpicture}[
  node distance=8mm,
  konsepbox/.style={draw=biru,fill=birumuda,rounded corners,align=center,inner sep=5pt,font=\small\bfseries},
  subbox/.style={draw=hijau,fill=hijaumuda,rounded corners,align=center,inner sep=4pt,font=\footnotesize},
  garis/.style={-Stealth,biru,thick}]
  \node[konsepbox] (akar) {DISTRIBUSI BINOMIAL};
  \node[subbox,below left=12mm and 16mm of akar] (ber) {Percobaan\\Bernoulli};
  \node[subbox,below=12mm of akar] (rum) {Rumus\\$P(X{=}k)$};
  \node[subbox,below right=12mm and 16mm of akar] (ev) {$E(X)$ \&\\Var$(X)$};
  \node[subbox,below=10mm of rum] (apl) {Aplikasi};
  \draw[garis] (akar)--(ber); \draw[garis] (akar)--(rum); \draw[garis] (akar)--(ev);
  \draw[garis] (rum)--(apl);
\end{tikzpicture}\\
\small \textbf{Peta konsep.} Dari satu percobaan Bernoulli menuju $n$ percobaan (binomial), lalu nilai harapan, variansi, dan aplikasinya.
\end{center}

\section{Percobaan Bernoulli dan Distribusi Binomial}
\begin{definisi}
\textbf{Percobaan Bernoulli} adalah percobaan dengan tepat dua hasil: "sukses" (peluang $p$) dan "gagal" (peluang $q=1-p$). \textbf{Percobaan binomial} adalah $n$ percobaan Bernoulli yang \emph{bebas} dengan $p$ tetap. Jika $X$ menyatakan banyaknya sukses:
\[
P(X=k)=\binom nk p^k q^{\,n-k},\qquad k=0,1,2,\dots,n.
\]
\end{definisi}

\begin{konsep}{Empat Syarat Percobaan Binomial}
\begin{enumerate}[leftmargin=*,topsep=2pt,label=(\arabic*)]
  \item banyaknya percobaan $n$ tetap;
  \item tiap percobaan hanya dua hasil (sukses/gagal);
  \item peluang sukses $p$ tetap di setiap percobaan;
  \item antarpercobaan saling bebas.
\end{enumerate}
\end{konsep}

\begin{contoh}{Peluang Binomial}
Koin adil dilempar $3$ kali. Tentukan peluang muncul tepat $2$ gambar.

\jawab $n=3$, $p=\tfrac12$, $k=2$: $P(X=2)=\binom32\left(\tfrac12\right)^2\left(\tfrac12\right)^1=3\cdot\tfrac18=\tfrac38$.
\end{contoh}

\section{Nilai Harapan dan Variansi}
\begin{konsep}{Rata-rata dan Sebaran}
Untuk $X\sim B(n,p)$:
\[
E(X)=np,\qquad \operatorname{Var}(X)=npq,\qquad \sigma=\sqrt{npq}.
\]
$E(X)$ adalah rata-rata banyaknya sukses dalam jangka panjang; $\sigma$ mengukur sebarannya.
\end{konsep}

\begin{contoh}{Nilai Harapan dan Variansi}
Sebuah dadu dilempar $60$ kali; "sukses" $=$ muncul mata $6$. Tentukan $E(X)$ dan $\operatorname{Var}(X)$.

\jawab $p=\tfrac16$, $q=\tfrac56$, $n=60$. $E(X)=60\cdot\tfrac16=10$; $\operatorname{Var}(X)=60\cdot\tfrac16\cdot\tfrac56=\tfrac{300}{36}\approx8{,}33$.
\end{contoh}

\begin{kesalahan}
Distribusi binomial mensyaratkan $p$ \emph{tetap} dan percobaan \emph{bebas}. Pengambilan \emph{tanpa pengembalian} dari populasi kecil melanggar syarat ini (peluang berubah tiap pengambilan) --- modelnya hipergeometrik, bukan binomial.
\end{kesalahan}

\section{Aplikasi Distribusi Binomial}
\begin{contoh}{Aplikasi Kendali Mutu}
Sebuah mesin menghasilkan $5\%$ produk cacat. Dari sampel $20$ produk, tentukan peluang tepat $1$ produk cacat.

\jawab $n=20$, $p=0{,}05$, $k=1$: $P(X=1)=\binom{20}{1}(0{,}05)^1(0{,}95)^{19}=20\cdot0{,}05\cdot(0{,}95)^{19}\approx0{,}377$.
\end{contoh}

\begin{dunianyata}
Distribusi binomial memodelkan kendali mutu, uji klinis (banyak pasien sembuh), survei (banyak responden setuju), keandalan sistem, dan A/B testing pada produk digital. Ia juga fondasi untuk uji hipotesis proporsi.
\end{dunianyata}

\section*{Ringkasan Bab}
\addcontentsline{toc}{section}{Ringkasan Bab}
\begin{itemize}[leftmargin=*,topsep=2pt]
  \item Bernoulli: dua hasil, sukses $p$, gagal $q=1-p$.
  \item Binomial $X\sim B(n,p)$: $P(X=k)=\binom nk p^k q^{n-k}$.
  \item Syarat: $n$ tetap, dua hasil, $p$ tetap, percobaan bebas.
  \item $E(X)=np$, $\operatorname{Var}(X)=npq$, $\sigma=\sqrt{npq}$.
\end{itemize}

\begin{refleksi}
\begin{itemize}[leftmargin=*,topsep=2pt]
  \item[\cek] Saya dapat mengenali percobaan binomial dan syaratnya.
  \item[\cek] Saya dapat menghitung $P(X=k)$ dengan rumus binomial.
  \item[\cek] Saya dapat menghitung $E(X)$, $\operatorname{Var}(X)$, dan $\sigma$.
  \item[\cek] Saya dapat menerapkan distribusi binomial pada masalah nyata.
\end{itemize}
\end{refleksi}

\section*{Soal Akhir Bab \thechapter}
\addcontentsline{toc}{section}{Soal Akhir Bab \thechapter}

\bagiansoal{A}{(Fundamental --- setara SNBT Dasar)}
\begin{enumerate}[leftmargin=*]
  \item Pada percobaan Bernoulli dengan $p=0{,}3$, tentukan $q$.
  \item Tentukan $P(X=2)$ untuk $n=4$, $p=0{,}5$.
  \item Hitung $E(X)$ untuk $n=10$, $p=0{,}2$.
  \item Hitung $\operatorname{Var}(X)$ untuk $n=10$, $p=0{,}2$.
  \item Koin adil dilempar $3$ kali. Tentukan peluang tepat $2$ gambar.
  \item Tentukan $\sigma$ untuk $n=100$, $p=0{,}5$.
  \item Tentukan $P(X=0)$ untuk $n=5$, $p=0{,}2$.
  \item Sebutkan empat syarat percobaan binomial.
  \item Dadu dilempar $60$ kali; sukses $=$ muncul $6$. Hitung $E(X)$.
  \item Tuliskan rumus $P(X=k)$ distribusi binomial.
\end{enumerate}

\bagiansoal{B}{(Intermediate --- setara SNBT Sulit)}
\begin{enumerate}[leftmargin=*]
  \item Dadu dilempar $5$ kali. Tentukan peluang muncul angka $6$ tepat $2$ kali.
  \item Peluang sukses $0{,}4$, $n=6$. Tentukan $P(X\ge1)$.
  \item Ujian $10$ soal pilihan ganda ($4$ opsi) ditebak acak. Tentukan nilai harapan banyak jawaban benar.
  \item Tentukan $P(X=4)$ untuk $n=8$, $p=0{,}5$.
  \item Tentukan $\operatorname{Var}(X)$ dan $\sigma$ untuk $n=50$, $p=0{,}4$.
\end{enumerate}

\bagiansoal{C}{(Advanced --- setara Olimpiade SMA, gabungan antarbab)}
\begin{enumerate}[leftmargin=*]
  \item \textbf{(Binomial + Komplemen).} Peluang lahir anak laki-laki $0{,}5$. Pada keluarga $5$ anak, tentukan peluang minimal $1$ anak perempuan.
  \item \textbf{(Binomial + Pencacahan).} Tunjukkan bahwa $\displaystyle\sum_{k=0}^{n}P(X=k)=1$.
  \item \textbf{(Binomial + Aplikasi).} Mesin menghasilkan $5\%$ cacat; diambil $20$ produk. Tentukan peluang tepat $1$ cacat.
  \item \textbf{(Binomial + Bersyarat).} Empat koin adil dilempar. Diketahui muncul minimal $1$ gambar. Tentukan peluang muncul tepat $2$ gambar.
  \item \textbf{(Binomial + Nilai Harapan).} Tiga koin dilempar; pemain dibayar Rp$1000$ tiap gambar. Tentukan nilai harapan pembayaran.
\end{enumerate}

\bagiansoal{D}{(Expert --- setara tahun pertama universitas, sintesis \& pembuktian)}
\begin{enumerate}[leftmargin=*]
  \item \textbf{(Pembuktian).} Buktikan $E(X)=np$ untuk $X\sim B(n,p)$.
  \item \textbf{(Pembuktian).} Buktikan $\operatorname{Var}(X)=npq$ (mis. dengan menulis $X$ sebagai jumlah peubah Bernoulli bebas).
  \item \textbf{(Aproksimasi Poisson).} Untuk $n=1000$, $p=0{,}001$, gunakan $\lambda=np$ untuk menaksir $P(X=0)$ via $P(X=k)\approx e^{-\lambda}\dfrac{\lambda^k}{k!}$.
  \item \textbf{(Aproksimasi Normal).} Untuk $n=100$, $p=0{,}5$, taksir $P(X=50)$ menggunakan puncak kurva normal $\dfrac{1}{\sigma\sqrt{2\pi}}$.
  \item \textbf{(Modus).} Tentukan nilai $k$ yang memaksimalkan $P(X=k)$ untuk $n=10$, $p=0{,}3$.
\end{enumerate}

\begin{challenge}
\begin{enumerate}[leftmargin=*,topsep=2pt]
  \item Berapa kali minimal sebuah koin adil harus dilempar agar peluang muncul \emph{minimal satu} gambar mencapai sedikitnya $0{,}99$?
  \item Pada keluarga dengan $n$ anak ($n$ genap, peluang laki/perempuan sama), tentukan peluang banyak anak laki-laki sama dengan banyak anak perempuan.
  \item Buktikan $\dfrac{P(X=k+1)}{P(X=k)}=\dfrac{n-k}{k+1}\cdot\dfrac{p}{q}$, lalu jelaskan bagaimana rasio ini menentukan letak modus distribusi.
\end{enumerate}
\end{challenge}

% ============================================================
\chapter{Distribusi Normal dan Uji Hipotesis}
% ============================================================
\begin{tujuan}
Setelah mempelajari bab ini, peserta didik mampu:
\begin{itemize}[leftmargin=*,topsep=2pt]
  \item mengenali ciri kurva normal serta peran mean dan simpangan baku;
  \item menghitung dan memaknai skor-$z$ (z-score) serta membaca luas di bawah kurva;
  \item menafsirkan kurva normal menggunakan aturan empiris $68$--$95$--$99{,}7$;
  \item merumuskan hipotesis nol dan alternatif;
  \item menghitung statistik uji dan mengambil keputusan uji hipotesis.
\end{itemize}
\textbf{Capaian Pembelajaran (Fase F).} Peserta didik dapat menggunakan distribusi normal untuk menafsirkan data dan melakukan uji hipotesis sederhana.
\end{tujuan}

\begin{pemantik}
\begin{itemize}[leftmargin=*,topsep=2pt]
  \item Mengapa banyak sekali data alami (tinggi badan, nilai ujian, kesalahan pengukuran) membentuk kurva "lonceng" yang mirip?
  \item Bagaimana membandingkan nilai $80$ pada ujian Matematika dengan $80$ pada ujian Fisika yang tingkat kesulitannya berbeda?
  \item Ketika sebuah pabrik mengklaim "isi rata-rata $500$ mL", bagaimana data sampel dapat membuktikan atau membantah klaim itu?
\end{itemize}
\end{pemantik}

\begin{studikasus}{Membandingkan Prestasi Lewat Skor-$z$}
Dina memperoleh $80$ pada ujian Matematika ($\mu=70,\ \sigma=5$) dan $85$ pada Fisika ($\mu=80,\ \sigma=10$). Pada ujian mana ia relatif lebih unggul?
\smallskip

\textbf{Amati dan duga.} Nilai mentah saja menyesatkan. Hitung posisi relatif: Matematika $z=\tfrac{80-70}{5}=2$; Fisika $z=\tfrac{85-80}{10}=0{,}5$. Ternyata $80$ di Matematika ($2$ simpangan baku di atas rata-rata) jauh lebih istimewa daripada $85$ di Fisika. \emph{Skor-$z$} menstandarkan data agar dapat dibandingkan --- ini kunci membaca distribusi normal.
\end{studikasus}

\section*{Peta Konsep Bab}
\addcontentsline{toc}{section}{Peta Konsep Bab}
\begin{center}
\begin{tikzpicture}[
  node distance=8mm,
  konsepbox/.style={draw=biru,fill=birumuda,rounded corners,align=center,inner sep=5pt,font=\small\bfseries},
  subbox/.style={draw=hijau,fill=hijaumuda,rounded corners,align=center,inner sep=4pt,font=\footnotesize},
  garis/.style={-Stealth,biru,thick}]
  \node[konsepbox] (akar) {NORMAL \&\ UJI HIPOTESIS};
  \node[subbox,below left=12mm and 20mm of akar] (kn) {Kurva normal\\\&\ skor-$z$};
  \node[subbox,below=12mm of akar] (in) {Interpretasi\\\&\ aturan empiris};
  \node[subbox,below right=12mm and 20mm of akar] (uh) {Uji hipotesis};
  \node[subbox,below=10mm of uh] (kep) {Statistik uji\\\&\ keputusan};
  \draw[garis] (akar)--(kn); \draw[garis] (akar)--(in); \draw[garis] (akar)--(uh);
  \draw[garis] (uh)--(kep);
\end{tikzpicture}\\
\small \textbf{Peta konsep.} Kurva normal dan skor-$z$ menjadi dasar interpretasi data hingga pengambilan keputusan uji hipotesis.
\end{center}

\section{Kurva Normal dan Skor-$z$}
\begin{konsep}{Distribusi Normal}
Distribusi normal $N(\mu,\sigma^2)$ berbentuk lonceng simetris terhadap \textbf{mean} $\mu$ (yang sekaligus median dan modus); lebar sebarannya ditentukan \textbf{simpangan baku} $\sigma$. Luas total di bawah kurva $=1$. \textbf{Distribusi normal baku} adalah $N(0,1)$, dilambangkan $Z$.
\end{konsep}

\begin{konsep}{Skor-$z$ (Standarisasi)}
Setiap nilai $x$ diubah ke skor baku:
\[
z=\frac{x-\mu}{\sigma}.
\]
Skor-$z$ menyatakan "berapa simpangan baku" $x$ dari mean. $z>0$ di atas rata-rata, $z<0$ di bawahnya. Luas $P(Z<z)$ dibaca dari tabel normal baku.
\end{konsep}

\begin{center}
\begin{tikzpicture}
\begin{axis}[axis lines=middle,width=11cm,height=5cm,
  xlabel={$z$},ymin=0,ymax=0.45,xmin=-3.6,xmax=3.6,
  xtick={-3,-2,-1,0,1,2,3},ytick=\empty,samples=100,domain=-3.6:3.6]
\addplot[very thick,biru]{exp(-x^2/2)/sqrt(2*pi)};
\addplot[draw=none,fill=birumuda,opacity=0.6,domain=-1:1]{exp(-x^2/2)/sqrt(2*pi)}\closedcycle;
\node at (axis cs:0,0.16){\scriptsize $68\%$};
\end{axis}
\end{tikzpicture}\\
\small Kurva normal baku. Daerah $\pm1\sigma$ memuat $\approx68\%$ data.
\end{center}

\section{Interpretasi Kurva: Aturan Empiris}
\begin{konsep}{Aturan $68$--$95$--$99{,}7$}
Pada distribusi normal:
\[
P(\mu-\sigma<X<\mu+\sigma)\approx68\%,\quad
P(\mu-2\sigma<X<\mu+2\sigma)\approx95\%,\quad
P(\mu-3\sigma<X<\mu+3\sigma)\approx99{,}7\%.
\]
Nilai tabel yang sering dipakai: $P(Z<1)=0{,}8413$, $P(Z<1{,}96)=0{,}975$, $P(Z<2)=0{,}9772$.
\end{konsep}

\begin{contoh}{Membaca Luas}
Nilai ujian berdistribusi normal $\mu=75$, $\sigma=10$. Berapa persen siswa bernilai di atas $85$?

\jawab $z=\tfrac{85-75}{10}=1$. $P(X>85)=P(Z>1)=1-0{,}8413=0{,}1587$, yaitu $\approx15{,}87\%$.
\end{contoh}

\section{Uji Hipotesis}
\begin{konsep}{Kerangka Uji Hipotesis}
\begin{itemize}[leftmargin=*,topsep=2pt]
  \item \textbf{Hipotesis nol} $H_0$: pernyataan "status quo" (mis. $\mu=\mu_0$).
  \item \textbf{Hipotesis alternatif} $H_1$: klaim yang ingin dibuktikan ($\mu\neq\mu_0$, $\mu>\mu_0$, atau $\mu<\mu_0$).
  \item \textbf{Statistik uji} (untuk rata-rata, $\sigma$ diketahui): $\displaystyle z=\frac{\bar x-\mu_0}{\sigma/\sqrt n}$.
  \item \textbf{Taraf nyata} $\alpha$ (mis. $0{,}05$) menentukan nilai kritis: dua arah $\pm1{,}96$; satu arah $1{,}645$.
  \item \textbf{Keputusan:} jika statistik uji jatuh di daerah kritis (atau $p\text{-value}<\alpha$), \emph{tolak} $H_0$; selain itu \emph{gagal menolak} $H_0$.
\end{itemize}
\end{konsep}

\begin{contoh}{Uji Hipotesis Dua Arah}
Pabrik mengklaim isi rata-rata botol $500$ mL. Sampel $n=25$ memberi $\bar x=495$ mL, $\sigma=10$, $\alpha=0{,}05$. Ujilah.

\jawab $H_0:\mu=500$, $H_1:\mu\neq500$. Statistik uji $z=\dfrac{495-500}{10/\sqrt{25}}=\dfrac{-5}{2}=-2{,}5$. Karena $|-2{,}5|>1{,}96$, $z$ jatuh di daerah kritis $\Rightarrow$ \textbf{tolak} $H_0$: isi rata-rata berbeda signifikan dari $500$ mL.
\end{contoh}

\begin{kesalahan}
"Gagal menolak $H_0$" bukan berarti "$H_0$ terbukti benar" --- hanya bukti belum cukup untuk menolaknya. Selain itu, $\alpha=0{,}05$ berarti masih ada $5\%$ peluang menolak $H_0$ yang sebenarnya benar (galat tipe I).
\end{kesalahan}

\begin{dunianyata}
Distribusi normal dan uji hipotesis adalah tulang punggung penelitian ilmiah, uji obat (uji klinis), kendali mutu industri, A/B testing produk digital, dan penilaian standar (skor terstandar seperti TOEFL/SAT). "Signifikansi statistik" dalam berita riset berakar pada konsep ini.
\end{dunianyata}

\section*{Ringkasan Bab}
\addcontentsline{toc}{section}{Ringkasan Bab}
\begin{itemize}[leftmargin=*,topsep=2pt]
  \item Kurva normal: lonceng simetris terhadap $\mu$, luas total $1$; baku $N(0,1)$.
  \item Skor-$z$: $z=\dfrac{x-\mu}{\sigma}$ untuk membandingkan dan membaca tabel.
  \item Aturan empiris: $68$--$95$--$99{,}7$ untuk $\pm1,\pm2,\pm3\,\sigma$.
  \item Uji hipotesis: $H_0$ vs $H_1$, statistik uji $z=\dfrac{\bar x-\mu_0}{\sigma/\sqrt n}$, bandingkan nilai kritis ($\pm1{,}96$ dua arah; $1{,}645$ satu arah).
\end{itemize}

\begin{refleksi}
\begin{itemize}[leftmargin=*,topsep=2pt]
  \item[\cek] Saya memahami ciri kurva normal serta peran $\mu$ dan $\sigma$.
  \item[\cek] Saya dapat menghitung dan menafsirkan skor-$z$.
  \item[\cek] Saya dapat menggunakan aturan empiris dan tabel normal.
  \item[\cek] Saya dapat merumuskan hipotesis, menghitung statistik uji, dan mengambil keputusan.
\end{itemize}
\end{refleksi}

\section*{Soal Akhir Bab \thechapter}
\addcontentsline{toc}{section}{Soal Akhir Bab \thechapter}

\bagiansoal{A}{(Fundamental --- setara SNBT Dasar)}
\begin{enumerate}[leftmargin=*]
  \item Berapa luas total di bawah kurva normal?
  \item Hitung skor-$z$ untuk $x=80$, $\mu=70$, $\sigma=5$.
  \item Menurut aturan empiris, berapa persen data dalam $\pm1\sigma$ dari mean?
  \item Berapa persen data dalam $\pm2\sigma$?
  \item Sebutkan mean dan simpangan baku distribusi normal baku.
  \item Berapa skor-$z$ untuk nilai yang sama dengan mean?
  \item Kurva normal simetris terhadap apa?
  \item Jika $x<\mu$, bagaimana tanda skor-$z$?
  \item Tentukan $P(Z<0)$.
  \item Tuliskan rumus skor-$z$.
\end{enumerate}

\bagiansoal{B}{(Intermediate --- setara SNBT Sulit)}
\begin{enumerate}[leftmargin=*]
  \item Nilai ujian $\mu=70$, $\sigma=8$. Hitung skor-$z$ untuk $x=86$.
  \item Tinggi badan $\mu=160$, $\sigma=6$. Berapa persen orang dengan tinggi antara $154$ dan $166$ cm?
  \item Diketahui $P(Z<1)=0{,}8413$. Tentukan $P(Z>1)$.
  \item Tentukan $P(0<Z<1)$.
  \item Skor terstandar $\mu=500$, $\sigma=100$. Hitung skor-$z$ untuk skor $650$.
\end{enumerate}

\bagiansoal{C}{(Advanced --- setara Olimpiade SMA, gabungan antarbab)}
\begin{enumerate}[leftmargin=*]
  \item \textbf{(Normal + Interpretasi).} Nilai $\mu=75$, $\sigma=10$. Tentukan persentase siswa bernilai di atas $85$.
  \item \textbf{(Normal + Binomial).} Koin adil dilempar $100$ kali. Dengan pendekatan normal, taksir $P(X\ge60)$.
  \item \textbf{(Normal + Persentil).} Nilai $\mu=70$, $\sigma=10$. Tentukan batas nilai persentil ke-$90$ (gunakan $z\approx1{,}28$).
  \item \textbf{(Uji Hipotesis --- Perumusan).} Sebuah klaim menyatakan rata-rata berat produk $=500$ g; peneliti menduga berat sebenarnya \emph{kurang}. Rumuskan $H_0$ dan $H_1$.
  \item \textbf{(Uji Hipotesis --- Statistik).} Sampel $n=36$, $\bar x=52$, $\mu_0=50$, $\sigma=6$. Hitung statistik uji $z$.
\end{enumerate}

\bagiansoal{D}{(Expert --- setara tahun pertama universitas, sintesis \& pembuktian)}
\begin{enumerate}[leftmargin=*]
  \item \textbf{(Uji Dua Arah Lengkap).} Pabrik klaim isi rata-rata $500$ mL. Sampel $n=25$, $\bar x=495$, $\sigma=10$, $\alpha=0{,}05$. Hitung $z$, bandingkan $\pm1{,}96$, dan ambil keputusan.
  \item \textbf{(Standarisasi).} Buktikan bahwa jika $X\sim N(\mu,\sigma^2)$ maka $Z=\dfrac{X-\mu}{\sigma}$ memiliki mean $0$ dan variansi $1$.
  \item \textbf{(Uji Satu Arah + $p$-value).} Metode baru diklaim meningkatkan skor dari $\mu_0=70$. Sampel $n=49$, $\bar x=72$, $\sigma=7$, $\alpha=0{,}05$. Hitung $z$, $p$-value, dan simpulkan.
  \item \textbf{(Galat).} Jelaskan galat tipe I dan tipe II dalam uji hipotesis serta kaitannya dengan $\alpha$.
  \item \textbf{(Interval Kepercayaan).} Sampel $n=100$, $\bar x=50$, $\sigma=8$. Susun interval kepercayaan $95\%$ untuk $\mu$.
\end{enumerate}

\begin{challenge}
\begin{enumerate}[leftmargin=*,topsep=2pt]
  \item Menurut aturan empiris, berapa persen data berada \emph{di luar} $\pm3\sigma$?
  \item Skor IQ berdistribusi normal $\mu=100$, $\sigma=15$. Berapa IQ minimum untuk masuk $2\%$ teratas (gunakan $z\approx2{,}05$)?
  \item Jelaskan mengapa binomial dengan $n$ besar dapat dihampiri $N(np,npq)$, lalu taksir $P(40\le X\le60)$ untuk $n=100$, $p=0{,}5$.
\end{enumerate}
\end{challenge}

% ============================================================
\chapter{Dimensi Tiga}
% ============================================================
\begin{tujuan}
Setelah mempelajari bab ini, peserta didik mampu:
\begin{itemize}[leftmargin=*,topsep=2pt]
  \item mengenali kedudukan titik, garis, dan bidang dalam ruang;
  \item menghitung jarak titik ke titik, titik ke garis, dan titik ke bidang;
  \item menentukan sudut antara dua garis, garis dan bidang, serta dua bidang;
  \item menerapkan teorema Pythagoras dan trigonometri pada bangun ruang.
\end{itemize}
\textbf{Capaian Pembelajaran (Fase F).} Peserta didik dapat menyelesaikan masalah jarak dan sudut dalam bangun ruang dimensi tiga.
\end{tujuan}

\begin{pemantik}
\begin{itemize}[leftmargin=*,topsep=2pt]
  \item Bagaimana mengukur "jarak" sebuah lampu di langit-langit ke salah satu rusuk lantai?
  \item Mengapa dua garis dalam ruang bisa tidak sejajar \emph{dan} tidak berpotongan sekaligus?
  \item Saat memasang rak miring, sudut apa yang sebenarnya kita atur --- sudut garis-bidang atau bidang-bidang?
\end{itemize}
\end{pemantik}

\begin{studikasus}{Kubus sebagai Laboratorium Ruang}
Hampir semua masalah dimensi tiga di sekolah dapat "dipasang" pada sebuah kubus $ABCD.EFGH$. Dengan menempatkan kubus pada sistem koordinat ($A$ di titik asal), setiap titik punya koordinat, sehingga jarak dan sudut dapat dihitung dengan Pythagoras dan vektor.
\begin{center}
\begin{tikzpicture}[scale=1.4,line join=round]
\coordinate (A) at (0,0); \coordinate (B) at (2,0);
\coordinate (C) at (2.8,0.8); \coordinate (D) at (0.8,0.8);
\coordinate (E) at (0,2); \coordinate (F) at (2,2);
\coordinate (G) at (2.8,2.8); \coordinate (H) at (0.8,2.8);
\draw[thick] (A)--(B)--(F)--(E)--cycle;
\draw[thick] (B)--(C)--(G)--(F);
\draw[thick] (E)--(H)--(G);
\draw[dashed] (A)--(D)--(C);
\draw[dashed] (D)--(H);
\foreach \p/\pos in {A/below left,B/below right,C/right,D/above,E/left,F/right,G/right,H/above}
  \fill (\p) circle (1.2pt) node[\pos,font=\scriptsize]{$\p$};
\end{tikzpicture}
\end{center}
\textbf{Amati dan duga.} Pada kubus rusuk $a$: berapa panjang diagonal sisi $AC$? Berapa diagonal ruang $AG$? Kedua nilai ($a\sqrt2$ dan $a\sqrt3$) menjadi "alat ukur" utama sepanjang bab ini.
\end{studikasus}

\section*{Peta Konsep Bab}
\addcontentsline{toc}{section}{Peta Konsep Bab}
\begin{center}
\begin{tikzpicture}[
  node distance=8mm,
  konsepbox/.style={draw=biru,fill=birumuda,rounded corners,align=center,inner sep=5pt,font=\small\bfseries},
  subbox/.style={draw=hijau,fill=hijaumuda,rounded corners,align=center,inner sep=4pt,font=\footnotesize},
  garis/.style={-Stealth,biru,thick}]
  \node[konsepbox] (akar) {DIMENSI TIGA};
  \node[subbox,below left=12mm and 18mm of akar] (ked) {Titik, garis,\\bidang \&\ kedudukan};
  \node[subbox,below=12mm of akar] (jar) {Jarak};
  \node[subbox,below right=12mm and 18mm of akar] (sud) {Sudut};
  \node[subbox,below=10mm of jar] (kub) {Model kubus\\\&\ koordinat};
  \draw[garis] (akar)--(ked); \draw[garis] (akar)--(jar); \draw[garis] (akar)--(sud);
  \draw[garis] (jar)--(kub);
\end{tikzpicture}\\
\small \textbf{Peta konsep.} Kedudukan unsur ruang menuntun perhitungan jarak dan sudut, dengan kubus sebagai model utama.
\end{center}

\begin{catatan}
\textbf{Pengingat Materi Kelas 10 (MAE Bab 3 --- Vektor).} Dimensi tiga menggunakan konsep vektor secara eksplisit. Pastikan kamu ingat:
\begin{itemize}[leftmargin=*,topsep=2pt]
  \item Representasi vektor sebagai komponen: $\vec a=(a_1,a_2)$ di 2D, yang akan dikembangkan ke $\vec a=(a_1,a_2,a_3)$ di 3D.
  \item Panjang vektor dari Pythagoras: $|\vec a|=\sqrt{a_1^2+a_2^2}$ (di 2D).
  \item Operasi vektor: penjumlahan, pengurangan, perkalian skalar, dan sifat-sifatnya.
  \item Dot product: $\vec a\cdot\vec b=|\vec a||\vec b|\cos\theta$ --- digunakan langsung untuk menghitung sudut antara dua ruas garis (atau bidang) di Dimensi Tiga.
\end{itemize}
\end{catatan}

\section{Titik, Garis, Bidang, dan Kedudukannya}
\begin{konsep}{Tiga Unsur Dasar: Titik, Garis, dan Bidang}
\begin{itemize}[leftmargin=*,topsep=2pt]
  \item \textbf{Titik} --- menyatakan \emph{posisi} tanpa panjang, lebar, atau tinggi (tak berukuran). Pada kubus: titik-titik sudut $A,B,\dots,H$, tetapi juga titik tengah rusuk, pusat bidang, dan pusat kubus.
  \item \textbf{Garis} --- himpunan titik yang memanjang lurus tak terbatas ke \emph{dua} arah. Bagian garis di antara dua titik disebut \textbf{ruas garis}. Setiap \emph{rusuk} kubus adalah ruas garis.
  \item \textbf{Bidang} --- permukaan datar yang meluas tak terbatas ke segala arah. Setiap \emph{sisi} kubus hanyalah \emph{bagian} dari sebuah bidang.
\end{itemize}
Melalui dua titik berbeda terdapat tepat satu garis; melalui tiga titik tak segaris terdapat tepat satu bidang.
\end{konsep}

\begin{center}
\begin{tikzpicture}[scale=1.8,line join=round,>=Stealth]
\coordinate (A) at (0,0); \coordinate (B) at (2,0);
\coordinate (C) at (2.8,0.8); \coordinate (D) at (0.8,0.8);
\coordinate (E) at (0,2); \coordinate (F) at (2,2);
\coordinate (G) at (2.8,2.8); \coordinate (H) at (0.8,2.8);
\fill[hijau,opacity=0.15] (A)--(B)--(F)--(E)--cycle;
\draw[thick] (A)--(B)--(F)--(E)--cycle;
\draw[thick] (B)--(C)--(G)--(F);
\draw[thick] (E)--(H)--(G);
\draw[dashed] (A)--(D)--(C);
\draw[dashed] (D)--(H);
\draw[ultra thick,oranye] (A)--(B);
\fill[merah] (H) circle (1.8pt);
\foreach \p/\pos in {A/below left,B/below,C/right,D/above left,E/left,F/right,G/right,H/above}
  \node[\pos,font=\scriptsize] at (\p) {$\p$};
% anotasi
\node[merah,font=\scriptsize,align=center] at (2.0,3.0) {\textbf{titik} (sudut $H$)};
\draw[->,merah,thin] (1.55,3.0)--(H);
\node[oranye,font=\scriptsize,align=center] at (1.0,-0.55) {\textbf{garis/rusuk} $AB$ (ruas garis)};
\draw[->,oranye,thin] (1.0,-0.32)--(1.0,-0.02);
\node[hijau!55!black,font=\scriptsize,align=center] at (3.7,1.4) {\textbf{bidang} $ABFE$\\(sisi depan)};
\draw[->,hijau!55!black,thin] (3.05,1.4)--(1.3,1.0);
\end{tikzpicture}\\
\small \textbf{Tiga unsur dasar pada kubus.} Titik (sudut $H$), ruas garis (rusuk $AB$), dan bidang (sisi $ABFE$). Bidang dan garis sejatinya tak terbatas; sisi dan rusuk hanyalah bagian yang tampak.
\end{center}

\begin{konsep}{Kedudukan Unsur Ruang}
\textbf{Dua garis} dapat: berpotongan (sebidang, satu titik sekutu), sejajar (sebidang, tanpa titik sekutu), atau \textbf{bersilangan} (tidak sebidang). \textbf{Garis terhadap bidang}: terletak pada bidang, sejajar bidang, atau menembus (memotong). \textbf{Dua bidang}: berpotongan (pada sebuah garis) atau sejajar. Melalui tiga titik tak segaris terdapat \emph{tepat satu} bidang.
\end{konsep}

\smallskip\noindent\textbf{Visualisasi kedudukan dua garis.}
\begin{center}
\begin{tabular}{ccc}
% (a) berpotongan
\begin{tikzpicture}[scale=0.72,line join=round]
\coordinate (A) at (0,0); \coordinate (B) at (2,0); \coordinate (C) at (2.8,0.8); \coordinate (D) at (0.8,0.8);
\coordinate (E) at (0,2); \coordinate (F) at (2,2); \coordinate (G) at (2.8,2.8); \coordinate (H) at (0.8,2.8);
\draw[thick] (A)--(B)--(F)--(E)--cycle; \draw[thick] (B)--(C)--(G)--(F); \draw[thick] (E)--(H)--(G);
\draw[dashed] (A)--(D)--(C); \draw[dashed] (D)--(H);
\draw[ultra thick,biru] (A)--(B); \draw[ultra thick,oranye] (B)--(C);
\fill[merah] (B) circle (2.4pt);
\node[font=\scriptsize,below left] at (A){$A$}; \node[font=\scriptsize,below right] at (B){$B$}; \node[font=\scriptsize,right] at (C){$C$};
\end{tikzpicture}
&
% (b) sejajar
\begin{tikzpicture}[scale=0.72,line join=round]
\coordinate (A) at (0,0); \coordinate (B) at (2,0); \coordinate (C) at (2.8,0.8); \coordinate (D) at (0.8,0.8);
\coordinate (E) at (0,2); \coordinate (F) at (2,2); \coordinate (G) at (2.8,2.8); \coordinate (H) at (0.8,2.8);
\draw[thick] (A)--(B)--(F)--(E)--cycle; \draw[thick] (B)--(C)--(G)--(F); \draw[thick] (E)--(H)--(G);
\draw[dashed] (A)--(D)--(C); \draw[dashed] (D)--(H);
\draw[ultra thick,biru] (A)--(B); \draw[ultra thick,oranye] (E)--(F);
\node[font=\scriptsize,below left] at (A){$A$}; \node[font=\scriptsize,below right] at (B){$B$}; \node[font=\scriptsize,above left] at (E){$E$}; \node[font=\scriptsize,above right] at (F){$F$};
\end{tikzpicture}
&
% (c) bersilangan
\begin{tikzpicture}[scale=0.72,line join=round]
\coordinate (A) at (0,0); \coordinate (B) at (2,0); \coordinate (C) at (2.8,0.8); \coordinate (D) at (0.8,0.8);
\coordinate (E) at (0,2); \coordinate (F) at (2,2); \coordinate (G) at (2.8,2.8); \coordinate (H) at (0.8,2.8);
\draw[thick] (A)--(B)--(F)--(E)--cycle; \draw[thick] (B)--(C)--(G)--(F); \draw[thick] (E)--(H)--(G);
\draw[dashed] (A)--(D)--(C); \draw[dashed] (D)--(H);
\draw[ultra thick,biru] (A)--(B); \draw[ultra thick,oranye] (C)--(G);
\node[font=\scriptsize,below left] at (A){$A$}; \node[font=\scriptsize,below] at (B){$B$}; \node[font=\scriptsize,below right] at (C){$C$}; \node[font=\scriptsize,right] at (G){$G$};
\end{tikzpicture}
\\[-1pt]
\footnotesize (a) \textbf{Berpotongan}: $AB$ \&\ $BC$ & \footnotesize (b) \textbf{Sejajar}: $AB \parallel EF$ & \footnotesize (c) \textbf{Bersilangan}: $AB$ \&\ $CG$\\
\footnotesize sebidang, bertemu di $B$ & \footnotesize sebidang, tak bertemu & \footnotesize \emph{tidak} sebidang, tak bertemu\\
\end{tabular}
\end{center}

\smallskip\noindent\textbf{Visualisasi kedudukan garis terhadap bidang.}
\begin{center}
\begin{tabular}{ccc}
% garis pada bidang
\begin{tikzpicture}[scale=0.72,line join=round]
\coordinate (A) at (0,0); \coordinate (B) at (2,0); \coordinate (C) at (2.8,0.8); \coordinate (D) at (0.8,0.8);
\coordinate (E) at (0,2); \coordinate (F) at (2,2); \coordinate (G) at (2.8,2.8); \coordinate (H) at (0.8,2.8);
\fill[biru,opacity=0.15] (A)--(B)--(C)--(D)--cycle;
\draw[thick] (A)--(B)--(F)--(E)--cycle; \draw[thick] (B)--(C)--(G)--(F); \draw[thick] (E)--(H)--(G);
\draw[dashed] (A)--(D)--(C); \draw[dashed] (D)--(H);
\draw[ultra thick,oranye] (A)--(B);
\node[font=\scriptsize,below left] at (A){$A$}; \node[font=\scriptsize,below right] at (B){$B$};
\end{tikzpicture}
&
% garis sejajar bidang
\begin{tikzpicture}[scale=0.72,line join=round]
\coordinate (A) at (0,0); \coordinate (B) at (2,0); \coordinate (C) at (2.8,0.8); \coordinate (D) at (0.8,0.8);
\coordinate (E) at (0,2); \coordinate (F) at (2,2); \coordinate (G) at (2.8,2.8); \coordinate (H) at (0.8,2.8);
\fill[biru,opacity=0.15] (A)--(B)--(C)--(D)--cycle;
\draw[thick] (A)--(B)--(F)--(E)--cycle; \draw[thick] (B)--(C)--(G)--(F); \draw[thick] (E)--(H)--(G);
\draw[dashed] (A)--(D)--(C); \draw[dashed] (D)--(H);
\draw[ultra thick,oranye] (E)--(F);
\node[font=\scriptsize,above left] at (E){$E$}; \node[font=\scriptsize,above right] at (F){$F$};
\end{tikzpicture}
&
% garis menembus bidang
\begin{tikzpicture}[scale=0.72,line join=round]
\coordinate (A) at (0,0); \coordinate (B) at (2,0); \coordinate (C) at (2.8,0.8); \coordinate (D) at (0.8,0.8);
\coordinate (E) at (0,2); \coordinate (F) at (2,2); \coordinate (G) at (2.8,2.8); \coordinate (H) at (0.8,2.8);
\fill[hijau,opacity=0.18] (B)--(D)--(E)--cycle;
\draw[thick] (A)--(B)--(F)--(E)--cycle; \draw[thick] (B)--(C)--(G)--(F); \draw[thick] (E)--(H)--(G);
\draw[dashed] (A)--(D)--(C); \draw[dashed] (D)--(H);
\draw[ultra thick,oranye] (A)--(G);
\fill[merah] (0.933,0.933) circle (2.2pt);
\node[font=\scriptsize,below left] at (A){$A$}; \node[font=\scriptsize,right] at (G){$G$};
\node[font=\scriptsize,merah,right] at (0.95,0.78){\,titik tembus};
\end{tikzpicture}
\\[-1pt]
\footnotesize \textbf{Terletak pada bidang} & \footnotesize \textbf{Sejajar bidang} & \footnotesize \textbf{Menembus bidang}\\
\footnotesize $AB$ pada alas $ABCD$ & \footnotesize $EF \parallel$ alas $ABCD$ & \footnotesize $AG$ menembus bidang $BDE$\\
\end{tabular}
\end{center}

\smallskip\noindent\textbf{Visualisasi kedudukan dua bidang.}
\begin{center}
\begin{tabular}{cc}
% berpotongan
\begin{tikzpicture}[scale=0.78,line join=round]
\coordinate (A) at (0,0); \coordinate (B) at (2,0); \coordinate (C) at (2.8,0.8); \coordinate (D) at (0.8,0.8);
\coordinate (E) at (0,2); \coordinate (F) at (2,2); \coordinate (G) at (2.8,2.8); \coordinate (H) at (0.8,2.8);
\fill[hijau,opacity=0.18] (A)--(B)--(F)--(E)--cycle;
\fill[biru,opacity=0.18] (B)--(C)--(G)--(F)--cycle;
\draw[thick] (A)--(B)--(F)--(E)--cycle; \draw[thick] (B)--(C)--(G)--(F); \draw[thick] (E)--(H)--(G);
\draw[dashed] (A)--(D)--(C); \draw[dashed] (D)--(H);
\draw[ultra thick,merah] (B)--(F);
\node[font=\scriptsize,below] at (B){$B$}; \node[font=\scriptsize,right] at (F){$F$};
\end{tikzpicture}
&
% sejajar
\begin{tikzpicture}[scale=0.78,line join=round]
\coordinate (A) at (0,0); \coordinate (B) at (2,0); \coordinate (C) at (2.8,0.8); \coordinate (D) at (0.8,0.8);
\coordinate (E) at (0,2); \coordinate (F) at (2,2); \coordinate (G) at (2.8,2.8); \coordinate (H) at (0.8,2.8);
\fill[hijau,opacity=0.18] (A)--(B)--(F)--(E)--cycle;
\fill[biru,opacity=0.18] (D)--(C)--(G)--(H)--cycle;
\draw[thick] (A)--(B)--(F)--(E)--cycle; \draw[thick] (B)--(C)--(G)--(F); \draw[thick] (E)--(H)--(G);
\draw[dashed] (A)--(D)--(C); \draw[dashed] (D)--(H);
\end{tikzpicture}
\\[-1pt]
\footnotesize \textbf{Berpotongan}: $ABFE$ \&\ $BCGF$ pada garis $BF$ & \footnotesize \textbf{Sejajar}: $ABFE \parallel DCGH$\\
\end{tabular}
\end{center}

\smallskip\noindent\textbf{Jenis-jenis titik penting pada kubus.}
\begin{center}
\begin{tikzpicture}[scale=1.7,line join=round,>=Stealth]
\coordinate (A) at (0,0); \coordinate (B) at (2,0); \coordinate (C) at (2.8,0.8); \coordinate (D) at (0.8,0.8);
\coordinate (E) at (0,2); \coordinate (F) at (2,2); \coordinate (G) at (2.8,2.8); \coordinate (H) at (0.8,2.8);
\draw[thick] (A)--(B)--(F)--(E)--cycle; \draw[thick] (B)--(C)--(G)--(F); \draw[thick] (E)--(H)--(G);
\draw[dashed] (A)--(D)--(C); \draw[dashed] (D)--(H);
\draw[dashed,gray] (G)--(C); % proyeksi G ke alas
% titik-titik
\fill[merah] (A) circle (1.8pt);
\fill[oranye] (1,0) circle (1.8pt);      % tengah AB
\fill[hijau!60!black] (1.4,2.4) circle (1.8pt); % pusat sisi atas EFGH
\fill[ungu] (1.4,1.4) circle (1.8pt);    % pusat kubus
\fill[tosca] (C) circle (1.8pt);         % proyeksi G pada alas
\node[font=\scriptsize,below left] at (A){$A$};
\node[merah,font=\scriptsize,align=center] at (-0.55,-0.25){titik\\sudut};
\draw[->,merah,thin] (-0.35,-0.05)--(A);
\node[oranye,font=\scriptsize,align=center] at (1,-0.5){tengah\\rusuk $AB$};
\draw[->,oranye,thin] (1,-0.28)--(1,-0.02);
\node[hijau!55!black,font=\scriptsize,align=center] at (2.7,3.05){pusat\\bidang $EFGH$};
\draw[->,hijau!55!black,thin] (2.2,3.0)--(1.45,2.45);
\node[ungu,font=\scriptsize,align=center] at (-0.5,1.6){pusat\\kubus};
\draw[->,ungu,thin] (-0.1,1.5)--(1.35,1.42);
\node[tosca,font=\scriptsize,align=center] at (3.7,0.7){proyeksi $G$\\pada alas $=C$};
\draw[->,tosca,thin] (3.1,0.7)--(2.85,0.8);
\node[font=\scriptsize,right] at (G){$G$};
\end{tikzpicture}\\
\small \textbf{Ragam titik.} Selain titik sudut, sering muncul: titik tengah rusuk, pusat bidang (perpotongan diagonal sisi), pusat kubus (perpotongan diagonal ruang), \emph{titik tembus} (garis $\cap$ bidang), dan \emph{titik proyeksi} (kaki garis tegak lurus, mis. proyeksi $G$ ke alas adalah $C$).
\end{center}

\begin{konsep}{Bahasa Alternatif yang Sering Mengecoh di Soal}
\begin{itemize}[leftmargin=*,topsep=2pt]
  \item \textbf{"rusuk $AB$"} = ruas garis; \textbf{"garis $AB$"} = garis tak hingga yang memuatnya (boleh \emph{diperpanjang}). Soal yang menyebut "perpanjangan $AB$" memakai bagian di luar rusuk.
  \item \textbf{"bidang alas"} $=ABCD$, \textbf{"bidang atas/tutup"} $=EFGH$, \textbf{"bidang sisi tegak"} $=ABFE,\,BCGF,\,\dots$
  \item \textbf{"bidang diagonal"} = bidang yang memuat dua rusuk berhadapan, mis. $ACGE$ atau $BDHF$ (berbentuk persegi panjang $a\times a\sqrt2$). Jangan dikira sisi biasa.
  \item \textbf{"diagonal sisi"} (mis. $AC,\,AF$) berbeda dari \textbf{"diagonal ruang"} (mis. $AG,\,BH$). Bedakan baik-baik.
  \item \textbf{"jarak"} titik ke garis/bidang \emph{selalu} berarti ruas \textbf{tegak lurus} (terpendek), bukan ke sembarang titik acuan.
  \item \textbf{"proyeksi $P$ pada bidang"} = kaki garis tegak lurus dari $P$; \textbf{"proyeksi garis pada bidang"} = bayangannya.
  \item \textbf{"sudut garis dan bidang"} diukur ke \emph{proyeksi} garis, \textbf{bukan} ke garis normal (lihat bagian Sudut).
  \item \textbf{"titik tembus"} = titik potong sebuah garis dengan sebuah bidang.
\end{itemize}
\end{konsep}

\begin{konsep}{Diagonal pada Kubus Rusuk $a$}
\textbf{Diagonal sisi} (mis. $AC$) $=a\sqrt2$ (dari Pythagoras $\sqrt{a^2+a^2}$). \textbf{Diagonal ruang} (mis. $AG$) $=a\sqrt3$ (dari $\sqrt{(a\sqrt2)^2+a^2}$). Sebuah kubus punya $12$ diagonal sisi dan $4$ diagonal ruang.
\end{konsep}

\section{Jarak dalam Ruang}
\begin{konsep}{Tiga Jenis Jarak}
\begin{itemize}[leftmargin=*,topsep=2pt]
  \item \textbf{Titik ke titik:} $d=\sqrt{(x_2-x_1)^2+(y_2-y_1)^2+(z_2-z_1)^2}$.
  \item \textbf{Titik ke garis:} panjang ruas \emph{tegak lurus} dari titik ke garis.
  \item \textbf{Titik ke bidang:} panjang ruas tegak lurus dari titik ke bidang. Untuk bidang $ax+by+cz+d=0$:
  \[
  d=\frac{|ax_0+by_0+cz_0+d|}{\sqrt{a^2+b^2+c^2}}.
  \]
\end{itemize}
\end{konsep}

\begin{contoh}{Jarak Titik ke Titik dan Diagonal}
Kubus rusuk $6$ cm. Tentukan panjang diagonal sisi dan diagonal ruangnya.

\jawab Diagonal sisi $=6\sqrt2$ cm; diagonal ruang $=6\sqrt3$ cm.
\end{contoh}

\begin{contoh}{Jarak Titik ke Bidang}
Tentukan jarak titik $P(1,1,1)$ ke bidang $2x+y+2z=6$.

\jawab $d=\dfrac{|2(1)+1(1)+2(1)-6|}{\sqrt{2^2+1^2+2^2}}=\dfrac{|-1|}{3}=\dfrac13$.
\end{contoh}

\section{Sudut dalam Ruang}
\begin{konsep}{Tiga Jenis Sudut}
\begin{itemize}[leftmargin=*,topsep=2pt]
  \item \textbf{Antara dua garis:} sudut yang dibentuk vektor arah keduanya, $\cos\theta=\dfrac{|\vec u\cdot\vec v|}{|\vec u|\,|\vec v|}$.
  \item \textbf{Antara garis dan bidang:} sudut antara garis dan \emph{proyeksinya} pada bidang (komplemen sudut garis dengan normal bidang).
  \item \textbf{Antara dua bidang (sudut dihedral):} sudut antara dua garis pada masing-masing bidang yang tegak lurus garis potong; setara sudut antar vektor normal.
\end{itemize}
\end{konsep}

\smallskip\noindent\textbf{Visualisasi tiga jenis sudut.}
\begin{center}
\begin{tabular}{ccc}
% (A) dua garis
\begin{tikzpicture}[scale=0.78,line join=round]
\coordinate (A) at (0,0); \coordinate (B) at (2,0); \coordinate (C) at (2.8,0.8); \coordinate (D) at (0.8,0.8);
\coordinate (E) at (0,2); \coordinate (F) at (2,2); \coordinate (G) at (2.8,2.8); \coordinate (H) at (0.8,2.8);
\draw[thick] (A)--(B)--(F)--(E)--cycle; \draw[thick] (B)--(C)--(G)--(F); \draw[thick] (E)--(H)--(G);
\draw[dashed] (A)--(D)--(C); \draw[dashed] (D)--(H);
\draw[ultra thick,biru] (A)--(B); \draw[ultra thick,oranye] (A)--(G);
\draw[merah,thick] (0.5,0) arc(0:45:0.5);
\node[merah,font=\scriptsize] at (0.72,0.26){$\theta$};
\node[font=\scriptsize,below left] at (A){$A$}; \node[font=\scriptsize,below right] at (B){$B$}; \node[font=\scriptsize,right] at (G){$G$};
\end{tikzpicture}
&
% (B) garis-bidang
\begin{tikzpicture}[scale=0.78,line join=round]
\coordinate (A) at (0,0); \coordinate (B) at (2,0); \coordinate (C) at (2.8,0.8); \coordinate (D) at (0.8,0.8);
\coordinate (E) at (0,2); \coordinate (F) at (2,2); \coordinate (G) at (2.8,2.8); \coordinate (H) at (0.8,2.8);
\fill[biru,opacity=0.13] (A)--(B)--(C)--(D)--cycle;
\draw[thick] (A)--(B)--(F)--(E)--cycle; \draw[thick] (B)--(C)--(G)--(F); \draw[thick] (E)--(H)--(G);
\draw[dashed] (A)--(D)--(C); \draw[dashed] (D)--(H);
\draw[ultra thick,oranye] (A)--(G);
\draw[ultra thick,hijau!60!black] (A)--(C);
\draw[thin,gray] (G)--(C);
\draw[thin] (2.627,0.7505)--(2.627,0.9305)--(2.8,0.98);
\draw[merah,thick] (0.6,0.171) arc(16:45:0.62);
\node[merah,font=\scriptsize] at (0.78,0.35){$\theta$};
\node[font=\scriptsize,below left] at (A){$A$}; \node[font=\scriptsize,right] at (C){$C$}; \node[font=\scriptsize,right] at (G){$G$};
\end{tikzpicture}
&
% (C) dua bidang (dihedral)
\begin{tikzpicture}[scale=0.78,line join=round]
\coordinate (A) at (0,0); \coordinate (B) at (2,0); \coordinate (C) at (2.8,0.8); \coordinate (D) at (0.8,0.8);
\coordinate (E) at (0,2); \coordinate (F) at (2,2); \coordinate (G) at (2.8,2.8); \coordinate (H) at (0.8,2.8);
\coordinate (O) at (1.4,0.4);
\fill[biru,opacity=0.13] (A)--(B)--(C)--(D)--cycle;
\fill[hijau,opacity=0.18] (B)--(D)--(G)--cycle;
\draw[thick] (A)--(B)--(F)--(E)--cycle; \draw[thick] (B)--(C)--(G)--(F); \draw[thick] (E)--(H)--(G);
\draw[dashed] (A)--(D)--(C); \draw[dashed] (D)--(H);
\draw[ultra thick,merah] (B)--(D);
\draw[thin,oranye] (O)--(G); \draw[thin,hijau!60!black] (O)--(C);
\draw[merah,thick] (1.88,0.537) arc(16:60:0.5);
\node[merah,font=\scriptsize] at (2.0,0.82){$\theta$};
\node[font=\scriptsize,below] at (B){$B$}; \node[font=\scriptsize,above left] at (D){$D$}; \node[font=\scriptsize,right] at (G){$G$}; \node[font=\scriptsize,below right] at (C){$C$};
\end{tikzpicture}
\\[-1pt]
\footnotesize (A) \textbf{Dua garis}: $\angle(AB,AG)$ & \footnotesize (B) \textbf{Garis--bidang}: $AG$ \&\ proyeksinya $AC$ & \footnotesize (C) \textbf{Dua bidang}: dihedral $BDG$--alas\\
\footnotesize arc di titik temu & \footnotesize $GC\perp$ alas; $\theta=\angle GAC$ & \footnotesize $\theta=\angle GOC$, $OG,OC\perp BD$\\
\end{tabular}
\end{center}

\begin{contoh}{Sudut Garis dan Bidang}
Pada kubus rusuk $a$, tentukan sudut antara diagonal ruang $AG$ dan bidang alas $ABCD$.

\jawab Proyeksi $AG$ pada alas adalah $AC$. Komponen tegak $=a$ (tinggi), $AC=a\sqrt2$. Maka $\tan\theta=\dfrac{a}{a\sqrt2}=\dfrac{1}{\sqrt2}$, sehingga $\theta=\arctan\tfrac{1}{\sqrt2}\approx35{,}26^\circ$.
\end{contoh}

\begin{kesalahan}
Sudut garis--bidang adalah sudut terhadap \emph{proyeksi} garis di bidang, \emph{bukan} terhadap garis normal. Keduanya saling berkomplemen ($90^\circ$); tertukar akan memberi jawaban yang salah.
\end{kesalahan}

\begin{dunianyata}
Geometri ruang mendasari arsitektur dan konstruksi (rangka atap, kemiringan tangga), grafika komputer dan game 3D (proyeksi, kamera), navigasi GPS (posisi $3$D), serta kristalografi dan kimia (sudut ikatan molekul, mis. $\approx109{,}5^\circ$ pada metana).
\end{dunianyata}

\section*{Ringkasan Bab}
\addcontentsline{toc}{section}{Ringkasan Bab}
\begin{itemize}[leftmargin=*,topsep=2pt]
  \item Kedudukan garis: berpotongan, sejajar, atau bersilangan; bidang: berpotongan atau sejajar.
  \item Kubus rusuk $a$: diagonal sisi $a\sqrt2$, diagonal ruang $a\sqrt3$.
  \item Jarak: titik--titik (Pythagoras $3$D), titik--garis \& titik--bidang (ruas tegak lurus); rumus titik--bidang $\dfrac{|ax_0+by_0+cz_0+d|}{\sqrt{a^2+b^2+c^2}}$.
  \item Sudut: antar garis ($\cos\theta$ dari vektor arah), garis--bidang (terhadap proyeksi), antar bidang (dihedral/normal).
\end{itemize}

\begin{refleksi}
\begin{itemize}[leftmargin=*,topsep=2pt]
  \item[\cek] Saya dapat menentukan kedudukan titik, garis, dan bidang.
  \item[\cek] Saya dapat menghitung jarak titik ke titik, garis, dan bidang.
  \item[\cek] Saya dapat menentukan sudut antar garis, garis--bidang, dan antar bidang.
  \item[\cek] Saya dapat memakai kubus dan koordinat sebagai alat bantu.
\end{itemize}
\end{refleksi}

\section*{Soal Akhir Bab \thechapter}
\addcontentsline{toc}{section}{Soal Akhir Bab \thechapter}

\bagiansoal{A}{(Fundamental --- setara SNBT Dasar)}
\begin{enumerate}[leftmargin=*]
  \item Kubus rusuk $6$ cm. Tentukan panjang diagonal sisinya.
  \item Tentukan panjang diagonal ruang kubus rusuk $6$ cm.
  \item Tentukan jarak titik $A(1,2,3)$ ke $B(4,6,3)$.
  \item Berapa banyak diagonal ruang pada sebuah kubus?
  \item Dua garis yang tidak sebidang dan tidak berpotongan disebut \dots
  \item Tentukan jarak titik $O(0,0,0)$ ke $P(2,3,6)$.
  \item Pada kubus $ABCD.EFGH$, bagaimana kedudukan rusuk $AB$ dan $DC$?
  \item Tentukan panjang rusuk kubus jika diagonal ruangnya $9\sqrt3$.
  \item Melalui tiga titik tak segaris terdapat berapa bidang?
  \item Tentukan jarak titik $(1,1,1)$ ke $(1,1,5)$.
\end{enumerate}

\bagiansoal{B}{(Intermediate --- setara SNBT Sulit)}
\begin{enumerate}[leftmargin=*]
  \item Kubus rusuk $4$. Tentukan jarak $A$ ke $C$.
  \item Kubus rusuk $4$. Tentukan jarak $A$ ke $G$.
  \item Kubus rusuk $6$. Tentukan jarak titik $A$ ke bidang $BDHF$.
  \item Pada kubus, tentukan sudut antara diagonal ruang $AG$ dan rusuk $AB$.
  \item Pada kubus, tentukan sudut antara diagonal ruang $AG$ dan bidang alas $ABCD$.
\end{enumerate}

\bagiansoal{C}{(Advanced --- setara Olimpiade SMA, gabungan antarbab)}
\begin{enumerate}[leftmargin=*]
  \item \textbf{(Dim Tiga + Trigonometri).} Pada kubus rusuk $a$, tentukan sudut antara bidang $BDG$ dan bidang alas $ABCD$.
  \item \textbf{(Dim Tiga + Vektor).} Kubus rusuk $6$. Tentukan jarak titik $B$ ke diagonal ruang $AG$.
  \item \textbf{(Dim Tiga + Geometri Analitik).} Tentukan jarak titik $P(1,1,1)$ ke bidang $2x+y+2z=6$.
  \item \textbf{(Dim Tiga + Pythagoras).} Limas $T.ABCD$ beralas persegi sisi $6$ dengan tinggi $TO=4$ ($O$ pusat alas). Tentukan panjang rusuk tegak $TA$.
  \item \textbf{(Dim Tiga + Trigonometri).} Pada limas soal sebelumnya, tentukan sudut antara rusuk $TA$ dan bidang alas.
\end{enumerate}

\bagiansoal{D}{(Expert --- setara tahun pertama universitas, sintesis \& pembuktian)}
\begin{enumerate}[leftmargin=*]
  \item \textbf{(Pembuktian).} Buktikan bahwa diagonal ruang kubus rusuk $a$ panjangnya $a\sqrt3$.
  \item \textbf{(Penurunan Rumus).} Turunkan rumus jarak titik $P(x_0,y_0,z_0)$ ke bidang $ax+by+cz+d=0$.
  \item \textbf{(Sudut Antar Bidang via Normal).} Tentukan sudut antara bidang $x+y+z=1$ dan bidang $x=0$.
  \item \textbf{(Volume).} Pada kubus rusuk $a$, tentukan volume bidang empat (tetrahedron) bertitik sudut $A,B,D,E$.
  \item \textbf{(Jarak Garis Bersilangan).} Pada kubus rusuk $a$, tentukan jarak antara garis $AC$ dan garis $BH$.
\end{enumerate}

\begin{challenge}
\begin{enumerate}[leftmargin=*,topsep=2pt]
  \item Pada kubus rusuk $a$, tentukan sudut antara dua diagonal ruang $AG$ dan $CE$.
  \item Pada kubus rusuk $a$, buktikan diagonal ruang $AG$ tegak lurus bidang $BDE$, lalu tentukan jarak titik $A$ ke bidang $BDE$.
  \item Tentukan jari-jari bola yang menyinggung semua rusuk sebuah kubus rusuk $a$.
\end{enumerate}
\end{challenge}

% ============================================================
\part*{\centering KUNCI JAWABAN DAN PEMBAHASAN}
\addcontentsline{toc}{part}{KUNCI JAWABAN DAN PEMBAHASAN}
% ============================================================
\noindent\textit{Setiap pembahasan disusun dengan tiga bagian: \textbf{Soal} (pernyataan ringkas), \textbf{Analisis} (strategi dan konsep kunci), dan \textbf{Pembahasan} (langkah penyelesaian). Gunakan kunci ini untuk mengoreksi diri, bukan sekadar mencocokkan hasil.}

% ============================================================
\chapter{Pembahasan Bab 1: Komposisi dan Invers}
% ============================================================

\kbagian{Bagian A}

\knomor{A.1}
\ksoal $f(x)=2x+1$, $g(x)=x-3$. Tentukan $(f\circ g)(x)$.
\kanalisis Substitusikan $g(x)$ ke dalam $f$.
\kbahas $(f\circ g)(x)=f(x-3)=2(x-3)+1=2x-5.$

\knomor{A.2}
\ksoal $f(x)=x^2$, $g(x)=x+2$. Hitung $(g\circ f)(2)$.
\kanalisis Kerjakan $f$ dahulu, lalu $g$.
\kbahas $f(2)=4$, maka $(g\circ f)(2)=g(4)=4+2=6.$

\knomor{A.3}
\ksoal Tentukan invers dari $f(x)=3x-6$.
\kanalisis Tulis $y=3x-6$, lalu nyatakan $x$ dalam $y$.
\kbahas $y=3x-6\Rightarrow x=\dfrac{y+6}{3}$. Jadi $f^{-1}(x)=\dfrac{x+6}{3}.$

\knomor{A.4}
\ksoal $f(x)=\dfrac{x+4}{2}$. Tentukan $f^{-1}(x)$.
\kanalisis Selesaikan $y=\tfrac{x+4}{2}$ untuk $x$.
\kbahas $2y=x+4\Rightarrow x=2y-4$. Jadi $f^{-1}(x)=2x-4.$

\knomor{A.5}
\ksoal $f(x)=2x$, $g(x)=x+5$. Hitung $(f\circ g)(3)$.
\kanalisis Kerjakan $g$ dahulu, lalu $f$.
\kbahas $g(3)=8$, maka $(f\circ g)(3)=f(8)=16.$

\kbagian{Bagian B}

\knomor{B.1}
\ksoal $f(x)=2x+1$, $g(x)=x^2-1$. Tentukan $(f\circ g)(x)$ dan $(g\circ f)(x)$.
\kanalisis Hitung kedua komposisi lalu bandingkan.
\kbahas $(f\circ g)(x)=2(x^2-1)+1=2x^2-1$. \quad $(g\circ f)(x)=(2x+1)^2-1=4x^2+4x$.
Karena $2x^2-1\neq 4x^2+4x$, terbukti $f\circ g\neq g\circ f$.

\knomor{B.2}
\ksoal $(f\circ g)(x)=4x+5$ dan $g(x)=2x+1$. Tentukan $f(x)$.
\kanalisis Misalkan $u=g(x)=2x+1$, nyatakan $x$ dalam $u$, lalu substitusi.
\kbahas Dari $u=2x+1$ diperoleh $x=\tfrac{u-1}{2}$. Maka $f(u)=4\cdot\tfrac{u-1}{2}+5=2(u-1)+5=2u+3$. Jadi $f(x)=2x+3.$

\knomor{B.3}
\ksoal Tentukan invers dari $f(x)=\dfrac{2x-1}{x+3}$.
\kanalisis Kalikan silang, kelompokkan suku ber-$x$.
\kbahas $y(x+3)=2x-1\Rightarrow x(y-2)=-1-3y\Rightarrow x=\dfrac{3y+1}{2-y}$. Jadi $f^{-1}(x)=\dfrac{3x+1}{2-x}$, dengan $x\neq2$.

\knomor{B.4}
\ksoal $f(x)=\sqrt{x-2}$. Tentukan domain $f$ dan $f^{-1}(x)$.
\kanalisis Syarat akar $\ge0$; untuk invers, kuadratkan.
\kbahas Domain $f$: $x\ge2$. Dari $y=\sqrt{x-2}$ ($y\ge0$) diperoleh $y^2=x-2\Rightarrow x=y^2+2$. Jadi $f^{-1}(x)=x^2+2$ dengan $x\ge0$.

\knomor{B.5}
\ksoal $f(x)=3x-2$, $g(x)=x+a$. Jika $(f\circ g)(2)=10$, tentukan $a$.
\kanalisis Susun $(f\circ g)(x)$ lalu masukkan $x=2$.
\kbahas $(f\circ g)(x)=3(x+a)-2=3x+3a-2$. Untuk $x=2$: $6+3a-2=10\Rightarrow 3a=6\Rightarrow a=2.$

\knomor{B.6}
\ksoal $f(x)=x+1$ dan $(f\circ g)(x)=x^2+2$. Tentukan $g(x)$.
\kanalisis $(f\circ g)(x)=g(x)+1$.
\kbahas $g(x)+1=x^2+2\Rightarrow g(x)=x^2+1.$

\knomor{B.7}
\ksoal Tentukan invers dari $f(x)=2^{x+1}$.
\kanalisis Gunakan logaritma sebagai invers eksponen.
\kbahas $y=2^{x+1}\Rightarrow x+1={}^{2}\!\log y\Rightarrow x={}^{2}\!\log y-1$. Jadi $f^{-1}(x)={}^{2}\!\log x-1$, dengan $x>0$.

\kbagian{Bagian C}

\knomor{C.1}
\ksoal $f(x)=2^{x}$, $g(x)=x+3$. Selesaikan $(f\circ g)(x)=32$.
\kanalisis Samakan basis $2$.
\kbahas $(f\circ g)(x)=2^{x+3}=32=2^5\Rightarrow x+3=5\Rightarrow x=2.$

\knomor{C.2}
\ksoal $f(x)={}^{10}\!\log(x-1)$. Tentukan $f^{-1}(x)$ dan domainnya.
\kanalisis Invers logaritma adalah eksponen.
\kbahas $y={}^{10}\!\log(x-1)\Rightarrow x-1=10^{y}\Rightarrow x=10^{y}+1$. Jadi $f^{-1}(x)=10^{x}+1$, dengan domain semua bilangan real (karena $10^x$ terdefinisi untuk semua $x$).

\knomor{C.3}
\ksoal $(f\circ g)(x)=x^2-4x+7$ dan $f(x)=x+3$. Tentukan $g(x)$ dan titik minimumnya.
\kanalisis $(f\circ g)(x)=g(x)+3$.
\kbahas $g(x)+3=x^2-4x+7\Rightarrow g(x)=x^2-4x+4=(x-2)^2$. Titik minimum di $(2,0)$.

\knomor{C.4}
\ksoal $f(x)=\dfrac{x}{x+1}$. Tentukan $(f\circ f\circ f)(x)$.
\kanalisis Hitung bertahap; perhatikan pola penyebut.
\kbahas $(f\circ f)(x)=\dfrac{\frac{x}{x+1}}{\frac{x}{x+1}+1}=\dfrac{x}{2x+1}$. Lalu $(f\circ f\circ f)(x)=\dfrac{\frac{x}{2x+1}}{\frac{x}{2x+1}+1}=\dfrac{x}{3x+1}.$

\knomor{C.5}
\ksoal $f(x)=\dfrac{x+1}{x-1}$. Buktikan $f^{-1}(x)=f(x)$.
\kanalisis Cukup tunjukkan $f$ involusi: $f\big(f(x)\big)=x$.
\kbahas $f\big(f(x)\big)=\dfrac{\frac{x+1}{x-1}+1}{\frac{x+1}{x-1}-1}=\dfrac{\frac{(x+1)+(x-1)}{x-1}}{\frac{(x+1)-(x-1)}{x-1}}=\dfrac{2x}{2}=x.$ Karena $f\circ f=I$, maka $f^{-1}=f$.

\knomor{C.6}
\ksoal $f(x)=\sin x$, $g(x)=2x$. Bandingkan $(f\circ g)$ dan $(g\circ f)$ beserta periodenya.
\kanalisis Susun masing-masing komposisi.
\kbahas $(f\circ g)(x)=\sin(2x)$ berperiode $\pi$; \quad $(g\circ f)(x)=2\sin x$ berperiode $2\pi$. Keduanya berbeda, baik bentuk maupun periodenya.

\knomor{C.7}
\ksoal $x_0=5$, $x_{n+1}=2x_n-3$. Tentukan rumus $x_n$.
\kanalisis Cari titik tetap $x^*$ dari $x=2x-3$, lalu tinjau $x_n-x^*$.
\kbahas Titik tetap: $x^*=3$. Maka $x_{n}-3=2(x_{n-1}-3)$, sehingga $x_n-3=2^{n}(x_0-3)=2^{n}\cdot2=2^{n+1}$. Jadi $x_n=3+2^{n+1}.$

\kbagian{Bagian D}

\knomor{D.1}
\ksoal Buktikan komposisi dua fungsi injektif bersifat injektif.
\kanalisis Pakai definisi injektif: $h(a)=h(b)\Rightarrow a=b$.
\kbahas Misalkan $f,g$ injektif dan $(f\circ g)(a)=(f\circ g)(b)$, yaitu $f\big(g(a)\big)=f\big(g(b)\big)$. Karena $f$ injektif, $g(a)=g(b)$. Karena $g$ injektif, $a=b$. Jadi $f\circ g$ injektif. $\blacksquare$

\knomor{D.2}
\ksoal $f(x)=\dfrac{e^{x}-e^{-x}}{2}$. Tentukan $f^{-1}(x)$.
\kanalisis Tulis $y=\tfrac{e^x-e^{-x}}{2}$, kalikan $2e^x$ agar menjadi kuadrat dalam $e^x$.
\kbahas $2y=e^x-e^{-x}\Rightarrow 2ye^x=e^{2x}-1\Rightarrow (e^x)^2-2y\,e^x-1=0$. Dengan rumus kuadrat, $e^x=\dfrac{2y\pm\sqrt{4y^2+4}}{2}=y\pm\sqrt{y^2+1}$. Karena $e^x>0$, ambil tanda $+$: $e^x=y+\sqrt{y^2+1}$, sehingga $x=\ln\!\big(y+\sqrt{y^2+1}\big)$. Jadi $f^{-1}(x)=\ln\!\big(x+\sqrt{x^2+1}\big).$

\knomor{D.3}
\ksoal $f(x)=\dfrac{1}{1-x}$. Tunjukkan $(f\circ f\circ f)(x)=x$.
\kanalisis Hitung dua komposisi pertama, lalu komposisi ketiga.
\kbahas $(f\circ f)(x)=\dfrac{1}{1-\frac{1}{1-x}}=\dfrac{1}{\frac{(1-x)-1}{1-x}}=\dfrac{1-x}{-x}=\dfrac{x-1}{x}$. Lalu $(f\circ f\circ f)(x)=f\!\left(\dfrac{x-1}{x}\right)=\dfrac{1}{1-\frac{x-1}{x}}=\dfrac{1}{\frac{x-(x-1)}{x}}=\dfrac{x}{1}=x.$ $\blacksquare$

\knomor{D.4}
\ksoal $z=\dfrac{x-\mu}{\sigma}$. Tentukan transformasi inversnya dan maknanya.
\kanalisis Selesaikan untuk $x$.
\kbahas $z\sigma=x-\mu\Rightarrow x=\mu+\sigma z$. Inversnya $x=\mu+\sigma z$ mengembalikan skor baku menjadi nilai aslinya: kalikan dengan simpangan baku lalu tambah rata-rata. Maknanya, baku ($z$) menyatakan "berapa simpangan baku suatu data dari rata-rata", dan invers mengembalikannya ke skala asli.

\knomor{D.5}
\ksoal Buktikan: (i) $g$ genap $\Rightarrow f\circ g$ genap; (ii) $f,g$ ganjil $\Rightarrow f\circ g$ ganjil.
\kanalisis Pakai definisi genap $g(-x)=g(x)$ dan ganjil $g(-x)=-g(x)$.
\kbahas (i) $(f\circ g)(-x)=f\big(g(-x)\big)=f\big(g(x)\big)=(f\circ g)(x)$, jadi genap. \\
(ii) $(f\circ g)(-x)=f\big(g(-x)\big)=f\big(-g(x)\big)=-f\big(g(x)\big)=-(f\circ g)(x)$, jadi ganjil. $\blacksquare$

\kbagian{Challenge Corner}

\knomor{1}
\ksoal Tentukan semua $f(x)=ax+b$ dengan $(f\circ f)(x)=x$.
\kanalisis Susun $(f\circ f)(x)$ lalu samakan dengan $x$ untuk semua $x$.
\kbahas $(f\circ f)(x)=a(ax+b)+b=a^2x+(ab+b)$. Agar $=x$: $a^2=1$ dan $ab+b=0$.
\begin{itemize}[leftmargin=*,topsep=2pt]
  \item $a=1$: maka $2b=0\Rightarrow b=0$, memberi $f(x)=x$ (identitas).
  \item $a=-1$: maka $ab+b=-b+b=0$ otomatis, sehingga $b$ bebas; $f(x)=-x+b$.
\end{itemize}
Jadi $f(x)=x$ atau $f(x)=-x+b$ untuk sembarang $b$.

\knomor{2}
\ksoal $f(x)+2f\!\left(\tfrac1x\right)=3x$. Tentukan $f(x)$.
\kanalisis Substitusi $x\to\tfrac1x$ untuk memperoleh persamaan kedua, lalu eliminasi.
\kbahas Persamaan (1): $f(x)+2f\!\left(\tfrac1x\right)=3x$. Substitusi $x\to\tfrac1x$ memberi (2): $f\!\left(\tfrac1x\right)+2f(x)=\dfrac3x$. Kalikan (2) dengan $2$: $2f\!\left(\tfrac1x\right)+4f(x)=\dfrac6x$. Kurangi (1): $3f(x)=\dfrac6x-3x$, sehingga $f(x)=\dfrac2x-x.$

\knomor{3}
\ksoal $f(x)=\dfrac{x}{\sqrt{1+x^2}}$. Buktikan $f^{(n)}(x)=\dfrac{x}{\sqrt{1+nx^2}}$.
\kanalisis Induksi: basis $n=1$ benar; asumsikan benar untuk $n$, buktikan untuk $n+1$.
\kbahas \emph{Basis:} $n=1$ jelas benar. \emph{Langkah induksi:} andaikan $f^{(n)}(x)=\dfrac{x}{\sqrt{1+nx^2}}$. Maka
\[
f^{(n+1)}(x)=f\!\big(f^{(n)}(x)\big)=\frac{\frac{x}{\sqrt{1+nx^2}}}{\sqrt{1+\frac{x^2}{1+nx^2}}}
=\frac{\frac{x}{\sqrt{1+nx^2}}}{\sqrt{\frac{1+nx^2+x^2}{1+nx^2}}}
=\frac{x}{\sqrt{1+(n+1)x^2}}.
\]
Dengan prinsip induksi, rumus berlaku untuk setiap bilangan asli $n$. $\blacksquare$

% ============================================================
\chapter{Pembahasan Bab 2: Lingkaran}
% ============================================================

\kbagian{Bagian A}

\knomor{A.1}
\ksoal Persamaan lingkaran berpusat $(0,0)$, jari-jari $5$.
\kanalisis Bentuk baku dengan $a=b=0$.
\kbahas $x^2+y^2=25.$

\knomor{A.2}
\ksoal Persamaan lingkaran berpusat $(3,-2)$, jari-jari $4$.
\kanalisis Substitusi ke $(x-a)^2+(y-b)^2=r^2$.
\kbahas $(x-3)^2+(y+2)^2=16.$

\knomor{A.3}
\ksoal Pusat dan jari-jari $x^2+y^2-6x+4y-12=0$.
\kanalisis $A=-6,B=4,C=-12$; pusat $\left(-\tfrac A2,-\tfrac B2\right)$.
\kbahas Pusat $(3,-2)$, $r=\sqrt{9+4+12}=\sqrt{25}=5.$

\knomor{A.4}
\ksoal Panjang busur, $r=6$ cm, sudut $90^\circ$.
\kanalisis Busur $=\tfrac{\theta}{360^\circ}\cdot2\pi r$.
\kbahas $\tfrac{90}{360}\cdot2\pi(6)=\tfrac14\cdot12\pi=3\pi$ cm.

\knomor{A.5}
\ksoal Luas juring, $r=10$ cm, sudut $72^\circ$.
\kanalisis Juring $=\tfrac{\theta}{360^\circ}\cdot\pi r^2$.
\kbahas $\tfrac{72}{360}\cdot\pi(100)=\tfrac15\cdot100\pi=20\pi$ cm$^2$.

\kbagian{Bagian B}

\knomor{B.1}
\ksoal Kedudukan $(1,2)$ terhadap $x^2+y^2=9$.
\kanalisis Hitung $K=x_0^2+y_0^2-9$.
\kbahas $1+4-9=-4<0$, jadi titik \textbf{di dalam} lingkaran.

\knomor{B.2}
\ksoal Lingkaran berpusat $(2,3)$ melalui $(5,7)$.
\kanalisis $r=$ jarak pusat ke titik.
\kbahas $r=\sqrt{(5-2)^2+(7-3)^2}=\sqrt{9+16}=5$. Persamaan: $(x-2)^2+(y-3)^2=25.$

\knomor{B.3}
\ksoal Garis singgung $x^2+y^2=25$ di $(3,4)$.
\kanalisis Gunakan $x_1x+y_1y=r^2$.
\kbahas $3x+4y=25.$

\knomor{B.4}
\ksoal Garis singgung $x^2+y^2=20$ bergradien $2$.
\kanalisis Rumus $y=mx\pm r\sqrt{1+m^2}$ dengan $r=\sqrt{20}$, $m=2$.
\kbahas $y=2x\pm\sqrt{20}\cdot\sqrt{1+4}=2x\pm\sqrt{100}=2x\pm10.$

\knomor{B.5}
\ksoal Kedudukan garis $y=x+1$ terhadap $x^2+y^2=4$.
\kanalisis Substitusi, tinjau diskriminan.
\kbahas $x^2+(x+1)^2=4\Rightarrow 2x^2+2x-3=0$. $D=4+24=28>0$, jadi garis \textbf{memotong} lingkaran di dua titik.

\knomor{B.6}
\ksoal Kedudukan $(5,1)$ terhadap $x^2+y^2-4x-2y-4=0$.
\kanalisis Pusat $(2,1)$, $r=3$; bandingkan jarak dengan $r$.
\kbahas Jarak $=\sqrt{(5-2)^2+(1-1)^2}=3=r$, jadi titik \textbf{terletak pada} lingkaran.

\knomor{B.7}
\ksoal Panjang tali busur, $r=5$, jarak ke pusat $3$.
\kanalisis Setengah tali busur $=\sqrt{r^2-d^2}$.
\kbahas Panjang $=2\sqrt{25-9}=2\cdot4=8.$

\kbagian{Bagian C}

\knomor{C.1}
\ksoal Panjang garis singgung dari $(7,1)$ ke $x^2+y^2=9$.
\kanalisis Panjang $=\sqrt{x_P^2+y_P^2-r^2}$.
\kbahas $\sqrt{49+1-9}=\sqrt{41}.$

\knomor{C.2}
\ksoal Daya titik $P(6,0)$ terhadap $x^2+y^2=4$.
\kanalisis Daya $=x_P^2+y_P^2-r^2$.
\kbahas $36+0-4=32>0$. Daya positif menandakan $P$ \textbf{di luar} lingkaran; $\sqrt{32}=4\sqrt2$ adalah panjang garis singgung dari $P$.

\knomor{C.3}
\ksoal Tunjukkan $3x+4y=25$ menyinggung $x^2+y^2=25$ dan cari titik singgung.
\kanalisis Bandingkan jarak pusat ke garis dengan $r=5$.
\kbahas $d=\dfrac{|{-25}|}{\sqrt{3^2+4^2}}=\dfrac{25}{5}=5=r$, jadi menyinggung. Titik singgung adalah kaki tegak lurus dari $O$: $\left(\tfrac{3\cdot25}{25},\tfrac{4\cdot25}{25}\right)=(3,4).$

\knomor{C.4}
\ksoal $y=mx$ menyinggung $(x-4)^2+y^2=4$. Tentukan $m$.
\kanalisis Jarak pusat $(4,0)$ ke garis $mx-y=0$ sama dengan $r=2$.
\kbahas $\dfrac{|4m|}{\sqrt{m^2+1}}=2\Rightarrow16m^2=4(m^2+1)\Rightarrow12m^2=4\Rightarrow m^2=\tfrac13$. Jadi $m=\pm\dfrac{1}{\sqrt3}=\pm\dfrac{\sqrt3}{3}.$

\knomor{C.5}
\ksoal Lingkaran melalui $A(0,0)$, $B(6,0)$, $C(0,8)$.
\kanalisis Pakai bentuk umum $x^2+y^2+Ax+By+C=0$.
\kbahas $A(0,0)\Rightarrow C=0$. $B(6,0)\Rightarrow36+6A=0\Rightarrow A=-6$. $C(0,8)\Rightarrow64+8B=0\Rightarrow B=-8$. Persamaan: $x^2+y^2-6x-8y=0$, yaitu $(x-3)^2+(y-4)^2=25.$

\knomor{C.6}
\ksoal $T$ pada $x^2+y^2=16$ bersudut $60^\circ$; cari $T$ dan busur dari $(4,0)$.
\kanalisis Parametrik $(r\cos\theta,r\sin\theta)$; busur $=r\theta$ (radian).
\kbahas $T=(4\cos60^\circ,4\sin60^\circ)=(2,\,2\sqrt3)$. Busur $=4\cdot\tfrac{\pi}{3}=\dfrac{4\pi}{3}.$

\kbagian{Bagian D}

\knomor{D.1}
\ksoal Buktikan garis singgung $x_1x+y_1y=r^2$ tegak lurus jari-jari di $(x_1,y_1)$.
\kanalisis Bandingkan gradien garis singgung dan gradien jari-jari $O(x_1,y_1)$.
\kbahas Gradien jari-jari $=\dfrac{y_1}{x_1}$. Garis singgung $x_1x+y_1y=r^2$ bergradien $-\dfrac{x_1}{y_1}$. Hasil kali gradien $=\dfrac{y_1}{x_1}\cdot\left(-\dfrac{x_1}{y_1}\right)=-1$, jadi keduanya tegak lurus. (Jika $y_1=0$, jari-jari mendatar dan garis singgung tegak --- tetap tegak lurus.) $\blacksquare$

\knomor{D.2}
\ksoal Titik pada $x^2+y^2=1$ yang terdekat/terjauh dari $(3,4)$.
\kanalisis Titik ekstrem terletak pada garis yang menghubungkan pusat $O$ ke $(3,4)$; jarak $O$ ke $(3,4)$ adalah $5$.
\kbahas Vektor satuan arah $(3,4)$ adalah $\left(\tfrac35,\tfrac45\right)$. Titik \textbf{terdekat} $\left(\tfrac35,\tfrac45\right)$ berjarak $5-1=4$; titik \textbf{terjauh} $\left(-\tfrac35,-\tfrac45\right)$ berjarak $5+1=6$.

\knomor{D.3}
\ksoal Garis kuasa dari $x^2+y^2=4$ dan $x^2+y^2-6x+5=0$.
\kanalisis Kurangkan kedua persamaan; suku kuadrat saling hapus.
\kbahas $(x^2+y^2-4)-(x^2+y^2-6x+5)=0\Rightarrow6x-9=0\Rightarrow x=\tfrac32$. Garis $x=\tfrac32$ adalah \textbf{sumbu radikal}: tempat kedudukan titik yang dayanya terhadap kedua lingkaran sama.

\knomor{D.4}
\ksoal Lingkaran melalui $(0,0)$, $(4,0)$, $(0,6)$.
\kanalisis Bentuk umum, substitusi tiap titik.
\kbahas $C=0$; $(4,0)\Rightarrow16+4A=0\Rightarrow A=-4$; $(0,6)\Rightarrow36+6B=0\Rightarrow B=-6$. Jadi $x^2+y^2-4x-6y=0$, yaitu $(x-2)^2+(y-3)^2=13$ (pusat $(2,3)$, $r=\sqrt{13}$).

\knomor{D.5}
\ksoal Buktikan $\big(a+r\cos\theta,\,b+r\sin\theta\big)$ pada $(x-a)^2+(y-b)^2=r^2$.
\kanalisis Substitusi langsung dan pakai identitas $\sin^2+\cos^2=1$.
\kbahas $(a+r\cos\theta-a)^2+(b+r\sin\theta-b)^2=r^2\cos^2\theta+r^2\sin^2\theta=r^2(\cos^2\theta+\sin^2\theta)=r^2.$ Terbukti untuk setiap $\theta$. $\blacksquare$

\kbagian{Challenge Corner}

\knomor{1}
\ksoal Panjang tali busur singgung $AB$ dari $P(7,1)$ ke $x^2+y^2=25$.
\kanalisis Tali busur singgung (chord of contact) berpersamaan $x_Px+y_Py=r^2$; hitung jarak pusat ke garis lalu setengah tali busur.
\kbahas Garis $AB$: $7x+y=25$. Jarak pusat $d=\dfrac{25}{\sqrt{49+1}}=\dfrac{25}{\sqrt{50}}=\dfrac{5}{\sqrt2}$. Maka $AB=2\sqrt{r^2-d^2}=2\sqrt{25-\tfrac{25}{2}}=2\sqrt{\tfrac{25}{2}}=2\cdot\dfrac{5}{\sqrt2}=5\sqrt2.$

\knomor{2}
\ksoal Lokus titik dengan $PA=2PB$, $A(0,0)$, $B(6,0)$.
\kanalisis Tulis $\sqrt{x^2+y^2}=2\sqrt{(x-6)^2+y^2}$ lalu kuadratkan.
\kbahas $x^2+y^2=4\big[(x-6)^2+y^2\big]=4x^2-48x+144+4y^2$. Maka $3x^2+3y^2-48x+144=0\Rightarrow x^2+y^2-16x+48=0$, yaitu $(x-8)^2+y^2=16$. Lokusnya \textbf{lingkaran} berpusat $(8,0)$ dengan jari-jari $4$ (lingkaran Apollonius).

% ============================================================
\chapter{Pembahasan Bab 3: Statistika Bivariat}
% ============================================================

\kbagian{Bagian A}

\knomor{A.1}
\ksoal Data $x:2,4,6$, $y:1,3,5$. Hitung $\bar x,\bar y$.
\kanalisis Rata-rata $=$ jumlah dibagi banyak data.
\kbahas $\bar x=\dfrac{2+4+6}{3}=4$, \quad $\bar y=\dfrac{1+3+5}{3}=3.$

\knomor{A.2}
\ksoal Titik cenderung turun dari kiri-atas ke kanan-bawah. Jenis korelasi?
\kanalisis Arah menurun menandakan hubungan berlawanan.
\kbahas Korelasi \textbf{negatif}.

\knomor{A.3}
\ksoal Tafsirkan $r=0{,}85$.
\kanalisis Tanda positif, $|r|\ge0{,}7$.
\kbahas Korelasi \textbf{positif} dan \textbf{kuat}: ketika satu variabel naik, yang lain cenderung naik dengan pola linear yang erat.

\knomor{A.4}
\ksoal $\hat y=3+2x$, prediksi pada $x=5$.
\kanalisis Substitusi langsung.
\kbahas $\hat y=3+2(5)=13.$

\knomor{A.5}
\ksoal $\hat y=10-0{,}5x$. Perubahan $\hat y$ jika $x$ naik $1$.
\kanalisis Gradien $b$ menyatakan perubahan $\hat y$ per satu satuan $x$.
\kbahas $\hat y$ \textbf{turun} $0{,}5$ satuan.

\kbagian{Bagian B}
\textit{Untuk data} $x:1,2,3,4,5$, $y:2,3,5,4,6$: $\bar x=3$, $\bar y=4$, $\sum x_iy_i=69$, $\sum x_i^2=55$, $\sum y_i^2=90$.

\knomor{B.1}
\ksoal Hitung $S_{xx}$ dan $S_{xy}$.
\kanalisis $S_{xx}=\sum x_i^2-n\bar x^2$, $S_{xy}=\sum x_iy_i-n\bar x\bar y$.
\kbahas $S_{xx}=55-5(9)=10$, \quad $S_{xy}=69-5(3)(4)=9.$

\knomor{B.2}
\ksoal Tentukan garis regresi.
\kanalisis $b=S_{xy}/S_{xx}$, $a=\bar y-b\bar x$.
\kbahas $b=\dfrac{9}{10}=0{,}9$, $a=4-0{,}9(3)=1{,}3$. Jadi $\hat y=1{,}3+0{,}9x.$

\knomor{B.3}
\ksoal Hitung $r$.
\kanalisis $r=\dfrac{S_{xy}}{\sqrt{S_{xx}S_{yy}}}$ dengan $S_{yy}=90-80=10$.
\kbahas $r=\dfrac{9}{\sqrt{10\cdot10}}=0{,}9$ (positif kuat).

\knomor{B.4}
\ksoal Hitung kovarians (populasi).
\kanalisis $\operatorname{cov}=\tfrac1n S_{xy}$.
\kbahas $\operatorname{cov}=\dfrac{9}{5}=1{,}8.$

\knomor{B.5}
\ksoal Prediksi $y$ pada $x=6$ dan residual jika $y$ nyata $=7$.
\kanalisis Pakai garis regresi lalu $e=y-\hat y$.
\kbahas $\hat y=1{,}3+0{,}9(6)=6{,}7$. Residual $e=7-6{,}7=0{,}3.$

\knomor{B.6}
\ksoal Tafsirkan $r=-0{,}95$; apakah membuktikan kausalitas?
\kanalisis Tanda negatif, $|r|$ sangat besar; korelasi $\neq$ sebab-akibat.
\kbahas Korelasi \textbf{negatif sangat kuat}: makin banyak jam bermain game, nilai ujian cenderung makin rendah. Namun ini \textbf{tidak} membuktikan game menyebabkan nilai turun --- mungkin ada faktor perancu (mis. kurang tidur, kurang waktu belajar) yang memengaruhi keduanya.

\knomor{B.7}
\ksoal Titik $(10,2)$ menyimpang jauh. Istilah dan pengaruhnya?
\kanalisis Titik ekstrem disebut outlier.
\kbahas Titik itu adalah \textbf{outlier} (pencilan). Outlier "menarik" garis regresi ke arahnya sehingga gradien dan perpotongan berubah, serta dapat menurunkan nilai $|r|$, membuat model kurang mewakili pola mayoritas data.

\kbagian{Bagian C}

\knomor{C.1}
\ksoal Buktikan garis regresi melalui $(\bar x,\bar y)$.
\kanalisis Substitusi $x=\bar x$ ke $\hat y=a+bx$ dengan $a=\bar y-b\bar x$.
\kbahas $\hat y(\bar x)=a+b\bar x=(\bar y-b\bar x)+b\bar x=\bar y.$ Jadi titik $(\bar x,\bar y)$ selalu dilalui. $\blacksquare$

\knomor{C.2}
\ksoal Pengaruh penggandaan $x\to2x$ terhadap $r$.
\kanalisis Tinjau perubahan $S_{xy}$ dan $\sqrt{S_{xx}S_{yy}}$.
\kbahas $S_{xy}\to2S_{xy}$, $S_{xx}\to4S_{xx}$, sehingga $r\to\dfrac{2S_{xy}}{\sqrt{4S_{xx}\cdot S_{yy}}}=\dfrac{2S_{xy}}{2\sqrt{S_{xx}S_{yy}}}=r$. Korelasi \textbf{tidak berubah}: $r$ tak terpengaruh oleh penskalaan linear positif.

\knomor{C.3}
\ksoal $S_{xy}=24$, $S_{xx}=16$, $\bar x=5$, $\bar y=20$. Regresi dan prediksi di $x=8$.
\kanalisis Hitung $b,a$ lalu substitusi.
\kbahas $b=\dfrac{24}{16}=1{,}5$, $a=20-1{,}5(5)=12{,}5$. $\hat y=12{,}5+1{,}5x$. Di $x=8$: $\hat y=12{,}5+12=24{,}5.$

\knomor{C.4}
\ksoal $S_{xy}=18$, $S_{xx}=25$, $S_{yy}=16$. Hitung $r$ dan $b$.
\kanalisis $r=\dfrac{S_{xy}}{\sqrt{S_{xx}S_{yy}}}$, $b=\dfrac{S_{xy}}{S_{xx}}$.
\kbahas $r=\dfrac{18}{\sqrt{25\cdot16}}=\dfrac{18}{20}=0{,}9$; \quad $b=\dfrac{18}{25}=0{,}72.$

\knomor{C.5}
\ksoal $\hat y=2+0{,}5x$, $r=0{,}6$. Tentukan gradien regresi $x$ atas $y$.
\kanalisis Gunakan $b_{yx}\cdot b_{xy}=r^2$.
\kbahas $b_{xy}=\dfrac{r^2}{b_{yx}}=\dfrac{0{,}36}{0{,}5}=0{,}72.$

\knomor{C.6}
\ksoal Es krim vs kasus tenggelam: mengapa bukan sebab-akibat?
\kanalisis Identifikasi variabel perancu.
\kbahas Keduanya dipicu oleh \textbf{cuaca panas}: saat panas, penjualan es krim naik \emph{dan} lebih banyak orang berenang sehingga kasus tenggelam naik. Cuaca adalah variabel perancu; es krim tidak menyebabkan orang tenggelam.

\kbagian{Bagian D}

\knomor{D.1}
\ksoal Buktikan $b=r\cdot\dfrac{s_y}{s_x}$.
\kanalisis Tulis $r,s_x,s_y$ dalam $S_{xx},S_{yy},S_{xy}$.
\kbahas $s_x=\sqrt{S_{xx}/n}$, $s_y=\sqrt{S_{yy}/n}$, $r=\dfrac{S_{xy}}{\sqrt{S_{xx}S_{yy}}}$. Maka
$r\cdot\dfrac{s_y}{s_x}=\dfrac{S_{xy}}{\sqrt{S_{xx}S_{yy}}}\cdot\dfrac{\sqrt{S_{yy}/n}}{\sqrt{S_{xx}/n}}=\dfrac{S_{xy}}{\sqrt{S_{xx}S_{yy}}}\cdot\sqrt{\dfrac{S_{yy}}{S_{xx}}}=\dfrac{S_{xy}}{S_{xx}}=b.$ $\blacksquare$

\knomor{D.2}
\ksoal Buktikan $|r|\le1$ dengan Cauchy--Schwarz.
\kanalisis Terapkan ketaksamaan pada $a_i=x_i-\bar x$, $b_i=y_i-\bar y$.
\kbahas Cauchy--Schwarz: $\left(\sum a_ib_i\right)^2\le\left(\sum a_i^2\right)\left(\sum b_i^2\right)$, yaitu $S_{xy}^2\le S_{xx}S_{yy}$. Maka $r^2=\dfrac{S_{xy}^2}{S_{xx}S_{yy}}\le1$, sehingga $|r|\le1$. $\blacksquare$

\knomor{D.3}
\ksoal Turunkan $b$ dan $a$ dengan kuadrat terkecil.
\kanalisis Minimumkan $Q(a,b)=\sum(y_i-a-bx_i)^2$; set turunan parsial nol.
\kbahas $\dfrac{\partial Q}{\partial a}=-2\sum(y_i-a-bx_i)=0\Rightarrow \sum y_i=na+b\sum x_i\Rightarrow \bar y=a+b\bar x\Rightarrow a=\bar y-b\bar x$.
$\dfrac{\partial Q}{\partial b}=-2\sum x_i(y_i-a-bx_i)=0\Rightarrow \sum x_iy_i=a\sum x_i+b\sum x_i^2$. Substitusi $a=\bar y-b\bar x$ dan sederhanakan memberi $S_{xy}=b\,S_{xx}$, jadi $b=\dfrac{S_{xy}}{S_{xx}}$. $\blacksquare$

\knomor{D.4}
\ksoal $y=c\,x^k$: tunjukkan $\log y$ linear terhadap $\log x$.
\kanalisis Ambil logaritma kedua ruas.
\kbahas $\log y=\log c+k\log x$. Dengan $Y=\log y$, $X=\log x$, diperoleh $Y=\log c+kX$ --- garis lurus bergradien $k$ dan perpotongan $\log c$. Maka regresi linear pada pasangan $(\log x,\log y)$ menaksir $k$ sebagai gradien dan $c=10^{(\text{perpotongan})}$.

\knomor{D.5}
\ksoal Jika $r=0$, pastikah tak ada hubungan? Beri contoh tandingan.
\kanalisis $r$ hanya mengukur hubungan \emph{linear}.
\kbahas Tidak. Ambil $x:-2,-1,0,1,2$ dan $y=x^2:4,1,0,1,4$. Di sini $\bar x=0$, sehingga $S_{xy}=\sum x_iy_i=(-2)(4)+(-1)(1)+0+(1)(1)+(2)(4)=0$, jadi $r=0$. Padahal $y$ \emph{sepenuhnya} ditentukan $x$ melalui hubungan kuadrat. Korelasi nol berarti tak ada hubungan \emph{linear}, bukan tak ada hubungan sama sekali.

\kbagian{Challenge Corner}

\knomor{1}
\ksoal Data $x:1,2,3,4$, $y:1,2,3,10$. Hitung $r$, lalu tanpa titik terakhir.
\kanalisis Hitung statistik ringkas, lalu ulangi untuk tiga titik pertama.
\kbahas $\bar x=2{,}5$, $\bar y=4$. $\sum x_iy_i=54$, $\sum x_i^2=30$, $\sum y_i^2=114$. $S_{xx}=30-25=5$, $S_{xy}=54-40=14$, $S_{yy}=114-64=50$. Maka $r=\dfrac{14}{\sqrt{5\cdot50}}=\dfrac{14}{\sqrt{250}}\approx0{,}885$. \\
Tanpa titik $(4,10)$: data $x:1,2,3$, $y:1,2,3$ tepat segaris, sehingga $r=1$. Jadi satu outlier menurunkan korelasi dari $1$ menjadi $\approx0{,}885$ --- menegaskan pentingnya memeriksa scatter plot.

\knomor{2}
\ksoal Buktikan $\sum(y_i-\hat y_i)=0$ pada regresi kuadrat terkecil.
\kanalisis Gunakan $\hat y_i=a+bx_i$ dengan $a=\bar y-b\bar x$.
\kbahas $\sum(y_i-\hat y_i)=\sum(y_i-a-bx_i)=\sum y_i-na-b\sum x_i=n\bar y-na-bn\bar x=n(\bar y-a-b\bar x)$. Karena $a=\bar y-b\bar x$, maka $\bar y-a-b\bar x=0$, sehingga jumlah seluruh residual $=0$. $\blacksquare$

% ============================================================
\chapter{Pembahasan Bab 4: Bilangan Kompleks}
% ============================================================

\kbagian{Bagian A}

\knomor{A.1}
\ksoal Hitung $i^2,i^3,i^4,i^{10}$.
\kanalisis Pangkat $i$ berulang periode $4$; tinjau sisa bagi $4$.
\kbahas $i^2=-1$, $i^3=-i$, $i^4=1$, dan $i^{10}=i^{4\cdot2+2}=i^2=-1.$

\knomor{A.2}
\ksoal $(3+2i)+(1-4i)$.
\kanalisis Jumlahkan bagian real dan imajiner terpisah.
\kbahas $(3+1)+(2-4)i=4-2i.$

\knomor{A.3}
\ksoal $(2+i)(3-i)$.
\kanalisis Kalikan distributif, ganti $i^2=-1$.
\kbahas $6-2i+3i-i^2=6+i+1=7+i.$

\knomor{A.4}
\ksoal Konjugat dan modulus $z=3-4i$.
\kanalisis $\bar z$ membalik tanda imajiner; $|z|=\sqrt{a^2+b^2}$.
\kbahas $\bar z=3+4i$, $|z|=\sqrt{9+16}=5.$

\knomor{A.5}
\ksoal Ubah $z=1+i$ ke polar.
\kanalisis $r=\sqrt{a^2+b^2}$, $\theta$ dari kuadran.
\kbahas $r=\sqrt2$, $\theta=45^\circ$, jadi $z=\sqrt2\,\mathrm{cis}\,45^\circ.$

\kbagian{Bagian B}

\knomor{B.1}
\ksoal $(5+3i)-(2-i)$.
\kanalisis Kurangkan bagian sejenis.
\kbahas $(5-2)+(3-(-1))i=3+4i.$

\knomor{B.2}
\ksoal $\dfrac{3+2i}{1-i}$.
\kanalisis Kalikan konjugat penyebut $1+i$.
\kbahas $\dfrac{(3+2i)(1+i)}{(1-i)(1+i)}=\dfrac{3+3i+2i+2i^2}{2}=\dfrac{1+5i}{2}=\dfrac12+\dfrac52 i.$

\knomor{B.3}
\ksoal Modulus dan argumen $z=-1+\sqrt3\,i$.
\kanalisis $r=\sqrt{a^2+b^2}$; $z$ di kuadran II.
\kbahas $r=\sqrt{1+3}=2$. $\tan\theta=\dfrac{\sqrt3}{-1}=-\sqrt3$ dengan $z$ di kuadran II, maka $\theta=120^\circ$.

\knomor{B.4}
\ksoal $i^{2025}$.
\kanalisis $2025\bmod4$.
\kbahas $2025=4\cdot506+1$, jadi $i^{2025}=i^1=i.$

\knomor{B.5}
\ksoal $z=2+i$: hitung $z\bar z$ dan bandingkan $|z|^2$.
\kanalisis $z\bar z=(a+bi)(a-bi)=a^2+b^2$.
\kbahas $z\bar z=(2+i)(2-i)=4-i^2=5$; $|z|^2=2^2+1^2=5$. Sama. $\checkmark$

\knomor{B.6}
\ksoal Selesaikan $z^2=-9$.
\kanalisis Akar dari bilangan negatif memakai $i$.
\kbahas $z=\pm\sqrt{-9}=\pm3i.$

\knomor{B.7}
\ksoal $z=2(\cos60^\circ+i\sin60^\circ)$ ke kartesius.
\kanalisis Substitusi nilai $\cos60^\circ,\sin60^\circ$.
\kbahas $z=2\left(\tfrac12+i\tfrac{\sqrt3}{2}\right)=1+\sqrt3\,i.$

\kbagian{Bagian C}

\knomor{C.1}
\ksoal $(1+i)^8$ dengan polar.
\kanalisis $1+i=\sqrt2\,\mathrm{cis}\,45^\circ$; pakai De Moivre.
\kbahas $(1+i)^8=(\sqrt2)^8\,\mathrm{cis}(8\cdot45^\circ)=16\,\mathrm{cis}\,360^\circ=16.$

\knomor{C.2}
\ksoal $x^2-4x+13=0$.
\kanalisis Rumus kuadrat; diskriminan negatif.
\kbahas $x=\dfrac{4\pm\sqrt{16-52}}{2}=\dfrac{4\pm\sqrt{-36}}{2}=\dfrac{4\pm6i}{2}=2\pm3i.$

\knomor{C.3}
\ksoal Semua akar pangkat tiga dari $8$.
\kanalisis $8=8\,\mathrm{cis}\,0^\circ$; rumus akar ke-$n$.
\kbahas $\sqrt[3]{8}=2$, sudut $0^\circ,120^\circ,240^\circ$: $\ w_0=2,\ w_1=-1+\sqrt3\,i,\ w_2=-1-\sqrt3\,i.$

\knomor{C.4}
\ksoal Turunkan $\cos2\theta,\sin2\theta$.
\kanalisis Kuadratkan $\cos\theta+i\sin\theta$ dua cara.
\kbahas $(\cos\theta+i\sin\theta)^2=\cos^2\theta-\sin^2\theta+i(2\sin\theta\cos\theta)$. Menurut De Moivre ini juga $\cos2\theta+i\sin2\theta$. Samakan bagian real dan imajiner: $\cos2\theta=\cos^2\theta-\sin^2\theta$, $\sin2\theta=2\sin\theta\cos\theta.$

\knomor{C.5}
\ksoal $z=3+4i$ dikali $i$ (rotasi $90^\circ$).
\kanalisis Kalikan dengan $i$.
\kbahas $i(3+4i)=3i+4i^2=-4+3i.$ (Modulus tetap $5$, sudut bertambah $90^\circ$.)

\knomor{C.6}
\ksoal Buktikan $|z_1z_2|=|z_1||z_2|$.
\kanalisis Tulis dalam bentuk polar.
\kbahas Misal $z_1=r_1\mathrm{cis}\,\theta_1$, $z_2=r_2\mathrm{cis}\,\theta_2$. Maka $z_1z_2=r_1r_2\,\mathrm{cis}(\theta_1+\theta_2)$, sehingga $|z_1z_2|=r_1r_2=|z_1||z_2|$. $\blacksquare$

\kbagian{Bagian D}

\knomor{D.1}
\ksoal Buktikan De Moivre dengan induksi.
\kanalisis Basis $n=1$; langkah induksi pakai rumus jumlah sudut.
\kbahas \emph{Basis:} $n=1$ jelas. \emph{Induksi:} andaikan $(\cos\theta+i\sin\theta)^n=\cos n\theta+i\sin n\theta$. Maka
$(\cos\theta+i\sin\theta)^{n+1}=(\cos n\theta+i\sin n\theta)(\cos\theta+i\sin\theta)$
$=(\cos n\theta\cos\theta-\sin n\theta\sin\theta)+i(\sin n\theta\cos\theta+\cos n\theta\sin\theta)$
$=\cos(n+1)\theta+i\sin(n+1)\theta.$ Dengan induksi terbukti untuk semua $n\in\mathbb{N}$. $\blacksquare$

\knomor{D.2}
\ksoal Semua akar $z^4=1$ dan jumlahnya.
\kanalisis Akar ke-$4$ dari $1=\mathrm{cis}\,0^\circ$: sudut $0^\circ,90^\circ,180^\circ,270^\circ$.
\kbahas Akar: $1,\,i,\,-1,\,-i$. Jumlah $=1+i-1-i=0.$

\knomor{D.3}
\ksoal Jumlah $1+i+i^2+\cdots+i^{2025}$.
\kanalisis Setiap $4$ suku berurutan berjumlah $0$; ada $2026$ suku.
\kbahas $1+i+i^2+i^3=0$. Banyak suku $2026=4\cdot506+2$, jadi $506$ kelompok berjumlah $0$ dan sisa dua suku $i^{2024}+i^{2025}=i^0+i^1=1+i$. Jumlah $=1+i.$

\knomor{D.4}
\ksoal Buktikan $\cos3\theta=4\cos^3\theta-3\cos\theta$.
\kanalisis Ambil bagian real dari $(\cos\theta+i\sin\theta)^3$.
\kbahas Bagian real dari $(\cos\theta+i\sin\theta)^3$ adalah $\cos^3\theta-3\cos\theta\sin^2\theta$, dan menurut De Moivre sama dengan $\cos3\theta$. Maka
$\cos3\theta=\cos^3\theta-3\cos\theta(1-\cos^2\theta)=4\cos^3\theta-3\cos\theta.$ $\blacksquare$

\knomor{D.5}
\ksoal Akar ke-$n$ dari $1$: segi-$n$ beraturan dan jumlah nol.
\kanalisis Akar $\zeta_k=\mathrm{cis}\,\tfrac{360^\circ k}{n}$; gunakan deret geometri.
\kbahas Akar-akarnya $\zeta_k=\mathrm{cis}\,\tfrac{360^\circ k}{n}$, $k=0,\dots,n-1$, berjarak sama ($\tfrac{360^\circ}{n}$) pada lingkaran satuan --- yaitu titik sudut segi-$n$ beraturan. Jumlahnya adalah deret geometri dengan rasio $\zeta=\mathrm{cis}\,\tfrac{360^\circ}{n}\neq1$:
$\sum_{k=0}^{n-1}\zeta^k=\dfrac{\zeta^n-1}{\zeta-1}=\dfrac{1-1}{\zeta-1}=0$ untuk $n\ge2$. $\blacksquare$

\kbagian{Challenge Corner}

\knomor{1}
\ksoal $\big(1+i\sqrt3\big)^{12}$.
\kanalisis Ubah ke polar lalu De Moivre.
\kbahas $1+i\sqrt3=2\,\mathrm{cis}\,60^\circ$. Maka $\big(1+i\sqrt3\big)^{12}=2^{12}\,\mathrm{cis}(12\cdot60^\circ)=4096\,\mathrm{cis}\,720^\circ=4096.$

\knomor{2}
\ksoal Semua akar pangkat enam dari $-64$.
\kanalisis $-64=64\,\mathrm{cis}\,180^\circ$; $\sqrt[6]{64}=2$, sudut $\tfrac{180^\circ+360^\circ k}{6}$.
\kbahas Untuk $k=0,\dots,5$ sudutnya $30^\circ,90^\circ,150^\circ,210^\circ,270^\circ,330^\circ$, sehingga akarnya
$2\,\mathrm{cis}\,30^\circ,\ 2\,\mathrm{cis}\,90^\circ,\ 2\,\mathrm{cis}\,150^\circ,\ 2\,\mathrm{cis}\,210^\circ,\ 2\,\mathrm{cis}\,270^\circ,\ 2\,\mathrm{cis}\,330^\circ$
(yaitu $\sqrt3+i,\ 2i,\ -\sqrt3+i,\ -\sqrt3-i,\ -2i,\ \sqrt3-i$).

\knomor{3}
\ksoal Buktikan $z+\tfrac1z=2\cos\theta\Rightarrow z^n+\tfrac1{z^n}=2\cos n\theta$.
\kanalisis Tunjukkan $z=\mathrm{cis}(\pm\theta)$, lalu pakai De Moivre.
\kbahas Dari $z+\tfrac1z=2\cos\theta$ diperoleh $z^2-2\cos\theta\,z+1=0$, sehingga $z=\cos\theta\pm i\sin\theta=\mathrm{cis}(\pm\theta)$. Ambil $z=\mathrm{cis}\,\theta$; maka $z^n=\mathrm{cis}\,n\theta$ dan $\tfrac1{z^n}=\mathrm{cis}(-n\theta)$. Jadi
$z^n+\dfrac1{z^n}=\mathrm{cis}\,n\theta+\mathrm{cis}(-n\theta)=2\cos n\theta.$ (Kasus $z=\mathrm{cis}(-\theta)$ memberi hasil sama.) $\blacksquare$

% ============================================================
\chapter{Pembahasan Bab 5: Polinomial}
% ============================================================

\kbagian{Bagian A}

\knomor{A.1}
\ksoal Sisa $P(x)=2x^3+x^2-7$ dibagi $(x-2)$.
\kanalisis Teorema sisa: sisa $=P(2)$.
\kbahas $P(2)=2(8)+4-7=16+4-7=13.$

\knomor{A.2}
\ksoal Tunjukkan $(x-1)$ faktor $x^3-1$.
\kanalisis Teorema faktor: cek $P(1)$.
\kbahas $P(1)=1-1=0$, jadi $(x-1)$ faktor. (Memang $x^3-1=(x-1)(x^2+x+1)$.)

\knomor{A.3}
\ksoal Faktorkan $x^3+2x^2-5x-6$ ($x=-1$ akar).
\kanalisis Bagi dengan $(x+1)$ lalu faktorkan kuadrat.
\kbahas Horner dengan $k=-1$ memberi $x^3+2x^2-5x-6=(x+1)(x^2+x-6)=(x+1)(x+3)(x-2).$

\knomor{A.4}
\ksoal Akar $x^2-7x+10=0$: jumlah dan hasil kali.
\kanalisis Vieta: $-\tfrac ba$ dan $\tfrac ca$.
\kbahas $x_1+x_2=7$, $x_1x_2=10.$

\knomor{A.5}
\ksoal Kandidat akar rasional $x^3-4x+1$.
\kanalisis $\tfrac pq$ dengan $p\mid1$, $q\mid1$.
\kbahas Kandidat: $\pm1$.

\kbagian{Bagian B}

\knomor{B.1}
\ksoal $(x+2)$ faktor $x^3+ax^2-x+6$; tentukan $a$.
\kanalisis Teorema faktor: $P(-2)=0$.
\kbahas $P(-2)=-8+4a+2+6=4a=0\Rightarrow a=0.$

\knomor{B.2}
\ksoal Semua akar $2x^3-3x^2-3x+2$.
\kanalisis Uji kandidat $\pm1,\pm2,\pm\tfrac12$.
\kbahas $x=2$: $16-12-6+2=0$. Horner: $(x-2)(2x^2+x-1)=(x-2)(2x-1)(x+1)$. Akar: $x=2,\ \tfrac12,\ -1.$

\knomor{B.3}
\ksoal Akar $x^2-5x+6=0$: hitung $x_1^2+x_2^2$ dan $\tfrac1{x_1}+\tfrac1{x_2}$.
\kanalisis Vieta: $x_1+x_2=5$, $x_1x_2=6$.
\kbahas $x_1^2+x_2^2=(x_1+x_2)^2-2x_1x_2=25-12=13$. $\dfrac1{x_1}+\dfrac1{x_2}=\dfrac{x_1+x_2}{x_1x_2}=\dfrac{5}{6}.$

\knomor{B.4}
\ksoal Akar $x^3-6x^2+11x-6$.
\kanalisis Kandidat $\pm1,\pm2,\pm3,\pm6$.
\kbahas $P(1)=0$, lalu $(x-1)(x^2-5x+6)=(x-1)(x-2)(x-3)$. Akar: $1,2,3.$

\knomor{B.5}
\ksoal $P(x)$ bersisa $3$ dibagi $(x-1)$ dan $-3$ dibagi $(x+2)$; sisa dibagi $(x-1)(x+2)$.
\kanalisis Sisa berbentuk $px+q$; pakai $P(1)=3$, $P(-2)=-3$.
\kbahas Tulis sisa $=px+q$. $P(1)=p+q=3$ dan $P(-2)=-2p+q=-3$. Kurangkan: $3p=6\Rightarrow p=2$, $q=1$. Sisa $=2x+1.$

\knomor{B.6}
\ksoal $x=3$ akar multiplisitas $\ge2$ dari $x^3-9x^2+kx-27$.
\kanalisis Syarat $P(3)=0$ dan $P'(3)=0$.
\kbahas $P(3)=27-81+3k-27=3k-81=0\Rightarrow k=27$. Cek $P'(x)=3x^2-18x+k$, $P'(3)=27-54+27=0$. $\checkmark$ (Bahkan $P(x)=(x-3)^3$.) Jadi $k=27.$

\kbagian{Bagian C}

\knomor{C.1}
\ksoal Akar $x^3-6x^2+11x-6=0$ aritmetika; tentukan beda.
\kanalisis Akar $1,2,3$ dari B.4.
\kbahas Akar $1,2,3$ membentuk barisan aritmetika dengan beda $1$. (Verifikasi Vieta: jumlah $6$, hasil kali $6$, sesuai.)

\knomor{C.2}
\ksoal Selesaikan $8^x-7\cdot4^x+14\cdot2^x-8=0$.
\kanalisis Substitusi $u=2^x$ ($8^x=u^3$, $4^x=u^2$).
\kbahas $u^3-7u^2+14u-8=0\Rightarrow(u-1)(u-2)(u-4)=0\Rightarrow u=1,2,4$. Maka $2^x=1,2,4$ memberi $x=0,1,2.$

\knomor{C.3}
\ksoal $x_1,x_2,x_3$ akar $x^3+px+q=0$; nyatakan $\sum x_i^2$.
\kanalisis Vieta: $\sum x_i=0$, $\sum x_ix_j=p$.
\kbahas $\sum x_i^2=\left(\sum x_i\right)^2-2\sum_{i<j}x_ix_j=0-2p=-2p.$

\knomor{C.4}
\ksoal Buktikan polinomial bulat ganjil di $0$ dan $1$ tak punya akar bulat.
\kanalisis Tinjau paritas $P(r)$ menurut paritas $r$.
\kbahas Untuk akar bulat $r$: jika $r$ genap maka $P(r)\equiv P(0)\pmod 2$ (semua suku ber-$r$ genap), yang ganjil $\neq0$. Jika $r$ ganjil maka $r^k\equiv1\pmod2$, sehingga $P(r)\equiv P(1)\pmod2$, juga ganjil $\neq0$. Karena $P(r)$ selalu ganjil, tak mungkin $P(r)=0$. $\blacksquare$

\knomor{C.5}
\ksoal Perilaku $P(x)=(x-1)^2(x+2)$.
\kanalisis Tinjau akar, multiplisitas, derajat, koefisien utama.
\kbahas Akar $x=1$ multiplisitas $2$ (kurva \emph{menyinggung} sumbu-$x$ lalu berbalik) dan $x=-2$ multiplisitas $1$ (kurva \emph{memotong}). Derajat $3$ dengan koefisien utama positif: $x\to-\infty\Rightarrow P\to-\infty$, $x\to+\infty\Rightarrow P\to+\infty$.

\knomor{C.6}
\ksoal $x^3-3x^2+kx-1=0$ akar $a,b,c$; jika $a+b+c=ab+bc+ca$, cari $k$.
\kanalisis Vieta: $a+b+c=3$, $ab+bc+ca=k$.
\kbahas Syarat $3=k$, jadi $k=3.$

\kbagian{Bagian D}

\knomor{D.1}
\ksoal Buktikan Vieta untuk kuadrat.
\kanalisis Jabarkan $(x-x_1)(x-x_2)$.
\kbahas $(x-x_1)(x-x_2)=x^2-(x_1+x_2)x+x_1x_2$. Samakan dengan $x^2+bx+c$: $-(x_1+x_2)=b\Rightarrow x_1+x_2=-b$, dan $x_1x_2=c$. $\blacksquare$

\knomor{D.2}
\ksoal Akar tak-real berpasangan konjugat; susun kubik monik real berakar $2$ dan $1+i$.
\kanalisis Jika koefisien real, $P(\bar z)=\overline{P(z)}$, sehingga $z$ akar $\Rightarrow\bar z$ akar. Maka $1-i$ juga akar.
\kbahas Akar: $2,\ 1+i,\ 1-i$. $P(x)=(x-2)\big((x-1)^2+1\big)=(x-2)(x^2-2x+2)=x^3-4x^2+6x-4.$

\knomor{D.3}
\ksoal $x_1,x_2,x_3$ akar $x^3+px+q=0$; hitung $\sum x_i^3$.
\kanalisis Tiap akar memenuhi $x_i^3=-px_i-q$.
\kbahas Jumlahkan: $\sum x_i^3=-p\sum x_i-3q=-p(0)-3q=-3q.$

\knomor{D.4}
\ksoal Buktikan $k$ akar multiplisitas $\ge2$ $\Leftrightarrow$ $P(k)=0$ dan $P'(k)=0$.
\kanalisis Tulis $P(x)=(x-k)^mQ(x)$ dengan $Q(k)\neq0$ dan turunkan.
\kbahas ($\Rightarrow$) Jika $m\ge2$: $P(k)=0$, dan $P'(x)=m(x-k)^{m-1}Q(x)+(x-k)^mQ'(x)$, sehingga $P'(k)=0$ (faktor $(x-k)^{m-1}$ dengan $m-1\ge1$). \\
($\Leftarrow$) Jika $m=1$: $P(x)=(x-k)Q(x)$, $P'(k)=Q(k)\neq0$. Jadi $P'(k)=0$ memaksa $m\ge2$. $\blacksquare$

\knomor{D.5}
\ksoal Buktikan polinomial derajat $n$ punya paling banyak $n$ akar berbeda.
\kanalisis Induksi memakai teorema faktor.
\kbahas Setiap akar $k$ memberi faktor $(x-k)$ (teorema faktor), menurunkan derajat sebesar $1$. Jika ada $n+1$ akar berbeda, $P(x)$ habis dibagi hasil kali $n+1$ faktor linear berbeda, sehingga derajatnya $\ge n+1>n$ --- kontradiksi (untuk $P\not\equiv0$). Maka paling banyak $n$ akar. $\blacksquare$

\kbagian{Challenge Corner}

\knomor{1}
\ksoal $x+\tfrac1x=3$; cari $x^3+\tfrac1{x^3}$.
\kanalisis Pakai identitas $a^3+b^3=(a+b)^3-3ab(a+b)$ dengan $a=x,b=\tfrac1x$.
\kbahas $x^3+\dfrac1{x^3}=\left(x+\dfrac1x\right)^3-3\left(x+\dfrac1x\right)=27-9=18.$

\knomor{2}
\ksoal $P$ monik derajat $4$, $P(k)=k^2$ untuk $k=1,2,3,4$. Cari $P(5)$.
\kanalisis Tinjau $Q(x)=P(x)-x^2$, monik derajat $4$ dengan akar $1,2,3,4$.
\kbahas $Q(x)=(x-1)(x-2)(x-3)(x-4)$, jadi $P(x)=x^2+(x-1)(x-2)(x-3)(x-4)$. Maka $P(5)=25+(4)(3)(2)(1)=25+24=49.$

\knomor{3}
\ksoal Buktikan jumlah kuadrat akar $x^n+a_{n-1}x^{n-1}+\cdots+a_0=0$ adalah $a_{n-1}^2-2a_{n-2}$.
\kanalisis Vieta: $\sum x_i=-a_{n-1}$, $\sum_{i<j}x_ix_j=a_{n-2}$.
\kbahas $\sum x_i^2=\left(\sum x_i\right)^2-2\sum_{i<j}x_ix_j=(-a_{n-1})^2-2a_{n-2}=a_{n-1}^2-2a_{n-2}.$ $\blacksquare$

% ============================================================
\chapter{Pembahasan Bab 6: Matriks}
% ============================================================

\kbagian{Bagian A}

\knomor{A.1}
\ksoal $A+B$ untuk $A=\begin{pmatrix}1&2\\3&4\end{pmatrix}$, $B=\begin{pmatrix}0&1\\-1&2\end{pmatrix}$.
\kanalisis Jumlahkan elemen seletak.
\kbahas $A+B=\begin{pmatrix}1&3\\2&6\end{pmatrix}.$

\knomor{A.2}
\ksoal $\det\begin{pmatrix}3&2\\1&4\end{pmatrix}$.
\kanalisis $ad-bc$.
\kbahas $12-2=10.$

\knomor{A.3}
\ksoal Invers $\begin{pmatrix}1&2\\3&5\end{pmatrix}$.
\kanalisis $\det=-1$; pakai rumus invers ordo $2$.
\kbahas $A^{-1}=\dfrac{1}{-1}\begin{pmatrix}5&-2\\-3&1\end{pmatrix}=\begin{pmatrix}-5&2\\3&-1\end{pmatrix}.$

\knomor{A.4}
\ksoal $\begin{pmatrix}1&2\\3&4\end{pmatrix}\begin{pmatrix}1\\1\end{pmatrix}$.
\kanalisis Perkalian matriks $\times$ vektor.
\kbahas $\begin{pmatrix}1+2\\3+4\end{pmatrix}=\begin{pmatrix}3\\7\end{pmatrix}.$

\knomor{A.5}
\ksoal Jenis matriks $\begin{pmatrix}2&0&0\\0&5&0\\0&0&3\end{pmatrix}$.
\kanalisis Elemen luar diagonal nol.
\kbahas Matriks \textbf{diagonal}.

\kbagian{Bagian B}

\knomor{B.1}
\ksoal $\det\begin{pmatrix}1&2&3\\0&1&4\\5&6&0\end{pmatrix}$ (Sarrus).
\kanalisis $aei+bfg+cdh-ceg-afh-bdi$.
\kbahas $=0+40+0-15-24-0=1.$

\knomor{B.2}
\ksoal $A=\begin{pmatrix}2&1\\1&3\end{pmatrix}$ simetris; hitung $A^2$.
\kanalisis $A^\top=A$ karena $a_{12}=a_{21}=1$.
\kbahas Simetris. $A^2=\begin{pmatrix}2&1\\1&3\end{pmatrix}\begin{pmatrix}2&1\\1&3\end{pmatrix}=\begin{pmatrix}5&5\\5&10\end{pmatrix}.$

\knomor{B.3}
\ksoal Selesaikan sistem dengan invers.
\kanalisis $A=\begin{pmatrix}1&2\\3&5\end{pmatrix}$, $\det=-1$, $A^{-1}=\begin{pmatrix}-5&2\\3&-1\end{pmatrix}$.
\kbahas $\mathbf x=A^{-1}\begin{pmatrix}4\\7\end{pmatrix}=\begin{pmatrix}-20+14\\12-7\end{pmatrix}=\begin{pmatrix}-6\\5\end{pmatrix}$. Jadi $x=-6,\ y=5.$

\knomor{B.4}
\ksoal Sistem yang sama dengan Cramer.
\kanalisis $x=\tfrac{\det A_1}{\det A}$, $y=\tfrac{\det A_2}{\det A}$.
\kbahas $\det A=-1$, $\det A_1=\det\begin{pmatrix}4&2\\7&5\end{pmatrix}=6$, $\det A_2=\det\begin{pmatrix}1&4\\3&7\end{pmatrix}=-5$. Maka $x=\tfrac{6}{-1}=-6$, $y=\tfrac{-5}{-1}=5.$

\knomor{B.5}
\ksoal $k$ agar $\begin{pmatrix}k&2\\8&k\end{pmatrix}$ singular.
\kanalisis Singular $\Leftrightarrow\det=0$.
\kbahas $k^2-16=0\Rightarrow k=\pm4.$

\knomor{B.6}
\ksoal $\det A=3$, ordo $2$. Hitung $\det(2A)$ dan $\det(A^{-1})$.
\kanalisis $\det(kA)=k^n\det A$; $\det(A^{-1})=1/\det A$.
\kbahas $\det(2A)=2^2\cdot3=12$; $\det(A^{-1})=\tfrac13.$

\knomor{B.7}
\ksoal $AB=I$, $A=\begin{pmatrix}2&1\\3&2\end{pmatrix}$. Tentukan $B$.
\kanalisis $B=A^{-1}$; $\det A=1$.
\kbahas $B=A^{-1}=\dfrac11\begin{pmatrix}2&-1\\-3&2\end{pmatrix}=\begin{pmatrix}2&-1\\-3&2\end{pmatrix}.$

\kbagian{Bagian C}

\knomor{C.1}
\ksoal $R=\begin{pmatrix}0&-1\\1&0\end{pmatrix}$; hitung $R^4$.
\kanalisis $R$ rotasi $90^\circ$; $R^2$ rotasi $180^\circ$.
\kbahas $R^2=\begin{pmatrix}-1&0\\0&-1\end{pmatrix}=-I$, maka $R^4=(-I)^2=I$. Artinya empat kali rotasi $90^\circ$ menghasilkan rotasi $360^\circ$, yaitu kembali ke posisi semula (identitas).

\knomor{C.2}
\ksoal Determinan koefisien dan $x$ (Cramer) untuk sistem 3 variabel.
\kanalisis Sarrus untuk $\det A$ dan $\det A_1$.
\kbahas $A=\begin{pmatrix}1&1&1\\2&-1&1\\1&2&-1\end{pmatrix}$, $\det A=7$. Ganti kolom pertama dengan $\begin{pmatrix}6\\3\\2\end{pmatrix}$: $\det A_1=7$. Maka $x=\tfrac{7}{7}=1.$ (Lengkapnya $x=1,y=2,z=3$.)

\knomor{C.3}
\ksoal Buktikan $\det(AB)=\det A\det B$ (ordo $2$).
\kanalisis Hitung langsung dengan $A=\begin{pmatrix}a&b\\c&d\end{pmatrix}$, $B=\begin{pmatrix}e&f\\g&h\end{pmatrix}$.
\kbahas $AB=\begin{pmatrix}ae+bg&af+bh\\ce+dg&cf+dh\end{pmatrix}$. Setelah dijabarkan,
$\det(AB)=(ae+bg)(cf+dh)-(af+bh)(ce+dg)=(ad-bc)(eh-fg)=\det A\cdot\det B.$ $\blacksquare$

\knomor{C.4}
\ksoal $A=\begin{pmatrix}1&1\\0&1\end{pmatrix}$; tentukan $A^n$.
\kanalisis Hitung $A^2,A^3$, duga pola, buktikan induksi.
\kbahas Dugaan $A^n=\begin{pmatrix}1&n\\0&1\end{pmatrix}$. \emph{Basis} $n=1$ benar. \emph{Induksi:}
\[
A^{n+1}=A^nA=\begin{pmatrix}1&n\\0&1\end{pmatrix}\begin{pmatrix}1&1\\0&1\end{pmatrix}=\begin{pmatrix}1&n+1\\0&1\end{pmatrix}.\quad\blacksquare
\]

\knomor{C.5}
\ksoal Pendapatan tiap cabang dari tabel penjualan dan harga.
\kanalisis Kalikan vektor harga (baris) dengan matriks penjualan.
\kbahas $\begin{pmatrix}20&35\end{pmatrix}\begin{pmatrix}12&9&15\\7&11&8\end{pmatrix}=\begin{pmatrix}485&565&580\end{pmatrix}$ (ribu rupiah). Jadi pendapatan cabang $1,2,3$ berturut-turut Rp$485$rb, Rp$565$rb, Rp$580$rb.

\knomor{C.6}
\ksoal Luas segitiga via determinan; hitung untuk $(0,0),(4,0),(1,3)$.
\kanalisis Vektor sisi $(x_2-x_1,y_2-y_1)$ dan $(x_3-x_1,y_3-y_1)$; luas $=\tfrac12|\det|$.
\kbahas $\det\begin{pmatrix}4&1\\0&3\end{pmatrix}=12$, sehingga luas $=\tfrac12|12|=6$ satuan luas.

\kbagian{Bagian D}

\knomor{D.1}
\ksoal Buktikan $A^2-(\operatorname{tr}A)A+(\det A)I=O$ (ordo $2$).
\kanalisis Hitung langsung dengan $A=\begin{pmatrix}a&b\\c&d\end{pmatrix}$.
\kbahas $A^2=\begin{pmatrix}a^2+bc&b(a+d)\\c(a+d)&d^2+bc\end{pmatrix}$, $(\operatorname{tr}A)A=(a+d)\begin{pmatrix}a&b\\c&d\end{pmatrix}$. Maka
$A^2-(a+d)A=\begin{pmatrix}bc-ad&0\\0&bc-ad\end{pmatrix}=-(ad-bc)I=-(\det A)I.$
Jadi $A^2-(\operatorname{tr}A)A+(\det A)I=O$. $\blacksquare$

\knomor{D.2}
\ksoal Buktikan $(AB)^{-1}=B^{-1}A^{-1}$.
\kanalisis Tunjukkan hasil kali dengan $AB$ menghasilkan $I$.
\kbahas $(AB)(B^{-1}A^{-1})=A(BB^{-1})A^{-1}=AIA^{-1}=AA^{-1}=I$, dan serupa dari kiri. Karena invers tunggal, $(AB)^{-1}=B^{-1}A^{-1}$. $\blacksquare$

\knomor{D.3}
\ksoal Buktikan $\det(A^{-1})=\dfrac1{\det A}$.
\kanalisis Terapkan $\det(AB)=\det A\det B$ pada $AA^{-1}=I$.
\kbahas $\det(AA^{-1})=\det A\cdot\det(A^{-1})$, dan $\det(I)=1$. Maka $\det A\cdot\det(A^{-1})=1$, sehingga $\det(A^{-1})=\dfrac1{\det A}$ (untuk $\det A\neq0$). $\blacksquare$

\knomor{D.4}
\ksoal $A=\begin{pmatrix}2&0\\0&3\end{pmatrix}$; tentukan $A^n$.
\kanalisis Perkalian matriks diagonal memangkatkan tiap elemen diagonal.
\kbahas $A^n=\begin{pmatrix}2^n&0\\0&3^n\end{pmatrix}$. Matriks diagonal mudah dipangkatkan karena tidak ada "pencampuran" antarbaris/kolom --- tiap arah berdiri sendiri.

\knomor{D.5}
\ksoal Tunjukkan $\det A=0\Rightarrow A\mathbf x=\mathbf b$ tak bersolusi tunggal.
\kanalisis $\det A=0$ berarti baris bergantung linear.
\kbahas Jika $\det A=0$, $A$ tak punya invers sehingga solusi tak tunggal. Contoh $A=\begin{pmatrix}1&2\\2&4\end{pmatrix}$: untuk $\mathbf b=\begin{pmatrix}1\\3\end{pmatrix}$ sistem \emph{tak konsisten} (tak ada solusi, karena $2(x+2y)=2\neq3$); untuk $\mathbf b=\begin{pmatrix}1\\2\end{pmatrix}$ kedua persamaan identik sehingga ada \emph{tak hingga} solusi. $\blacksquare$

\kbagian{Challenge Corner}

\knomor{1}
\ksoal $F=\begin{pmatrix}1&1\\1&0\end{pmatrix}$: buktikan $F^n=\begin{pmatrix}F_{n+1}&F_n\\F_n&F_{n-1}\end{pmatrix}$ dan simpulkan $F_{n+1}F_{n-1}-F_n^2=(-1)^n$.
\kanalisis Induksi pada $n$; lalu ambil determinan kedua ruas.
\kbahas \emph{Basis} $n=1$: $F^1=\begin{pmatrix}F_2&F_1\\F_1&F_0\end{pmatrix}=\begin{pmatrix}1&1\\1&0\end{pmatrix}$. \emph{Induksi:}
\[
F^{n+1}=F^nF=\begin{pmatrix}F_{n+1}&F_n\\F_n&F_{n-1}\end{pmatrix}\begin{pmatrix}1&1\\1&0\end{pmatrix}=\begin{pmatrix}F_{n+1}+F_n&F_{n+1}\\F_n+F_{n-1}&F_n\end{pmatrix}=\begin{pmatrix}F_{n+2}&F_{n+1}\\F_{n+1}&F_n\end{pmatrix}.
\]
Ambil determinan: $\det(F^n)=(\det F)^n=(-1)^n$, sedangkan $\det\begin{pmatrix}F_{n+1}&F_n\\F_n&F_{n-1}\end{pmatrix}=F_{n+1}F_{n-1}-F_n^2$. Jadi $F_{n+1}F_{n-1}-F_n^2=(-1)^n$ (identitas Cassini). $\blacksquare$

\knomor{2}
\ksoal Semua $A$ ordo $2$ dengan $A^2=I$, $A\neq\pm I$.
\kanalisis $A^2=I$ berarti nilai eigen $\pm1$; selain $\pm I$, haruslah $\operatorname{tr}A=0$ dan $\det A=-1$.
\kbahas Tulis $A=\begin{pmatrix}a&b\\c&d\end{pmatrix}$. Dari Cayley--Hamilton, $A^2=(\operatorname{tr}A)A-(\det A)I$. Agar $A^2=I$ dengan $A\neq\pm I$ haruslah $\operatorname{tr}A=a+d=0$ dan $\det A=-1$, yakni $d=-a$ dan $-a^2-bc=-1$, sehingga $a^2+bc=1$. Maka semua matriks berbentuk $A=\begin{pmatrix}a&b\\c&-a\end{pmatrix}$ dengan $a^2+bc=1$ (misalnya $\begin{pmatrix}0&1\\1&0\end{pmatrix}$, $\begin{pmatrix}1&0\\0&-1\end{pmatrix}$). $\blacksquare$

% ============================================================
\chapter{Pembahasan Bab 7: Geometri Transformasi}
% ============================================================

\kbagian{Bagian A}

\knomor{A.1}
\ksoal Bayangan $A(2,5)$ oleh translasi $\begin{psmallmatrix}3\\-1\end{psmallmatrix}$.
\kanalisis Tambahkan komponen translasi.
\kbahas $A'=(2+3,\,5-1)=(5,4).$

\knomor{A.2}
\ksoal Bayangan $B(4,-2)$ oleh refleksi sumbu-$x$.
\kanalisis $(x,y)\mapsto(x,-y)$.
\kbahas $B'=(4,2).$

\knomor{A.3}
\ksoal Bayangan $C(1,3)$ oleh rotasi $90^\circ$ ccw.
\kanalisis $(x,y)\mapsto(-y,x)$.
\kbahas $C'=(-3,1).$

\knomor{A.4}
\ksoal Bayangan $D(2,-3)$ oleh dilatasi $[O,2]$.
\kanalisis Kalikan koordinat dengan $2$.
\kbahas $D'=(4,-6).$

\knomor{A.5}
\ksoal Bayangan $E(3,4)$ oleh refleksi garis $y=x$.
\kanalisis $(x,y)\mapsto(y,x)$.
\kbahas $E'=(4,3).$

\kbagian{Bagian B}

\knomor{B.1}
\ksoal $P(2,-3)$: refleksi sumbu-$x$ lalu dilatasi $[O,2]$.
\kanalisis Kerjakan berurutan.
\kbahas Refleksi: $(2,3)$. Dilatasi: $(4,6)$. Jadi $P'=(4,6).$

\knomor{B.2}
\ksoal Bayangan $Q(1,2)$ oleh rotasi $180^\circ$.
\kanalisis $(x,y)\mapsto(-x,-y)$.
\kbahas $Q'=(-1,-2).$

\knomor{B.3}
\ksoal Matriks gabungan "refleksi sumbu-$y$ lalu rotasi $90^\circ$"; terapkan pada $(3,1)$.
\kanalisis $M=M_2M_1$ dengan $M_1$ refleksi, $M_2$ rotasi.
\kbahas $M=\begin{pmatrix}0&-1\\1&0\end{pmatrix}\begin{pmatrix}-1&0\\0&1\end{pmatrix}=\begin{pmatrix}0&-1\\-1&0\end{pmatrix}$. Terapkan: $\begin{pmatrix}0&-1\\-1&0\end{pmatrix}\begin{pmatrix}3\\1\end{pmatrix}=\begin{pmatrix}-1\\-3\end{pmatrix}$, jadi bayangannya $(-1,-3).$

\knomor{B.4}
\ksoal $(5,2)$ ditranslasi $\begin{psmallmatrix}-2\\4\end{psmallmatrix}$ lalu refleksi sumbu-$x$.
\kanalisis Translasi dulu, lalu refleksi.
\kbahas Translasi: $(3,6)$. Refleksi sumbu-$x$: $(3,-6)$. Jadi bayangannya $(3,-6).$

\knomor{B.5}
\ksoal Luas bayangan segitiga (luas $6$) oleh dilatasi $[O,3]$.
\kanalisis Luas dikali $k^2$.
\kbahas $6\times3^2=54$ satuan luas.

\knomor{B.6}
\ksoal Bayangan $(2,0)$ oleh rotasi $60^\circ$.
\kanalisis Pakai matriks rotasi.
\kbahas $\begin{pmatrix}\cos60^\circ&-\sin60^\circ\\ \sin60^\circ&\cos60^\circ\end{pmatrix}\begin{pmatrix}2\\0\end{pmatrix}=\begin{pmatrix}2\cos60^\circ\\2\sin60^\circ\end{pmatrix}=\begin{pmatrix}1\\\sqrt3\end{pmatrix}.$ Jadi $(1,\sqrt3).$

\knomor{B.7}
\ksoal Bayangan garis $y=2x+1$ oleh refleksi sumbu-$x$.
\kanalisis $(x,y)\mapsto(x,-y)$; ganti $y\to-y$.
\kbahas $-y=2x+1\Rightarrow y=-2x-1.$

\kbagian{Bagian C}

\knomor{C.1}
\ksoal Tunjukkan rotasi $90^\circ$ dua kali $=$ rotasi $180^\circ$.
\kanalisis Kuadratkan matriks rotasi $90^\circ$.
\kbahas $\begin{pmatrix}0&-1\\1&0\end{pmatrix}^2=\begin{pmatrix}-1&0\\0&-1\end{pmatrix}$, yaitu matriks rotasi $180^\circ$. $\blacksquare$

\knomor{C.2}
\ksoal Buktikan rotasi $\alpha$ lalu $\beta$ $=$ rotasi $\alpha+\beta$; turunkan rumus jumlah sudut.
\kanalisis Kalikan $R(\beta)R(\alpha)$ dan samakan dengan $R(\alpha+\beta)$.
\kbahas Kalikan kedua matriks rotasi:
\[
R(\beta)R(\alpha)=\begin{pmatrix}\cos\beta&-\sin\beta\\ \sin\beta&\cos\beta\end{pmatrix}\begin{pmatrix}\cos\alpha&-\sin\alpha\\ \sin\alpha&\cos\alpha\end{pmatrix}.
\]
Entri pojok kiri-atas hasilnya $\cos\beta\cos\alpha-\sin\beta\sin\alpha$ dan entri kiri-bawah $\sin\beta\cos\alpha+\cos\beta\sin\alpha$. Bandingkan dengan kolom pertama $R(\alpha+\beta)$, yaitu $\begin{psmallmatrix}\cos(\alpha+\beta)\\ \sin(\alpha+\beta)\end{psmallmatrix}$: diperoleh
\[
\cos(\alpha+\beta)=\cos\alpha\cos\beta-\sin\alpha\sin\beta,\qquad
\sin(\alpha+\beta)=\sin\alpha\cos\beta+\cos\alpha\sin\beta. \quad\blacksquare
\]

\knomor{C.3}
\ksoal Lingkaran $x^2+y^2=4$ didilatasi $[O,3]$.
\kanalisis $(x,y)\mapsto(3x,3y)$; substitusi balik.
\kbahas Misal bayangan $(x',y')=(3x,3y)$, maka $x=\tfrac{x'}3$, $y=\tfrac{y'}3$. Substitusi: $\tfrac{x'^2}{9}+\tfrac{y'^2}{9}=4\Rightarrow x'^2+y'^2=36$. Bayangan lingkaran $x^2+y^2=36$ (jari-jari $6$). Luas menjadi $9$ kali ($=k^2$).

\knomor{C.4}
\ksoal Refleksi sumbu-$x$ lalu sumbu-$y$ menghasilkan apa?
\kanalisis Kalikan kedua matriks refleksi.
\kbahas $\begin{pmatrix}-1&0\\0&1\end{pmatrix}\begin{pmatrix}1&0\\0&-1\end{pmatrix}=\begin{pmatrix}-1&0\\0&-1\end{pmatrix}$, yaitu \textbf{rotasi $180^\circ$} (pencerminan titik terhadap $O$). $\blacksquare$

\knomor{C.5}
\ksoal Luas bayangan persegi satuan oleh $M=\begin{psmallmatrix}2&1\\0&3\end{psmallmatrix}$.
\kanalisis Faktor luas $=|\det M|$.
\kbahas $\det M=6$, jadi luas bayangan $=|6|\times1=6$ satuan luas.

\knomor{C.6}
\ksoal Bayangan $y=x^2$ oleh dilatasi $[O,2]$.
\kanalisis $(x,y)\mapsto(2x,2y)$; substitusi balik.
\kbahas $x=\tfrac{x'}2$, $y=\tfrac{y'}2$. Dari $y=x^2$: $\tfrac{y'}2=\left(\tfrac{x'}2\right)^2=\tfrac{x'^2}4\Rightarrow y'=\tfrac{x'^2}2$. Bayangan: $y=\tfrac12x^2.$

\kbagian{Bagian D}

\knomor{D.1}
\ksoal Buktikan $R(\alpha)R(\beta)=R(\alpha+\beta)$ dan $R(\alpha)^{-1}=R(-\alpha)$.
\kanalisis Gunakan hasil C.2; untuk invers ambil $\beta=-\alpha$.
\kbahas Dari C.2, $R(\alpha)R(\beta)=R(\alpha+\beta)$. Maka $R(\alpha)R(-\alpha)=R(0)=I$, sehingga $R(\alpha)^{-1}=R(-\alpha)$. $\blacksquare$

\knomor{D.2}
\ksoal Matriks refleksi terhadap garis $y=(\tan\theta)x$; tunjukkan $\det=-1$.
\kanalisis Refleksi terhadap garis bersudut $\theta$ memetakan sudut $\phi\mapsto2\theta-\phi$.
\kbahas Matriksnya $\begin{pmatrix}\cos2\theta&\sin2\theta\\ \sin2\theta&-\cos2\theta\end{pmatrix}$. Determinannya $=-\cos^2 2\theta-\sin^2 2\theta=-1$. (Determinan $-1$ menandai pembalikan orientasi --- ciri refleksi.) $\blacksquare$

\knomor{D.3}
\ksoal Buktikan rotasi/refleksi isometri; dilatasi mengalikan jarak $|k|$.
\kanalisis Tinjau panjang vektor selisih $\|M\mathbf u-M\mathbf v\|=\|M(\mathbf u-\mathbf v)\|$.
\kbahas Untuk matriks rotasi/refleksi $M$ berlaku $M^\top M=I$ (ortogonal), sehingga $\|M\mathbf w\|^2=\mathbf w^\top M^\top M\mathbf w=\mathbf w^\top\mathbf w=\|\mathbf w\|^2$; jarak terjaga. Untuk dilatasi $M=kI$, $\|M\mathbf w\|=|k|\,\|\mathbf w\|$, jadi jarak dikali $|k|$. $\blacksquare$

\knomor{D.4}
\ksoal Buktikan dua refleksi terhadap garis bersudut $\varphi$ $=$ rotasi $2\varphi$.
\kanalisis Kalikan dua matriks refleksi pada sudut $0$ dan $\varphi$.
\kbahas Refleksi terhadap garis bersudut $\theta$ bermatriks $F(\theta)=\begin{psmallmatrix}\cos2\theta&\sin2\theta\\ \sin2\theta&-\cos2\theta\end{psmallmatrix}$. Maka $F(\varphi)F(0)$ menghasilkan $\begin{psmallmatrix}\cos2\varphi&-\sin2\varphi\\ \sin2\varphi&\cos2\varphi\end{psmallmatrix}=R(2\varphi)$. Jadi komposisinya rotasi sebesar $2\varphi$. $\blacksquare$

\knomor{D.5}
\ksoal Translasi sebagai matriks $3\times3$ (koordinat homogen).
\kanalisis Tambahkan koordinat ketiga bernilai $1$.
\kbahas Tulis titik sebagai $\begin{psmallmatrix}x\\y\\1\end{psmallmatrix}$. Translasi $\begin{psmallmatrix}a\\b\end{psmallmatrix}$ menjadi
\[
\begin{pmatrix}1&0&a\\0&1&b\\0&0&1\end{pmatrix}\begin{pmatrix}x\\y\\1\end{pmatrix}=\begin{pmatrix}x+a\\y+b\\1\end{pmatrix}.
\]
Dengan trik ini, translasi pun menjadi perkalian matriks sehingga dapat digabung dengan rotasi/dilatasi dalam satu rantai.

\kbagian{Challenge Corner}

\knomor{1}
\ksoal Rotasi $40^\circ$ sebanyak $9$ kali.
\kanalisis Jumlahkan sudut.
\kbahas Total sudut $9\times40^\circ=360^\circ$, yaitu rotasi penuh $=$ identitas. Titik kembali ke posisi semula.

\knomor{2}
\ksoal Matriks "dilatasi $[O,2]$ lalu rotasi $90^\circ$"; samakah dengan urutan kebalikan?
\kanalisis Hitung $M_2M_1$ dan $M_1M_2$.
\kbahas $M_1=\begin{psmallmatrix}2&0\\0&2\end{psmallmatrix}$, $M_2=\begin{psmallmatrix}0&-1\\1&0\end{psmallmatrix}$. $M_2M_1=\begin{psmallmatrix}0&-2\\2&0\end{psmallmatrix}$, dan $M_1M_2=\begin{psmallmatrix}0&-2\\2&0\end{psmallmatrix}$ pula. Keduanya \textbf{sama}, karena dilatasi (kelipatan matriks identitas) komutatif dengan setiap matriks. $\blacksquare$

% ============================================================
\chapter{Pembahasan Bab 8: Analisis Grafik}
% ============================================================

\kbagian{Bagian A}

\knomor{A.1}
\ksoal Domain $f(x)=\sqrt{x-3}$.
\kanalisis Syarat akar $\ge0$.
\kbahas $x-3\ge0\Rightarrow x\ge3$. Domain $[3,\infty)$.

\knomor{A.2}
\ksoal Domain $g(x)=\dfrac{1}{x+2}$.
\kanalisis Penyebut $\neq0$.
\kbahas $x\neq-2$.

\knomor{A.3}
\ksoal Asimtot tegak $f(x)=\dfrac1{x-1}$.
\kanalisis Penyebut nol.
\kbahas $x=1$.

\knomor{A.4}
\ksoal Geser $y=x^2$ ke atas $3$.
\kanalisis $f(x)+k$ menggeser ke atas.
\kbahas $y=x^2+3$.

\knomor{A.5}
\ksoal Range $f(x)=x^2+1$.
\kanalisis $x^2\ge0$.
\kbahas $y\ge1$, range $[1,\infty)$.

\kbagian{Bagian B}

\knomor{B.1}
\ksoal Domain dan range $f(x)=\sqrt{4-x^2}$.
\kanalisis $4-x^2\ge0$; nilai maksimum saat $x=0$.
\kbahas Domain $[-2,2]$; nilai dari $0$ (di $x=\pm2$) sampai $2$ (di $x=0$), range $[0,2]$.

\knomor{B.2}
\ksoal Geser $2$ ke kanan, $1$ ke bawah.
\kanalisis $f(x-h)+k$ dengan $h=2$, $k=-1$.
\kbahas $y=f(x-2)-1$.

\knomor{B.3}
\ksoal Asimtot datar $f(x)=\dfrac{2x+1}{x-3}$.
\kanalisis Rasio koefisien pangkat tertinggi.
\kbahas $y=\dfrac{2}{1}=2$.

\knomor{B.4}
\ksoal Uraikan $y=-2f(x-1)$ dari $f(x)=x^2$.
\kanalisis Gabungan geser, regang, cermin.
\kbahas Geser $1$ ke kanan ($(x-1)^2$), regang vertikal $2\times$, cermin terhadap sumbu-$x$: $y=-2(x-1)^2$. Puncak $(1,0)$, terbuka ke bawah.

\knomor{B.5}
\ksoal Domain $f(x)=\dfrac{\sqrt{x-1}}{x-4}$.
\kanalisis $x-1\ge0$ dan $x-4\neq0$.
\kbahas $x\ge1$ dan $x\neq4$: domain $[1,4)\cup(4,\infty)$.

\knomor{B.6}
\ksoal $f(-2),f(0),f(3)$ untuk fungsi piecewise.
\kanalisis Pilih aturan sesuai interval.
\kbahas $f(-2)=-2+1=-1$; $f(0)=2(0)=0$; $f(3)=2(3)=6$.

\kbagian{Bagian C}

\knomor{C.1}
\ksoal Puncak dan sumbu simetri $y=2(x-3)^2+4$.
\kanalisis Bentuk puncak $a(x-h)^2+k$.
\kbahas Puncak $(3,4)$, sumbu simetri $x=3$.

\knomor{C.2}
\ksoal Transformasi $y=2^x\to2^{x-1}+3$ dan asimtotnya.
\kanalisis $x-1$ geser kanan, $+3$ geser atas.
\kbahas Geser $1$ ke kanan dan $3$ ke atas. Asimtot datar bergeser dari $y=0$ menjadi $y=3$.

\knomor{C.3}
\ksoal Domain dan asimtot $f(x)=\dfrac{x+2}{x^2-x-6}$.
\kanalisis Faktorkan penyebut.
\kbahas $x^2-x-6=(x-3)(x+2)$, jadi $f(x)=\dfrac{x+2}{(x-3)(x+2)}=\dfrac{1}{x-3}$ (dengan lubang di $x=-2$). Domain $x\neq3$ dan $x\neq-2$. Asimtot tegak $x=3$, asimtot datar $y=0$.

\knomor{C.4}
\ksoal Domain $f(x)={}^{10}\!\log(x^2-4)$.
\kanalisis Argumen logaritma $>0$.
\kbahas $x^2-4>0\Rightarrow x<-2$ atau $x>2$. Domain $(-\infty,-2)\cup(2,\infty)$.

\knomor{C.5}
\ksoal Paritas $f(x)=x^3-x$.
\kanalisis Hitung $f(-x)$.
\kbahas $f(-x)=-x^3+x=-(x^3-x)=-f(x)$, jadi \textbf{ganjil}.

\knomor{C.6}
\ksoal Domain $(f\circ g)(x)$ dengan $f=\sqrt x$, $g=x^2-9$.
\kanalisis Syarat $g(x)\ge0$.
\kbahas $x^2-9\ge0\Rightarrow x\le-3$ atau $x\ge3$. Domain $(-\infty,-3]\cup[3,\infty)$.

\kbagian{Bagian D}

\knomor{D.1}
\ksoal Buktikan genap$+$genap genap; ganjil$+$ganjil ganjil.
\kanalisis Pakai definisi pada $h=f+g$.
\kbahas Genap: $h(-x)=f(-x)+g(-x)=f(x)+g(x)=h(x)$. Ganjil: $h(-x)=-f(x)-g(x)=-h(x)$. $\blacksquare$

\knomor{D.2}
\ksoal Asimtot $f(x)=\dfrac{ax+b}{cx+d}$, $c\neq0$.
\kanalisis Penyebut nol; perilaku $x\to\infty$.
\kbahas Asimtot tegak $x=-\dfrac dc$; asimtot datar $y=\dfrac ac$ (rasio koefisien $x$).

\knomor{D.3}
\ksoal Tulis $ax^2+bx+c=a(x-h)^2+k$.
\kanalisis Melengkapkan kuadrat.
\kbahas $ax^2+bx+c=a\left(x+\dfrac{b}{2a}\right)^2+c-\dfrac{b^2}{4a}$, jadi $h=-\dfrac{b}{2a}$, $k=\dfrac{4ac-b^2}{4a}$.

\knomor{D.4}
\ksoal Buktikan $f$ = (bagian genap) $+$ (bagian ganjil).
\kanalisis Definisikan $g(x)=\tfrac12[f(x)+f(-x)]$, $h(x)=\tfrac12[f(x)-f(-x)]$.
\kbahas $g(-x)=\tfrac12[f(-x)+f(x)]=g(x)$ (genap); $h(-x)=\tfrac12[f(-x)-f(x)]=-h(x)$ (ganjil); dan $g(x)+h(x)=f(x)$. $\blacksquare$

\kbagian{Challenge Corner}

\knomor{1}
\ksoal Kuadrat berpuncak $(2,5)$ melalui $(0,1)$.
\kanalisis Bentuk puncak $a(x-2)^2+5$.
\kbahas $f(0)=4a+5=1\Rightarrow a=-1$. Jadi $f(x)=-(x-2)^2+5.$

\knomor{2}
\ksoal Sketsa $y=|x^2-4|$.
\kanalisis Mulai dari $y=x^2-4$, lalu cerminkan bagian negatif.
\kbahas $x^2-4<0$ pada $-2<x<2$; di sana grafik $y=x^2-4$ (yang berada di bawah sumbu-$x$) dicerminkan ke atas. Hasilnya berbentuk "W-melengkung" dengan akar di $x=\pm2$ dan puncak lokal $(0,4)$; di luar $[-2,2]$ grafik sama dengan $x^2-4$.

% ============================================================
\chapter{Pembahasan Bab 9: Fungsi Rasional}
% ============================================================

\kbagian{Bagian A}

\knomor{A.1}
\ksoal Tentukan domain $f(x)=\dfrac{1}{x-3}$.
\kanalisis Domain rasional = semua real kecuali pembuat nol penyebut.
\kbahas $x-3=0\Rightarrow x=3$. Domain $=\{x\mid x\neq 3\}$.

\knomor{A.2}
\ksoal Tentukan domain $f(x)=\dfrac{x+1}{x^2-4}$.
\kanalisis Cari pembuat nol penyebut $x^2-4$.
\kbahas $x^2-4=0\Rightarrow x=\pm2$. Domain $=\{x\mid x\neq -2,\ x\neq 2\}$.

\knomor{A.3}
\ksoal Tentukan asimtot tegak $f(x)=\dfrac{5}{x+2}$.
\kanalisis Asimtot tegak di pembuat nol penyebut (pembilang tak nol di situ).
\kbahas $x+2=0\Rightarrow$ asimtot tegak $x=-2$.

\knomor{A.4}
\ksoal Tentukan asimtot datar $f(x)=\dfrac{3x+1}{x-4}$.
\kanalisis Derajat pembilang $=$ derajat penyebut, jadi rasio koefisien utama.
\kbahas $y=\dfrac{3}{1}=3$.

\knomor{A.5}
\ksoal Tentukan asimtot datar $f(x)=\dfrac{2x+1}{x^2-1}$.
\kanalisis Derajat pembilang $(1)<$ derajat penyebut $(2)$.
\kbahas Asimtot datar $y=0$.

\knomor{A.6}
\ksoal Selesaikan $\dfrac{2}{x}=4$.
\kanalisis Kalikan kedua ruas dengan $x$ ($x\neq0$).
\kbahas $2=4x\Rightarrow x=\dfrac12$.

\knomor{A.7}
\ksoal Selesaikan $\dfrac{x+1}{x-2}=0$.
\kanalisis Pecahan nol jika pembilang nol dan penyebut tak nol.
\kbahas $x+1=0\Rightarrow x=-1$ (memenuhi $x\neq2$). HP $=\{-1\}$.

\knomor{A.8}
\ksoal Tentukan semua asimtot tegak $f(x)=\dfrac{x+1}{x^2-x}$.
\kanalisis Faktorkan penyebut; pastikan pembilang tak ikut nol.
\kbahas $x^2-x=x(x-1)=0\Rightarrow x=0$ atau $x=1$. Pembilang $x+1$ tak nol di situ, jadi asimtot tegak $x=0$ dan $x=1$.

\knomor{A.9}
\ksoal Hitung $f(2)$ untuk $f(x)=\dfrac{x^2-1}{x+1}$.
\kanalisis Sederhanakan dulu: $\dfrac{(x-1)(x+1)}{x+1}=x-1$ untuk $x\neq-1$.
\kbahas $f(2)=2-1=1$.

\knomor{A.10}
\ksoal Tentukan titik potong sumbu-$x$ dari $f(x)=\dfrac{2x-6}{x+1}$.
\kanalisis Titik potong sumbu-$x$ saat pembilang nol.
\kbahas $2x-6=0\Rightarrow x=3$. Titiknya $(3,0)$.

\kbagian{Bagian B}

\knomor{B.1}
\ksoal Tentukan asimtot tegak dan datar $f(x)=\dfrac{2x^2+3}{x^2-1}$.
\kanalisis Penyebut nol $\to$ asimtot tegak; derajat sama $\to$ rasio koefisien utama.
\kbahas $x^2-1=0\Rightarrow x=\pm1$ (asimtot tegak). Derajat sama $(2,2)$: asimtot datar $y=\dfrac{2}{1}=2$.

\knomor{B.2}
\ksoal Tentukan asimtot miring $f(x)=\dfrac{x^2+1}{x-1}$.
\kanalisis $\deg P=\deg Q+1$, lakukan pembagian polinomial.
\kbahas $x^2+1=(x-1)(x+1)+2$, jadi $f(x)=x+1+\dfrac{2}{x-1}$. Asimtot miring $y=x+1$.

\knomor{B.3}
\ksoal Selesaikan $\dfrac{1}{x-1}+\dfrac{1}{x+1}=1$.
\kanalisis Samakan penyebut $(x-1)(x+1)=x^2-1$, lalu selesaikan kuadrat; cek $x\neq\pm1$.
\kbahas $\dfrac{2x}{x^2-1}=1\Rightarrow 2x=x^2-1\Rightarrow x^2-2x-1=0\Rightarrow x=1\pm\sqrt2$. Keduanya sah. HP $=\{1-\sqrt2,\ 1+\sqrt2\}$.

\knomor{B.4}
\ksoal Selesaikan $\dfrac{x-2}{x+1}\ge 0$.
\kanalisis Pembuat nol: $x=2$ (tutup), $x=-1$ (buka). Uji tanda garis bilangan.
\kbahas Tanda: $x<-1\Rightarrow(-)/(-)=+$; $-1<x<2\Rightarrow(-)/(+)=-$; $x>2\Rightarrow+$. Maka HP $=\{x\mid x<-1\ \text{atau}\ x\ge2\}$.

\knomor{B.5}
\ksoal Selesaikan $\dfrac{x+3}{x-2}<1$.
\kanalisis Pindahkan ke satu ruas, jangan kali silang.
\kbahas $\dfrac{x+3}{x-2}-1=\dfrac{5}{x-2}<0$. Karena $5>0$, perlu $x-2<0\Rightarrow x<2$. HP $=\{x\mid x<2\}$.

\kbagian{Bagian C}

\knomor{C.1}
\ksoal Sederhanakan dan jelaskan grafik $f(x)=\dfrac{x^2-4}{x-2}$.
\kanalisis Faktor sekutu $\Rightarrow$ lubang, bukan asimtot tegak.
\kbahas $\dfrac{(x-2)(x+2)}{x-2}=x+2$ untuk $x\neq2$. Grafiknya garis $y=x+2$ dengan \textbf{lubang} di $(2,4)$; tidak ada asimtot tegak.

\knomor{C.2}
\ksoal Hitung $\displaystyle\lim_{x\to\infty}\dfrac{3x^2-x+1}{2x^2+5}$.
\kanalisis Bagi pembilang dan penyebut dengan $x^2$ (derajat tertinggi).
\kbahas $\lim\dfrac{3-\frac1x+\frac1{x^2}}{2+\frac5{x^2}}=\dfrac{3}{2}$. Nilai ini sekaligus asimtot datar $y=\tfrac32$.

\knomor{C.3}
\ksoal Tentukan $(f\circ f)(x)$ dan domainnya untuk $f(x)=\dfrac1x$.
\kanalisis Substitusikan $f$ ke dirinya; perhatikan syarat domain.
\kbahas $(f\circ f)(x)=f\!\left(\tfrac1x\right)=\dfrac{1}{1/x}=x$. Walau hasil akhir $x$, syarat $x\neq0$ tetap berlaku. Domain $=\{x\mid x\neq0\}$.

\knomor{C.4}
\ksoal Selesaikan $\dfrac{x^2-1}{x^2-4}\le 0$.
\kanalisis Faktorkan: $\dfrac{(x-1)(x+1)}{(x-2)(x+2)}$. Nol di $x=\pm1$ (tutup), tak terdefinisi di $x=\pm2$ (buka).
\kbahas Uji tanda: $(-\infty,-2):+$; $(-2,-1):-$; $(-1,1):+$; $(1,2):-$; $(2,\infty):+$. Bernilai $\le0$ pada $(-2,-1]\cup[1,2)$. HP $=\{x\mid -2<x\le-1\ \text{atau}\ 1\le x<2\}$.

\knomor{C.5}
\ksoal Untuk $x>0$, tentukan minimum $f(x)=x+\dfrac1x$.
\kanalisis Gunakan AM--GM pada $x$ dan $\tfrac1x$.
\kbahas $x+\dfrac1x\ge2\sqrt{x\cdot\tfrac1x}=2$, kesamaan saat $x=\tfrac1x\Rightarrow x=1$. Minimum $=2$ di $x=1$.

\kbagian{Bagian D}

\knomor{D.1}
\ksoal Tuliskan $\dfrac{3x+5}{(x+1)(x+2)}=\dfrac{A}{x+1}+\dfrac{B}{x+2}$.
\kanalisis Kalikan silang lalu substitusi nilai $x$ yang melenyapkan salah satu faktor.
\kbahas $3x+5=A(x+2)+B(x+1)$. $x=-1:\ 2=A\Rightarrow A=2$. $x=-2:\ -1=-B\Rightarrow B=1$. Jadi $\dfrac{2}{x+1}+\dfrac{1}{x+2}$.

\knomor{D.2}
\ksoal Buktikan asimtot miring $f(x)=\dfrac{ax^2+bx+c}{x+d}$ adalah $y=ax+(b-ad)$.
\kanalisis Bagi pembilang dengan $(x+d)$.
\kbahas $ax^2+bx+c=(x+d)\big(ax+(b-ad)\big)+\big(c-d(b-ad)\big)$, sehingga $f(x)=ax+(b-ad)+\dfrac{c-d(b-ad)}{x+d}$. Saat $x\to\pm\infty$ suku terakhir $\to0$, jadi asimtot miring $y=ax+(b-ad)$. $\blacksquare$

\knomor{D.3}
\ksoal Buktikan turunan $f(x)=\dfrac1x$ adalah $-\dfrac{1}{x^2}$ (definisi).
\kanalisis Gunakan $f'(x)=\lim_{h\to0}\dfrac{f(x+h)-f(x)}{h}$.
\kbahas $\dfrac{\frac{1}{x+h}-\frac1x}{h}=\dfrac{\frac{x-(x+h)}{x(x+h)}}{h}=\dfrac{-h}{h\,x(x+h)}=\dfrac{-1}{x(x+h)}$. Saat $h\to0$ diperoleh $-\dfrac{1}{x^2}$. $\blacksquare$

\knomor{D.4}
\ksoal Buktikan $x+\dfrac1x\ge2$ untuk $x>0$, dan kapan kesamaan.
\kanalisis Mulai dari kuadrat yang selalu tak negatif.
\kbahas $\left(\sqrt{x}-\dfrac{1}{\sqrt{x}}\right)^2\ge0\Rightarrow x-2+\dfrac1x\ge0\Rightarrow x+\dfrac1x\ge2$. Kesamaan saat $\sqrt{x}=\tfrac{1}{\sqrt{x}}$, yaitu $x=1$. $\blacksquare$

\knomor{D.5}
\ksoal Tentukan $f^{-1}(x)$ untuk $f(x)=\dfrac{x+1}{x-1}$ dan tunjukkan $f=f^{-1}$.
\kanalisis Tukar $x\leftrightarrow y$ lalu selesaikan untuk $y$.
\kbahas $y=\dfrac{x+1}{x-1}\Rightarrow y(x-1)=x+1\Rightarrow x(y-1)=y+1\Rightarrow x=\dfrac{y+1}{y-1}$. Maka $f^{-1}(x)=\dfrac{x+1}{x-1}=f(x)$, jadi $f$ adalah inversnya sendiri (involusi). $\blacksquare$

\kbagian{Challenge Corner}

\knomor{1}
\ksoal Selesaikan $\dfrac{x^2-3x+2}{x^2-5x+6}>0$.
\kanalisis Faktorkan: $\dfrac{(x-1)(x-2)}{(x-2)(x-3)}$. Faktor $(x-2)$ memberi lubang ($x\neq2$).
\kbahas Untuk $x\neq2$, ekspresi menjadi $\dfrac{x-1}{x-3}>0\Rightarrow x<1$ atau $x>3$. Nilai $x=2$ tidak termasuk selang itu, jadi HP $=\{x\mid x<1\ \text{atau}\ x>3\}$.

\knomor{2}
\ksoal Tentukan maksimum, minimum, dan range $f(x)=\dfrac{x}{x^2+1}$.
\kanalisis Pakai turunan, atau syarat diskriminan dari $y(x^2+1)=x$.
\kbahas $yx^2-x+y=0$ punya solusi real bila $1-4y^2\ge0\Rightarrow -\tfrac12\le y\le\tfrac12$. Jadi range $\left[-\tfrac12,\tfrac12\right]$; maksimum $\tfrac12$ di $x=1$ dan minimum $-\tfrac12$ di $x=-1$.

\knomor{3}
\ksoal Buktikan $f(x)=\dfrac{x^2+x+1}{x^2-x+1}$ hanya bernilai pada $\left[\tfrac13,3\right]$.
\kanalisis Penyebut selalu positif (diskriminan $<0$). Susun sebagai kuadrat dalam $x$ dan syaratkan diskriminan $\ge0$.
\kbahas $y(x^2-x+1)=x^2+x+1\Rightarrow(y-1)x^2-(y+1)x+(y-1)=0$. Untuk $y=1$, $x=0$ (tercapai). Untuk $y\neq1$, syarat real: $(y+1)^2-4(y-1)^2\ge0$. Faktorkan $\big[(y+1)-2(y-1)\big]\big[(y+1)+2(y-1)\big]=(3-y)(3y-1)\ge0\Rightarrow\tfrac13\le y\le3$. Jadi range $=\left[\tfrac13,3\right]$. $\blacksquare$

% ============================================================
\chapter{Pembahasan Bab 10: Nilai Mutlak dan Bentuk Irasional}
% ============================================================

\kbagian{Bagian A}

\knomor{A.1}
\ksoal Hitung $|-7|$.
\kanalisis Nilai mutlak = jarak dari $0$, selalu tak negatif.
\kbahas $|-7|=7$.

\knomor{A.2}
\ksoal Selesaikan $|x|=5$.
\kanalisis $|x|=a\iff x=\pm a$.
\kbahas $x=5$ atau $x=-5$. HP $=\{-5,5\}$.

\knomor{A.3}
\ksoal Selesaikan $|x-3|=2$.
\kanalisis $|x-3|=2\iff x-3=\pm2$.
\kbahas $x=5$ atau $x=1$. HP $=\{1,5\}$.

\knomor{A.4}
\ksoal Selesaikan $|2x|=8$.
\kanalisis $|2x|=2|x|=8\Rightarrow|x|=4$.
\kbahas $x=\pm4$. HP $=\{-4,4\}$.

\knomor{A.5}
\ksoal Sederhanakan $\sqrt{50}$.
\kanalisis Keluarkan faktor kuadrat sempurna.
\kbahas $\sqrt{50}=\sqrt{25\cdot2}=5\sqrt2$.

\knomor{A.6}
\ksoal Sederhanakan $\sqrt{12}+\sqrt{27}$.
\kanalisis Sederhanakan tiap akar menjadi sejenis.
\kbahas $2\sqrt3+3\sqrt3=5\sqrt3$.

\knomor{A.7}
\ksoal Rasionalkan $\dfrac{1}{\sqrt3}$.
\kanalisis Kalikan dengan $\dfrac{\sqrt3}{\sqrt3}$.
\kbahas $\dfrac{1}{\sqrt3}=\dfrac{\sqrt3}{3}$.

\knomor{A.8}
\ksoal Rasionalkan $\dfrac{2}{\sqrt5}$.
\kanalisis Kalikan dengan $\dfrac{\sqrt5}{\sqrt5}$.
\kbahas $\dfrac{2}{\sqrt5}=\dfrac{2\sqrt5}{5}$.

\knomor{A.9}
\ksoal Selesaikan $\sqrt{x}=4$.
\kanalisis Kuadratkan; syarat $x\ge0$ otomatis.
\kbahas $x=16$.

\knomor{A.10}
\ksoal Hitung $|3-8|+|-2|$.
\kanalisis Hitung tiap nilai mutlak lalu jumlahkan.
\kbahas $|-5|+|-2|=5+2=7$.

\kbagian{Bagian B}

\knomor{B.1}
\ksoal Selesaikan $|x-2|=|2x+1|$.
\kanalisis $|f|=|g|\iff f=g$ atau $f=-g$.
\kbahas $x-2=2x+1\Rightarrow x=-3$; atau $x-2=-(2x+1)\Rightarrow 3x=1\Rightarrow x=\tfrac13$. HP $=\left\{-3,\tfrac13\right\}$.

\knomor{B.2}
\ksoal Selesaikan $|x-1|<3$.
\kanalisis $|u|<a\iff -a<u<a$.
\kbahas $-3<x-1<3\Rightarrow -2<x<4$. HP $=\{x\mid -2<x<4\}$.

\knomor{B.3}
\ksoal Selesaikan $|2x+1|\ge5$.
\kanalisis $|u|\ge a\iff u\ge a$ atau $u\le-a$.
\kbahas $2x+1\ge5\Rightarrow x\ge2$; atau $2x+1\le-5\Rightarrow x\le-3$. HP $=\{x\mid x\le-3\ \text{atau}\ x\ge2\}$.

\knomor{B.4}
\ksoal Rasionalkan $\dfrac{3}{\sqrt5-\sqrt2}$.
\kanalisis Kalikan dengan sekawan $\sqrt5+\sqrt2$.
\kbahas $\dfrac{3(\sqrt5+\sqrt2)}{5-2}=\dfrac{3(\sqrt5+\sqrt2)}{3}=\sqrt5+\sqrt2$.

\knomor{B.5}
\ksoal Selesaikan $\sqrt{2x+3}=x$.
\kanalisis Syarat $x\ge0$; kuadratkan lalu cek akar palsu.
\kbahas $2x+3=x^2\Rightarrow x^2-2x-3=0\Rightarrow(x-3)(x+1)=0$. $x=-1$ ditolak ($x\ge0$); $x=3$ memenuhi ($\sqrt9=3$). HP $=\{3\}$.

\kbagian{Bagian C}

\knomor{C.1}
\ksoal Peroleh grafik $y=|x-2|+1$ dari $y=|x|$; tentukan verteks.
\kanalisis $|x-h|+k$ menggeser grafik $h$ kanan, $k$ atas.
\kbahas Geser $y=|x|$ sejauh $2$ ke kanan lalu $1$ ke atas; verteks $(2,1)$, kurva "V" terbuka ke atas.

\knomor{C.2}
\ksoal Selesaikan $|x^2-4|=3$.
\kanalisis $x^2-4=\pm3$.
\kbahas $x^2-4=3\Rightarrow x^2=7\Rightarrow x=\pm\sqrt7$; $x^2-4=-3\Rightarrow x^2=1\Rightarrow x=\pm1$. HP $=\{-\sqrt7,-1,1,\sqrt7\}$.

\knomor{C.3}
\ksoal Selesaikan $\sqrt{x-1}<2$.
\kanalisis Syarat $x-1\ge0$; kuadratkan ($a=2>0$).
\kbahas $0\le x-1<4\Rightarrow 1\le x<5$. HP $=\{x\mid 1\le x<5\}$.

\knomor{C.4}
\ksoal Selesaikan $|x|<|x-2|$.
\kanalisis Kuadratkan (kedua ruas tak negatif).
\kbahas $x^2<x^2-4x+4\Rightarrow 4x<4\Rightarrow x<1$. HP $=\{x\mid x<1\}$. (Secara geometris: jarak ke $0$ lebih kecil daripada jarak ke $2$.)

\knomor{C.5}
\ksoal Sederhanakan $\sqrt{7+2\sqrt{10}}$.
\kanalisis Cari $a,b$ dengan $a+b=7$, $ab=10$ agar $7+2\sqrt{10}=(\sqrt a+\sqrt b)^2$.
\kbahas $a=5,b=2$: $(\sqrt5+\sqrt2)^2=7+2\sqrt{10}$. Maka $\sqrt{7+2\sqrt{10}}=\sqrt5+\sqrt2$.

\kbagian{Bagian D}

\knomor{D.1}
\ksoal Buktikan $|a+b|\le|a|+|b|$.
\kanalisis Pakai $-|x|\le x\le|x|$ lalu jumlahkan.
\kbahas Dari $-|a|\le a\le|a|$ dan $-|b|\le b\le|b|$, jumlahkan: $-(|a|+|b|)\le a+b\le|a|+|b|$. Ini setara $|a+b|\le|a|+|b|$. $\blacksquare$

\knomor{D.2}
\ksoal Hitung $\displaystyle\lim_{x\to0}\dfrac{\sqrt{x+1}-1}{x}$.
\kanalisis Kalikan sekawan $\sqrt{x+1}+1$.
\kbahas $\dfrac{(\sqrt{x+1}-1)(\sqrt{x+1}+1)}{x(\sqrt{x+1}+1)}=\dfrac{x}{x(\sqrt{x+1}+1)}=\dfrac{1}{\sqrt{x+1}+1}\xrightarrow[x\to0]{}\dfrac12$.

\knomor{D.3}
\ksoal Tunjukkan $f(x)=|x|$ tak terdiferensialkan di $x=0$.
\kanalisis Bandingkan limit turunan dari kiri dan kanan.
\kbahas Untuk $x>0$, $f(x)=x$ sehingga laju $=+1$; untuk $x<0$, $f(x)=-x$ sehingga laju $=-1$. Limit hasil bagi selisih dari kanan $=+1$ dan dari kiri $=-1$, tidak sama, jadi turunan di $0$ tidak ada. $\blacksquare$

\knomor{D.4}
\ksoal Buktikan $\sqrt2$ irasional.
\kanalisis Pengandaian (kontradiksi): andaikan $\sqrt2=\tfrac pq$ pecahan paling sederhana.
\kbahas Andaikan $\sqrt2=\dfrac pq$ dengan $\gcd(p,q)=1$. Maka $p^2=2q^2$, jadi $p^2$ genap $\Rightarrow p$ genap, tulis $p=2k$. Maka $4k^2=2q^2\Rightarrow q^2=2k^2$, jadi $q$ juga genap. Maka $p,q$ sama-sama genap --- bertentangan dengan $\gcd(p,q)=1$. Jadi $\sqrt2$ irasional. $\blacksquare$

\knomor{D.5}
\ksoal Selesaikan sistem $|x+y|=6$ dan $|x-y|=2$.
\kanalisis $x+y=\pm6$ dan $x-y=\pm2$; padukan keempat kombinasi.
\kbahas Dari $x+y=\pm6$, $x-y=\pm2$ diperoleh $(x,y)\in\{(4,2),(2,4),(-2,-4),(-4,-2)\}$.

\kbagian{Challenge Corner}

\knomor{1}
\ksoal Selesaikan $|x-1|+|x-3|=4$.
\kanalisis Tinjau tiga selang: $x<1$, $1\le x\le3$, $x>3$.
\kbahas $x<1$: $(1-x)+(3-x)=4-2x=4\Rightarrow x=0$ (memenuhi). $1\le x\le3$: jumlah $=2\neq4$ (tak ada). $x>3$: $(x-1)+(x-3)=2x-4=4\Rightarrow x=4$ (memenuhi). HP $=\{0,4\}$.

\knomor{2}
\ksoal Selesaikan $\sqrt{x+4}-\sqrt{x-1}=1$.
\kanalisis Misal $a=\sqrt{x+4},\ b=\sqrt{x-1}$; pakai $a^2-b^2=(a-b)(a+b)$.
\kbahas $a-b=1$ dan $a^2-b^2=(x+4)-(x-1)=5$, jadi $a+b=5$. Maka $a=3,\ b=2$, sehingga $x+4=9\Rightarrow x=5$ (cek: $\sqrt9-\sqrt4=3-2=1$ \checkmark). HP $=\{5\}$.

\knomor{3}
\ksoal Minimum $f(x)=|x-1|+|x-2|+|x-3|$.
\kanalisis Jumlah jarak ke tiga titik diminimumkan di \emph{median}.
\kbahas Median dari $1,2,3$ adalah $2$. $f(2)=1+0+1=2$. Minimum $=2$ tercapai di $x=2$.

% ============================================================
\chapter{Pembahasan Bab 11: Pemodelan Matematika}
% ============================================================

\kbagian{Bagian A}

\knomor{A.1}
\ksoal $C(x)=5000+2000x$, hitung $C(10)$.
\kanalisis Substitusi $x=10$.
\kbahas $C(10)=5000+20000=25000$ rupiah.

\knomor{A.2}
\ksoal Model bakteri menggandakan tiap jam, awal $100$.
\kanalisis Rasio $2$ per jam.
\kbahas $P(t)=100\cdot2^{t}$.

\knomor{A.3}
\ksoal Mobil menyusut $10\%$/tahun dari Rp$200$jt.
\kanalisis Peluruhan $A(1-p)^t$.
\kbahas $V(t)=200(0{,}9)^t$ juta; $V(1)=180$ juta.

\knomor{A.4}
\ksoal $h(t)=20t-5t^2$, hitung $h(2)$.
\kanalisis Substitusi.
\kbahas $h(2)=40-20=20$.

\knomor{A.5}
\ksoal Suhu awal $T(t)=25+60(0{,}95)^t$.
\kanalisis $t=0\Rightarrow(0{,}95)^0=1$.
\kbahas $T(0)=25+60=85$ derajat.

\kbagian{Bagian B}

\knomor{B.1}
\ksoal $h(t)=20t-5t^2$: waktu mencapai tanah dan tinggi maksimum.
\kanalisis Akar $h=0$; puncak parabola.
\kbahas $5t(4-t)=0\Rightarrow t=0$ atau $t=4$ (tiba di tanah saat $t=4$ s). Puncak di $t=2$: $h(2)=20$ (tinggi maksimum).

\knomor{B.2}
\ksoal Bunga majemuk $5\%$, awal Rp$1$jt, setelah $3$ tahun.
\kanalisis $A(1+p)^t$.
\kbahas $M(t)=1\cdot(1{,}05)^t$ juta; $M(3)=(1{,}05)^3=1{,}157625$ juta $\approx$ Rp$1.157.625$.

\knomor{B.3}
\ksoal $P(t)=100\cdot2^t=1600$.
\kanalisis Samakan basis $2$.
\kbahas $2^t=16=2^4\Rightarrow t=4$ jam.

\knomor{B.4}
\ksoal Waktu paruh $5$ tahun; fraksi setelah $15$ tahun.
\kanalisis $N(t)=N_0(\tfrac12)^{t/5}$.
\kbahas Model $N(t)=N_0\left(\tfrac12\right)^{t/5}$. Setelah $15$ tahun: $\left(\tfrac12\right)^{3}=\tfrac18$ tersisa.

\knomor{B.5}
\ksoal Model linear melalui $(0,3)$ dan $(4,11)$.
\kanalisis Gradien $=\tfrac{\Delta y}{\Delta x}$, perpotongan dari $(0,3)$.
\kbahas $b=\dfrac{11-3}{4}=2$, $a=3$. Jadi $y=2x+3$.

\knomor{B.6}
\ksoal $R(x)=x(100-x)$, harga pemaksimal.
\kanalisis Parabola terbuka ke bawah, puncak di tengah akar.
\kbahas $R(x)=-x^2+100x$, puncak $x=50$. Pendapatan maksimum $R(50)=2500$.

\kbagian{Bagian C}

\knomor{C.1}
\ksoal Populasi $3\%$/tahun, $500.000\to1.000.000$.
\kanalisis $1{,}03^t=2$; pakai logaritma.
\kbahas $500000(1{,}03)^t=1000000\Rightarrow1{,}03^t=2\Rightarrow t=\dfrac{\log 2}{\log 1{,}03}\approx23{,}4$ tahun.

\knomor{C.2}
\ksoal Tali $40$ m menjadi persegi panjang berluas maksimum.
\kanalisis Keliling $2(p+l)=40\Rightarrow p+l=20$; luas maksimum saat $p=l$.
\kbahas Luas $L=p(20-p)=-p^2+20p$ maksimum di $p=10$, jadi persegi $10\times10$ dengan luas $100$ m$^2$.

\knomor{C.3}
\ksoal $M={}^{10}\!\log(I/I_0)$, $I=1000I_0$.
\kanalisis $\log 1000=3$.
\kbahas $M={}^{10}\!\log(1000)=3$.

\knomor{C.4}
\ksoal $P(t)=\dfrac{1000}{1+9e^{-t}}$: populasi awal dan batas.
\kanalisis $t=0$ dan $t\to\infty$.
\kbahas $P(0)=\dfrac{1000}{1+9}=100$. Saat $t\to\infty$, $e^{-t}\to0$, sehingga $P\to1000$ (kapasitas tampung).

\knomor{C.5}
\ksoal Data lurus pada skala $\log$: model dan penaksiran parameter.
\kanalisis $\log y$ linear terhadap $t$ menandakan eksponensial.
\kbahas Model \textbf{eksponensial} $y=A r^t$. Karena $\log y=\log A+(\log r)t$, lakukan regresi linear pada $(t,\log y)$: gradiennya menaksir $\log r$ dan perpotongannya menaksir $\log A$.

\knomor{C.6}
\ksoal $P(x)=-2x^2+80x-200$: $x$ pemaksimal dan keuntungan maksimum.
\kanalisis Puncak parabola $x=-\tfrac{b}{2a}$.
\kbahas $x=-\dfrac{80}{2(-2)}=20$; $P(20)=-800+1600-200=600$.

\kbagian{Bagian D}

\knomor{D.1}
\ksoal Tunjukkan $\dfrac{dP}{dt}=kP$ untuk $P=P_0e^{kt}$.
\kanalisis Turunkan terhadap $t$.
\kbahas $\dfrac{dP}{dt}=P_0\cdot k\,e^{kt}=k\big(P_0e^{kt}\big)=kP.$ Laju sebanding dengan populasi. $\blacksquare$

\knomor{D.2}
\ksoal Turunkan $T=\dfrac{\ln2}{\lambda}$ dari $N=N_0e^{-\lambda t}$.
\kanalisis Saat $N=\tfrac12N_0$.
\kbahas $N_0e^{-\lambda T}=\tfrac12N_0\Rightarrow e^{-\lambda T}=\tfrac12\Rightarrow-\lambda T=-\ln2\Rightarrow T=\dfrac{\ln2}{\lambda}.$ $\blacksquare$

\knomor{D.3}
\ksoal Mengapa eksponensial gagal; bagaimana logistik memperbaiki.
\kanalisis Tinjau sumber daya terbatas.
\kbahas Model eksponensial tumbuh tanpa batas, padahal sumber daya (makanan, ruang) terbatas. Model logistik menambahkan \emph{kapasitas tampung} $K$: laju $\tfrac{dP}{dt}=rP\left(1-\tfrac{P}{K}\right)$ melambat saat $P$ mendekati $K$, sehingga kurva mendatar --- lebih realistis.

\knomor{D.4}
\ksoal Model penyebaran penyakit menggandakan tiap $3$ hari.
\kanalisis Eksponensial berbasis penggandaan.
\kbahas Jika kasus awal $C_0$, maka $C(t)=C_0\cdot2^{t/3}$ ($t$ dalam hari). \emph{Keterbatasan:} mengasumsikan laju tetap dan populasi tak terbatas; mengabaikan kekebalan, intervensi, dan kejenuhan. Untuk jangka panjang lebih cocok model logistik atau SIR.

\kbagian{Challenge Corner}

\knomor{1}
\ksoal Kapan $B(t)=100(1{,}2)^t$ melampaui $A(t)=100+50t$?
\kanalisis Bandingkan nilai untuk beberapa $t$.
\kbahas $t=9$: $B\approx515{,}98$, $A=550$ ($B<A$). $t=10$: $B\approx619{,}17$, $A=600$ ($B>A$). Jadi $B$ mulai melampaui $A$ pada sekitar \textbf{bulan ke-10}.

\knomor{2}
\ksoal Tunjukkan pertumbuhan logistik tercepat saat $P=\tfrac K2$.
\kanalisis Laju logistik $\tfrac{dP}{dt}=rP\left(1-\tfrac PK\right)$ sebagai fungsi $P$.
\kbahas Laju $g(P)=rP\left(1-\tfrac PK\right)=rP-\tfrac{r}{K}P^2$ adalah parabola terbuka ke bawah dalam $P$, maksimum di $P=-\dfrac{r}{2(-r/K)}=\dfrac K2$. Jadi pertumbuhan tercepat (titik belok kurva $P(t)$) terjadi saat $P=\dfrac K2$. $\blacksquare$

% ============================================================
\chapter{Pembahasan Bab 12: Trigonometri Lanjutan}
% ============================================================

\kbagian{Bagian A}

\knomor{A.1}
\ksoal $\sin30^\circ+\cos60^\circ$.
\kanalisis Nilai sudut istimewa.
\kbahas $\tfrac12+\tfrac12=1.$

\knomor{A.2}
\ksoal $\sin\theta=\tfrac35$, $\theta$ lancip; cari $\cos\theta$.
\kanalisis $\cos\theta=\sqrt{1-\sin^2\theta}$.
\kbahas $\cos\theta=\sqrt{1-\tfrac{9}{25}}=\tfrac45.$

\knomor{A.3}
\ksoal $\sin2\theta$ jika $\sin\theta=\tfrac35,\cos\theta=\tfrac45$.
\kanalisis $\sin2\theta=2\sin\theta\cos\theta$.
\kbahas $2\cdot\tfrac35\cdot\tfrac45=\tfrac{24}{25}.$

\knomor{A.4}
\ksoal $\sin x=\tfrac12$, $0^\circ\le x<360^\circ$.
\kanalisis $x=\alpha$ atau $180^\circ-\alpha$.
\kbahas $x=30^\circ$ atau $x=150^\circ$.

\knomor{A.5}
\ksoal Amplitudo dan periode $y=3\sin(2x)$.
\kanalisis Amplitudo $=|A|$, periode $=\tfrac{360^\circ}{|\omega|}$.
\kbahas Amplitudo $3$, periode $\tfrac{360^\circ}{2}=180^\circ$.

\kbagian{Bagian B}

\knomor{B.1}
\ksoal Buktikan $\dfrac{1-\cos^2\theta}{\sin\theta}=\sin\theta$.
\kanalisis $1-\cos^2\theta=\sin^2\theta$.
\kbahas $\dfrac{\sin^2\theta}{\sin\theta}=\sin\theta.$ $\blacksquare$

\knomor{B.2}
\ksoal $\cos75^\circ$.
\kanalisis $75^\circ=45^\circ+30^\circ$.
\kbahas $\cos75^\circ=\cos45^\circ\cos30^\circ-\sin45^\circ\sin30^\circ=\tfrac{\sqrt2}{2}\cdot\tfrac{\sqrt3}{2}-\tfrac{\sqrt2}{2}\cdot\tfrac12=\dfrac{\sqrt6-\sqrt2}{4}.$

\knomor{B.3}
\ksoal $2\cos x+1=0$, $0^\circ\le x<360^\circ$.
\kanalisis $\cos x=-\tfrac12$.
\kbahas $x=120^\circ$ atau $x=240^\circ$.

\knomor{B.4}
\ksoal Aturan sinus: $a=10,A=30^\circ,B=45^\circ$, cari $b$.
\kanalisis $\dfrac{a}{\sin A}=\dfrac{b}{\sin B}$.
\kbahas $\dfrac{10}{\sin30^\circ}=\dfrac{b}{\sin45^\circ}\Rightarrow20=\dfrac{b}{\frac{\sqrt2}{2}}\Rightarrow b=10\sqrt2.$

\knomor{B.5}
\ksoal Buktikan $\cos2\theta=1-2\sin^2\theta$.
\kanalisis Mulai dari $\cos2\theta=\cos^2\theta-\sin^2\theta$ dan $\cos^2\theta=1-\sin^2\theta$.
\kbahas $\cos2\theta=(1-\sin^2\theta)-\sin^2\theta=1-2\sin^2\theta.$ $\blacksquare$

\knomor{B.6}
\ksoal Luas segitiga $a=8,b=6,C=60^\circ$.
\kanalisis $L=\tfrac12 ab\sin C$.
\kbahas $L=\tfrac12(8)(6)\sin60^\circ=24\cdot\tfrac{\sqrt3}{2}=12\sqrt3.$

\knomor{B.7}
\ksoal $\sin2x=\sin x$, $0^\circ\le x<360^\circ$.
\kanalisis $\sin2x=2\sin x\cos x$; faktorkan.
\kbahas $2\sin x\cos x-\sin x=0\Rightarrow\sin x(2\cos x-1)=0$. $\sin x=0\Rightarrow x=0^\circ,180^\circ$; $\cos x=\tfrac12\Rightarrow x=60^\circ,300^\circ$. HP $=\{0^\circ,60^\circ,180^\circ,300^\circ\}$.

\kbagian{Bagian C}

\knomor{C.1}
\ksoal $\sin\theta+\cos\theta=\tfrac75$; cari $\sin\theta\cos\theta$.
\kanalisis Kuadratkan kedua ruas.
\kbahas $(\sin\theta+\cos\theta)^2=1+2\sin\theta\cos\theta=\tfrac{49}{25}\Rightarrow2\sin\theta\cos\theta=\tfrac{24}{25}\Rightarrow\sin\theta\cos\theta=\tfrac{12}{25}.$

\knomor{C.2}
\ksoal $2\cos^2x-3\cos x+1=0$, $0^\circ\le x<360^\circ$.
\kanalisis Faktorkan dalam $\cos x$.
\kbahas $(2\cos x-1)(\cos x-1)=0\Rightarrow\cos x=\tfrac12$ atau $\cos x=1$. Maka $x=60^\circ,300^\circ$ atau $x=0^\circ$. HP $=\{0^\circ,60^\circ,300^\circ\}$.

\knomor{C.3}
\ksoal Buktikan $\sin3\theta=3\sin\theta-4\sin^3\theta$.
\kanalisis Ambil bagian imajiner $(\cos\theta+i\sin\theta)^3$.
\kbahas Bagian imajiner $=3\cos^2\theta\sin\theta-\sin^3\theta=3(1-\sin^2\theta)\sin\theta-\sin^3\theta=3\sin\theta-4\sin^3\theta$, yang sama dengan $\sin3\theta$ (De Moivre). $\blacksquare$

\knomor{C.4}
\ksoal $\sin x>\tfrac12$, $0^\circ\le x<360^\circ$.
\kanalisis Tinjau di mana grafik sin di atas $\tfrac12$.
\kbahas $\sin x=\tfrac12$ di $x=30^\circ,150^\circ$; di antaranya $\sin x>\tfrac12$. HP $=\{x\mid30^\circ<x<150^\circ\}$.

\knomor{C.5}
\ksoal Segitiga sisi $5,7,8$: kosinus sudut terbesar.
\kanalisis Sudut terbesar menghadap sisi $8$; aturan cosinus.
\kbahas $\cos\gamma=\dfrac{5^2+7^2-8^2}{2\cdot5\cdot7}=\dfrac{25+49-64}{70}=\dfrac{10}{70}=\dfrac17.$

\knomor{C.6}
\ksoal $h(t)=2+1{,}5\sin(30^\circ t)=2{,}75$; waktu pertama.
\kanalisis Selesaikan $\sin(30^\circ t)=\tfrac12$.
\kbahas $1{,}5\sin(30^\circ t)=0{,}75\Rightarrow\sin(30^\circ t)=\tfrac12$. Nilai terkecil positif: $30^\circ t=30^\circ\Rightarrow t=1$ jam.

\kbagian{Bagian D}

\knomor{D.1}
\ksoal Buktikan $\cos(\alpha-\beta)=\cos\alpha\cos\beta+\sin\alpha\sin\beta$ via dot product.
\kanalisis Dua vektor satuan bersudut $\alpha$ dan $\beta$.
\kbahas Ambil $\mathbf u=(\cos\alpha,\sin\alpha)$, $\mathbf v=(\cos\beta,\sin\beta)$. Maka $\mathbf u\cdot\mathbf v=\cos\alpha\cos\beta+\sin\alpha\sin\beta$. Di sisi lain $\mathbf u\cdot\mathbf v=|\mathbf u||\mathbf v|\cos(\alpha-\beta)=\cos(\alpha-\beta)$ (sudut antara keduanya $\alpha-\beta$). Jadi keduanya sama. $\blacksquare$

\knomor{D.2}
\ksoal Buktikan $\sin A+\sin B=2\sin\tfrac{A+B}{2}\cos\tfrac{A-B}{2}$.
\kanalisis Tulis $A=\tfrac{A+B}{2}+\tfrac{A-B}{2}$, $B=\tfrac{A+B}{2}-\tfrac{A-B}{2}$, lalu jabarkan.
\kbahas Misal $p=\tfrac{A+B}{2}$, $q=\tfrac{A-B}{2}$. Maka $\sin A=\sin(p+q)=\sin p\cos q+\cos p\sin q$ dan $\sin B=\sin(p-q)=\sin p\cos q-\cos p\sin q$. Jumlahnya $2\sin p\cos q=2\sin\tfrac{A+B}{2}\cos\tfrac{A-B}{2}$. $\blacksquare$

\knomor{D.3}
\ksoal Hitung $\lim_{x\to0}\dfrac{1-\cos x}{x^2}$.
\kanalisis Kalikan dengan $\dfrac{1+\cos x}{1+\cos x}$.
\kbahas $\dfrac{1-\cos x}{x^2}\cdot\dfrac{1+\cos x}{1+\cos x}=\dfrac{1-\cos^2x}{x^2(1+\cos x)}=\dfrac{\sin^2x}{x^2(1+\cos x)}=\left(\dfrac{\sin x}{x}\right)^2\dfrac{1}{1+\cos x}\to1\cdot\dfrac12=\dfrac12.$

\knomor{D.4}
\ksoal Buktikan rumus jumlah $\sum_{k=0}^{n-1}\sin(a+kd)$.
\kanalisis Kalikan $2\sin\tfrac d2$ agar teleskopik.
\kbahas $2\sin\tfrac d2\sin(a+kd)=\cos\!\big(a+kd-\tfrac d2\big)-\cos\!\big(a+kd+\tfrac d2\big)$. Saat dijumlahkan $k=0,\dots,n-1$, suku-suku saling menghapus (teleskop), menyisakan
$\cos\!\big(a-\tfrac d2\big)-\cos\!\big(a+(n-\tfrac12)d\big)=2\sin\tfrac{nd}{2}\sin\!\big(a+\tfrac{(n-1)d}{2}\big).$
Bagi dengan $2\sin\tfrac d2$ memberi hasil yang diminta. $\blacksquare$

\knomor{D.5}
\ksoal Periode $\sin x$; tulis $3\sin x+4\cos x=R\sin(x+\varphi)$.
\kanalisis $R\sin(x+\varphi)=R\cos\varphi\sin x+R\sin\varphi\cos x$; samakan koefisien.
\kbahas Karena $\sin(x+2\pi)=\sin x$, periodenya $2\pi$. Samakan: $R\cos\varphi=3$, $R\sin\varphi=4$, sehingga $R=\sqrt{3^2+4^2}=5$ dan $\tan\varphi=\tfrac43$, yakni $\varphi\approx53{,}13^\circ$. Jadi $3\sin x+4\cos x=5\sin(x+53{,}13^\circ)$.

\kbagian{Challenge Corner}

\knomor{1}
\ksoal Hitung $\cos20^\circ\cos40^\circ\cos80^\circ$.
\kanalisis Kalikan dengan $\sin20^\circ$ dan pakai $\sin2\theta=2\sin\theta\cos\theta$ berulang.
\kbahas $\sin20^\circ\cdot(\cos20^\circ\cos40^\circ\cos80^\circ)=\tfrac12\sin40^\circ\cos40^\circ\cos80^\circ=\tfrac14\sin80^\circ\cos80^\circ=\tfrac18\sin160^\circ=\tfrac18\sin20^\circ$. Karena $\sin20^\circ\neq0$, maka $\cos20^\circ\cos40^\circ\cos80^\circ=\dfrac18.$

\knomor{2}
\ksoal Buktikan $\tan1^\circ\tan2^\circ\cdots\tan89^\circ=1$.
\kanalisis Pasangkan $\theta$ dengan $90^\circ-\theta$.
\kbahas $\tan\theta\cdot\tan(90^\circ-\theta)=\tan\theta\cot\theta=1$. Pasangkan $1^\circ$ dengan $89^\circ$, $2^\circ$ dengan $88^\circ$, dan seterusnya; suku tengah $\tan45^\circ=1$. Seluruh hasil kali $=1$. $\blacksquare$

% ============================================================
\chapter{Pembahasan Bab 13: Logika Matematika}
% ============================================================

\kbagian{Bagian A}

\knomor{A.1}
\ksoal Negasi "Semua siswa rajin".
\kanalisis $\lnot(\forall x,P(x))\equiv\exists x,\lnot P(x)$.
\kbahas "Ada (sekurangnya satu) siswa yang tidak rajin".

\knomor{A.2}
\ksoal Nilai kebenaran "$2+2=4$ dan $3>5$".
\kanalisis Konjungsi benar hanya jika kedua komponen benar.
\kbahas $2+2=4$ (B), $3>5$ (S); B$\wedge$S $=$ \textbf{S}.

\knomor{A.3}
\ksoal Nilai kebenaran "$2>5$ atau $4=4$".
\kanalisis Disjungsi salah hanya jika keduanya salah.
\kbahas S$\vee$B $=$ \textbf{B}.

\knomor{A.4}
\ksoal Negasi "Hari ini hujan".
\kanalisis Balik nilai dengan $\lnot$.
\kbahas "Hari ini tidak hujan".

\knomor{A.5}
\ksoal Kontraposisi "Jika hujan maka jalan basah".
\kanalisis Kontraposisi $\lnot q\Rightarrow\lnot p$.
\kbahas "Jika jalan tidak basah maka tidak hujan".

\knomor{A.6}
\ksoal Nilai "Jika $2>3$ maka $5>1$".
\kanalisis Anteseden salah membuat implikasi benar.
\kbahas $2>3$ (S) $\Rightarrow$ implikasi \textbf{B}.

\knomor{A.7}
\ksoal Negasi $p\wedge q$.
\kanalisis De Morgan.
\kbahas $\lnot(p\wedge q)\equiv\lnot p\vee\lnot q$.

\knomor{A.8}
\ksoal Konvers "Jika $x>0$ maka $x^2>0$".
\kanalisis Konvers $q\Rightarrow p$ (tukar).
\kbahas "Jika $x^2>0$ maka $x>0$".

\knomor{A.9}
\ksoal Kapan $p\Leftrightarrow q$ benar?
\kanalisis Biimplikasi benar saat nilai keduanya sama.
\kbahas Saat $p$ dan $q$ keduanya benar atau keduanya salah.

\knomor{A.10}
\ksoal Negasi "$\exists x,\,P(x)$".
\kanalisis $\lnot\exists\equiv\forall\lnot$.
\kbahas $\forall x,\,\lnot P(x)$.

\kbagian{Bagian B}

\knomor{B.1}
\ksoal Tabel kebenaran $p\Rightarrow q$; kapan salah?
\kanalisis Tinjau keempat kombinasi.
\kbahas
\begin{center}\begin{tabular}{c c | c}\toprule $p$&$q$&$p\Rightarrow q$\\\midrule B&B&B\\B&S&S\\S&B&B\\S&S&B\\\bottomrule\end{tabular}\end{center}
Bernilai \textbf{salah hanya} saat $p$ benar dan $q$ salah.

\knomor{B.2}
\ksoal Negasi "Jika $n$ genap maka $n^2$ genap".
\kanalisis $\lnot(p\Rightarrow q)\equiv p\wedge\lnot q$.
\kbahas "$n$ genap dan $n^2$ tidak genap (ganjil)".

\knomor{B.3}
\ksoal Tunjukkan $p\Rightarrow q\equiv\lnot p\vee q$.
\kanalisis Bandingkan kolom tabel kebenaran.
\kbahas
\begin{center}\begin{tabular}{c c | c c}\toprule $p$&$q$&$p\Rightarrow q$&$\lnot p\vee q$\\\midrule B&B&B&B\\B&S&S&S\\S&B&B&B\\S&S&B&B\\\bottomrule\end{tabular}\end{center}
Kedua kolom identik, jadi ekuivalen. $\blacksquare$

\knomor{B.4}
\ksoal Negasi "$\forall x\in\mathbb{R},\ x^2\ge0$".
\kanalisis $\lnot\forall\equiv\exists\lnot$.
\kbahas "Ada bilangan real $x$ dengan $x^2<0$".

\knomor{B.5}
\ksoal Nilai $(p\wedge\lnot q)\Rightarrow r$ untuk $p$=B, $q$=S, $r$=B.
\kanalisis Hitung bertahap.
\kbahas $\lnot q=$ B; $p\wedge\lnot q=$ B$\wedge$B$=$ B; B$\Rightarrow$B $=$ \textbf{B}.

\kbagian{Bagian C}

\knomor{C.1}
\ksoal Buktikan "jika $n$ ganjil maka $n^2$ ganjil".
\kanalisis Tulis $n=2k+1$ lalu kuadratkan.
\kbahas $n=2k+1\Rightarrow n^2=4k^2+4k+1=2(2k^2+2k)+1$, berbentuk genap$+1$, jadi ganjil. $\blacksquare$

\knomor{C.2}
\ksoal Tunjukkan $[(p\Rightarrow q)\wedge p]\Rightarrow q$ tautologi.
\kanalisis Periksa seluruh baris bernilai B.
\kbahas
\begin{center}\begin{tabular}{c c | c c c}\toprule $p$&$q$&$p\Rightarrow q$&$(p\Rightarrow q)\wedge p$&$[\,]\Rightarrow q$\\\midrule B&B&B&B&B\\B&S&S&S&B\\S&B&B&S&B\\S&S&B&S&B\\\bottomrule\end{tabular}\end{center}
Kolom terakhir seluruhnya B $\Rightarrow$ tautologi. $\blacksquare$

\knomor{C.3}
\ksoal Negasikan "$\forall x\,\exists y,\ x+y=0$".
\kanalisis Tukar tiap kuantor dan negasikan inti.
\kbahas $\exists x\,\forall y,\ x+y\neq0$.

\knomor{C.4}
\ksoal Nilai kebenaran "$\forall x\in\mathbb{R}\,\exists y\in\mathbb{R},\ y^2=x$".
\kanalisis Cari penyangkal (kuantor universal salah bila ada satu $x$ gagal).
\kbahas \textbf{Salah}: untuk $x=-1$ tidak ada $y$ real dengan $y^2=-1$.

\knomor{C.5}
\ksoal Sederhanakan $\lnot(p\Rightarrow q)$ dan $\lnot(p\Leftrightarrow q)$.
\kanalisis Pakai negasi implikasi dan definisi biimplikasi.
\kbahas $\lnot(p\Rightarrow q)\equiv p\wedge\lnot q$; $\lnot(p\Leftrightarrow q)\equiv(p\wedge\lnot q)\vee(\lnot p\wedge q)$ (yaitu "$p$ dan $q$ berbeda nilai").

\kbagian{Bagian D}

\knomor{D.1}
\ksoal Buktikan bilangan prima tak hingga banyaknya.
\kanalisis Andaikan berhingga, bentuk bilangan baru (kontradiksi).
\kbahas Andaikan prima hanya $p_1,p_2,\dots,p_n$. Tinjau $N=p_1p_2\cdots p_n+1$. $N$ tak habis dibagi $p_i$ mana pun (sisa $1$), jadi $N$ memiliki faktor prima di luar daftar atau $N$ sendiri prima --- bertentangan dengan asumsi daftar lengkap. Maka prima tak hingga. $\blacksquare$

\knomor{D.2}
\ksoal Buktikan $[(p\Rightarrow q)\wedge(q\Rightarrow r)]\Rightarrow(p\Rightarrow r)$ tautologi.
\kanalisis Cukup tinjau kasus anteseden benar (premis kedua-duanya benar).
\kbahas Jika $p\Rightarrow q$ dan $q\Rightarrow r$ benar, ambil sembarang kasus $p$ benar; maka $q$ benar (premis 1), lalu $r$ benar (premis 2), sehingga $p\Rightarrow r$ benar. Bila anteseden gabungan salah, implikasi otomatis benar. Jadi selalu benar. $\blacksquare$

\knomor{D.3}
\ksoal Buktikan jumlah dua bilangan genap genap.
\kanalisis Pembuktian langsung dengan bentuk umum.
\kbahas Misal $a=2m$, $b=2n$. Maka $a+b=2m+2n=2(m+n)$, kelipatan $2$, jadi genap. $\blacksquare$

\knomor{D.4}
\ksoal Buktikan: jika $3n+2$ ganjil maka $n$ ganjil.
\kanalisis Kontraposisi: $n$ genap $\Rightarrow 3n+2$ genap.
\kbahas Misal $n=2k$. Maka $3n+2=6k+2=2(3k+1)$ genap. Kontraposisi benar, jadi pernyataan asal benar. $\blacksquare$

\knomor{D.5}
\ksoal Buktikan "$\forall x\in\mathbb{R},\ x^2-2x+2>0$".
\kanalisis Lengkapkan kuadrat.
\kbahas $x^2-2x+2=(x-1)^2+1\ge1>0$ untuk setiap $x$. $\blacksquare$

\kbagian{Challenge Corner}

\knomor{1}
\ksoal Penduduk $A$ berkata "Saya penipu". Tunjukkan ini mustahil.
\kanalisis Uji dua kemungkinan jenis $A$.
\kbahas Jika $A$ ksatria (berkata benar), maka pernyataannya benar $\Rightarrow A$ penipu --- kontradiksi. Jika $A$ penipu (berkata bohong), maka pernyataannya salah $\Rightarrow A$ bukan penipu (ksatria) --- kontradiksi. Kedua kasus mustahil, jadi tak ada penduduk yang dapat mengucapkannya. $\blacksquare$

\knomor{2}
\ksoal Buktikan dengan kontradiksi: jika $3n+2$ genap maka $n$ genap.
\kanalisis Andaikan $3n+2$ genap tetapi $n$ ganjil.
\kbahas Misal $n=2k+1$. Maka $3n+2=6k+3+2=6k+5=2(3k+2)+1$ ganjil --- bertentangan dengan asumsi $3n+2$ genap. Jadi $n$ pasti genap. $\blacksquare$

\knomor{3}
\ksoal Buktikan $\sqrt2+\sqrt3$ irasional.
\kanalisis Andaikan rasional, kuadratkan, isolasi $\sqrt6$.
\kbahas Andaikan $\sqrt2+\sqrt3=r$ rasional. Maka $r^2=2+2\sqrt6+3=5+2\sqrt6$, sehingga $\sqrt6=\dfrac{r^2-5}{2}$ rasional. Padahal $\sqrt6$ irasional (sebab $6$ bukan kuadrat sempurna) --- kontradiksi. Jadi $\sqrt2+\sqrt3$ irasional. $\blacksquare$

% ============================================================
\chapter{Pembahasan Bab 14: Induksi Matematika}
% ============================================================

\kbagian{Bagian A}

\knomor{A.1}
\ksoal Basis induksi $1+2+\cdots+n=\tfrac{n(n+1)}2$.
\kanalisis Substitusi $n=1$ ke kedua ruas.
\kbahas Ruas kiri $=1$; ruas kanan $=\tfrac{1\cdot2}{2}=1$. Sama, basis benar.

\knomor{A.2}
\ksoal Hipotesis induksi $P(k)$ untuk $1+3+\cdots+(2n-1)=n^2$.
\kanalisis Ganti $n$ dengan $k$.
\kbahas $P(k):\ 1+3+5+\cdots+(2k-1)=k^2$.

\knomor{A.3}
\ksoal Hitung $1+2+\cdots+10$.
\kanalisis Pakai rumus $\tfrac{n(n+1)}2$.
\kbahas $\tfrac{10\cdot11}{2}=55$.

\knomor{A.4}
\ksoal Hitung $1+3+\cdots+(2\cdot5-1)$.
\kanalisis Jumlah $n$ ganjil pertama $=n^2$, $n=5$.
\kbahas $5^2=25$.

\knomor{A.5}
\ksoal Periksa basis $n=1$ untuk $2^n>n$.
\kanalisis Substitusi $n=1$.
\kbahas $2^1=2>1$, benar.

\knomor{A.6}
\ksoal Suku yang ditambahkan dari $n=k$ ke $n=k+1$.
\kanalisis Suku berikutnya dalam $1+2+\cdots$.
\kbahas Suku $(k+1)$.

\knomor{A.7}
\ksoal Periksa $3\mid(n^3-n)$ untuk $n=2$.
\kanalisis Hitung $n^3-n$.
\kbahas $8-2=6$, dan $3\mid6$. Benar.

\knomor{A.8}
\ksoal Bentuk $P(k+1)$ dari $P(n):n^2$.
\kanalisis Ganti $n$ dengan $k+1$.
\kbahas $(k+1)^2$.

\knomor{A.9}
\ksoal Hitung $2+4+6+8$ dan bandingkan $n(n+1)$, $n=4$.
\kanalisis Jumlahkan langsung.
\kbahas $2+4+6+8=20$; $n(n+1)=4\cdot5=20$. Sama.

\knomor{A.10}
\ksoal Basis terkecil untuk $2^n>n^2$.
\kanalisis $2^n>n^2$ gagal di $n=2,3,4$, berlaku mulai $n=5$.
\kbahas $n_0=5$.

\kbagian{Bagian B}

\knomor{B.1}
\ksoal Buktikan $1+2+\cdots+n=\tfrac{n(n+1)}2$.
\kanalisis Basis $n=1$; langkah tambahkan $(k+1)$.
\kbahas Basis benar ($1=1$). Andaikan $1+\cdots+k=\tfrac{k(k+1)}2$. Tambah $(k+1)$: $\tfrac{k(k+1)}2+(k+1)=\tfrac{(k+1)(k+2)}2$, sesuai rumus untuk $k+1$. $\blacksquare$

\knomor{B.2}
\ksoal Buktikan $1+3+\cdots+(2n-1)=n^2$.
\kanalisis Suku ke-$(k+1)$ adalah $2(k+1)-1=2k+1$.
\kbahas Basis $n=1$: $1=1^2$. Andaikan jumlah $=k^2$. Tambah $2k+1$: $k^2+2k+1=(k+1)^2$. $\blacksquare$

\knomor{B.3}
\ksoal Buktikan $1^2+\cdots+n^2=\tfrac{n(n+1)(2n+1)}6$.
\kanalisis Tambahkan $(k+1)^2$ ke hipotesis lalu faktorkan.
\kbahas Basis $n=1$: $1=\tfrac{1\cdot2\cdot3}{6}=1$. Langkah: $\tfrac{k(k+1)(2k+1)}6+(k+1)^2=\tfrac{(k+1)[k(2k+1)+6(k+1)]}6=\tfrac{(k+1)(2k^2+7k+6)}6=\tfrac{(k+1)(k+2)(2k+3)}6$, sesuai rumus $k+1$. $\blacksquare$

\knomor{B.4}
\ksoal Buktikan $3\mid(n^3-n)$.
\kanalisis Uraikan $(k+1)^3-(k+1)$.
\kbahas Basis $n=1$: $0$ habis dibagi $3$. Langkah: $(k+1)^3-(k+1)=(k^3-k)+3(k^2+k)$. Suku pertama kelipatan $3$ (hipotesis), suku kedua kelipatan $3$. Jadi habis dibagi $3$. $\blacksquare$

\knomor{B.5}
\ksoal Buktikan $1+2+\cdots+2^{n-1}=2^n-1$.
\kanalisis Tambahkan $2^k$ ke hipotesis.
\kbahas Basis $n=1$: $2^0=1=2^1-1$. Langkah: $(2^k-1)+2^k=2\cdot2^k-1=2^{k+1}-1$. $\blacksquare$

\kbagian{Bagian C}

\knomor{C.1}
\ksoal Buktikan $a+ar+\cdots+ar^{n-1}=\dfrac{a(r^n-1)}{r-1}$, $r\neq1$.
\kanalisis Tambahkan $ar^k$ ke hipotesis dan satukan pecahan.
\kbahas Basis $n=1$: $a=\tfrac{a(r-1)}{r-1}=a$. Langkah: $\tfrac{a(r^k-1)}{r-1}+ar^k=\tfrac{a(r^k-1)+ar^k(r-1)}{r-1}=\tfrac{ar^{k+1}-a}{r-1}=\tfrac{a(r^{k+1}-1)}{r-1}$. $\blacksquare$

\knomor{C.2}
\ksoal Buktikan $2^n>n^2$ untuk $n\ge5$.
\kanalisis Basis $n=5$; langkah pakai $2k^2\ge(k+1)^2$.
\kbahas Basis: $32>25$. Andaikan $2^k>k^2$. Maka $2^{k+1}=2\cdot2^k>2k^2$. Karena $k\ge5$, $2k^2-(k+1)^2=k^2-2k-1=(k-1)^2-2>0$, jadi $2k^2>(k+1)^2$. Maka $2^{k+1}>(k+1)^2$. $\blacksquare$

\knomor{C.3}
\ksoal Buktikan $6\mid(n^3+5n)$.
\kanalisis Uraikan $(k+1)^3+5(k+1)$; tunjukkan tambahan kelipatan $6$.
\kbahas Basis $n=1$: $1+5=6$. Langkah: $(k+1)^3+5(k+1)=(k^3+5k)+3(k^2+k+2)$. Suku pertama kelipatan $6$ (hipotesis). Pada suku kedua, $k^2+k=k(k+1)$ genap, jadi $k^2+k+2$ genap; maka $3(k^2+k+2)$ kelipatan $6$. Jumlahnya kelipatan $6$. $\blacksquare$

\knomor{C.4}
\ksoal Buktikan $(1+x)^n\ge1+nx$ untuk $x\ge-1$ (Bernoulli).
\kanalisis Kalikan hipotesis dengan $(1+x)\ge0$.
\kbahas Basis $n=1$: $1+x\ge1+x$. Langkah: $(1+x)^{k+1}=(1+x)^k(1+x)\ge(1+kx)(1+x)=1+(k+1)x+kx^2\ge1+(k+1)x$ (karena $kx^2\ge0$). $\blacksquare$

\knomor{C.5}
\ksoal Buktikan $\begin{pmatrix}1&1\\0&1\end{pmatrix}^{\!n}=\begin{pmatrix}1&n\\0&1\end{pmatrix}$.
\kanalisis Kalikan hipotesis dengan matriks dasar.
\kbahas Basis $n=1$ jelas. Langkah: $\begin{pmatrix}1&k\\0&1\end{pmatrix}\begin{pmatrix}1&1\\0&1\end{pmatrix}=\begin{pmatrix}1&k+1\\0&1\end{pmatrix}$. $\blacksquare$

\kbagian{Bagian D}

\knomor{D.1}
\ksoal Buktikan setiap $n\ge2$ hasil kali bilangan prima.
\kanalisis Induksi kuat: pakai semua kasus $<k+1$.
\kbahas Basis $n=2$ prima. Andaikan benar untuk semua $2,\dots,k$. Untuk $k+1$: bila prima, selesai; bila tidak, $k+1=ab$ dengan $2\le a,b\le k$. Menurut hipotesis kuat, $a$ dan $b$ masing-masing hasil kali prima, jadi $k+1=ab$ juga. $\blacksquare$

\knomor{D.2}
\ksoal Buktikan $\begin{pmatrix}1&1\\1&0\end{pmatrix}^{\!n}=\begin{pmatrix}F_{n+1}&F_n\\F_n&F_{n-1}\end{pmatrix}$.
\kanalisis Gunakan $F_{k+1}+F_k=F_{k+2}$ saat mengalikan.
\kbahas Basis $n=1$: ruas kanan $\begin{pmatrix}F_2&F_1\\F_1&F_0\end{pmatrix}=\begin{pmatrix}1&1\\1&0\end{pmatrix}$. Langkah: $\begin{pmatrix}F_{k+1}&F_k\\F_k&F_{k-1}\end{pmatrix}\begin{pmatrix}1&1\\1&0\end{pmatrix}=\begin{pmatrix}F_{k+1}+F_k&F_{k+1}\\F_k+F_{k-1}&F_k\end{pmatrix}=\begin{pmatrix}F_{k+2}&F_{k+1}\\F_{k+1}&F_k\end{pmatrix}$. $\blacksquare$

\knomor{D.3}
\ksoal Buktikan $1+\tfrac12+\cdots+\tfrac{1}{2^n}\ge1+\tfrac n2$.
\kanalisis Dari $2^k$ ke $2^{k+1}$, tambahkan $2^k$ suku, masing-masing $\ge\tfrac{1}{2^{k+1}}$.
\kbahas Basis $n=1$: $1+\tfrac12=\tfrac32\ge\tfrac32$. Langkah: suku-suku baru $\tfrac{1}{2^k+1}+\cdots+\tfrac{1}{2^{k+1}}$ berjumlah $2^k$ buah, tiap $\ge\tfrac{1}{2^{k+1}}$, sehingga totalnya $\ge\tfrac{2^k}{2^{k+1}}=\tfrac12$. Maka jumlah baru $\ge\big(1+\tfrac k2\big)+\tfrac12=1+\tfrac{k+1}2$. $\blacksquare$

\knomor{D.4}
\ksoal Buktikan himpunan $n$ anggota punya $2^n$ himpunan bagian.
\kanalisis Tambahkan satu anggota; tiap subset lama menghasilkan dua subset.
\kbahas Basis $n=0$: himpunan kosong punya $1=2^0$ himpunan bagian. Langkah: himpunan $k+1$ anggota; setiap himpunan bagian entah memuat anggota baru atau tidak. Yang tidak memuat: $2^k$ (hipotesis); yang memuat: $2^k$ pula. Total $2\cdot2^k=2^{k+1}$. $\blacksquare$

\knomor{D.5}
\ksoal Buktikan De Moivre $(\cos\theta+i\sin\theta)^n=\cos n\theta+i\sin n\theta$.
\kanalisis Kalikan hipotesis dengan $(\cos\theta+i\sin\theta)$, pakai rumus jumlah sudut.
\kbahas Basis $n=1$ trivial. Langkah: $(\cos k\theta+i\sin k\theta)(\cos\theta+i\sin\theta)=(\cos k\theta\cos\theta-\sin k\theta\sin\theta)+i(\sin k\theta\cos\theta+\cos k\theta\sin\theta)=\cos(k+1)\theta+i\sin(k+1)\theta$. $\blacksquare$

\kbagian{Challenge Corner}

\knomor{1}
\ksoal Letak kesalahan "bukti semua kuda berwarna sama".
\kanalisis Periksa apakah langkah berlaku tepat saat $k=1\to k+1=2$.
\kbahas Langkah induksi menyimpulkan warna sama dengan menumpang-tindihkan dua kelompok $k$ kuda. Tumpang-tindih ini hanya ada bila $k\ge2$. Untuk $k=1\to2$, kedua kelompok berukuran $1$ \emph{tidak beririsan}, sehingga tak ada penghubung warna --- rantai induksi putus di $n=2$. Maka kesimpulannya tidak sah. $\blacksquare$

\knomor{2}
\ksoal Buktikan $n!>2^n$ untuk $n\ge4$.
\kanalisis Basis $n=4$; langkah kalikan dengan $(k+1)$.
\kbahas Basis: $4!=24>16=2^4$. Langkah: $(k+1)!=(k+1)\cdot k!>(k+1)\cdot2^k\ge2\cdot2^k=2^{k+1}$ karena $k+1\ge5>2$. $\blacksquare$

\knomor{3}
\ksoal Buktikan $F_1+F_2+\cdots+F_n=F_{n+2}-1$.
\kanalisis Tambahkan $F_{k+1}$ dan pakai $F_{k+2}+F_{k+1}=F_{k+3}$.
\kbahas Basis $n=1$: $F_1=1=F_3-1=2-1$. Langkah: $(F_{k+2}-1)+F_{k+1}=(F_{k+2}+F_{k+1})-1=F_{k+3}-1$. $\blacksquare$

% ============================================================
\chapter{Pembahasan Bab 15: Aturan Pencacahan}
% ============================================================

\kbagian{Bagian A}

\knomor{A.1}
\ksoal Hitung $5!$.
\kanalisis Kalikan $5\cdot4\cdot3\cdot2\cdot1$.
\kbahas $5!=120$.

\knomor{A.2}
\ksoal Hitung $0!$.
\kanalisis Definisi $0!=1$.
\kbahas $0!=1$.

\knomor{A.3}
\ksoal Hitung $P(5,2)$.
\kanalisis $P(n,r)=\dfrac{n!}{(n-r)!}$.
\kbahas $\dfrac{5!}{3!}=5\cdot4=20$.

\knomor{A.4}
\ksoal Hitung $C(5,2)$.
\kanalisis $\binom nr=\dfrac{n!}{r!(n-r)!}$.
\kbahas $\dfrac{5!}{2!\,3!}=\dfrac{20}{2}=10$.

\knomor{A.5}
\ksoal Pilih satu item dari $3$ baju atau $2$ celana.
\kanalisis "Atau" $\to$ aturan penjumlahan.
\kbahas $3+2=5$ pilihan.

\knomor{A.6}
\ksoal Pasang (baju, celana) dari $3$ baju, $2$ celana.
\kanalisis "Dan" $\to$ aturan perkalian.
\kbahas $3\times2=6$ pasang.

\knomor{A.7}
\ksoal Hitung $\dfrac{6!}{4!}$.
\kanalisis Coret faktor sekutu.
\kbahas $6\cdot5=30$.

\knomor{A.8}
\ksoal Hitung $\binom60+\binom66$.
\kanalisis $\binom n0=\binom nn=1$.
\kbahas $1+1=2$.

\knomor{A.9}
\ksoal Susunan huruf "PQRS".
\kanalisis Empat huruf berbeda $\to 4!$.
\kbahas $4!=24$.

\knomor{A.10}
\ksoal Hitung $\binom75$.
\kanalisis $\binom75=\binom72$.
\kbahas $\dfrac{7\cdot6}{2}=21$.

\kbagian{Bagian B}

\knomor{B.1}
\ksoal Bilangan $3$ angka berbeda dari $\{1,2,3,4,5\}$.
\kanalisis Urutan penting, tanpa pengulangan: $P(5,3)$.
\kbahas $P(5,3)=5\cdot4\cdot3=60$.

\knomor{B.2}
\ksoal $6$ orang di meja bundar.
\kanalisis Permutasi siklis $(n-1)!$.
\kbahas $(6-1)!=120$.

\knomor{B.3}
\ksoal Susunan huruf "MATEMATIKA".
\kanalisis $10$ huruf dengan A$\times3$, M$\times2$, T$\times2$.
\kbahas $\dfrac{10!}{3!\,2!\,2!}=\dfrac{3628800}{24}=151200$.

\knomor{B.4}
\ksoal Pilih $3$ dari $10$ orang (tanpa jabatan).
\kanalisis Urutan tak penting $\to$ kombinasi.
\kbahas $\binom{10}{3}=120$.

\knomor{B.5}
\ksoal Koefisien $x^2$ pada $(1+x)^5$.
\kanalisis Suku $\binom5k x^k$, ambil $k=2$.
\kbahas $\binom52=10$.

\kbagian{Bagian C}

\knomor{C.1}
\ksoal Banyak cara memilih $5$ kartu dari $52$.
\kanalisis Urutan tak penting: $\binom{52}{5}$.
\kbahas $\binom{52}{5}=2\,598\,960$.

\knomor{C.2}
\ksoal Bilangan genap $3$ angka berbeda dari $\{1,2,3,4,5\}$.
\kanalisis Angka satuan harus genap ($2$ atau $4$); isi dua posisi sisa dari $4$ angka tersisa.
\kbahas Satuan: $2$ pilihan. Ratusan--puluhan: $4\cdot3=12$. Total $2\times12=24$.

\knomor{C.3}
\ksoal Suku ke-$4$ dari $(2x+3)^5$.
\kanalisis Suku ke-$(k+1)=\binom5k(2x)^{5-k}3^k$; ke-$4\Rightarrow k=3$.
\kbahas $\binom53(2x)^2 3^3=10\cdot4x^2\cdot27=1080x^2$.

\knomor{C.4}
\ksoal Susunan "DUTA" dengan vokal berdampingan.
\kanalisis Gabung vokal $\{U,A\}$ menjadi satu blok.
\kbahas Unit: $\{UA\}, D, T \Rightarrow 3!$ susunan; di dalam blok $2!$. Total $3!\cdot2!=12$.

\knomor{C.5}
\ksoal Buktikan $\binom nk=\binom{n}{n-k}$ secara kombinatorial.
\kanalisis Memilih $k$ untuk diambil setara memilih $n-k$ untuk ditinggalkan.
\kbahas Setiap cara memilih $k$ objek menentukan secara unik $n-k$ objek yang tidak dipilih, dan sebaliknya. Korespondensi satu-satu ini memberi $\binom nk=\binom{n}{n-k}$. $\blacksquare$

\kbagian{Bagian D}

\knomor{D.1}
\ksoal Buktikan $\displaystyle\sum_{k=0}^n\binom nk=2^n$.
\kanalisis Substitusi $a=b=1$ pada teorema binomial, atau hitung subset.
\kbahas $(1+1)^n=\sum_{k=0}^n\binom nk 1^{n-k}1^k=\sum\binom nk$, dan ruas kiri $=2^n$. (Tafsiran: total himpunan bagian dari $n$ objek.) $\blacksquare$

\knomor{D.2}
\ksoal Buktikan $\binom nk=\binom{n-1}{k-1}+\binom{n-1}{k}$.
\kanalisis Tetapkan satu objek khusus; pilihan memuatnya atau tidak.
\kbahas Untuk memilih $k$ dari $n$ objek, sorot satu objek $x$. Jika $x$ terpilih, sisanya $\binom{n-1}{k-1}$ cara; jika tidak, $\binom{n-1}{k}$ cara. Keduanya saling lepas, jumlahnya $\binom nk$. $\blacksquare$

\knomor{D.3}
\ksoal Solusi bulat tak negatif $x_1+x_2+x_3=10$.
\kanalisis Bintang dan batang: $10$ bintang, $2$ batang.
\kbahas $\binom{10+3-1}{3-1}=\binom{12}{2}=66$.

\knomor{D.4}
\ksoal Bilangan $1$--$100$ yang habis dibagi $2$ atau $3$.
\kanalisis Inklusi--eksklusi $|A\cup B|=|A|+|B|-|A\cap B|$.
\kbahas $|A|=50$ (kelipatan $2$), $|B|=33$ (kelipatan $3$), $|A\cap B|=16$ (kelipatan $6$). Maka $50+33-16=67$.

\knomor{D.5}
\ksoal Buktikan $\displaystyle\sum_{k=0}^n k\binom nk=n\,2^{n-1}$.
\kanalisis Pakai identitas $k\binom nk=n\binom{n-1}{k-1}$.
\kbahas $\sum_{k=0}^n k\binom nk=\sum_{k=1}^n n\binom{n-1}{k-1}=n\sum_{j=0}^{n-1}\binom{n-1}{j}=n\cdot2^{n-1}$. $\blacksquare$

\kbagian{Challenge Corner}

\knomor{1}
\ksoal Banyak diagonal segi-$n$.
\kanalisis Pilih $2$ titik sudut, kurangi $n$ sisi.
\kbahas $\binom n2-n=\dfrac{n(n-1)}2-n=\dfrac{n(n-3)}2$.

\knomor{2}
\ksoal Jalur terpendek $(0,0)\to(m,n)$ (kanan/atas).
\kanalisis Susun urutan $m$ langkah "kanan" dan $n$ langkah "atas".
\kbahas Jumlah cara menyusun barisan $m+n$ langkah dengan $m$ di antaranya "kanan": $\dbinom{m+n}{m}$.

\knomor{3}
\ksoal Buktikan $\displaystyle\sum_{k=0}^n\binom nk^2=\binom{2n}{n}$.
\kanalisis Hitung dua kali: pilih $n$ orang dari gabungan dua kelompok berukuran $n$.
\kbahas Pilih $n$ orang dari $2n$ ($n$ "putra" $+$ $n$ "putri"): $\binom{2n}{n}$. Di sisi lain, jika $k$ diambil dari putra dan $n-k$ dari putri, banyaknya $\binom nk\binom{n}{n-k}=\binom nk^2$. Jumlahkan atas $k$: $\sum_{k=0}^n\binom nk^2=\binom{2n}{n}$. $\blacksquare$

% ============================================================
\chapter{Pembahasan Bab 16: Peluang}
% ============================================================

\kbagian{Bagian A}

\knomor{A.1}
\ksoal Ruang sampel satu dadu.
\kanalisis Daftar semua hasil mungkin.
\kbahas $S=\{1,2,3,4,5,6\}$.

\knomor{A.2}
\ksoal Peluang mata genap.
\kanalisis $A=\{2,4,6\}$.
\kbahas $P=\dfrac{3}{6}=\dfrac12$.

\knomor{A.3}
\ksoal Peluang mata $>4$.
\kanalisis $A=\{5,6\}$.
\kbahas $P=\dfrac{2}{6}=\dfrac13$.

\knomor{A.4}
\ksoal Peluang gambar pada satu koin.
\kanalisis $S=\{A,G\}$.
\kbahas $P=\dfrac12$.

\knomor{A.5}
\ksoal Peluang kartu hati dari $52$.
\kanalisis Ada $13$ kartu hati.
\kbahas $P=\dfrac{13}{52}=\dfrac14$.

\knomor{A.6}
\ksoal $P(A^c)$ jika $P(A)=0{,}3$.
\kanalisis Komplemen.
\kbahas $1-0{,}3=0{,}7$.

\knomor{A.7}
\ksoal Banyak anggota ruang sampel dua koin.
\kanalisis $\{AA,AG,GA,GG\}$.
\kbahas $4$.

\knomor{A.8}
\ksoal Peluang dua gambar pada dua koin.
\kanalisis Hanya $\{GG\}$ dari $4$ hasil.
\kbahas $P=\dfrac14$.

\knomor{A.9}
\ksoal $P(A\cup B)$ saling lepas, $P(A)=0{,}5$, $P(B)=0{,}3$.
\kanalisis Saling lepas $\Rightarrow$ jumlahkan.
\kbahas $0{,}5+0{,}3=0{,}8$.

\knomor{A.10}
\ksoal Peluang bola merah ($3$ merah, $2$ biru).
\kanalisis $n(A)=3$, $n(S)=5$.
\kbahas $P=\dfrac35$.

\kbagian{Bagian B}

\knomor{B.1}
\ksoal Peluang jumlah dua dadu $=7$.
\kanalisis Pasangan berjumlah $7$ ada $6$, dari $36$.
\kbahas $\{(1,6),(2,5),(3,4),(4,3),(5,2),(6,1)\}\Rightarrow P=\dfrac{6}{36}=\dfrac16$.

\knomor{B.2}
\ksoal $P(A\cup B)$, $P(A)=0{,}6$, $P(B)=0{,}5$, $P(A\cap B)=0{,}2$.
\kanalisis Rumus gabungan.
\kbahas $0{,}6+0{,}5-0{,}2=0{,}9$.

\knomor{B.3}
\ksoal $P(A\mid B)$, $P(A\cap B)=0{,}12$, $P(B)=0{,}3$.
\kanalisis $P(A\mid B)=\dfrac{P(A\cap B)}{P(B)}$.
\kbahas $\dfrac{0{,}12}{0{,}3}=0{,}4$.

\knomor{B.4}
\ksoal Peluang King atau hati dari $52$.
\kanalisis Gabungan dengan irisan King hati $=1$.
\kbahas $\dfrac{4+13-1}{52}=\dfrac{16}{52}=\dfrac{4}{13}$.

\knomor{B.5}
\ksoal Peluang dua bola merah ($5$ merah, $3$ putih), tanpa pengembalian.
\kanalisis Kalikan peluang berurutan.
\kbahas $\dfrac58\cdot\dfrac47=\dfrac{20}{56}=\dfrac{5}{14}$.

\kbagian{Bagian C}

\knomor{C.1}
\ksoal Peluang Andi terpilih (pilih $3$ dari $10$).
\kanalisis Hitung kasus memuat Andi $\div$ total, atau langsung $\tfrac{3}{10}$.
\kbahas $\dfrac{\binom92}{\binom{10}3}=\dfrac{36}{120}=\dfrac{3}{10}$.

\knomor{C.2}
\ksoal $A,B$ independen, $P(A)=0{,}4$, $P(B)=0{,}5$.
\kanalisis $P(A\cap B)=P(A)P(B)$.
\kbahas $P(A\cap B)=0{,}2$; $P(A\cup B)=0{,}4+0{,}5-0{,}2=0{,}7$.

\knomor{C.3}
\ksoal $P(\text{hasil}=3\mid \text{ganjil})$.
\kanalisis Bersyarat pada $\{1,3,5\}$.
\kbahas $\dfrac{P(3)}{P(\text{ganjil})}=\dfrac{1/6}{3/6}=\dfrac13$.

\knomor{C.4}
\ksoal Peluang minimal satu gambar (tiga koin).
\kanalisis Pakai komplemen "tak ada gambar".
\kbahas $1-\left(\tfrac12\right)^3=1-\dfrac18=\dfrac78$.

\knomor{C.5}
\ksoal Bayes: $P(A\mid \text{cacat})$.
\kanalisis Hukum peluang total untuk penyebut.
\kbahas $P(\text{cacat})=0{,}6\cdot0{,}02+0{,}4\cdot0{,}05=0{,}032$; $P(A\mid\text{cacat})=\dfrac{0{,}012}{0{,}032}=0{,}375$.

\kbagian{Bagian D}

\knomor{D.1}
\ksoal Tes medis: $P(\text{sakit}\mid +)$.
\kanalisis Bayes; prevalensi $1\%$, sensitivitas $99\%$, positif palsu $5\%$.
\kbahas $P(+)=0{,}01\cdot0{,}99+0{,}99\cdot0{,}05=0{,}0099+0{,}0495=0{,}0594$. Maka $P(\text{sakit}\mid+)=\dfrac{0{,}0099}{0{,}0594}\approx0{,}167$ (sekitar $16{,}7\%$).

\knomor{D.2}
\ksoal Buktikan $P(A^c)=1-P(A)$.
\kanalisis $A$ dan $A^c$ saling lepas dan menutupi $S$.
\kbahas $A\cup A^c=S$ dan $A\cap A^c=\varnothing$, jadi $P(A)+P(A^c)=P(S)=1$, sehingga $P(A^c)=1-P(A)$. $\blacksquare$

\knomor{D.3}
\ksoal Buktikan $A$ independen $B\Rightarrow A$ independen $B^c$.
\kanalisis Pisahkan $A=(A\cap B)\cup(A\cap B^c)$.
\kbahas $P(A\cap B^c)=P(A)-P(A\cap B)=P(A)-P(A)P(B)=P(A)\big(1-P(B)\big)=P(A)P(B^c)$. Jadi $A$ dan $B^c$ independen. $\blacksquare$

\knomor{D.4}
\ksoal Nilai harapan pembayaran (mata dadu, ribuan).
\kanalisis $E=\sum x_i P(x_i)$.
\kbahas $E=\dfrac{1+2+3+4+5+6}{6}=\dfrac{21}{6}=3{,}5$ (ribu rupiah).

\knomor{D.5}
\ksoal Hukum peluang total: peluang bola merah.
\kanalisis Bobot tiap kotak $\tfrac12$.
\kbahas $P(\text{merah})=\tfrac12\cdot\tfrac25+\tfrac12\cdot\tfrac45=\tfrac15+\tfrac25=\dfrac35$.

\kbagian{Challenge Corner}

\knomor{1}
\ksoal Monty Hall: sebaiknya pindah?
\kanalisis Bandingkan peluang menang strategi "tetap" vs "pindah".
\kbahas Peluang pilihan awal benar $=\tfrac13$ (strategi tetap menang $\tfrac13$). Maka strategi \textbf{pindah} menang bila pilihan awal salah, yaitu $\tfrac23$. Jadi sebaiknya berpindah; peluang menang $\dfrac23$.

\knomor{2}
\ksoal Masalah ulang tahun ($23$ orang).
\kanalisis Pakai komplemen "semua tanggal berbeda".
\kbahas $P(\text{ada sama})=1-\dfrac{365\cdot364\cdots343}{365^{23}}=1-\dfrac{365!}{342!\,\cdot365^{23}}\approx0{,}507$. Walau $23$ terasa sedikit, banyaknya \emph{pasangan} $\binom{23}{2}=253$ membuat peluang kecocokan melampaui $\tfrac12$.

\knomor{3}
\ksoal Peluang tepat dua dari tiga penembak ($0{,}6;0{,}7;0{,}8$) mengenai.
\kanalisis Jumlahkan tiga kemungkinan "tepat dua".
\kbahas $0{,}6\cdot0{,}7\cdot0{,}2+0{,}6\cdot0{,}3\cdot0{,}8+0{,}4\cdot0{,}7\cdot0{,}8=0{,}084+0{,}144+0{,}224=0{,}452$.

% ============================================================
\chapter{Pembahasan Bab 17: Irisan Kerucut}
% ============================================================

\kbagian{Bagian A}

\knomor{A.1}
\ksoal Fokus parabola $y^2=12x$.
\kanalisis $4p=12$.
\kbahas $p=3$, fokus $(3,0)$.

\knomor{A.2}
\ksoal Puncak elips $\dfrac{x^2}{25}+\dfrac{y^2}{9}=1$.
\kanalisis $a=5$, $b=3$.
\kbahas Puncak $(\pm5,0)$ dan $(0,\pm3)$.

\knomor{A.3}
\ksoal Asimtot hiperbola $\dfrac{x^2}{16}-\dfrac{y^2}{9}=1$.
\kanalisis $y=\pm\tfrac ba x$ dengan $a=4,b=3$.
\kbahas $y=\pm\tfrac34 x$.

\knomor{A.4}
\ksoal Parabola puncak $O$, fokus $(2,0)$.
\kanalisis $p=2$, bentuk $y^2=4px$.
\kbahas $y^2=8x$.

\knomor{A.5}
\ksoal Jenis irisan kerucut $e=1$.
\kanalisis Eksentrisitas $1$.
\kbahas Parabola.

\kbagian{Bagian B}

\knomor{B.1}
\ksoal Fokus dan eksentrisitas elips $\dfrac{x^2}{169}+\dfrac{y^2}{144}=1$.
\kanalisis $c^2=a^2-b^2$, $e=\tfrac ca$.
\kbahas $a=13$, $b=12$, $c=\sqrt{169-144}=5$. Fokus $(\pm5,0)$, $e=\tfrac{5}{13}$.

\knomor{B.2}
\ksoal Fokus hiperbola $\dfrac{x^2}{9}-\dfrac{y^2}{16}=1$.
\kanalisis $c^2=a^2+b^2$.
\kbahas $c=\sqrt{9+16}=5$, fokus $(\pm5,0)$.

\knomor{B.3}
\ksoal Elips puncak $(\pm5,0)$, fokus $(\pm3,0)$.
\kanalisis $a=5$, $c=3$, $b^2=a^2-c^2$.
\kbahas $b^2=25-9=16$. Persamaan $\dfrac{x^2}{25}+\dfrac{y^2}{16}=1$.

\knomor{B.4}
\ksoal Bentuk baku $x^2+4y^2-4x+8y+4=0$.
\kanalisis Melengkapkan kuadrat.
\kbahas $(x-2)^2+4(y+1)^2=4\Rightarrow\dfrac{(x-2)^2}{4}+(y+1)^2=1$. Pusat $(2,-1)$.

\knomor{B.5}
\ksoal Puncak dan fokus $(y-1)^2=8(x-2)$.
\kanalisis Bentuk $(y-k)^2=4p(x-h)$.
\kbahas Puncak $(2,1)$, $4p=8\Rightarrow p=2$, fokus $(2+2,1)=(4,1)$.

\knomor{B.6}
\ksoal Direktriks parabola $y^2=-16x$.
\kanalisis $4p=-16\Rightarrow p=-4$; direktriks $x=-p$.
\kbahas $x=4$.

\kbagian{Bagian C}

\knomor{C.1}
\ksoal Buktikan jumlah jarak ke dua fokus pada $\tfrac{x^2}{25}+\tfrac{y^2}{9}=1$ adalah $10$.
\kanalisis Definisi elips: jumlah jarak $=2a$.
\kbahas $a=5$, fokus $(\pm4,0)$ ($c=4$). Menurut definisi elips, untuk setiap titik di elips jumlah jaraknya ke kedua fokus $=2a=10$. (Sebagai cek: di puncak $(5,0)$, jarak $1+9=10$.) $\blacksquare$

\knomor{C.2}
\ksoal Titik potong $y=x+1$ dengan $y^2=4x$.
\kanalisis Substitusi.
\kbahas $(x+1)^2=4x\Rightarrow x^2-2x+1=0\Rightarrow(x-1)^2=0\Rightarrow x=1$, $y=2$. Garis menyinggung di $(1,2)$.

\knomor{C.3}
\ksoal Garis singgung $y^2=8x$ di $(2,4)$.
\kanalisis Rumus singgung $yy_1=2p(x+x_1)$ dengan $4p=8$.
\kbahas $2p=4$, jadi $4y=4(x+2)\Rightarrow y=x+2$.

\knomor{C.4}
\ksoal Elips sumbu mayor $10$, $e=0{,}6$.
\kanalisis $a=5$, $c=ea$, $b^2=a^2-c^2$.
\kbahas $c=0{,}6\cdot5=3$, $b^2=25-9=16$. Persamaan $\dfrac{x^2}{25}+\dfrac{y^2}{16}=1$.

\knomor{C.5}
\ksoal Hiperbola asimtot $y=\pm\tfrac34x$, melalui $(4,0)$.
\kanalisis $(4,0)$ adalah puncak $\Rightarrow a=4$; $\tfrac ba=\tfrac34$.
\kbahas $a=4$, $b=3$. Persamaan $\dfrac{x^2}{16}-\dfrac{y^2}{9}=1$.

\knomor{C.6}
\ksoal Orbit: perihelion $2$ SA, aphelion $8$ SA.
\kanalisis $a-c=2$, $a+c=8$.
\kbahas Jumlahkan: $2a=10\Rightarrow a=5$; selisih: $2c=6\Rightarrow c=3$. Eksentrisitas $e=\tfrac ca=\tfrac35=0{,}6$.

\kbagian{Bagian D}

\knomor{D.1}
\ksoal Turunkan persamaan elips dari definisi jumlah jarak $=2a$.
\kanalisis Tulis $\sqrt{(x+c)^2+y^2}+\sqrt{(x-c)^2+y^2}=2a$, isolasi dan kuadratkan dua kali.
\kbahas Pindahkan satu akar lalu kuadratkan: $\sqrt{(x+c)^2+y^2}=2a-\sqrt{(x-c)^2+y^2}$. Setelah dikuadratkan dan disederhanakan diperoleh $a\sqrt{(x-c)^2+y^2}=a^2-cx$; kuadratkan lagi: $a^2[(x-c)^2+y^2]=(a^2-cx)^2$, yang menyederhana menjadi $(a^2-c^2)x^2+a^2y^2=a^2(a^2-c^2)$. Dengan $b^2=a^2-c^2$, bagi $a^2b^2$: $\dfrac{x^2}{a^2}+\dfrac{y^2}{b^2}=1$. $\blacksquare$

\knomor{D.2}
\ksoal Mengapa sinar sejajar sumbu parabola dipantulkan ke fokus.
\kanalisis Sifat geometris garis singgung parabola dan hukum pemantulan.
\kbahas Pada tiap titik parabola, garis singgung membentuk sudut sama terhadap (i) garis sejajar sumbu dan (ii) garis ke fokus. Menurut hukum pemantulan (sudut datang $=$ sudut pantul), sinar yang datang sejajar sumbu dipantulkan tepat menuju fokus. Inilah dasar antena dan reflektor parabola. $\blacksquare$

\knomor{D.3}
\ksoal Parametrisasi $(a\cos t,b\sin t)$ dan luas elips.
\kanalisis Substitusi ke persamaan; luas via integral.
\kbahas $\dfrac{(a\cos t)^2}{a^2}+\dfrac{(b\sin t)^2}{b^2}=\cos^2 t+\sin^2 t=1$, jadi titik selalu di elips. Luasnya $\displaystyle\int$ memberi $L=\pi ab$ (elips adalah lingkaran satuan yang diregang $a$ dan $b$ kali, sehingga luas $\pi\cdot a\cdot b$).

\knomor{D.4}
\ksoal Buktikan rasio fokus--direktriks $=e\in(0,1)$ menghasilkan elips.
\kanalisis Tulis jarak ke fokus $=e\times$ jarak ke direktriks, lalu kuadratkan.
\kbahas Ambil fokus $(c,0)$ dan direktriks $x=\tfrac ae$ ($a=\tfrac ce$). Syarat $\sqrt{(x-c)^2+y^2}=e\left|x-\tfrac ae\right|$. Kuadratkan dan sederhanakan; suku $x^2(1-e^2)$ muncul, dan karena $0<e<1$ maka $1-e^2>0$, menghasilkan persamaan berbentuk $\dfrac{x^2}{a^2}+\dfrac{y^2}{a^2(1-e^2)}=1$ --- sebuah elips dengan $b^2=a^2(1-e^2)$. $\blacksquare$

\knomor{D.5}
\ksoal Peran $B^2-4AC$ pada $Ax^2+Bxy+Cy^2+Dx+Ey+F=0$.
\kanalisis Diskriminan irisan kerucut.
\kbahas Tanda $B^2-4AC$ menentukan jenis: $B^2-4AC<0$ \textbf{elips} (atau lingkaran bila $A=C,B=0$); $=0$ \textbf{parabola}; $>0$ \textbf{hiperbola}. Suku $Bxy$ menandakan rotasi sumbu.

\kbagian{Challenge Corner}

\knomor{1}
\ksoal Lintasan komet parabola, jarak terdekat $r_0$.
\kanalisis Puncak parabola di perihelion; jarak puncak ke fokus $=p$.
\kbahas Letakkan puncak di $O$ dengan fokus (Matahari) di $(r_0,0)$, sehingga $p=r_0$ dan lintasannya $y^2=4r_0x$. Karena eksentrisitas parabola $e=1$, kurva terbuka tak berhingga; komet hanya melintas sekali dan \textbf{tak pernah kembali} (tidak ada orbit tertutup seperti elips).

\knomor{2}
\ksoal Panjang latus rectum elips $\tfrac{x^2}{a^2}+\tfrac{y^2}{b^2}=1$.
\kanalisis Substitusi $x=c$ (di fokus) ke persamaan.
\kbahas Di $x=c$: $\dfrac{c^2}{a^2}+\dfrac{y^2}{b^2}=1\Rightarrow y^2=b^2\left(1-\tfrac{c^2}{a^2}\right)=b^2\cdot\dfrac{a^2-c^2}{a^2}=\dfrac{b^4}{a^2}$ (karena $a^2-c^2=b^2$). Maka $y=\pm\dfrac{b^2}{a}$, sehingga panjang latus rectum $=\dfrac{2b^2}{a}$. $\blacksquare$

% ============================================================
\chapter{Pembahasan Bab 18: Limit Fungsi}
% ============================================================

\kbagian{Bagian A}

\knomor{A.1}
\ksoal $\lim_{x\to2}(3x+1)$.
\kanalisis Fungsi kontinu, substitusi langsung.
\kbahas $3(2)+1=7$.

\knomor{A.2}
\ksoal $\lim_{x\to3}\dfrac{x^2-9}{x-3}$.
\kanalisis Faktorkan ($\tfrac00$).
\kbahas $\dfrac{(x-3)(x+3)}{x-3}=x+3\to6$.

\knomor{A.3}
\ksoal $\lim_{x\to\infty}\dfrac{2x+1}{x-3}$.
\kanalisis Bagi pangkat tertinggi.
\kbahas $\dfrac{2+1/x}{1-3/x}\to2$.

\knomor{A.4}
\ksoal $\lim_{x\to1}(x^2+2x)$.
\kanalisis Substitusi.
\kbahas $1+2=3$.

\knomor{A.5}
\ksoal $\lim_{x\to0}\dfrac{\sin x}{x}$.
\kanalisis Limit istimewa.
\kbahas $1$.

\kbagian{Bagian B}

\knomor{B.1}
\ksoal $\lim_{x\to2}\dfrac{x^2-4}{x-2}$.
\kanalisis Faktorkan.
\kbahas $x+2\to4$.

\knomor{B.2}
\ksoal $\lim_{x\to\infty}\dfrac{3x^2+2x}{x^2-1}$.
\kanalisis Bagi $x^2$.
\kbahas $\dfrac{3+2/x}{1-1/x^2}\to3$.

\knomor{B.3}
\ksoal $\lim_{x\to0}\dfrac{\sqrt{x+4}-2}{x}$.
\kanalisis Kalikan sekawan.
\kbahas $\dfrac{x}{x(\sqrt{x+4}+2)}=\dfrac{1}{\sqrt{x+4}+2}\to\dfrac14$.

\knomor{B.4}
\ksoal $\lim_{x\to\infty}\left(\sqrt{x^2+x}-x\right)$.
\kanalisis Kalikan sekawan, bagi $x$.
\kbahas $\dfrac{x}{\sqrt{x^2+x}+x}=\dfrac{1}{\sqrt{1+1/x}+1}\to\dfrac12$.

\knomor{B.5}
\ksoal $\lim_{x\to3}\dfrac{x-3}{|x-3|}$.
\kanalisis Tinjau limit kiri dan kanan.
\kbahas Kanan ($x>3$): $\tfrac{x-3}{x-3}=1$; kiri ($x<3$): $\tfrac{x-3}{-(x-3)}=-1$. Karena berbeda, limit \textbf{tidak ada}.

\knomor{B.6}
\ksoal $\lim_{x\to0}\dfrac{\sin3x}{2x}$.
\kanalisis Sesuaikan ke bentuk $\tfrac{\sin u}{u}$.
\kbahas $\dfrac{3}{2}\cdot\dfrac{\sin3x}{3x}\to\dfrac32$.

\knomor{B.7}
\ksoal $\lim_{x\to3}\dfrac{x^2-9}{x-3}$ dengan L'Hôpital.
\kanalisis Bentuk $\tfrac00$; turunkan pembilang dan penyebut terpisah.
\kbahas $\lim_{x\to3}\dfrac{2x}{1}=6$.

\knomor{B.8}
\ksoal $\lim_{x\to0}\dfrac{\sin x}{x}$ dengan L'Hôpital.
\kanalisis Bentuk $\tfrac00$.
\kbahas $\lim_{x\to0}\dfrac{\cos x}{1}=1$.

\kbagian{Bagian C}

\knomor{C.1}
\ksoal $\lim_{x\to1}\dfrac{x^3-1}{x-1}$.
\kanalisis Faktorkan $x^3-1=(x-1)(x^2+x+1)$.
\kbahas $x^2+x+1\to3$.

\knomor{C.2}
\ksoal $\lim_{x\to0}\dfrac{1-\cos x}{x^2}$.
\kanalisis Kalikan $\tfrac{1+\cos x}{1+\cos x}$.
\kbahas $\dfrac{\sin^2x}{x^2(1+\cos x)}=\left(\dfrac{\sin x}{x}\right)^2\dfrac{1}{1+\cos x}\to1\cdot\dfrac12=\dfrac12$.

\knomor{C.3}
\ksoal $\lim_{x\to4}\dfrac{x-4}{\sqrt{x}-2}$.
\kanalisis Kalikan sekawan $\sqrt x+2$ (atau faktorkan $x-4=(\sqrt x-2)(\sqrt x+2)$).
\kbahas $\sqrt{x}+2\to4$.

\knomor{C.4}
\ksoal $\lim_{x\to\infty}\left(\sqrt{4x^2+x}-2x\right)$.
\kanalisis Sekawan, bagi $x$.
\kbahas $\dfrac{x}{\sqrt{4x^2+x}+2x}=\dfrac{1}{\sqrt{4+1/x}+2}\to\dfrac{1}{4}$.

\knomor{C.5}
\ksoal $\lim_{h\to0}\dfrac{(2+h)^2-4}{h}$.
\kanalisis Jabarkan; ini turunan $x^2$ di $x=2$.
\kbahas $\dfrac{4+4h+h^2-4}{h}=4+h\to4$. Maknanya: gradien garis singgung $y=x^2$ di $x=2$ adalah $4$.

\knomor{C.6}
\ksoal Tentukan $a$ agar $f$ kontinu di $x=1$.
\kanalisis Kontinu $\Leftrightarrow a=\lim_{x\to1}\tfrac{x^2-1}{x-1}$.
\kbahas $\lim_{x\to1}\dfrac{(x-1)(x+1)}{x-1}=\lim(x+1)=2$. Jadi $a=2$.

\knomor{C.7}
\ksoal $\lim_{x\to0}\dfrac{x-\sin x}{x^3}$ dengan L'Hôpital.
\kanalisis Bentuk $\tfrac00$; terapkan L'Hôpital berulang.
\kbahas $\dfrac{x-\sin x}{x^3}\xrightarrow{\text{L'H}}\dfrac{1-\cos x}{3x^2}\xrightarrow{\text{L'H}}\dfrac{\sin x}{6x}\xrightarrow{\text{L'H}}\dfrac{\cos x}{6}\to\dfrac16$.

\knomor{C.8}
\ksoal $\lim_{x\to\infty}\dfrac{\ln x}{x}$ dengan L'Hôpital.
\kanalisis Bentuk $\tfrac\infty\infty$.
\kbahas $\lim_{x\to\infty}\dfrac{1/x}{1}=\lim_{x\to\infty}\dfrac1x=0$.

\kbagian{Bagian D}

\knomor{D.1}
\ksoal Buktikan $\lim_{x\to0}x^2\sin\tfrac1x=0$.
\kanalisis Teorema apit dengan $-1\le\sin\tfrac1x\le1$.
\kbahas $-x^2\le x^2\sin\tfrac1x\le x^2$. Karena $\lim_{x\to0}(-x^2)=\lim_{x\to0}x^2=0$, menurut teorema apit $\lim_{x\to0}x^2\sin\tfrac1x=0$. $\blacksquare$

\knomor{D.2}
\ksoal $\lim_{x\to\infty}(1+\tfrac1x)^x=e$; tentukan $\lim_{x\to\infty}(1+\tfrac2x)^x$.
\kanalisis Definisi $e$; tulis ulang eksponennya.
\kbahas $(1+\tfrac1x)^x\to e$ adalah definisi bilangan $e$. Untuk yang kedua: $\left(1+\tfrac2x\right)^x=\left[\left(1+\tfrac{2}{x}\right)^{x/2}\right]^2\to e^2$.

\knomor{D.3}
\ksoal $\lim_{x\to0}\dfrac{\tan x-\sin x}{x^3}$.
\kanalisis Tulis $\tan x-\sin x=\sin x\cdot\dfrac{1-\cos x}{\cos x}$.
\kbahas $\dfrac{\sin x(1-\cos x)}{x^3\cos x}=\dfrac{\sin x}{x}\cdot\dfrac{1-\cos x}{x^2}\cdot\dfrac{1}{\cos x}\to1\cdot\dfrac12\cdot1=\dfrac12$.

\knomor{D.4}
\ksoal Buktikan $\lim_{x\to0}\dfrac{\sin x}{x}=1$ dengan luas.
\kanalisis Bandingkan luas segitiga kecil, juring, dan segitiga besar pada lingkaran satuan.
\kbahas Untuk $0<x<\tfrac\pi2$, luas segitiga $\le$ luas juring $\le$ luas segitiga luar memberi $\tfrac12\sin x\le\tfrac12 x\le\tfrac12\tan x$. Bagi $\tfrac12\sin x$: $1\le\dfrac{x}{\sin x}\le\dfrac{1}{\cos x}$. Saat $x\to0$, $\cos x\to1$, sehingga menurut teorema apit $\dfrac{\sin x}{x}\to1$. $\blacksquare$

\knomor{D.5}
\ksoal $\lim_{x\to0^+}x\ln x$ (bentuk $0\cdot\infty$) dengan L'Hôpital.
\kanalisis Ubah ke $\tfrac\infty\infty$: tulis $x\ln x=\dfrac{\ln x}{1/x}$.
\kbahas $\lim_{x\to0^+}\dfrac{\ln x}{1/x}\xrightarrow{\text{L'H}}\lim_{x\to0^+}\dfrac{1/x}{-1/x^2}=\lim_{x\to0^+}(-x)=0$.

\kbagian{Challenge Corner}

\knomor{1}
\ksoal $\lim_{x\to0}\dfrac{\sqrt{1+x}-\sqrt{1-x}}{x}$.
\kanalisis Kalikan sekawan.
\kbahas $\dfrac{(1+x)-(1-x)}{x(\sqrt{1+x}+\sqrt{1-x})}=\dfrac{2x}{x(\sqrt{1+x}+\sqrt{1-x})}=\dfrac{2}{\sqrt{1+x}+\sqrt{1-x}}\to\dfrac{2}{2}=1$.

\knomor{2}
\ksoal Tentukan $a,b$ agar $\lim_{x\to\infty}\left(\dfrac{x^2+1}{x+1}-ax-b\right)=0$.
\kanalisis Bagi bersusun $\dfrac{x^2+1}{x+1}$.
\kbahas $\dfrac{x^2+1}{x+1}=x-1+\dfrac{2}{x+1}$. Maka $\dfrac{x^2+1}{x+1}-ax-b=(1-a)x-(1+b)+\dfrac{2}{x+1}$. Agar limitnya $0$: $a=1$ dan $b=-1$.

% ============================================================
\chapter{Pembahasan Bab 19: Turunan Fungsi}
% ============================================================

\kbagian{Bagian A}

\knomor{A.1}
\ksoal $f(x)=3x^2-4x+1$.
\kanalisis Aturan pangkat suku demi suku.
\kbahas $f'(x)=6x-4$.

\knomor{A.2}
\ksoal $f(x)=x^3-2x^2+5$.
\kanalisis Aturan pangkat.
\kbahas $f'(x)=3x^2-4x$.

\knomor{A.3}
\ksoal $f(x)=5x^4$.
\kanalisis $\tfrac{d}{dx}x^n=nx^{n-1}$.
\kbahas $f'(x)=20x^3$.

\knomor{A.4}
\ksoal Gradien $y=x^2-3x$ di $x=2$.
\kanalisis $f'(2)$.
\kbahas $f'(x)=2x-3$, $f'(2)=1$.

\knomor{A.5}
\ksoal $f(x)=\sqrt{x}$.
\kanalisis $\sqrt x=x^{1/2}$.
\kbahas $f'(x)=\tfrac12x^{-1/2}=\dfrac{1}{2\sqrt{x}}$.

\kbagian{Bagian B}

\knomor{B.1}
\ksoal $f(x)=(x^2+1)(x-3)$.
\kanalisis Aturan hasil kali.
\kbahas $f'(x)=2x(x-3)+(x^2+1)=3x^2-6x+1$.

\knomor{B.2}
\ksoal $f(x)=\dfrac{x}{x+1}$.
\kanalisis Aturan hasil bagi.
\kbahas $f'(x)=\dfrac{(1)(x+1)-x(1)}{(x+1)^2}=\dfrac{1}{(x+1)^2}$.

\knomor{B.3}
\ksoal $f(x)=(2x+1)^5$.
\kanalisis Aturan rantai.
\kbahas $f'(x)=5(2x+1)^4\cdot2=10(2x+1)^4$.

\knomor{B.4}
\ksoal $f(x)=\sin(3x)$.
\kanalisis Rantai dengan $(\sin u)'=\cos u$.
\kbahas $f'(x)=3\cos(3x)$.

\knomor{B.5}
\ksoal $f(x)=e^{2x}$.
\kanalisis Rantai dengan $(e^u)'=e^u$.
\kbahas $f'(x)=2e^{2x}$.

\knomor{B.6}
\ksoal $f(x)=\ln(x^2+1)$.
\kanalisis Rantai dengan $(\ln u)'=\tfrac{u'}{u}$.
\kbahas $f'(x)=\dfrac{2x}{x^2+1}$.

\knomor{B.7}
\ksoal Titik stasioner $f(x)=x^3-3x$.
\kanalisis $f'(x)=0$.
\kbahas $f'(x)=3x^2-3=0\Rightarrow x=\pm1$. Titik $(1,-2)$ dan $(-1,2)$.

\kbagian{Bagian C}

\knomor{C.1}
\ksoal Kawat $20$ cm, persegi panjang luas maksimum.
\kanalisis Keliling $2(p+l)=20$; nyatakan luas dalam satu peubah.
\kbahas $p+l=10$, $L=p(10-p)$, $L'(p)=10-2p=0\Rightarrow p=5$. Jadi persegi $5\times5$ dengan luas maksimum $25$ cm$^2$.

\knomor{C.2}
\ksoal Garis singgung $y=x^2+1$ di $x=1$.
\kanalisis Gradien $f'(1)$, titik $(1,2)$.
\kbahas $f'(x)=2x$, $f'(1)=2$. Garis: $y-2=2(x-1)\Rightarrow y=2x$.

\knomor{C.3}
\ksoal $s(t)=t^3-6t^2+9t$: kecepatan dan saat berhenti.
\kanalisis $v=s'(t)$; berhenti saat $v=0$.
\kbahas $v(t)=3t^2-12t+9=3(t-1)(t-3)$. $v=0$ saat $t=1$ s dan $t=3$ s.

\knomor{C.4}
\ksoal $f(x)=\sin^2 x$.
\kanalisis Rantai: $(\sin x)^2$.
\kbahas $f'(x)=2\sin x\cos x=\sin 2x$.

\knomor{C.5}
\ksoal Maksimum $f(x)=-x^2+4x+1$.
\kanalisis $f'(x)=0$; parabola terbuka ke bawah.
\kbahas $f'(x)=-2x+4=0\Rightarrow x=2$; $f(2)=5$ (maksimum).

\knomor{C.6}
\ksoal $f(x)=\sqrt{x^2+1}$.
\kanalisis Rantai.
\kbahas $f'(x)=\dfrac{2x}{2\sqrt{x^2+1}}=\dfrac{x}{\sqrt{x^2+1}}$.

\kbagian{Bagian D}

\knomor{D.1}
\ksoal Buktikan turunan $x^2$ adalah $2x$ dari definisi.
\kanalisis Pakai $f'(x)=\lim_{h\to0}\tfrac{f(x+h)-f(x)}{h}$.
\kbahas $\dfrac{(x+h)^2-x^2}{h}=\dfrac{2xh+h^2}{h}=2x+h\to2x$. $\blacksquare$

\knomor{D.2}
\ksoal Buktikan $(fg)'=f'g+fg'$.
\kanalisis Tambah-kurang suku $f(x+h)g(x)$ pada pembilang.
\kbahas
\[
\dfrac{f(x+h)g(x+h)-f(x)g(x)}{h}=\dfrac{[f(x+h)-f(x)]g(x+h)}{h}+\dfrac{f(x)[g(x+h)-g(x)]}{h}.
\]
Saat $h\to0$ menjadi $f'(x)g(x)+f(x)g'(x)$. $\blacksquare$

\knomor{D.3}
\ksoal Silinder volume tetap, luas minimum $\Rightarrow h=2r$.
\kanalisis $V=\pi r^2h$ tetap; $S=2\pi r^2+2\pi rh$; nyatakan $S(r)$ lalu turunkan.
\kbahas Dari $V=\pi r^2h$, $h=\tfrac{V}{\pi r^2}$. Maka $S=2\pi r^2+\dfrac{2V}{r}$, $S'(r)=4\pi r-\dfrac{2V}{r^2}=0\Rightarrow4\pi r^3=2V=2\pi r^2h\Rightarrow 2r=h$. Jadi tinggi optimal $=2\times$ jari-jari. $\blacksquare$

\knomor{D.4}
\ksoal Buktikan $(\sin x)'=\cos x$ dari definisi.
\kanalisis Pakai rumus selisih dan dua limit istimewa.
\kbahas
\[
\dfrac{\sin(x+h)-\sin x}{h}=\dfrac{\sin x\cos h+\cos x\sin h-\sin x}{h}=\sin x\cdot\dfrac{\cos h-1}{h}+\cos x\cdot\dfrac{\sin h}{h}.
\]
Saat $h\to0$: $\sin x\cdot0+\cos x\cdot1=\cos x$. $\blacksquare$

\knomor{D.5}
\ksoal Mengapa $f'(x)=0$ pada interval $\Rightarrow f$ konstan.
\kanalisis Teorema Nilai Rata-Rata (MVT).
\kbahas Untuk dua titik $a<b$ pada interval, MVT memberi $c$ dengan $f(b)-f(a)=f'(c)(b-a)$. Karena $f'(c)=0$, maka $f(b)=f(a)$ untuk setiap pasangan titik, sehingga $f$ konstan. $\blacksquare$

\kbagian{Challenge Corner}

\knomor{1}
\ksoal Turunan ke-$n$ dari $f(x)=\dfrac1x$.
\kanalisis Turunkan beberapa kali, amati pola.
\kbahas $f'=-x^{-2}$, $f''=2x^{-3}$, $f'''=-6x^{-4}$, \dots Polanya $f^{(n)}(x)=\dfrac{(-1)^n\,n!}{x^{n+1}}$.

\knomor{2}
\ksoal Minimum $f(x)=x+\dfrac1x$ ($x>0$).
\kanalisis $f'(x)=0$; bandingkan AM--GM.
\kbahas $f'(x)=1-\dfrac{1}{x^2}=0\Rightarrow x=1$ (untuk $x>0$), $f(1)=2$ minimum. Sesuai AM--GM: $x+\dfrac1x\ge2\sqrt{x\cdot\tfrac1x}=2$, dengan kesamaan saat $x=1$. $\blacksquare$

% ============================================================
\chapter{Pembahasan Bab 20: Integral}
% ============================================================

\kbagian{Bagian A}

\knomor{A.1}
\ksoal $\int x^2\,dx$.
\kanalisis Aturan pangkat integral.
\kbahas $\dfrac{x^3}{3}+C$.

\knomor{A.2}
\ksoal $\int(2x+1)\,dx$.
\kanalisis Integralkan tiap suku.
\kbahas $x^2+x+C$.

\knomor{A.3}
\ksoal $\int_0^2(2x+1)\,dx$.
\kanalisis $[x^2+x]_0^2$.
\kbahas $4+2-0=6$.

\knomor{A.4}
\ksoal $\int_1^2 x^2\,dx$.
\kanalisis $[\tfrac{x^3}{3}]_1^2$.
\kbahas $\tfrac83-\tfrac13=\tfrac73$.

\knomor{A.5}
\ksoal $\int 3x^2\,dx$.
\kanalisis Pangkat.
\kbahas $x^3+C$.

\kbagian{Bagian B}

\knomor{B.1}
\ksoal $\int(x^3-2x^2+5)\,dx$.
\kanalisis Integralkan tiap suku.
\kbahas $\dfrac{x^4}{4}-\dfrac{2x^3}{3}+5x+C$.

\knomor{B.2}
\ksoal $\int_0^1(3x^2+2x)\,dx$.
\kanalisis $[x^3+x^2]_0^1$.
\kbahas $1+1=2$.

\knomor{B.3}
\ksoal $\int\dfrac{x^2+1}{\sqrt{x}}\,dx$.
\kanalisis Tulis $x^{3/2}+x^{-1/2}$.
\kbahas $\int(x^{3/2}+x^{-1/2})\,dx=\dfrac{2}{5}x^{5/2}+2x^{1/2}+C$.

\knomor{B.4}
\ksoal $\int_0^{\pi}\sin x\,dx$.
\kanalisis Antiturunan $-\cos x$.
\kbahas $[-\cos x]_0^\pi=-\cos\pi+\cos0=1+1=2$.

\knomor{B.5}
\ksoal $\int 2x(x^2+1)^3\,dx$.
\kanalisis Substitusi $u=x^2+1$, $du=2x\,dx$.
\kbahas $\int u^3\,du=\dfrac{u^4}{4}+C=\dfrac{(x^2+1)^4}{4}+C$.

\knomor{B.6}
\ksoal Luas di bawah $y=x^2$ dari $0$ sampai $3$.
\kanalisis $\int_0^3 x^2\,dx$.
\kbahas $\left[\tfrac{x^3}{3}\right]_0^3=9$.

\kbagian{Bagian C}

\knomor{C.1}
\ksoal Luas antara $y=x^2$ dan $y=x$, $0\le x\le1$.
\kanalisis $\int_0^1(\text{atas}-\text{bawah})$; di sini $x\ge x^2$.
\kbahas $\int_0^1(x-x^2)\,dx=\left[\tfrac{x^2}{2}-\tfrac{x^3}{3}\right]_0^1=\tfrac12-\tfrac13=\dfrac16$.

\knomor{C.2}
\ksoal $\int x\,e^{x^2}\,dx$.
\kanalisis Substitusi $u=x^2$, $du=2x\,dx$.
\kbahas $\tfrac12\int e^u\,du=\tfrac12e^{x^2}+C$.

\knomor{C.3}
\ksoal $\dfrac{d}{dx}\int_0^x\sin(t^2)\,dt$.
\kanalisis Teorema Dasar Kalkulus bagian 1.
\kbahas $\sin(x^2)$.

\knomor{C.4}
\ksoal Volume putar $y=\sqrt x$, $0\le x\le4$, terhadap sumbu-$x$.
\kanalisis $V=\pi\int_0^4 y^2\,dx$.
\kbahas $V=\pi\int_0^4 x\,dx=\pi\left[\tfrac{x^2}{2}\right]_0^4=8\pi$.

\knomor{C.5}
\ksoal $\int_{-2}^2 x^3\,dx$.
\kanalisis $x^3$ fungsi ganjil pada selang simetris.
\kbahas Karena $x^3$ ganjil, integral pada $[-2,2]$ saling menghapus: hasilnya $0$. (Cek: $[\tfrac{x^4}{4}]_{-2}^2=4-4=0$.)

\knomor{C.6}
\ksoal $v(t)=3t^2-6t$, $s(0)=2$; tentukan $s(t)$.
\kanalisis $s(t)=\int v\,dt$, tentukan $C$ dari syarat awal.
\kbahas $s(t)=t^3-3t^2+C$; $s(0)=C=2$. Jadi $s(t)=t^3-3t^2+2$.

\kbagian{Bagian D}

\knomor{D.1}
\ksoal Nyatakan kedua bagian Teorema Dasar Kalkulus.
\kanalisis Hubungan turunan dan integral.
\kbahas \textbf{Bagian 1:} jika $G(x)=\int_a^x f(t)\,dt$ maka $G'(x)=f(x)$ --- menurunkan integral mengembalikan integran. \textbf{Bagian 2:} $\int_a^b f(x)\,dx=F(b)-F(a)$ untuk sembarang antiturunan $F$. Keduanya menunjukkan integral adalah kebalikan turunan.

\knomor{D.2}
\ksoal $\int_0^{\pi/2}\sin x\cos x\,dx$.
\kanalisis Substitusi $u=\sin x$.
\kbahas $du=\cos x\,dx$; batas $x:0\to\tfrac\pi2$ menjadi $u:0\to1$. $\int_0^1 u\,du=\left[\tfrac{u^2}{2}\right]_0^1=\dfrac12$.

\knomor{D.3}
\ksoal Buktikan $\int_a^b f=-\int_b^a f$ dan $\int_a^a f=0$.
\kanalisis Pakai $\int_a^b f=F(b)-F(a)$.
\kbahas $\int_b^a f=F(a)-F(b)=-(F(b)-F(a))=-\int_a^b f$. Dan $\int_a^a f=F(a)-F(a)=0$. $\blacksquare$

\knomor{D.4}
\ksoal Turunkan volume bola $\tfrac43\pi R^3$ dengan integral putar.
\kanalisis Putar setengah lingkaran $y=\sqrt{R^2-x^2}$ terhadap sumbu-$x$.
\kbahas $V=\pi\int_{-R}^{R}(R^2-x^2)\,dx=\pi\left[R^2x-\tfrac{x^3}{3}\right]_{-R}^{R}=\pi\left(2R^3-\tfrac{2R^3}{3}\right)=\dfrac{4}{3}\pi R^3$. $\blacksquare$

\kbagian{Challenge Corner}

\knomor{1}
\ksoal $\int_0^1\dfrac{x}{1+x^2}\,dx$.
\kanalisis Substitusi $u=1+x^2$.
\kbahas $du=2x\,dx$; batas $u:1\to2$. $\tfrac12\int_1^2\dfrac{du}{u}=\tfrac12[\ln u]_1^2=\dfrac12\ln2$.

\knomor{2}
\ksoal Buktikan $\int_0^a f(x)\,dx=\int_0^a f(a-x)\,dx$; hitung $\int_0^{\pi/2}\dfrac{\sin x}{\sin x+\cos x}\,dx$.
\kanalisis Substitusi $u=a-x$ untuk identitas; terapkan pada integral.
\kbahas Dengan $u=a-x$, $du=-dx$, batas terbalik memberi $\int_0^a f(a-x)\,dx$ --- terbukti. Misal $I=\int_0^{\pi/2}\dfrac{\sin x}{\sin x+\cos x}\,dx$. Pakai $x\to\tfrac\pi2-x$: $\sin\to\cos$, sehingga $I=\int_0^{\pi/2}\dfrac{\cos x}{\cos x+\sin x}\,dx$. Jumlahkan kedua bentuk: $2I=\int_0^{\pi/2}1\,dx=\dfrac\pi2$, jadi $I=\dfrac\pi4$. $\blacksquare$

% ============================================================
\chapter{Pembahasan Bab 21: Distribusi Binomial}
% ============================================================

\kbagian{Bagian A}

\knomor{A.1}
\ksoal $q$ jika $p=0{,}3$.
\kanalisis $q=1-p$.
\kbahas $q=0{,}7$.

\knomor{A.2}
\ksoal $P(X=2)$, $n=4$, $p=0{,}5$.
\kanalisis $\binom42(0{,}5)^2(0{,}5)^2$.
\kbahas $6\cdot(0{,}5)^4=\dfrac{6}{16}=0{,}375$.

\knomor{A.3}
\ksoal $E(X)$, $n=10$, $p=0{,}2$.
\kanalisis $E=np$.
\kbahas $10\cdot0{,}2=2$.

\knomor{A.4}
\ksoal $\operatorname{Var}(X)$, $n=10$, $p=0{,}2$.
\kanalisis $\operatorname{Var}=npq$.
\kbahas $10\cdot0{,}2\cdot0{,}8=1{,}6$.

\knomor{A.5}
\ksoal Peluang tepat $2$ gambar dari $3$ lemparan.
\kanalisis $\binom32(0{,}5)^3$.
\kbahas $3\cdot\tfrac18=\dfrac38=0{,}375$.

\knomor{A.6}
\ksoal $\sigma$, $n=100$, $p=0{,}5$.
\kanalisis $\sigma=\sqrt{npq}$.
\kbahas $\sqrt{100\cdot0{,}25}=\sqrt{25}=5$.

\knomor{A.7}
\ksoal $P(X=0)$, $n=5$, $p=0{,}2$.
\kanalisis $\binom50 p^0 q^5=q^5$.
\kbahas $(0{,}8)^5=0{,}32768$.

\knomor{A.8}
\ksoal Empat syarat percobaan binomial.
\kanalisis Sebutkan ciri.
\kbahas (1) $n$ tetap; (2) dua hasil (sukses/gagal); (3) $p$ tetap; (4) percobaan saling bebas.

\knomor{A.9}
\ksoal $E(X)$, dadu $60$ kali, sukses $=6$.
\kanalisis $p=\tfrac16$, $E=np$.
\kbahas $60\cdot\tfrac16=10$.

\knomor{A.10}
\ksoal Rumus $P(X=k)$.
\kanalisis Distribusi binomial.
\kbahas $P(X=k)=\binom nk p^k(1-p)^{n-k}$.

\kbagian{Bagian B}

\knomor{B.1}
\ksoal Peluang muncul $6$ tepat $2$ kali dari $5$ lemparan dadu.
\kanalisis $n=5$, $p=\tfrac16$, $k=2$.
\kbahas $\binom52\left(\tfrac16\right)^2\left(\tfrac56\right)^3=10\cdot\dfrac{1}{36}\cdot\dfrac{125}{216}=\dfrac{1250}{7776}\approx0{,}161$.

\knomor{B.2}
\ksoal $P(X\ge1)$, $p=0{,}4$, $n=6$.
\kanalisis Komplemen $1-P(X=0)$.
\kbahas $1-(0{,}6)^6=1-0{,}046656=0{,}9533$.

\knomor{B.3}
\ksoal $E$ jawaban benar, $10$ soal, tebak acak $4$ opsi.
\kanalisis $p=0{,}25$, $E=np$.
\kbahas $10\cdot0{,}25=2{,}5$.

\knomor{B.4}
\ksoal $P(X=4)$, $n=8$, $p=0{,}5$.
\kanalisis $\binom84(0{,}5)^8$.
\kbahas $\dfrac{70}{256}\approx0{,}273$.

\knomor{B.5}
\ksoal $\operatorname{Var}$ dan $\sigma$, $n=50$, $p=0{,}4$.
\kanalisis $\operatorname{Var}=npq$.
\kbahas $\operatorname{Var}=50\cdot0{,}4\cdot0{,}6=12$; $\sigma=\sqrt{12}\approx3{,}46$.

\kbagian{Bagian C}

\knomor{C.1}
\ksoal Peluang minimal $1$ perempuan dari $5$ anak.
\kanalisis Komplemen "semua laki-laki".
\kbahas $1-(0{,}5)^5=1-\dfrac{1}{32}=\dfrac{31}{32}$.

\knomor{C.2}
\ksoal Tunjukkan $\sum_{k=0}^n P(X=k)=1$.
\kanalisis Teorema binomial pada $(p+q)^n$.
\kbahas $\sum_{k=0}^n\binom nk p^k q^{n-k}=(p+q)^n=1^n=1$. $\blacksquare$

\knomor{C.3}
\ksoal Peluang tepat $1$ cacat dari $20$ ($p=0{,}05$).
\kanalisis $\binom{20}{1}(0{,}05)(0{,}95)^{19}$.
\kbahas $20\cdot0{,}05\cdot(0{,}95)^{19}=(0{,}95)^{19}\approx0{,}377$.

\knomor{C.4}
\ksoal Peluang tepat $2$ gambar, diketahui minimal $1$ gambar (4 koin).
\kanalisis Peluang bersyarat $\dfrac{P(X=2)}{P(X\ge1)}$.
\kbahas $P(X=2)=\dfrac{6}{16}$, $P(X\ge1)=1-\dfrac{1}{16}=\dfrac{15}{16}$. Maka $\dfrac{6/16}{15/16}=\dfrac{6}{15}=\dfrac25$.

\knomor{C.5}
\ksoal Nilai harapan pembayaran ($3$ koin, Rp$1000$/gambar).
\kanalisis $E(\text{gambar})=np$.
\kbahas $E(\text{gambar})=3\cdot\tfrac12=1{,}5$; pembayaran $=1{,}5\times1000=\text{Rp}1500$.

\kbagian{Bagian D}

\knomor{D.1}
\ksoal Buktikan $E(X)=np$.
\kanalisis Pakai $k\binom nk=n\binom{n-1}{k-1}$.
\kbahas $E(X)=\sum_{k=0}^n k\binom nk p^k q^{n-k}=\sum_{k=1}^n n\binom{n-1}{k-1}p^k q^{n-k}=np\sum_{j=0}^{n-1}\binom{n-1}{j}p^{j}q^{(n-1)-j}=np(p+q)^{n-1}=np$. $\blacksquare$

\knomor{D.2}
\ksoal Buktikan $\operatorname{Var}(X)=npq$.
\kanalisis Tulis $X=X_1+\cdots+X_n$, $X_i$ Bernoulli bebas.
\kbahas Tiap $X_i$: $E(X_i)=p$, $\operatorname{Var}(X_i)=E(X_i^2)-p^2=p-p^2=pq$. Karena $X_i$ saling bebas, variansi jumlah $=$ jumlah variansi: $\operatorname{Var}(X)=\sum_{i=1}^n pq=npq$. $\blacksquare$

\knomor{D.3}
\ksoal Taksir $P(X=0)$, $n=1000$, $p=0{,}001$ (Poisson).
\kanalisis $\lambda=np=1$; $P(X=0)\approx e^{-\lambda}$.
\kbahas $\lambda=1$, $P(X=0)\approx e^{-1}\approx0{,}368$.

\knomor{D.4}
\ksoal Taksir $P(X=50)$, $n=100$, $p=0{,}5$ (normal).
\kanalisis $\mu=50$, $\sigma=\sqrt{npq}=5$; puncak $\dfrac{1}{\sigma\sqrt{2\pi}}$.
\kbahas $P(X=50)\approx\dfrac{1}{5\sqrt{2\pi}}\approx\dfrac{1}{12{,}53}\approx0{,}0798$.

\knomor{D.5}
\ksoal Modus untuk $n=10$, $p=0{,}3$.
\kanalisis Modus $=\lfloor(n+1)p\rfloor$.
\kbahas $\lfloor 11\cdot0{,}3\rfloor=\lfloor3{,}3\rfloor=3$. Jadi $P(X=k)$ maksimum di $k=3$.

\kbagian{Challenge Corner}

\knomor{1}
\ksoal Minimal berapa lemparan agar $P(\text{minimal 1 gambar})\ge0{,}99$.
\kanalisis $1-(0{,}5)^n\ge0{,}99\Rightarrow(0{,}5)^n\le0{,}01$.
\kbahas $n\ge\dfrac{\log0{,}01}{\log0{,}5}=\dfrac{-2}{-0{,}30103}\approx6{,}64$. Karena $n$ bulat, $n=7$.

\knomor{2}
\ksoal Peluang banyak laki $=$ banyak perempuan ($n$ genap).
\kanalisis Banyak laki $=\tfrac n2$ pada binomial $p=\tfrac12$.
\kbahas $P\!\left(X=\tfrac n2\right)=\dbinom{n}{n/2}\left(\tfrac12\right)^{n}$. (Mis. $n=4$: $\tfrac{6}{16}=\tfrac38$.)

\knomor{3}
\ksoal Buktikan $\dfrac{P(X=k+1)}{P(X=k)}=\dfrac{n-k}{k+1}\cdot\dfrac pq$ dan kaitannya dengan modus.
\kanalisis Bandingkan dua suku berurutan.
\kbahas $\dfrac{\binom{n}{k+1}p^{k+1}q^{n-k-1}}{\binom nk p^k q^{n-k}}=\dfrac{\binom{n}{k+1}}{\binom nk}\cdot\dfrac pq=\dfrac{n-k}{k+1}\cdot\dfrac pq$. Selama rasio ini $>1$, peluang masih naik; saat $<1$, mulai turun. Titik balik (modus) terjadi ketika rasio melewati $1$, memberi modus $\lfloor(n+1)p\rfloor$. $\blacksquare$

% ============================================================
\chapter{Pembahasan Bab 22: Distribusi Normal dan Uji Hipotesis}
% ============================================================
\noindent\textit{Nilai tabel normal baku yang digunakan: $P(Z<1)=0{,}8413$, $P(Z<1{,}96)=0{,}975$, $P(Z<2)=0{,}9772$, $P(Z<2{,}05)\approx0{,}98$.}

\kbagian{Bagian A}

\knomor{A.1}
\ksoal Luas total di bawah kurva normal.
\kanalisis Sifat fungsi peluang.
\kbahas $1$.

\knomor{A.2}
\ksoal Skor-$z$, $x=80$, $\mu=70$, $\sigma=5$.
\kanalisis $z=\dfrac{x-\mu}{\sigma}$.
\kbahas $\dfrac{80-70}{5}=2$.

\knomor{A.3}
\ksoal Persen data dalam $\pm1\sigma$.
\kanalisis Aturan empiris.
\kbahas $\approx68\%$.

\knomor{A.4}
\ksoal Persen data dalam $\pm2\sigma$.
\kanalisis Aturan empiris.
\kbahas $\approx95\%$.

\knomor{A.5}
\ksoal Mean dan $\sigma$ normal baku.
\kanalisis Definisi $N(0,1)$.
\kbahas Mean $=0$, $\sigma=1$.

\knomor{A.6}
\ksoal Skor-$z$ untuk nilai $=$ mean.
\kanalisis $x=\mu\Rightarrow z=0$.
\kbahas $0$.

\knomor{A.7}
\ksoal Kurva normal simetris terhadap apa?
\kanalisis Pusat distribusi.
\kbahas Terhadap mean $\mu$.

\knomor{A.8}
\ksoal Tanda $z$ jika $x<\mu$.
\kanalisis Pembilang $x-\mu<0$.
\kbahas Negatif.

\knomor{A.9}
\ksoal $P(Z<0)$.
\kanalisis Simetri di $0$.
\kbahas $0{,}5$.

\knomor{A.10}
\ksoal Rumus skor-$z$.
\kanalisis Standarisasi.
\kbahas $z=\dfrac{x-\mu}{\sigma}$.

\kbagian{Bagian B}

\knomor{B.1}
\ksoal Skor-$z$, $x=86$, $\mu=70$, $\sigma=8$.
\kanalisis Substitusi.
\kbahas $\dfrac{86-70}{8}=2$.

\knomor{B.2}
\ksoal Persen tinggi antara $154$ dan $166$ ($\mu=160$, $\sigma=6$).
\kanalisis $154=\mu-\sigma$, $166=\mu+\sigma$.
\kbahas Rentang $\pm1\sigma\Rightarrow\approx68\%$.

\knomor{B.3}
\ksoal $P(Z>1)$.
\kanalisis Komplemen $1-P(Z<1)$.
\kbahas $1-0{,}8413=0{,}1587$.

\knomor{B.4}
\ksoal $P(0<Z<1)$.
\kanalisis $P(Z<1)-P(Z<0)$.
\kbahas $0{,}8413-0{,}5=0{,}3413$.

\knomor{B.5}
\ksoal Skor-$z$, skor $650$, $\mu=500$, $\sigma=100$.
\kanalisis Substitusi.
\kbahas $\dfrac{650-500}{100}=1{,}5$.

\kbagian{Bagian C}

\knomor{C.1}
\ksoal Persen siswa nilai $>85$ ($\mu=75$, $\sigma=10$).
\kanalisis $z=1$, lalu $P(Z>1)$.
\kbahas $z=1$; $P(Z>1)=1-0{,}8413=0{,}1587\approx15{,}87\%$.

\knomor{C.2}
\ksoal Taksir $P(X\ge60)$, koin $n=100$.
\kanalisis $\mu=np=50$, $\sigma=\sqrt{npq}=5$; standarkan.
\kbahas $z=\dfrac{60-50}{5}=2$; $P(Z\ge2)=1-0{,}9772=0{,}0228\approx2{,}28\%$.

\knomor{C.3}
\ksoal Batas persentil ke-$90$ ($\mu=70$, $\sigma=10$).
\kanalisis $x=\mu+z\sigma$ dengan $z=1{,}28$.
\kbahas $x=70+1{,}28\cdot10=82{,}8$.

\knomor{C.4}
\ksoal Rumuskan $H_0$, $H_1$ (berat diduga kurang dari $500$).
\kanalisis Klaim "kurang" $\Rightarrow$ uji satu arah kiri.
\kbahas $H_0:\mu=500$ \quad melawan \quad $H_1:\mu<500$.

\knomor{C.5}
\ksoal Statistik uji $z$, $n=36$, $\bar x=52$, $\mu_0=50$, $\sigma=6$.
\kanalisis $z=\dfrac{\bar x-\mu_0}{\sigma/\sqrt n}$.
\kbahas $z=\dfrac{52-50}{6/6}=\dfrac{2}{1}=2$.

\kbagian{Bagian D}

\knomor{D.1}
\ksoal Uji dua arah: $n=25$, $\bar x=495$, $\mu_0=500$, $\sigma=10$, $\alpha=0{,}05$.
\kanalisis Hitung $z$, bandingkan $\pm1{,}96$.
\kbahas $z=\dfrac{495-500}{10/\sqrt{25}}=\dfrac{-5}{2}=-2{,}5$. Karena $|-2{,}5|>1{,}96$, $z$ di daerah kritis $\Rightarrow$ \textbf{tolak} $H_0$: isi rata-rata berbeda signifikan dari $500$ mL.

\knomor{D.2}
\ksoal Buktikan $Z=\dfrac{X-\mu}{\sigma}$ punya mean $0$, variansi $1$.
\kanalisis Pakai sifat linear $E$ dan $\operatorname{Var}$.
\kbahas $E(Z)=\dfrac{E(X)-\mu}{\sigma}=\dfrac{\mu-\mu}{\sigma}=0$. $\operatorname{Var}(Z)=\dfrac{1}{\sigma^2}\operatorname{Var}(X)=\dfrac{\sigma^2}{\sigma^2}=1$. $\blacksquare$

\knomor{D.3}
\ksoal Uji satu arah: $n=49$, $\bar x=72$, $\mu_0=70$, $\sigma=7$, $\alpha=0{,}05$.
\kanalisis $z=\dfrac{\bar x-\mu_0}{\sigma/\sqrt n}$; bandingkan $1{,}645$ dan hitung $p$-value.
\kbahas $z=\dfrac{72-70}{7/7}=2$. Nilai kritis satu arah $1{,}645$; karena $2>1{,}645$ (dan $p\text{-value}=P(Z>2)=0{,}0228<0{,}05$), \textbf{tolak} $H_0$: ada peningkatan skor yang signifikan.

\knomor{D.4}
\ksoal Jelaskan galat tipe I dan II.
\kanalisis Empat kemungkinan keputusan vs kebenaran $H_0$.
\kbahas \textbf{Galat tipe I}: menolak $H_0$ padahal $H_0$ benar (peluangnya $=\alpha$). \textbf{Galat tipe II}: gagal menolak $H_0$ padahal $H_0$ salah (peluangnya $\beta$). Menurunkan $\alpha$ cenderung menaikkan $\beta$; keduanya ditekan dengan memperbesar ukuran sampel.

\knomor{D.5}
\ksoal Interval kepercayaan $95\%$, $n=100$, $\bar x=50$, $\sigma=8$.
\kanalisis $\bar x\pm1{,}96\cdot\dfrac{\sigma}{\sqrt n}$.
\kbahas Margin $=1{,}96\cdot\dfrac{8}{10}=1{,}568$. Interval $50\pm1{,}568=(48{,}43;\ 51{,}57)$.

\kbagian{Challenge Corner}

\knomor{1}
\ksoal Persen data di luar $\pm3\sigma$.
\kanalisis $100\%-99{,}7\%$.
\kbahas $\approx0{,}3\%$.

\knomor{2}
\ksoal IQ minimum untuk $2\%$ teratas ($\mu=100$, $\sigma=15$).
\kanalisis $x=\mu+z\sigma$, $z\approx2{,}05$.
\kbahas $x=100+2{,}05\cdot15=130{,}75\approx131$.

\knomor{3}
\ksoal Aproksimasi normal binomial; taksir $P(40\le X\le60)$, $n=100$, $p=0{,}5$.
\kanalisis Untuk $n$ besar, jumlah banyak peubah Bernoulli bebas mendekati normal (teorema limit pusat) dengan $\mu=np$, $\sigma=\sqrt{npq}$.
\kbahas $\mu=50$, $\sigma=5$. $z_1=\dfrac{40-50}{5}=-2$, $z_2=\dfrac{60-50}{5}=2$. $P(-2<Z<2)\approx0{,}9544$, yaitu $\approx95{,}44\%$. $\blacksquare$

% ============================================================
\chapter{Pembahasan Bab 23: Dimensi Tiga}
% ============================================================
\noindent\textit{Acuan koordinat kubus rusuk $a$: $A(0,0,0)$, $B(a,0,0)$, $C(a,a,0)$, $D(0,a,0)$, $E(0,0,a)$, $F(a,0,a)$, $G(a,a,a)$, $H(0,a,a)$.}

\kbagian{Bagian A}

\knomor{A.1}
\ksoal Diagonal sisi kubus rusuk $6$.
\kanalisis $a\sqrt2$.
\kbahas $6\sqrt2$ cm.

\knomor{A.2}
\ksoal Diagonal ruang kubus rusuk $6$.
\kanalisis $a\sqrt3$.
\kbahas $6\sqrt3$ cm.

\knomor{A.3}
\ksoal Jarak $A(1,2,3)$ ke $B(4,6,3)$.
\kanalisis Rumus jarak $3$D.
\kbahas $\sqrt{3^2+4^2+0^2}=\sqrt{25}=5$.

\knomor{A.4}
\ksoal Banyak diagonal ruang kubus.
\kanalisis Menghubungkan titik sudut berseberangan.
\kbahas $4$.

\knomor{A.5}
\ksoal Dua garis tak sebidang dan tak berpotongan.
\kanalisis Definisi kedudukan.
\kbahas Bersilangan.

\knomor{A.6}
\ksoal Jarak $O(0,0,0)$ ke $P(2,3,6)$.
\kanalisis Rumus jarak.
\kbahas $\sqrt{4+9+36}=\sqrt{49}=7$.

\knomor{A.7}
\ksoal Kedudukan $AB$ dan $DC$.
\kanalisis Rusuk berhadapan pada sisi alas.
\kbahas Sejajar.

\knomor{A.8}
\ksoal Rusuk kubus jika diagonal ruang $9\sqrt3$.
\kanalisis $a\sqrt3=9\sqrt3$.
\kbahas $a=9$.

\knomor{A.9}
\ksoal Banyak bidang melalui tiga titik tak segaris.
\kanalisis Aksioma bidang.
\kbahas Tepat satu.

\knomor{A.10}
\ksoal Jarak $(1,1,1)$ ke $(1,1,5)$.
\kanalisis Hanya beda pada $z$.
\kbahas $|5-1|=4$.

\kbagian{Bagian B}

\knomor{B.1}
\ksoal Jarak $A$ ke $C$, kubus rusuk $4$.
\kanalisis Diagonal sisi.
\kbahas $4\sqrt2$.

\knomor{B.2}
\ksoal Jarak $A$ ke $G$, kubus rusuk $4$.
\kanalisis Diagonal ruang.
\kbahas $4\sqrt3$.

\knomor{B.3}
\ksoal Jarak $A$ ke bidang $BDHF$, kubus rusuk $6$.
\kanalisis Bidang $BDHF$ memuat diagonal $BD$ dan tegak; jarak $A$ = jarak $A$ ke garis $BD$ di alas $=\tfrac12 AC$.
\kbahas $AC=6\sqrt2$, jarak $=\tfrac12\cdot6\sqrt2=3\sqrt2$.

\knomor{B.4}
\ksoal Sudut antara $AG$ dan $AB$.
\kanalisis $\vec{AG}=(a,a,a)$, $\vec{AB}=(a,0,0)$; $\cos\theta=\dfrac{|\vec{AG}\cdot\vec{AB}|}{|\vec{AG}||\vec{AB}|}$.
\kbahas $\cos\theta=\dfrac{a^2}{a\sqrt3\cdot a}=\dfrac{1}{\sqrt3}$, jadi $\theta=\arccos\tfrac{1}{\sqrt3}\approx54{,}74^\circ$.

\knomor{B.5}
\ksoal Sudut antara $AG$ dan bidang alas.
\kanalisis Proyeksi $AG$ ke alas $=AC$; $\tan\theta=\dfrac{a}{AC}$.
\kbahas $\tan\theta=\dfrac{a}{a\sqrt2}=\dfrac{1}{\sqrt2}$, jadi $\theta\approx35{,}26^\circ$.

\kbagian{Bagian C}

\knomor{C.1}
\ksoal Sudut antara bidang $BDG$ dan alas $ABCD$.
\kanalisis Garis potong $BD$; ambil titik $C$, proyeksi $G$ ke alas; $\tan\theta=\dfrac{CG}{\text{jarak }C\text{ ke }BD}$.
\kbahas Jarak $C$ ke $BD=\tfrac12 AC=\tfrac{a\sqrt2}{2}$, dan $CG=a$. $\tan\theta=\dfrac{a}{a\sqrt2/2}=\sqrt2$, jadi $\theta=\arctan\sqrt2\approx54{,}74^\circ$.

\knomor{C.2}
\ksoal Jarak $B$ ke diagonal ruang $AG$, kubus rusuk $6$.
\kanalisis $\vec{AB}=(a,0,0)$, arah $AG$: $\hat u=\tfrac1{\sqrt3}(1,1,1)$; $d^2=|\vec{AB}|^2-(\vec{AB}\cdot\hat u)^2$.
\kbahas $d^2=a^2-\left(\dfrac{a}{\sqrt3}\right)^2=a^2-\dfrac{a^2}{3}=\dfrac{2a^2}{3}$, jadi $d=a\sqrt{\tfrac23}=\dfrac{a\sqrt6}{3}$. Untuk $a=6$: $d=2\sqrt6$.

\knomor{C.3}
\ksoal Jarak $P(1,1,1)$ ke bidang $2x+y+2z=6$.
\kanalisis Rumus jarak titik--bidang.
\kbahas $d=\dfrac{|2+1+2-6|}{\sqrt{4+1+4}}=\dfrac{1}{3}$.

\knomor{C.4}
\ksoal Rusuk tegak $TA$, alas persegi sisi $6$, tinggi $TO=4$.
\kanalisis $OA=\tfrac12$ diagonal alas $=3\sqrt2$; $TA=\sqrt{TO^2+OA^2}$.
\kbahas $TA=\sqrt{4^2+(3\sqrt2)^2}=\sqrt{16+18}=\sqrt{34}$.

\knomor{C.5}
\ksoal Sudut antara $TA$ dan bidang alas.
\kanalisis $\tan\theta=\dfrac{TO}{OA}$.
\kbahas $\tan\theta=\dfrac{4}{3\sqrt2}=\dfrac{2\sqrt2}{3}\approx0{,}943$, jadi $\theta\approx43{,}3^\circ$.

\kbagian{Bagian D}

\knomor{D.1}
\ksoal Buktikan diagonal ruang kubus rusuk $a$ adalah $a\sqrt3$.
\kanalisis Dua kali Pythagoras: diagonal sisi lalu diagonal ruang.
\kbahas Diagonal sisi $AC=\sqrt{a^2+a^2}=a\sqrt2$. Diagonal ruang $AG$ tegak lurus terhadap $CG=a$: $AG=\sqrt{AC^2+CG^2}=\sqrt{2a^2+a^2}=a\sqrt3$. $\blacksquare$

\knomor{D.2}
\ksoal Turunkan jarak $P(x_0,y_0,z_0)$ ke bidang $ax+by+cz+d=0$.
\kanalisis Proyeksikan vektor dari titik bidang ke $P$ pada normal $\vec n=(a,b,c)$.
\kbahas Ambil $Q(x_1,y_1,z_1)$ pada bidang ($ax_1+by_1+cz_1+d=0$). Jarak $=$ proyeksi $\vec{QP}$ pada $\hat n$: $\dfrac{|\vec{QP}\cdot\vec n|}{|\vec n|}=\dfrac{|a(x_0-x_1)+b(y_0-y_1)+c(z_0-z_1)|}{\sqrt{a^2+b^2+c^2}}=\dfrac{|ax_0+by_0+cz_0+d|}{\sqrt{a^2+b^2+c^2}}$, memakai $-(ax_1+by_1+cz_1)=d$. $\blacksquare$

\knomor{D.3}
\ksoal Sudut antara bidang $x+y+z=1$ dan $x=0$.
\kanalisis Sudut antar bidang $=$ sudut antar normalnya.
\kbahas $\vec n_1=(1,1,1)$, $\vec n_2=(1,0,0)$. $\cos\theta=\dfrac{|1|}{\sqrt3\cdot1}=\dfrac{1}{\sqrt3}$, jadi $\theta\approx54{,}74^\circ$.

\knomor{D.4}
\ksoal Volume tetrahedron $A,B,D,E$ pada kubus rusuk $a$.
\kanalisis $V=\tfrac16|\det[\vec{AB},\vec{AD},\vec{AE}]|$.
\kbahas $\vec{AB}=(a,0,0)$, $\vec{AD}=(0,a,0)$, $\vec{AE}=(0,0,a)$; determinannya $a^3$. $V=\dfrac{a^3}{6}$.

\knomor{D.5}
\ksoal Jarak garis bersilangan $AC$ dan $BH$, kubus rusuk $a$.
\kanalisis $d=\dfrac{|\vec{AB}\cdot(\vec u\times\vec v)|}{|\vec u\times\vec v|}$ dengan $\vec u=\vec{AC}$, $\vec v=\vec{BH}$.
\kbahas $\vec u=(1,1,0)$, $\vec v=(-1,1,1)$; $\vec u\times\vec v=(1,-1,2)$, $|\cdot|=\sqrt6$. Dengan $\vec{AB}=(a,0,0)$: $d=\dfrac{|(a,0,0)\cdot(1,-1,2)|}{\sqrt6}=\dfrac{a}{\sqrt6}=\dfrac{a\sqrt6}{6}$.

\kbagian{Challenge Corner}

\knomor{1}
\ksoal Sudut antara dua diagonal ruang $AG$ dan $CE$.
\kanalisis $\vec{AG}=(1,1,1)$, $\vec{CE}=(-1,-1,1)$ (arah).
\kbahas $\cos\theta=\dfrac{\vec{AG}\cdot\vec{CE}}{|\vec{AG}||\vec{CE}|}=\dfrac{-1-1+1}{3}=-\dfrac13$, jadi $\theta=\arccos\!\left(-\tfrac13\right)\approx109{,}47^\circ$ (sudut lancipnya $\approx70{,}53^\circ$).

\knomor{2}
\ksoal Buktikan $AG\perp$ bidang $BDE$ dan jarak $A$ ke $BDE$.
\kanalisis Tentukan persamaan bidang $BDE$; bandingkan normal dengan arah $AG$.
\kbahas $B(a,0,0),D(0,a,0),E(0,0,a)$ memenuhi $x+y+z=a$, jadi normal bidang $(1,1,1)$ sejajar $\vec{AG}=(a,a,a)$ --- maka $AG\perp$ bidang $BDE$. Jarak $A(0,0,0)$: $\dfrac{|0+0+0-a|}{\sqrt3}=\dfrac{a}{\sqrt3}=\dfrac{a\sqrt3}{3}$. $\blacksquare$

\knomor{3}
\ksoal Jari-jari bola yang menyinggung semua rusuk kubus rusuk $a$.
\kanalisis Bola berpusat di pusat kubus; titik singgung di tengah tiap rusuk.
\kbahas Pusat $\left(\tfrac a2,\tfrac a2,\tfrac a2\right)$; titik tengah rusuk $AB$ adalah $\left(\tfrac a2,0,0\right)$. Jari-jari $=\sqrt{0+\left(\tfrac a2\right)^2+\left(\tfrac a2\right)^2}=\dfrac{a}{\sqrt2}=\dfrac{a\sqrt2}{2}$.

\end{document}
